% Options for packages loaded elsewhere
\PassOptionsToPackage{unicode}{hyperref}
\PassOptionsToPackage{hyphens}{url}
\PassOptionsToPackage{dvipsnames,svgnames,x11names}{xcolor}
%
\documentclass[
  10pt,
  twoside,openany]{book}
\usepackage{amsmath,amssymb}
\usepackage{iftex}
\ifPDFTeX
  \usepackage[T1]{fontenc}
  \usepackage[utf8]{inputenc}
  \usepackage{textcomp} % provide euro and other symbols
\else % if luatex or xetex
  \usepackage{unicode-math} % this also loads fontspec
  \defaultfontfeatures{Scale=MatchLowercase}
  \defaultfontfeatures[\rmfamily]{Ligatures=TeX,Scale=1}
\fi
\usepackage{lmodern}
\ifPDFTeX\else
  % xetex/luatex font selection
\fi
% Use upquote if available, for straight quotes in verbatim environments
\IfFileExists{upquote.sty}{\usepackage{upquote}}{}
\IfFileExists{microtype.sty}{% use microtype if available
  \usepackage[]{microtype}
  \UseMicrotypeSet[protrusion]{basicmath} % disable protrusion for tt fonts
}{}
\makeatletter
\@ifundefined{KOMAClassName}{% if non-KOMA class
  \IfFileExists{parskip.sty}{%
    \usepackage{parskip}
  }{% else
    \setlength{\parindent}{0pt}
    \setlength{\parskip}{6pt plus 2pt minus 1pt}}
}{% if KOMA class
  \KOMAoptions{parskip=half}}
\makeatother
\usepackage{xcolor}
\usepackage[a4paper,inner=25mm,outer=17mm,top=20mm,bottom=20mm,headsep=6mm,footskip=9mm]{geometry}
\usepackage{longtable,booktabs,array}
\usepackage{calc} % for calculating minipage widths
% Correct order of tables after \paragraph or \subparagraph
\usepackage{etoolbox}
\makeatletter
\patchcmd\longtable{\par}{\if@noskipsec\mbox{}\fi\par}{}{}
\makeatother
% Allow footnotes in longtable head/foot
\IfFileExists{footnotehyper.sty}{\usepackage{footnotehyper}}{\usepackage{footnote}}
\makesavenoteenv{longtable}
\setlength{\emergencystretch}{3em} % prevent overfull lines
\providecommand{\tightlist}{%
  \setlength{\itemsep}{0pt}\setlength{\parskip}{0pt}}
\setcounter{secnumdepth}{-\maxdimen} % remove section numbering
\ifLuaTeX
\usepackage[bidi=basic]{babel}
\else
\usepackage[bidi=default]{babel}
\fi
\babelprovide[main,import]{vietnamese}
% get rid of language-specific shorthands (see #6817):
\let\LanguageShortHands\languageshorthands
\def\languageshorthands#1{}
% ---- Discrete Mathematics Open-Book Toolkit: preamble bổ sung cho pandoc (XeLaTeX) ----
\usepackage{fontspec}
\setmainfont{Noto Serif}
\setsansfont{Noto Sans}
\setmonofont{DejaVu Sans Mono}[Scale=0.88]
\usepackage{amssymb}
\usepackage{newunicodechar}
% Ký tự Unicode trong văn bản thường -> ký hiệu toán (font chính có thể thiếu glyph)
\newunicodechar{→}{\ensuremath{\rightarrow}}
\newunicodechar{←}{\ensuremath{\leftarrow}}
\newunicodechar{⇒}{\ensuremath{\Rightarrow}}
\newunicodechar{⇔}{\ensuremath{\Leftrightarrow}}
\newunicodechar{↔}{\ensuremath{\leftrightarrow}}
\newunicodechar{✓}{\ensuremath{\checkmark}}
\newunicodechar{≤}{\ensuremath{\le}}
\newunicodechar{≥}{\ensuremath{\ge}}
\newunicodechar{≠}{\ensuremath{\ne}}
\newunicodechar{∞}{\ensuremath{\infty}}
\newunicodechar{×}{\ensuremath{\times}}
\newunicodechar{∴}{\ensuremath{\therefore}}
\newunicodechar{∀}{\ensuremath{\forall}}
\newunicodechar{∃}{\ensuremath{\exists}}
\newunicodechar{∈}{\ensuremath{\in}}
\newunicodechar{∧}{\ensuremath{\wedge}}
\newunicodechar{∨}{\ensuremath{\vee}}
\newunicodechar{¬}{\ensuremath{\neg}}
\newunicodechar{↦}{\ensuremath{\mapsto}}
\newunicodechar{−}{\ensuremath{-}}
\newunicodechar{∗}{\ensuremath{*}}
\newunicodechar{∩}{\ensuremath{\cap}}
\newunicodechar{≼}{\ensuremath{\preceq}}
\newunicodechar{∪}{\ensuremath{\cup}}
\newunicodechar{⊆}{\ensuremath{\subseteq}}
\newunicodechar{①}{\textcircled{\footnotesize 1}}
\newunicodechar{②}{\textcircled{\footnotesize 2}}
\newunicodechar{③}{\textcircled{\footnotesize 3}}
\newunicodechar{④}{\textcircled{\footnotesize 4}}
\newunicodechar{⑤}{\textcircled{\footnotesize 5}}
\newunicodechar{⑥}{\textcircled{\footnotesize 6}}
\newunicodechar{⑦}{\textcircled{\footnotesize 7}}
\newunicodechar{⑧}{\textcircled{\footnotesize 8}}
\newunicodechar{⑨}{\textcircled{\footnotesize 9}}
% Bố cục tiêu đề: mã ví dụ/dạng bài nổi bật, có kẻ ngang
\usepackage{xcolor}
\definecolor{idblue}{HTML}{1F4E79}
\usepackage{titlesec}
\titleformat{\chapter}[hang]{\normalfont\sffamily\huge\bfseries\color{idblue}}{}{0pt}{}[\vspace{2pt}\titlerule]
\titlespacing*{\chapter}{0pt}{0pt}{12pt}
\titleformat{\section}{\normalfont\sffamily\Large\bfseries\color{idblue}}{}{0pt}{}
\titleformat{\subsection}{\normalfont\sffamily\large\bfseries}{}{0pt}{}[\vspace{1pt}\titlerule]
\titlespacing*{\subsection}{0pt}{10pt plus 4pt}{4pt}
% Header/footer
\usepackage{fancyhdr}
\pagestyle{fancy}
\fancyhf{}
\fancyhead[LE,RO]{\sffamily\small\thepage}
\fancyhead[RE,LO]{\sffamily\small\nouppercase{\leftmark}}
\fancyfoot[C]{\sffamily\scriptsize Discrete Mathematics Open-Book Toolkit — Cấu trúc rời rạc}
\renewcommand{\chaptermark}[1]{\markboth{#1}{}}
\fancypagestyle{plain}{\fancyhf{}\fancyhead[LE,RO]{\sffamily\small\thepage}\fancyfoot[C]{\sffamily\scriptsize Discrete Mathematics Open-Book Toolkit — Cấu trúc rời rạc}}
% Bảng: chữ nhỏ hơn một chút, không cắt
\usepackage{etoolbox}
\AtBeginEnvironment{longtable}{\small}
\setlength{\LTpre}{4pt}\setlength{\LTpost}{4pt}
% Tránh ngắt trang ngay sau tiêu đề ví dụ; cho phép công thức dài ngắt dòng
\allowdisplaybreaks
\setlength{\emergencystretch}{3em}
\raggedbottom
\widowpenalty=10000 \clubpenalty=10000
\ifLuaTeX
  \usepackage{selnolig}  % disable illegal ligatures
\fi
\IfFileExists{bookmark.sty}{\usepackage{bookmark}}{\usepackage{hyperref}}
\IfFileExists{xurl.sty}{\usepackage{xurl}}{} % add URL line breaks if available
\urlstyle{same}
\hypersetup{
  pdftitle={Discrete Mathematics Open-Book Toolkit},
  pdflang={vi},
  colorlinks=true,
  linkcolor={idblue},
  filecolor={Maroon},
  citecolor={Blue},
  urlcolor={idblue},
  pdfcreator={LaTeX via pandoc}}

\title{Discrete Mathematics Open-Book Toolkit}
\usepackage{etoolbox}
\makeatletter
\providecommand{\subtitle}[1]{% add subtitle to \maketitle
  \apptocmd{\@title}{\par {\large #1 \par}}{}{}
}
\makeatother
\subtitle{Bộ công cụ tra cứu và giải bài Cấu trúc rời rạc (theo slide
Phạm Thế Bảo, 5 chương)}
\author{}
\date{Phiên bản 1.0 --- 2026-11-12}

\begin{document}
\frontmatter
\maketitle

{
\hypersetup{linkcolor=}
\setcounter{tocdepth}{0}
\tableofcontents
}
\mainmatter
\hypertarget{how-to-use}{%
\chapter{Cách dùng bộ tài liệu}\label{how-to-use}}

\textbf{Chuỗi tra cứu:} Đề bài → \textbf{Mục lục từ khóa} (trang đầu) →
mã dạng bài (ví dụ \texttt{REC-03}) → mục dạng bài (\emph{Kiểm tra
trước}, công thức F-, phương pháp M-) → ví dụ \texttt{REC-03-Exx} → lời
giải → kiểm tra \& bẫy \texttt{T-xx}.

\textbf{Hệ thống mã (ổn định, không phụ thuộc số trang):}

\begin{longtable}[]{@{}
  >{\raggedright\arraybackslash}p{(\columnwidth - 4\tabcolsep) * \real{0.1429}}
  >{\raggedright\arraybackslash}p{(\columnwidth - 4\tabcolsep) * \real{0.7143}}
  >{\raggedright\arraybackslash}p{(\columnwidth - 4\tabcolsep) * \real{0.1429}}@{}}
\toprule\noalign{}
\begin{minipage}[b]{\linewidth}\raggedright
Mẫu mã
\end{minipage} & \begin{minipage}[b]{\linewidth}\raggedright
Ý nghĩa
\end{minipage} & \begin{minipage}[b]{\linewidth}\raggedright
Ví dụ
\end{minipage} \\
\midrule\noalign{}
\endhead
\bottomrule\noalign{}
\endlastfoot
\texttt{TOP-NN} & Dạng bài (TOP = chủ đề: LOG, QNT, SET, FUN, IND, CNT,
REC, REL, POS, BOO, GRA) & \texttt{FUN-02} \\
\texttt{TOP-NN-Exx} & Ví dụ thứ xx của dạng bài & \texttt{FUN-02-E01} \\
\texttt{F-TOP-NN} & Công thức / định lý & \texttt{F-REC-03} \\
\texttt{M-NN} & Phương pháp giải & \texttt{M-17} \\
\texttt{T-NN} & Bẫy / lỗi thường gặp & \texttt{T-15} \\
\texttt{S\textless{}k\textgreater{}-p\textless{}n\textgreater{}} &
Nguồn: tệp slide \(k\) (1 = Logic, 2 = Phép đếm, 3 = Quan hệ, 4 = Đại số
Bool, 5 = Đồ thị), trang PDF \(n\) & \texttt{S2-p57} \\
\end{longtable}

\emph{Ghi chú:} spec mẫu dùng hậu tố \texttt{V02}; bộ tài liệu dùng
\texttt{Exx} cho số thứ tự ví dụ để không trùng với nhãn biến thể
V0--V4.

\textbf{Nhãn trong ví dụ:} \textbf{D1/D2/D3} = cơ bản/trung bình/nâng
cao; \textbf{V0} chuẩn, \textbf{V1} đổi dữ kiện, \textbf{V2} đổi cách
diễn đạt, \textbf{V3} kết hợp kiến thức, \textbf{V4} trường hợp đặc
biệt/bẫy. \textbf{{[}SV{]}} = đề lấy từ slide; \textbf{{[}AI-GENERATED
PRACTICE{]}} = đề do bộ tài liệu tạo thêm (không phải đề của giảng
viên); \textbf{{[}CC{]}} = lời giải đã kiểm tra độc lập bằng máy tính;
\textbf{{[}INFERRED{]}} = suy luận, nguồn không nói trực tiếp;
\textbf{{[}CONFLICT{]}} = nguồn mâu thuẫn.

\textbf{Phạm vi:} 5 chương slide \emph{Cấu trúc rời rạc} (Phạm Thế Bảo).
Chưa có đề cương thi/đề thi cũ ⇒ phạm vi thi là \textbf{giả định} (xem
phụ lục Khoảng trống \& mâu thuẫn). Không có thống kê tần suất dạng bài.

\textbf{Ký hiệu:} phủ định \(\neg p\) (= \(\bar p\) trong slide); bù
Bool \(\bar x\); thứ tự \(\preceq\); \(\mathbb N=\{0,1,2,\dots\}\).

\hypertarget{keyword-index}{%
\chapter{Mục lục từ khóa (Keyword Index)}\label{keyword-index}}

Tra theo \textbf{từ khóa trong đề} → mã dạng bài → trang. Cột ``Phân
biệt'' giúp tránh nhầm khi một từ khóa có thể dẫn tới dạng khác.

\begin{longtable}[]{@{}
  >{\raggedright\arraybackslash}p{(\columnwidth - 8\tabcolsep) * \real{0.3509}}
  >{\raggedright\arraybackslash}p{(\columnwidth - 8\tabcolsep) * \real{0.2105}}
  >{\raggedright\arraybackslash}p{(\columnwidth - 8\tabcolsep) * \real{0.1228}}
  >{\raggedright\arraybackslash}p{(\columnwidth - 8\tabcolsep) * \real{0.2456}}
  >{\raggedright\arraybackslash}p{(\columnwidth - 8\tabcolsep) * \real{0.0702}}@{}}
\toprule\noalign{}
\begin{minipage}[b]{\linewidth}\raggedright
Từ khóa đề bài
\end{minipage} & \begin{minipage}[b]{\linewidth}\raggedright
Dấu hiệu nhận diện
\end{minipage} & \begin{minipage}[b]{\linewidth}\raggedright
Mã dạng bài
\end{minipage} & \begin{minipage}[b]{\linewidth}\raggedright
Phân biệt
\end{minipage} & \begin{minipage}[b]{\linewidth}\raggedright
Trang
\end{minipage} \\
\midrule\noalign{}
\endhead
\bottomrule\noalign{}
\endlastfoot
\(\forall\), \(\exists\), với mọi, tồn tại --- đúng hay sai & vị từ trên
tập số & \protect\hyperlink{QNT-01}{QNT-01} & & \pageref{QNT-01} \\
\(P(X)\), tập các tập con, lực lượng & &
\protect\hyperlink{SET-02}{SET-02} & & \pageref{SET-02} \\
ánh xạ ngược, khả nghịch, \(f^{-1}\) & chỉ khi song ánh &
\protect\hyperlink{FUN-02}{FUN-02} & ảnh ngược của tập → FUN-01 &
\pageref{FUN-02} \\
bậc, bậc vào, bậc ra, số cạnh, định lý bắt tay & &
\protect\hyperlink{GRA-01}{GRA-01} & & \pageref{GRA-01} \\
bao nhiêu số \ldots{} chữ số (khác nhau), chẵn, chia hết cho 5 & tập chữ
số có 0 & \protect\hyperlink{CNT-01}{CNT-01} & & \pageref{CNT-01} \\
biểu đồ Hasse & & \protect\hyperlink{POS-03}{POS-03} & &
\pageref{POS-03} \\
chặn trên, chặn dưới, sup, inf, \(a\vee b\), \(a\wedge b\) & &
\protect\hyperlink{POS-04}{POS-04} & & \pageref{POS-04} \\
chia \(k\) vật giống nhau cho \(n\) người / hộp & &
\protect\hyperlink{CNT-06}{CNT-06} & vật khác nhau → \(n^k\) (CNT-04) &
\pageref{CNT-06} \\
chứng minh đẳng thức tập hợp & \(A\setminus(B\cup C)=\dots\) &
\protect\hyperlink{SET-01}{SET-01} & & \pageref{SET-01} \\
chứng minh suy luận, \(\therefore\), quy tắc suy diễn & tiền đề / vạch
ngang / kết luận & \protect\hyperlink{LOG-04}{LOG-04} & nghi sai →
LOG-05 & \pageref{LOG-04} \\
chứng minh theo trường hợp & & \protect\hyperlink{LOG-04}{LOG-04} & chia
hết theo \(n\bmod k\) → IND-01 & \pageref{LOG-04} \\
chứng minh tương đương logic, \(\Leftrightarrow\) & hai dạng mệnh đề &
\protect\hyperlink{LOG-03}{LOG-03} & & \pageref{LOG-03} \\
chuồng bồ câu, Dirichlet, luôn tồn tại hai\ldots{} & chứng minh tồn tại
& \protect\hyperlink{CNT-02}{CNT-02} & & \pageref{CNT-02} \\
có phải ánh xạ; ảnh \(f(A)\), ảnh ngược \(f^{-1}(B)\) & &
\protect\hyperlink{FUN-01}{FUN-01} & & \pageref{FUN-01} \\
cổng AND, OR, NOT, NAND, NOR; vẽ mạch, viết biểu thức của sơ đồ & &
\protect\hyperlink{BOO-04}{BOO-04} & & \pageref{BOO-04} \\
công thức đa thức tối tiểu, tối thiểu hóa hàm Bool, tế bào lớn & &
\protect\hyperlink{BOO-03}{BOO-03} & & \pageref{BOO-03} \\
đẳng cấu & & \protect\hyperlink{GRA-03}{GRA-03} & & \pageref{GRA-03} \\
dạng nối rời chính tắc, từ tối tiểu & &
\protect\hyperlink{BOO-01}{BOO-01} & & \pageref{BOO-01} \\
điều kiện cần, điều kiện đủ, nếu\ldots{} thì\ldots{} & đọc \(p\to q\) &
\protect\hyperlink{LOG-02}{LOG-02} & xem F-LOG-01 & \pageref{LOG-02} \\
đỉnh khớp, cầu & & \protect\hyperlink{GRA-04}{GRA-04} & Fleury → GRA-05
& \pageref{GRA-04} \\
đơn ánh, toàn ánh, song ánh & chú ý miền \& đích &
\protect\hyperlink{FUN-02}{FUN-02} & & \pageref{FUN-02} \\
đơn đồ thị, đa đồ thị, giả đồ thị, khuyên & &
\protect\hyperlink{GRA-01}{GRA-01} & & \pageref{GRA-01} \\
đồng dư, modulo, \(\mathbb Z_m\), phân hoạch & &
\protect\hyperlink{REL-03}{REL-03} & & \pageref{REL-03} \\
đường đi đơn/sơ cấp, chu trình sơ cấp, liên thông & &
\protect\hyperlink{GRA-04}{GRA-04} & & \pageref{GRA-04} \\
đường đi ngắn nhất, khoảng cách, Dijkstra & trọng số dương &
\protect\hyperlink{GRA-06}{GRA-06} & có trọng số âm → GRA-07 &
\pageref{GRA-06} \\
Euler, đi qua mỗi cạnh đúng một lần, vẽ một nét, Königsberg & &
\protect\hyperlink{GRA-05}{GRA-05} & & \pageref{GRA-05} \\
Fleury & & \protect\hyperlink{GRA-05}{GRA-05} & & \pageref{GRA-05} \\
Ford--Bellman, trọng số âm, mạch âm & &
\protect\hyperlink{GRA-07}{GRA-07} & & \pageref{GRA-07} \\
giải hệ thức đệ qui, vế phải = 0 & thuần nhất &
\protect\hyperlink{REC-02}{REC-02} & & \pageref{REC-02} \\
hàm Bool, bảng chân trị, bỏ phiếu & mô tả bằng lời &
\protect\hyperlink{BOO-01}{BOO-01} & & \pageref{BOO-01} \\
hằng đúng, hằng sai, mâu thuẫn, chân lý & &
\protect\hyperlink{LOG-02}{LOG-02} & & \pageref{LOG-02} \\
hệ quả logic \(\Rightarrow\) & & \protect\hyperlink{LOG-03}{LOG-03} & &
\pageref{LOG-03} \\
hoặc / ít nhất một trong / không chia hết cho & đếm hợp các tập &
\protect\hyperlink{CNT-03}{CNT-03} & & \pageref{CNT-03} \\
hoán vị, chỉnh hợp, tổ hợp, chọn ủy ban, xếp hàng & thứ tự? lặp? &
\protect\hyperlink{CNT-04}{CNT-04} & có vật giống nhau → CNT-05 &
\pageref{CNT-04} \\
hợp, giao, hiệu, bù, De Morgan tập hợp & &
\protect\hyperlink{SET-01}{SET-01} & & \pageref{SET-01} \\
lập bảng chân trị (dạng mệnh đề) & cho \(E(p,q,\dots)\) &
\protect\hyperlink{LOG-02}{LOG-02} & hàm Bool từ mô tả → BOO-01 &
\pageref{LOG-02} \\
liệt kê quan hệ, ma trận biểu diễn \(M_R\) & &
\protect\hyperlink{REL-01}{REL-01} & ma trận kề đồ thị → GRA-02 &
\pageref{REL-01} \\
lời văn: ``nếu\ldots{} thì\ldots, biết rằng\ldots, suy ra\ldots{}'' &
đặt biến mệnh đề & \protect\hyperlink{LOG-05}{LOG-05} & nếu chắc đúng →
LOG-04 & \pageref{LOG-05} \\
ma trận kề & & \protect\hyperlink{GRA-02}{GRA-02} & ma trận trọng số →
GRA-06 & \pageref{GRA-02} \\
mệnh đề? sơ cấp / phức hợp & danh sách câu nói &
\protect\hyperlink{LOG-01}{LOG-01} & & \pageref{LOG-01} \\
nguyên lý cộng / nhân, có bao nhiêu cách & các bước hoặc các phương án &
\protect\hyperlink{CNT-01}{CNT-01} & chọn nhóm → CNT-04 &
\pageref{CNT-01} \\
phản chứng, quy tắc mâu thuẫn & kết luận dạng \(a\to b\) &
\protect\hyperlink{LOG-04}{LOG-04} & & \pageref{LOG-04} \\
phản xạ, đối xứng, phản xứng, bắc cầu & &
\protect\hyperlink{REL-02}{REL-02} & & \pageref{REL-02} \\
phủ định mệnh đề lượng từ, \(\varepsilon\)--\(\delta\) & &
\protect\hyperlink{QNT-02}{QNT-02} & & \pageref{QNT-02} \\
phương trình đặc trưng, nghiệm kép & &
\protect\hyperlink{REC-02}{REC-02} & có vế phải → REC-03 &
\pageref{REC-02} \\
quan hệ thứ tự, poset, thứ tự toàn phần, so sánh được & &
\protect\hyperlink{POS-01}{POS-01} & & \pageref{POS-01} \\
quan hệ tương đương, lớp tương đương \([a]\) & &
\protect\hyperlink{REL-03}{REL-03} & & \pageref{REL-03} \\
quy nạp, \(P(n)\), với mọi \(n\ge n_0\) & tổng, bất đẳng thức, chia hết
& \protect\hyperlink{IND-01}{IND-01} & & \pageref{IND-01} \\
rút gọn / đơn giản dạng mệnh đề & & \protect\hyperlink{LOG-03}{LOG-03} &
& \pageref{LOG-03} \\
sắp xếp chữ cái của một từ, vật giống nhau & &
\protect\hyperlink{CNT-05}{CNT-05} & & \pageref{CNT-05} \\
sắp xếp topo, thứ tự thực hiện công việc & &
\protect\hyperlink{POS-05}{POS-05} & & \pageref{POS-05} \\
số nghiệm nguyên không âm \(x_1+\dots+x_n=k\) & ràng buộc \(\ge,\le,>\)
& \protect\hyperlink{CNT-06}{CNT-06} & & \pageref{CNT-06} \\
suy luận có đúng không, phản ví dụ & &
\protect\hyperlink{LOG-05}{LOG-05} & & \pageref{LOG-05} \\
thiết kế mạch cầu dao / công tắc / đèn & &
\protect\hyperlink{BOO-04}{BOO-04} & & \pageref{BOO-04} \\
thứ tự lượng từ \(\forall x\exists y\) / \(\exists y\forall x\) & &
\protect\hyperlink{QNT-01}{QNT-01} & & \pageref{QNT-01} \\
thứ tự từ điển, so sánh chuỗi & & \protect\hyperlink{POS-02}{POS-02} & &
\pageref{POS-02} \\
tích ánh xạ, hàm hợp \(g\circ f\) & & \protect\hyperlink{FUN-03}{FUN-03}
& & \pageref{FUN-03} \\
tích Đề-các \(A\times B\) & & \protect\hyperlink{SET-02}{SET-02} & &
\pageref{SET-02} \\
tìm hệ thức đệ qui, cầu thang, tháp Hà Nội, chuỗi bit & đặt \(x_n\) &
\protect\hyperlink{REC-01}{REC-01} & & \pageref{REC-01} \\
tổ hợp lặp, chọn \(k\) vật từ \(n\) loại & &
\protect\hyperlink{CNT-06}{CNT-06} & & \pageref{CNT-06} \\
tối đại, tối tiểu, lớn nhất, nhỏ nhất & &
\protect\hyperlink{POS-03}{POS-03} & của tập con \(A\) → POS-04 &
\pageref{POS-03} \\
vẽ biểu đồ Karnaugh; tế bào là kar của đơn thức nào & &
\protect\hyperlink{BOO-02}{BOO-02} & & \pageref{BOO-02} \\
vế phải đa thức / \(\beta^n P(n)\), nghiệm riêng & không thuần nhất &
\protect\hyperlink{REC-03}{REC-03} & & \pageref{REC-03} \\
vế phải là tổng nhiều dạng & nhiều \(\beta\) khác nhau &
\protect\hyperlink{REC-04}{REC-04} & & \pageref{REC-04} \\
\end{longtable}

\hypertarget{chapter-map}{%
\chapter{Bản đồ theo chương (Chapter Map)}\label{chapter-map}}

\begin{longtable}[]{@{}
  >{\raggedright\arraybackslash}p{(\columnwidth - 12\tabcolsep) * \real{0.1311}}
  >{\raggedright\arraybackslash}p{(\columnwidth - 12\tabcolsep) * \real{0.1639}}
  >{\raggedright\arraybackslash}p{(\columnwidth - 12\tabcolsep) * \real{0.2295}}
  >{\raggedright\arraybackslash}p{(\columnwidth - 12\tabcolsep) * \real{0.1967}}
  >{\raggedright\arraybackslash}p{(\columnwidth - 12\tabcolsep) * \real{0.1475}}
  >{\raggedright\arraybackslash}p{(\columnwidth - 12\tabcolsep) * \real{0.0656}}
  >{\raggedright\arraybackslash}p{(\columnwidth - 12\tabcolsep) * \real{0.0656}}@{}}
\toprule\noalign{}
\begin{minipage}[b]{\linewidth}\raggedright
Chương (nguồn)
\end{minipage} & \begin{minipage}[b]{\linewidth}\raggedright
Chủ đề chính
\end{minipage} & \begin{minipage}[b]{\linewidth}\raggedright
Định nghĩa quan trọng
\end{minipage} & \begin{minipage}[b]{\linewidth}\raggedright
Công thức / định lý
\end{minipage} & \begin{minipage}[b]{\linewidth}\raggedright
Dạng bài
\end{minipage} & \begin{minipage}[b]{\linewidth}\raggedright
Ví dụ (trang)
\end{minipage} & \begin{minipage}[b]{\linewidth}\raggedright
Nguồn
\end{minipage} \\
\midrule\noalign{}
\endhead
\bottomrule\noalign{}
\endlastfoot
\textbf{I. Cơ sở logic} (S1) & Mệnh đề, dạng mệnh đề, suy diễn & mệnh
đề, sơ cấp/phức hợp, hằng đúng/sai, tương đương, hệ quả &
\protect\hyperlink{F-LOG-01}{F-LOG-01}--\protect\hyperlink{F-LOG-04}{F-LOG-04}
& \protect\hyperlink{LOG-01}{LOG-01}--\protect\hyperlink{LOG-05}{LOG-05}
& \pageref{ex-ch1} & S1-p5--56 \\
& Vị từ, lượng từ & vị từ, \(\forall\), \(\exists\) &
\protect\hyperlink{F-QNT-01}{F-QNT-01},
\protect\hyperlink{F-QNT-02}{F-QNT-02} &
\protect\hyperlink{QNT-01}{QNT-01}, \protect\hyperlink{QNT-02}{QNT-02} &
\pageref{QNT-01-E01} & S1-p58--72 \\
& Tập hợp & lực lượng, \(P(X)\), \(A\times B\) &
\protect\hyperlink{F-SET-01}{F-SET-01},
\protect\hyperlink{F-SET-02}{F-SET-02} &
\protect\hyperlink{SET-01}{SET-01}, \protect\hyperlink{SET-02}{SET-02} &
\pageref{SET-01-E01} & S1-p73--82 \\
& Ánh xạ & ảnh, ảnh ngược, đơn/toàn/song ánh, ánh xạ ngược, tích &
\protect\hyperlink{F-FUN-01}{F-FUN-01}--\protect\hyperlink{F-FUN-03}{F-FUN-03}
& \protect\hyperlink{FUN-01}{FUN-01}--\protect\hyperlink{FUN-03}{FUN-03}
& \pageref{FUN-01-E01} & S1-p83--96 \\
& Qui nạp & bước cơ sở, bước quy nạp &
\protect\hyperlink{F-IND-01}{F-IND-01} &
\protect\hyperlink{IND-01}{IND-01} & \pageref{IND-01-E01} &
S1-p97--100 \\
\textbf{II. Phép đếm} (S2) & Nguyên lý, tổ hợp & cộng, nhân, Dirichlet,
bù trừ; \(P_n,A_n^k,C_n^k\) &
\protect\hyperlink{F-CNT-01}{F-CNT-01}--\protect\hyperlink{F-CNT-04}{F-CNT-04}
& \protect\hyperlink{CNT-01}{CNT-01}--\protect\hyperlink{CNT-04}{CNT-04}
& \pageref{ex-ch2a} & S2-p3--16 \\
& Hoán vị lặp, tổ hợp lặp & \(\frac{n!}{n_1!\cdots n_k!}\), \(K_n^k\) &
\protect\hyperlink{F-CNT-05}{F-CNT-05},
\protect\hyperlink{F-CNT-06}{F-CNT-06} &
\protect\hyperlink{CNT-05}{CNT-05}, \protect\hyperlink{CNT-06}{CNT-06} &
\pageref{CNT-05-E01} & S2-p17--25 \\
& Hệ thức đệ qui & tuyến tính cấp \(k\), PTĐT, nghiệm TQ/riêng &
\protect\hyperlink{F-REC-01}{F-REC-01}--\protect\hyperlink{F-REC-05}{F-REC-05}
& \protect\hyperlink{REC-01}{REC-01}--\protect\hyperlink{REC-04}{REC-04}
& \pageref{ex-ch2b} & S2-p26--67 \\
\textbf{III. Quan hệ} (S3) & Tính chất, biểu diễn, tương đương & 4 tính
chất, \(M_R\), lớp tương đương, đồng dư &
\protect\hyperlink{F-REL-01}{F-REL-01}--\protect\hyperlink{F-REL-03}{F-REL-03}
& \protect\hyperlink{REL-01}{REL-01}--\protect\hyperlink{REL-03}{REL-03}
& \pageref{ex-ch3} & S3-p4--27 \\
& Quan hệ thứ tự & poset, toàn phần, từ điển, Hasse, tối đại/tiểu, chặn,
topo &
\protect\hyperlink{F-POS-01}{F-POS-01}--\protect\hyperlink{F-POS-03}{F-POS-03}
& \protect\hyperlink{POS-01}{POS-01}--\protect\hyperlink{POS-05}{POS-05}
& \pageref{POS-01-E01} & S3-p28--56 \\
\textbf{IV. Đại số Bool} (S4) & Hàm Bool & từ đơn, đơn thức, từ tối
tiểu, dạng nối rời chính tắc & \protect\hyperlink{F-BOO-01}{F-BOO-01},
\protect\hyperlink{F-BOO-02}{F-BOO-02} &
\protect\hyperlink{BOO-01}{BOO-01} & \pageref{ex-ch4} & S4-p5--19 \\
& Karnaugh & tế bào, tế bào lớn, phủ tối tiểu &
\protect\hyperlink{F-BOO-03}{F-BOO-03},
\protect\hyperlink{F-BOO-04}{F-BOO-04} &
\protect\hyperlink{BOO-02}{BOO-02}, \protect\hyperlink{BOO-03}{BOO-03} &
\pageref{BOO-02-E01} & S4-p20--65 \\
& Mạng logic & 5 cổng & \protect\hyperlink{F-BOO-05}{F-BOO-05} &
\protect\hyperlink{BOO-04}{BOO-04} & \pageref{BOO-04-E01} &
S4-p66--78 \\
\textbf{V. Đồ thị} (S5) & Khái niệm, ma trận, đẳng cấu & bậc, đơn/đa/giả
đồ thị, ma trận kề, đẳng cấu &
\protect\hyperlink{F-GRA-01}{F-GRA-01}--\protect\hyperlink{F-GRA-03}{F-GRA-03}
& \protect\hyperlink{GRA-01}{GRA-01}--\protect\hyperlink{GRA-03}{GRA-03}
& \pageref{ex-ch5} & S5-p3--35 \\
& Liên thông, Euler & đường đi đơn/sơ cấp, đỉnh khớp, cầu, Euler &
\protect\hyperlink{F-GRA-04}{F-GRA-04} &
\protect\hyperlink{GRA-04}{GRA-04}, \protect\hyperlink{GRA-05}{GRA-05} &
\pageref{GRA-04-E01} & S5-p36--51 \\
& Đường đi ngắn nhất & trọng số, khoảng cách, mạch âm &
\protect\hyperlink{F-GRA-05}{F-GRA-05},
\protect\hyperlink{F-GRA-06}{F-GRA-06} &
\protect\hyperlink{GRA-06}{GRA-06}, \protect\hyperlink{GRA-07}{GRA-07} &
\pageref{GRA-06-E01} & S5-p52--86 \\
\end{longtable}

\hypertarget{method-index}{%
\chapter{Mục lục theo phương pháp giải (Method
Index)}\label{method-index}}

\begin{longtable}[]{@{}
  >{\raggedright\arraybackslash}p{(\columnwidth - 8\tabcolsep) * \real{0.0930}}
  >{\raggedright\arraybackslash}p{(\columnwidth - 8\tabcolsep) * \real{0.3256}}
  >{\raggedright\arraybackslash}p{(\columnwidth - 8\tabcolsep) * \real{0.2326}}
  >{\raggedright\arraybackslash}p{(\columnwidth - 8\tabcolsep) * \real{0.2791}}
  >{\raggedright\arraybackslash}p{(\columnwidth - 8\tabcolsep) * \real{0.0698}}@{}}
\toprule\noalign{}
\begin{minipage}[b]{\linewidth}\raggedright
Mã
\end{minipage} & \begin{minipage}[b]{\linewidth}\raggedright
Phương pháp
\end{minipage} & \begin{minipage}[b]{\linewidth}\raggedright
Dùng cho dạng bài
\end{minipage} & \begin{minipage}[b]{\linewidth}\raggedright
Ví dụ minh họa
\end{minipage} & \begin{minipage}[b]{\linewidth}\raggedright
Trang
\end{minipage} \\
\midrule\noalign{}
\endhead
\bottomrule\noalign{}
\endlastfoot
\protect\hyperlink{M-01}{M-01} & Lập bảng chân trị &
\protect\hyperlink{LOG-02}{LOG-02}, \protect\hyperlink{LOG-03}{LOG-03} &
\protect\hyperlink{LOG-02-E01}{LOG-02-E01} & \pageref{M-01} \\
\protect\hyperlink{M-02}{M-02} & Biến đổi tương đương bằng luật logic &
\protect\hyperlink{LOG-03}{LOG-03} &
\protect\hyperlink{LOG-03-E01}{LOG-03-E01} & \pageref{M-02} \\
\protect\hyperlink{M-03}{M-03} & Chuỗi quy tắc suy diễn (chứng minh trực
tiếp) & \protect\hyperlink{LOG-04}{LOG-04} &
\protect\hyperlink{LOG-04-E01}{LOG-04-E01} & \pageref{M-03} \\
\protect\hyperlink{M-04}{M-04} & Phản chứng (quy tắc mâu thuẫn) &
\protect\hyperlink{LOG-04}{LOG-04} &
\protect\hyperlink{LOG-04-E02}{LOG-04-E02},
\protect\hyperlink{LOG-04-E03}{LOG-04-E03} & \pageref{M-04} \\
\protect\hyperlink{M-05}{M-05} & Kiểm tra suy luận bằng hệ chân trị /
phản ví dụ & \protect\hyperlink{LOG-05}{LOG-05} &
\protect\hyperlink{LOG-05-E01}{LOG-05-E01},
\protect\hyperlink{LOG-05-E02}{LOG-05-E02} & \pageref{M-05} \\
\protect\hyperlink{M-06}{M-06} & Xét chân trị mệnh đề lượng từ &
\protect\hyperlink{QNT-01}{QNT-01} &
\protect\hyperlink{QNT-01-E01}{QNT-01-E01} & \pageref{M-06} \\
\protect\hyperlink{M-07}{M-07} & Phủ định mệnh đề lượng từ &
\protect\hyperlink{QNT-02}{QNT-02} &
\protect\hyperlink{QNT-02-E01}{QNT-02-E01} & \pageref{M-07} \\
\protect\hyperlink{M-08}{M-08} & Chứng minh đẳng thức tập hợp &
\protect\hyperlink{SET-01}{SET-01} &
\protect\hyperlink{SET-01-E01}{SET-01-E01} & \pageref{M-08} \\
\protect\hyperlink{M-09}{M-09} & Xét đơn/toàn/song ánh qua phương trình
\(f(x)=y\) & \protect\hyperlink{FUN-01}{FUN-01},
\protect\hyperlink{FUN-02}{FUN-02} &
\protect\hyperlink{FUN-02-E01}{FUN-02-E01},
\protect\hyperlink{FUN-02-E02}{FUN-02-E02} & \pageref{M-09} \\
\protect\hyperlink{M-10}{M-10} & Quy nạp toán học &
\protect\hyperlink{IND-01}{IND-01} &
\protect\hyperlink{IND-01-E01}{IND-01-E01},
\protect\hyperlink{IND-01-E03}{IND-01-E03} & \pageref{M-10} \\
\protect\hyperlink{M-11}{M-11} & Đếm bằng chia trường hợp + nguyên lý
nhân & \protect\hyperlink{CNT-01}{CNT-01} &
\protect\hyperlink{CNT-01-E01}{CNT-01-E01} & \pageref{M-11} \\
\protect\hyperlink{M-12}{M-12} & Nguyên lý chuồng bồ câu &
\protect\hyperlink{CNT-02}{CNT-02} &
\protect\hyperlink{CNT-02-E01}{CNT-02-E01} & \pageref{M-12} \\
\protect\hyperlink{M-13}{M-13} & Bù trừ / đếm phần bù &
\protect\hyperlink{CNT-03}{CNT-03}, \protect\hyperlink{CNT-06}{CNT-06} &
\protect\hyperlink{CNT-03-E02}{CNT-03-E02},
\protect\hyperlink{CNT-06-E02}{CNT-06-E02} & \pageref{M-13} \\
\protect\hyperlink{M-14}{M-14} & Chọn đúng công thức tổ hợp (cây quyết
định) & \protect\hyperlink{CNT-04}{CNT-04} &
\protect\hyperlink{CNT-04-E01}{CNT-04-E01} & \pageref{M-14} \\
\protect\hyperlink{M-15}{M-15} & Đổi biến cho nghiệm nguyên có ràng buộc
& \protect\hyperlink{CNT-06}{CNT-06} &
\protect\hyperlink{CNT-06-E02}{CNT-06-E02},
\protect\hyperlink{CNT-06-E03}{CNT-06-E03} & \pageref{M-15} \\
\protect\hyperlink{M-16}{M-16} & Lập hệ thức đệ qui &
\protect\hyperlink{REC-01}{REC-01} &
\protect\hyperlink{REC-01-E01}{REC-01-E01},
\protect\hyperlink{REC-01-E02}{REC-01-E02} & \pageref{M-16} \\
\protect\hyperlink{M-17}{M-17} & Giải hệ thức đệ qui tuyến tính (quy
trình đầy đủ) & \protect\hyperlink{REC-02}{REC-02},
\protect\hyperlink{REC-03}{REC-03}, \protect\hyperlink{REC-04}{REC-04} &
\protect\hyperlink{REC-02-E02}{REC-02-E02},
\protect\hyperlink{REC-03-E02}{REC-03-E02},
\protect\hyperlink{REC-04-E01}{REC-04-E01} & \pageref{M-17} \\
\protect\hyperlink{M-18}{M-18} & Kiểm tra tính chất quan hệ &
\protect\hyperlink{REL-02}{REL-02}, \protect\hyperlink{REL-03}{REL-03},
\protect\hyperlink{POS-01}{POS-01} &
\protect\hyperlink{POS-01-E01}{POS-01-E01},
\protect\hyperlink{REL-02-E01}{REL-02-E01},
\protect\hyperlink{REL-02-E02}{REL-02-E02} & \pageref{M-18} \\
\protect\hyperlink{M-19}{M-19} & Tìm lớp tương đương &
\protect\hyperlink{REL-03}{REL-03} &
\protect\hyperlink{REL-03-E01}{REL-03-E01},
\protect\hyperlink{REL-03-E02}{REL-03-E02} & \pageref{M-19} \\
\protect\hyperlink{M-20}{M-20} & Vẽ Hasse và đọc phần tử đặc biệt, chặn
& \protect\hyperlink{POS-03}{POS-03}, \protect\hyperlink{POS-04}{POS-04}
& \protect\hyperlink{POS-03-E01}{POS-03-E01},
\protect\hyperlink{POS-04-E01}{POS-04-E01} & \pageref{M-20} \\
\protect\hyperlink{M-21}{M-21} & Sắp xếp topo &
\protect\hyperlink{POS-05}{POS-05} &
\protect\hyperlink{POS-05-E01}{POS-05-E01} & \pageref{M-21} \\
\protect\hyperlink{M-22}{M-22} & Thuật toán Karnaugh 5 bước &
\protect\hyperlink{BOO-03}{BOO-03} &
\protect\hyperlink{BOO-03-E01}{BOO-03-E01},
\protect\hyperlink{BOO-03-E02}{BOO-03-E02} & \pageref{M-22} \\
\protect\hyperlink{M-23}{M-23} & Từ bảng chân trị / mô tả đến công thức
và mạch & \protect\hyperlink{BOO-01}{BOO-01},
\protect\hyperlink{BOO-04}{BOO-04} &
\protect\hyperlink{BOO-04-E01}{BOO-04-E01},
\protect\hyperlink{BOO-04-E02}{BOO-04-E02} & \pageref{M-23} \\
\protect\hyperlink{M-24}{M-24} & Kiểm tra đẳng cấu &
\protect\hyperlink{GRA-03}{GRA-03} &
\protect\hyperlink{GRA-03-E01}{GRA-03-E01} & \pageref{M-24} \\
\protect\hyperlink{M-25}{M-25} & Euler: kiểm tra điều kiện + Fleury &
\protect\hyperlink{GRA-05}{GRA-05} &
\protect\hyperlink{GRA-05-E01}{GRA-05-E01},
\protect\hyperlink{GRA-05-E02}{GRA-05-E02} & \pageref{M-25} \\
\protect\hyperlink{M-26}{M-26} & Dijkstra dạng bảng &
\protect\hyperlink{GRA-06}{GRA-06} &
\protect\hyperlink{GRA-06-E01}{GRA-06-E01} & \pageref{M-26} \\
\protect\hyperlink{M-27}{M-27} & Ford--Bellman dạng bảng &
\protect\hyperlink{GRA-07}{GRA-07} &
\protect\hyperlink{GRA-07-E01}{GRA-07-E01},
\protect\hyperlink{GRA-07-E02}{GRA-07-E02} & \pageref{M-27} \\
\end{longtable}

\hypertarget{taxonomy}{%
\chapter{Phân loại dạng bài (Problem Taxonomy)}\label{taxonomy}}

Mỗi dạng: \textbf{Từ khóa} (cách đề hỏi) · \textbf{Dấu hiệu} ·
\textbf{Kiểm tra trước} · \textbf{Kiến thức} (mã F-) · \textbf{Phương
pháp} (mã M-) · \textbf{Ví dụ} · \textbf{Nguồn}. Mã ví dụ dạng
\texttt{DẠNG-Exx}; nhãn độ khó D1--D3 và biến thể V0--V4 ghi trong từng
ví dụ.

\hypertarget{chux1b0ux1a1ng-i-cux1a1-sux1edf-logic}{%
\section{Chương I --- Cơ sở
logic}\label{chux1b0ux1a1ng-i-cux1a1-sux1edf-logic}}

\hypertarget{LOG-01}{%
\subsection{LOG-01 · Nhận biết mệnh đề, phân loại sơ cấp/phức
hợp}\label{LOG-01}}

\begin{itemize}
\tightlist
\item
  \textbf{Từ khóa:} ``có phải là mệnh đề không'', ``mệnh đề sơ cấp /
  phức hợp'', ``xác định chân trị''.
\item
  \textbf{Dấu hiệu:} cho danh sách câu nói tự nhiên.
\item
  \textbf{Kiểm tra trước:} câu hỏi, cảm thán, mệnh lệnh, câu chứa biến
  tự do (``n là số tự nhiên'') \textbf{không} là mệnh đề.
\item
  \textbf{Kiến thức:} \protect\hyperlink{F-LOG-01}{F-LOG-01}.
  \textbf{Phương pháp:} đối chiếu định nghĩa (S1-p6, p9).
\item
  \textbf{Ví dụ:} \protect\hyperlink{LOG-01-E01}{LOG-01-E01}.
  \textbf{Nguồn:} S1-p6--9.
\end{itemize}

\hypertarget{LOG-02}{%
\subsection{LOG-02 · Lập bảng chân trị; hằng đúng / hằng
sai}\label{LOG-02}}

\begin{itemize}
\tightlist
\item
  \textbf{Từ khóa:} ``lập bảng chân trị'', ``hằng đúng'', ``hằng sai'',
  ``mâu thuẫn'', ``chân lý''.
\item
  \textbf{Dấu hiệu:} cho dạng mệnh đề \(E(p,q,\dots)\).
\item
  \textbf{Kiểm tra trước:} số dòng \(=2^n\); thứ tự ưu tiên ngoặc.
\item
  \textbf{Kiến thức:} \protect\hyperlink{F-LOG-01}{F-LOG-01}.
  \textbf{Phương pháp:} \protect\hyperlink{M-01}{M-01}.
\item
  \textbf{Ví dụ:} \protect\hyperlink{LOG-02-E01}{LOG-02-E01},
  \protect\hyperlink{LOG-02-E02}{LOG-02-E02}. \textbf{Nguồn:}
  S1-p18--22.
\end{itemize}

\hypertarget{LOG-03}{%
\subsection{LOG-03 · Chứng minh tương đương logic / rút gọn dạng mệnh
đề}\label{LOG-03}}

\begin{itemize}
\tightlist
\item
  \textbf{Từ khóa:} ``chứng minh \(E \Leftrightarrow F\)'', ``tương
  đương logic'', ``rút gọn'', ``đơn giản biểu thức'', ``hệ quả logic''.
\item
  \textbf{Dấu hiệu:} hai dạng mệnh đề nối bằng \(\Leftrightarrow\) hoặc
  \(\Rightarrow\).
\item
  \textbf{Kiểm tra trước:} mỗi bước phải nêu tên luật;
  ``\(\Rightarrow\)'' (hệ quả) khác ``\(\Leftrightarrow\)''.
\item
  \textbf{Kiến thức:} \protect\hyperlink{F-LOG-02}{F-LOG-02},
  \protect\hyperlink{F-LOG-03}{F-LOG-03}. \textbf{Phương pháp:}
  \protect\hyperlink{M-02}{M-02} (hoặc \protect\hyperlink{M-01}{M-01} để
  đối chiếu).
\item
  \textbf{Ví dụ:} \protect\hyperlink{LOG-03-E01}{LOG-03-E01},
  \protect\hyperlink{LOG-03-E02}{LOG-03-E02}. \textbf{Nguồn:}
  S1-p23--30.
\end{itemize}

\hypertarget{LOG-04}{%
\subsection{LOG-04 · Chứng minh suy luận đúng (quy tắc suy diễn, phản
chứng)}\label{LOG-04}}

\begin{itemize}
\tightlist
\item
  \textbf{Từ khóa:} ``chứng minh suy luận'', ``kiểm tra suy luận'',
  ``\(\therefore\)'', ``suy ra'', ``tiền đề -- kết luận''.
\item
  \textbf{Dấu hiệu:} danh sách tiền đề trên vạch ngang, kết luận dưới có
  dấu \(\therefore\).
\item
  \textbf{Kiểm tra trước:} kết luận dạng kéo theo → có thể dùng phản
  chứng (thêm \(\neg h\)) hoặc đổi sang phản đảo.
\item
  \textbf{Kiến thức:} \protect\hyperlink{F-LOG-04}{F-LOG-04}.
  \textbf{Phương pháp:} \protect\hyperlink{M-03}{M-03},
  \protect\hyperlink{M-04}{M-04}.
\item
  \textbf{Ví dụ:} \protect\hyperlink{LOG-04-E01}{LOG-04-E01},
  \protect\hyperlink{LOG-04-E02}{LOG-04-E02},
  \protect\hyperlink{LOG-04-E03}{LOG-04-E03}. \textbf{Nguồn:}
  S1-p31--54.
\end{itemize}

\hypertarget{LOG-05}{%
\subsection{LOG-05 · Suy luận sai --- tìm phản ví dụ; bài toán lời
văn}\label{LOG-05}}

\begin{itemize}
\tightlist
\item
  \textbf{Từ khóa:} ``suy luận sau có đúng không'', ``đúng hay sai'',
  đoạn văn có ``nếu\ldots{} thì\ldots{}'', ``biết rằng''.
\item
  \textbf{Dấu hiệu:} không chắc kết luận đúng; bài toán lời văn cần đặt
  biến mệnh đề.
\item
  \textbf{Kiểm tra trước:} gán kết luận = 0, giải hệ ``mọi tiền đề =
  1''; có nghiệm ⇒ suy luận sai.
\item
  \textbf{Kiến thức:} \protect\hyperlink{F-LOG-04}{F-LOG-04} (quy tắc 9
  --- phản ví dụ). \textbf{Phương pháp:} \protect\hyperlink{M-05}{M-05}.
\item
  \textbf{Ví dụ:} \protect\hyperlink{LOG-05-E01}{LOG-05-E01},
  \protect\hyperlink{LOG-05-E02}{LOG-05-E02}. \textbf{Nguồn:}
  S1-p48--49, p55--56.
\end{itemize}

\hypertarget{QNT-01}{%
\subsection{QNT-01 · Xét chân trị mệnh đề có lượng từ}\label{QNT-01}}

\begin{itemize}
\tightlist
\item
  \textbf{Từ khóa:} ``∀'', ``∃'', ``với mọi'', ``tồn tại'', ``đúng hay
  sai'', ``có ít nhất một''.
\item
  \textbf{Dấu hiệu:} vị từ một hoặc hai biến trên tập số.
\item
  \textbf{Kiểm tra trước:} \textbf{thứ tự lượng từ}
  (\(\forall x\exists y\) ≠ \(\exists y\forall x\)); miền của biến.
\item
  \textbf{Kiến thức:} \protect\hyperlink{F-QNT-01}{F-QNT-01}.
  \textbf{Phương pháp:} \protect\hyperlink{M-06}{M-06}.
\item
  \textbf{Ví dụ:} \protect\hyperlink{QNT-01-E01}{QNT-01-E01},
  \protect\hyperlink{QNT-01-E02}{QNT-01-E02}. \textbf{Nguồn:}
  S1-p59--68.
\end{itemize}

\hypertarget{QNT-02}{%
\subsection{QNT-02 · Phủ định mệnh đề có lượng từ}\label{QNT-02}}

\begin{itemize}
\tightlist
\item
  \textbf{Từ khóa:} ``phủ định'', ``viết mệnh đề phủ định'',
  ``\(\overline{\forall x\dots}\)''.
\item
  \textbf{Dấu hiệu:} mệnh đề lượng từ hóa cần lấy phủ định.
\item
  \textbf{Kiểm tra trước:} đổi \(\forall\leftrightarrow\exists\), phủ
  định vị từ; phủ định của \(A\to B\) là \(A\wedge\neg B\); giữ nguyên
  miền.
\item
  \textbf{Kiến thức:} \protect\hyperlink{F-QNT-02}{F-QNT-02}.
  \textbf{Phương pháp:} \protect\hyperlink{M-07}{M-07}.
\item
  \textbf{Ví dụ:} \protect\hyperlink{QNT-02-E01}{QNT-02-E01}.
  \textbf{Nguồn:} S1-p69--71.
\end{itemize}

\hypertarget{SET-01}{%
\subsection{SET-01 · Phép toán tập hợp, chứng minh đẳng thức tập
hợp}\label{SET-01}}

\begin{itemize}
\tightlist
\item
  \textbf{Từ khóa:} ``hợp, giao, hiệu, bù'', ``chứng minh
  \(A\setminus(B\cup C)=\dots\)'', ``De Morgan''.
\item
  \textbf{Kiểm tra trước:} bù lấy trong tập nào (vũ trụ); hiệu không
  giao hoán.
\item
  \textbf{Kiến thức:} \protect\hyperlink{F-SET-01}{F-SET-01}.
  \textbf{Phương pháp:} \protect\hyperlink{M-08}{M-08}.
\item
  \textbf{Ví dụ:} \protect\hyperlink{SET-01-E01}{SET-01-E01}.
  \textbf{Nguồn:} S1-p73--79.
\end{itemize}

\hypertarget{SET-02}{%
\subsection{SET-02 · Tập lũy thừa, tích Đề-các, lực
lượng}\label{SET-02}}

\begin{itemize}
\tightlist
\item
  \textbf{Từ khóa:} ``\(P(X)\)'', ``tập các tập con'',
  ``\(A\times B\)'', ``\(|A\times B|\)'', ``lực lượng''.
\item
  \textbf{Kiểm tra trước:} \(\emptyset\in P(X)\), \(X\in P(X)\);
  \(A\times B\neq B\times A\).
\item
  \textbf{Kiến thức:} \protect\hyperlink{F-SET-02}{F-SET-02}. \textbf{Ví
  dụ:} \protect\hyperlink{SET-02-E01}{SET-02-E01}. \textbf{Nguồn:}
  S1-p80--82.
\end{itemize}

\hypertarget{FUN-01}{%
\subsection{FUN-01 · Ánh xạ: kiểm tra định nghĩa, ảnh và ảnh
ngược}\label{FUN-01}}

\begin{itemize}
\tightlist
\item
  \textbf{Từ khóa:} ``có phải ánh xạ'', ``\(f(A)\)'', ``\(f^{-1}(B)\)'',
  ``ảnh'', ``ảnh ngược / tạo ảnh''.
\item
  \textbf{Kiểm tra trước:} mỗi \(x\) có \textbf{đúng một} ảnh trong
  \(Y\); \(f^{-1}(B)\) tồn tại kể cả khi \(f\) không song ánh.
\item
  \textbf{Kiến thức:} \protect\hyperlink{F-FUN-01}{F-FUN-01}.
  \textbf{Phương pháp:} \protect\hyperlink{M-09}{M-09}.
\item
  \textbf{Ví dụ:} \protect\hyperlink{FUN-01-E01}{FUN-01-E01}.
  \textbf{Nguồn:} S1-p83--87.
\end{itemize}

\hypertarget{FUN-02}{%
\subsection{FUN-02 · Đơn ánh, toàn ánh, song ánh; ánh xạ
ngược}\label{FUN-02}}

\begin{itemize}
\tightlist
\item
  \textbf{Từ khóa:} ``đơn ánh'', ``toàn ánh'', ``song ánh'', ``ánh xạ
  ngược'', ``khả nghịch'', ``1-1'', ``lên''.
\item
  \textbf{Dấu hiệu:} cho \(f:X\to Y\) cụ thể (chú ý \(X\), \(Y\)).
\item
  \textbf{Kiểm tra trước:} \textbf{miền và tập đích} quyết định kết quả
  (\(x^2+1\) đơn ánh trên \(\mathbb N\) nhưng không trên \(\mathbb R\)).
\item
  \textbf{Kiến thức:} \protect\hyperlink{F-FUN-02}{F-FUN-02},
  \protect\hyperlink{F-FUN-03}{F-FUN-03}. \textbf{Phương pháp:}
  \protect\hyperlink{M-09}{M-09}.
\item
  \textbf{Ví dụ:} \protect\hyperlink{FUN-02-E01}{FUN-02-E01},
  \protect\hyperlink{FUN-02-E02}{FUN-02-E02}. \textbf{Nguồn:}
  S1-p88--94.
\end{itemize}

\hypertarget{FUN-03}{%
\subsection{\texorpdfstring{FUN-03 · Tích (hợp) ánh xạ \(g\circ f\),
\(f\circ g\)}{FUN-03 · Tích (hợp) ánh xạ g\textbackslash circ f, f\textbackslash circ g}}\label{FUN-03}}

\begin{itemize}
\tightlist
\item
  \textbf{Từ khóa:} ``tìm \(g\circ f\)'', ``\(f\circ g\)'', ``tích ánh
  xạ'', ``hàm hợp''.
\item
  \textbf{Kiểm tra trước:} \(g\circ f(x)=g(f(x))\) --- áp \(f\) trước;
  với hàm từng khúc, xét điều kiện trên \textbf{\(f(x)\)} chứ không trên
  \(x\).
\item
  \textbf{Kiến thức:} \protect\hyperlink{F-FUN-03}{F-FUN-03}. \textbf{Ví
  dụ:} \protect\hyperlink{FUN-03-E01}{FUN-03-E01}. \textbf{Nguồn:}
  S1-p95--96.
\end{itemize}

\hypertarget{IND-01}{%
\subsection{IND-01 · Chứng minh bằng quy nạp toán học}\label{IND-01}}

\begin{itemize}
\tightlist
\item
  \textbf{Từ khóa:} ``chứng minh với mọi \(n\ge n_0\)'', ``quy nạp'',
  ``\(P(n)\)'', tổng \(1+3+\dots\), ``chia hết'', ``\(2^n<n!\)''.
\item
  \textbf{Kiểm tra trước:} bước cơ sở đúng giá trị \(n_0\) của đề (có
  thể khác 1); bước quy nạp phải \textbf{dùng} giả thiết quy nạp.
\item
  \textbf{Kiến thức:} \protect\hyperlink{F-IND-01}{F-IND-01}.
  \textbf{Phương pháp:} \protect\hyperlink{M-10}{M-10}.
\item
  \textbf{Ví dụ:} \protect\hyperlink{IND-01-E01}{IND-01-E01} \ldots{}
  \protect\hyperlink{IND-01-E04}{IND-01-E04}. \textbf{Nguồn:}
  S1-p97--100, S1-p47.
\end{itemize}

\hypertarget{chux1b0ux1a1ng-ii-phuxe9p-ux111ux1ebfm}{%
\section{Chương II --- Phép
đếm}\label{chux1b0ux1a1ng-ii-phuxe9p-ux111ux1ebfm}}

\hypertarget{CNT-01}{%
\subsection{CNT-01 · Nguyên lý cộng, nhân (đếm số, chia trường
hợp)}\label{CNT-01}}

\begin{itemize}
\tightlist
\item
  \textbf{Từ khóa:} ``có bao nhiêu cách'', ``bao nhiêu số \ldots{} chữ
  số khác nhau'', ``chia hết cho 2/5'', ``chẵn''.
\item
  \textbf{Dấu hiệu:} lập số từ tập chữ số có chứa 0; điều kiện ở chữ số
  cuối.
\item
  \textbf{Kiểm tra trước:} chữ số đầu \(\ne0\); ``khác nhau'' hay không;
  xét chữ số bị ràng buộc \textbf{trước}.
\item
  \textbf{Kiến thức:} \protect\hyperlink{F-CNT-01}{F-CNT-01}.
  \textbf{Phương pháp:} \protect\hyperlink{M-11}{M-11}.
\item
  \textbf{Ví dụ:} \protect\hyperlink{CNT-01-E01}{CNT-01-E01},
  \protect\hyperlink{CNT-01-E02}{CNT-01-E02}. \textbf{Nguồn:} S2-p3--5.
\end{itemize}

\hypertarget{CNT-02}{%
\subsection{CNT-02 · Nguyên lý chuồng bồ câu (Dirichlet)}\label{CNT-02}}

\begin{itemize}
\tightlist
\item
  \textbf{Từ khóa:} ``chứng minh tồn tại'', ``ít nhất'', ``luôn có hai
  phần tử'', ``cùng ngày sinh''.
\item
  \textbf{Kiểm tra trước:} xác định rõ ``chim'' và ``chuồng''; dùng
  \(\lceil n/k\rceil\).
\item
  \textbf{Kiến thức:} \protect\hyperlink{F-CNT-02}{F-CNT-02}.
  \textbf{Phương pháp:} \protect\hyperlink{M-12}{M-12}.
\item
  \textbf{Ví dụ:} \protect\hyperlink{CNT-02-E01}{CNT-02-E01}.
  \textbf{Nguồn:} S2-p6--7.
\end{itemize}

\hypertarget{CNT-03}{%
\subsection{CNT-03 · Nguyên lý bù trừ}\label{CNT-03}}

\begin{itemize}
\tightlist
\item
  \textbf{Từ khóa:} ``học tiếng Anh hoặc tiếng Pháp'', ``chia hết cho 2
  hoặc 3 hoặc 5'', ``ít nhất một'', ``không chia hết cho''.
\item
  \textbf{Kiến thức:} \protect\hyperlink{F-CNT-03}{F-CNT-03}.
  \textbf{Phương pháp:} \protect\hyperlink{M-13}{M-13}.
\item
  \textbf{Ví dụ:} \protect\hyperlink{CNT-03-E01}{CNT-03-E01},
  \protect\hyperlink{CNT-03-E02}{CNT-03-E02}. \textbf{Nguồn:} S2-p8--10.
\end{itemize}

\hypertarget{CNT-04}{%
\subsection{CNT-04 · Hoán vị, chỉnh hợp, tổ hợp}\label{CNT-04}}

\begin{itemize}
\tightlist
\item
  \textbf{Từ khóa:} ``sắp xếp'', ``xếp hàng'' (hoán vị); ``chọn và sắp
  thứ tự'', ``lập số có \(k\) chữ số khác nhau'' (chỉnh hợp); ``chọn
  nhóm'', ``tập con \(k\) phần tử'', ``ủy ban'' (tổ hợp).
\item
  \textbf{Kiểm tra trước:} có \textbf{thứ tự} không? có \textbf{lặp}
  không? (cây quyết định \protect\hyperlink{M-14}{M-14}).
\item
  \textbf{Kiến thức:} \protect\hyperlink{F-CNT-04}{F-CNT-04}. \textbf{Ví
  dụ:} \protect\hyperlink{CNT-04-E01}{CNT-04-E01}. \textbf{Nguồn:}
  S2-p11--16.
\end{itemize}

\hypertarget{CNT-05}{%
\subsection{CNT-05 · Hoán vị lặp (sắp xếp có phần tử giống
nhau)}\label{CNT-05}}

\begin{itemize}
\tightlist
\item
  \textbf{Từ khóa:} ``sắp xếp các chữ cái của từ'', ``chuỗi khác nhau'',
  ``đồ vật giống nhau''.
\item
  \textbf{Kiến thức:} \protect\hyperlink{F-CNT-05}{F-CNT-05}. \textbf{Ví
  dụ:} \protect\hyperlink{CNT-05-E01}{CNT-05-E01}. \textbf{Nguồn:}
  S2-p17--18.
\end{itemize}

\hypertarget{CNT-06}{%
\subsection{CNT-06 · Tổ hợp lặp; số nghiệm nguyên không âm có ràng buộc;
chia vật giống nhau}\label{CNT-06}}

\begin{itemize}
\tightlist
\item
  \textbf{Từ khóa:} ``chọn \(k\) vật từ \(n\) loại (được lặp)'', ``số
  nghiệm nguyên không âm \(x_1+\dots+x_n=k\)'', ``chia \(k\) vật giống
  nhau vào \(n\) hộp'', ``ít nhất'', ``không quá''.
\item
  \textbf{Kiểm tra trước:} ràng buộc ``\(\ge a\)'' → đổi biến;
  ``\(\le b\)'' → \textbf{lấy phần bù} (trừ trường hợp ``\(\ge b+1\)'');
  ``\(>\)'' → đổi thành ``\(\ge\)'' với số nguyên.
\item
  \textbf{Kiến thức:} \protect\hyperlink{F-CNT-06}{F-CNT-06}.
  \textbf{Phương pháp:} \protect\hyperlink{M-15}{M-15},
  \protect\hyperlink{M-13}{M-13}.
\item
  \textbf{Ví dụ:} \protect\hyperlink{CNT-06-E01}{CNT-06-E01},
  \protect\hyperlink{CNT-06-E02}{CNT-06-E02},
  \protect\hyperlink{CNT-06-E03}{CNT-06-E03}. \textbf{Nguồn:}
  S2-p19--25.
\end{itemize}

\hypertarget{REC-01}{%
\subsection{REC-01 · Lập hệ thức đệ qui từ bài toán}\label{REC-01}}

\begin{itemize}
\tightlist
\item
  \textbf{Từ khóa:} ``tìm một hệ thức đệ qui cho \(x_n\)'', ``cầu
  thang'', ``tháp Hà Nội'', ``số cách'', ``chuỗi bit độ dài \(n\)''.
\item
  \textbf{Kiểm tra trước:} các trường hợp phải \textbf{rời nhau và đầy
  đủ}; tìm đủ điều kiện ban đầu (\(x_1\), \(x_2\)\ldots).
\item
  \textbf{Kiến thức:} \protect\hyperlink{F-CNT-01}{F-CNT-01} (nguyên lý
  cộng). \textbf{Phương pháp:} \protect\hyperlink{M-16}{M-16}.
\item
  \textbf{Ví dụ:} \protect\hyperlink{REC-01-E01}{REC-01-E01},
  \protect\hyperlink{REC-01-E02}{REC-01-E02},
  \protect\hyperlink{REC-01-E03}{REC-01-E03}. \textbf{Nguồn:}
  S2-p27--39.
\end{itemize}

\hypertarget{REC-02}{%
\subsection{REC-02 · Giải hệ thức đệ qui tuyến tính thuần nhất (cấp 1,
2)}\label{REC-02}}

\begin{itemize}
\tightlist
\item
  \textbf{Từ khóa:} ``giải hệ thức'', ``vế phải = 0'', ``tìm nghiệm tổng
  quát'', ``nghiệm riêng thỏa \(x_0, x_1\)''.
\item
  \textbf{Kiểm tra trước:} viết đúng PTĐT (bậc = cấp); nghiệm kép →
  \((C_1+nC_2)\lambda_0^n\).
\item
  \textbf{Kiến thức:} \protect\hyperlink{F-REC-01}{F-REC-01},
  \protect\hyperlink{F-REC-02}{F-REC-02}. \textbf{Phương pháp:}
  \protect\hyperlink{M-17}{M-17}.
\item
  \textbf{Ví dụ:} \protect\hyperlink{REC-02-E01}{REC-02-E01},
  \protect\hyperlink{REC-02-E02}{REC-02-E02},
  \protect\hyperlink{REC-02-E03}{REC-02-E03}. \textbf{Nguồn:}
  S2-p41--46.
\end{itemize}

\hypertarget{REC-03}{%
\subsection{\texorpdfstring{REC-03 · Giải hệ thức không thuần nhất, vế
phải
\(f_n=\beta^nP_r(n)\)}{REC-03 · Giải hệ thức không thuần nhất, vế phải f\_n=\textbackslash beta\^{}nP\_r(n)}}\label{REC-03}}

\begin{itemize}
\tightlist
\item
  \textbf{Từ khóa:} vế phải là đa thức theo \(n\), hoặc
  \(\beta^n\times\) đa thức; ``\((18n+12)3^n\)'', ``\(12n+8\)'', ``hằng
  số''.
\item
  \textbf{Kiểm tra trước:} \textbf{\(\beta\) có là nghiệm của PTĐT
  không, bội mấy?} → nhân \(n^0, n^1, n^2\); đa thức thuần
  (\(\beta=1\)!).
\item
  \textbf{Kiến thức:} \protect\hyperlink{F-REC-03}{F-REC-03},
  \protect\hyperlink{F-REC-05}{F-REC-05}. \textbf{Phương pháp:}
  \protect\hyperlink{M-17}{M-17}.
\item
  \textbf{Ví dụ:} \protect\hyperlink{REC-03-E01}{REC-03-E01} \ldots{}
  \protect\hyperlink{REC-03-E09}{REC-03-E09}. \textbf{Nguồn:}
  S2-p47--63, p67.
\end{itemize}

\hypertarget{REC-04}{%
\subsection{REC-04 · Vế phải là tổng nhiều dạng (chồng chất nghiệm
riêng)}\label{REC-04}}

\begin{itemize}
\tightlist
\item
  \textbf{Từ khóa:} vế phải có dạng ``\(20+(2-n)2^{n-2}+3\cdot4^n\)''.
\item
  \textbf{Kiểm tra trước:} tách thành các dạng 1 theo từng \(\beta\)
  khác nhau; xét bội của \textbf{từng} \(\beta\).
\item
  \textbf{Kiến thức:} \protect\hyperlink{F-REC-04}{F-REC-04}.
  \textbf{Phương pháp:} \protect\hyperlink{M-17}{M-17}.
\item
  \textbf{Ví dụ:} \protect\hyperlink{REC-04-E01}{REC-04-E01}.
  \textbf{Nguồn:} S2-p52, p64--66.
\end{itemize}

\hypertarget{chux1b0ux1a1ng-iii-quan-hux1ec7}{%
\section{Chương III --- Quan hệ}\label{chux1b0ux1a1ng-iii-quan-hux1ec7}}

\hypertarget{REL-01}{%
\subsection{REL-01 · Liệt kê quan hệ; ma trận biểu diễn}\label{REL-01}}

\begin{itemize}
\tightlist
\item
  \textbf{Từ khóa:} ``liệt kê các cặp'', ``ma trận biểu diễn \(M_R\)'',
  ``từ ma trận suy ra các cặp''.
\item
  \textbf{Kiểm tra trước:} hàng = phần tử của \(A\) (tập nguồn), cột =
  phần tử của \(B\); thứ tự liệt kê.
\item
  \textbf{Kiến thức:} \protect\hyperlink{F-REL-01}{F-REL-01}. \textbf{Ví
  dụ:} \protect\hyperlink{REL-01-E01}{REL-01-E01}. \textbf{Nguồn:}
  S3-p4--6, p12--15.
\end{itemize}

\hypertarget{REL-02}{%
\subsection{REL-02 · Kiểm tra tính chất phản xạ, đối xứng, phản xứng,
bắc cầu}\label{REL-02}}

\begin{itemize}
\tightlist
\item
  \textbf{Từ khóa:} ``quan hệ có tính chất gì'', ``phản xạ'', ``đối
  xứng'', ``phản xứng / phản đối xứng'', ``bắc cầu / truyền''.
\item
  \textbf{Kiểm tra trước:} chứng minh ``có'' phải tổng quát; ``không''
  chỉ cần \textbf{một phản ví dụ cụ thể}.
\item
  \textbf{Kiến thức:} \protect\hyperlink{F-REL-01}{F-REL-01}.
  \textbf{Phương pháp:} \protect\hyperlink{M-18}{M-18}.
\item
  \textbf{Ví dụ:} \protect\hyperlink{REL-02-E01}{REL-02-E01},
  \protect\hyperlink{REL-02-E02}{REL-02-E02}. \textbf{Nguồn:} S3-p7--11,
  p16--18.
\end{itemize}

\hypertarget{REL-03}{%
\subsection{REL-03 · Quan hệ tương đương, lớp tương đương, đồng
dư}\label{REL-03}}

\begin{itemize}
\tightlist
\item
  \textbf{Từ khóa:} ``chứng minh quan hệ tương đương'', ``tìm các lớp
  tương đương'', ``\([a]_R\)'', ``đồng dư modulo \(m\)'', ``phân
  hoạch'', ``\(\mathbb Z_m\)''.
\item
  \textbf{Kiến thức:} \protect\hyperlink{F-REL-02}{F-REL-02},
  \protect\hyperlink{F-REL-03}{F-REL-03}. \textbf{Phương pháp:}
  \protect\hyperlink{M-18}{M-18}, \protect\hyperlink{M-19}{M-19}.
\item
  \textbf{Ví dụ:} \protect\hyperlink{REL-03-E01}{REL-03-E01},
  \protect\hyperlink{REL-03-E02}{REL-03-E02}. \textbf{Nguồn:}
  S3-p19--27.
\end{itemize}

\hypertarget{POS-01}{%
\subsection{POS-01 · Chứng minh quan hệ thứ tự; thứ tự toàn
phần}\label{POS-01}}

\begin{itemize}
\tightlist
\item
  \textbf{Từ khóa:} ``có phải poset'', ``quan hệ thứ tự'', ``thứ tự toàn
  phần / tuyến tính'', ``so sánh được''.
\item
  \textbf{Kiểm tra trước:} cần \textbf{phản xạ + phản xứng + bắc cầu};
  tập nguồn (\(\mathbb Z\) hay \(\mathbb Z^+\)) quyết định phản xứng của
  ``\textbar{}''.
\item
  \textbf{Kiến thức:} \protect\hyperlink{F-POS-01}{F-POS-01}.
  \textbf{Phương pháp:} \protect\hyperlink{M-18}{M-18}.
\item
  \textbf{Ví dụ:} \protect\hyperlink{POS-01-E01}{POS-01-E01}.
  \textbf{Nguồn:} S3-p28--35.
\end{itemize}

\hypertarget{POS-02}{%
\subsection{POS-02 · Thứ tự từ điển}\label{POS-02}}

\begin{itemize}
\tightlist
\item
  \textbf{Từ khóa:} ``thứ tự từ điển'', ``so sánh chuỗi'', ``chuỗi
  bit'', ``\(\Sigma^*\)''.
\item
  \textbf{Kiểm tra trước:} chuỗi là \textbf{tiền tố} của chuỗi kia thì
  nhỏ hơn; so sánh tại vị trí khác đầu tiên.
\item
  \textbf{Kiến thức:} \protect\hyperlink{F-POS-02}{F-POS-02}. \textbf{Ví
  dụ:} \protect\hyperlink{POS-02-E01}{POS-02-E01}. \textbf{Nguồn:}
  S3-p36--41.
\end{itemize}

\hypertarget{POS-03}{%
\subsection{POS-03 · Biểu đồ Hasse; tối đại, tối tiểu, lớn nhất, nhỏ
nhất}\label{POS-03}}

\begin{itemize}
\tightlist
\item
  \textbf{Từ khóa:} ``vẽ biểu đồ Hasse'', ``phần tử tối đại / tối
  tiểu'', ``lớn nhất / nhỏ nhất''.
\item
  \textbf{Kiểm tra trước:} chỉ vẽ cung \textbf{trội trực tiếp}; tối đại
  ≠ lớn nhất (có thể nhiều tối đại, không có lớn nhất).
\item
  \textbf{Kiến thức:} \protect\hyperlink{F-POS-03}{F-POS-03}.
  \textbf{Phương pháp:} \protect\hyperlink{M-20}{M-20}.
\item
  \textbf{Ví dụ:} \protect\hyperlink{POS-03-E01}{POS-03-E01}.
  \textbf{Nguồn:} S3-p42--49.
\end{itemize}

\hypertarget{POS-04}{%
\subsection{POS-04 · Chặn trên, chặn dưới, sup, inf}\label{POS-04}}

\begin{itemize}
\tightlist
\item
  \textbf{Từ khóa:} ``chặn trên / chặn dưới của \(A\)'', ``chặn trên nhỏ
  nhất'', ``chặn dưới lớn nhất'', ``sup'', ``inf'', ``\(a\vee b\),
  \(a\wedge b\)''.
\item
  \textbf{Kiểm tra trước:} chặn có thể nằm ngoài \(A\); sup có thể không
  tồn tại.
\item
  \textbf{Kiến thức:} \protect\hyperlink{F-POS-03}{F-POS-03}.
  \textbf{Phương pháp:} \protect\hyperlink{M-20}{M-20}.
\item
  \textbf{Ví dụ:} \protect\hyperlink{POS-04-E01}{POS-04-E01}.
  \textbf{Nguồn:} S3-p50--53.
\end{itemize}

\hypertarget{POS-05}{%
\subsection{POS-05 · Sắp xếp topo}\label{POS-05}}

\begin{itemize}
\tightlist
\item
  \textbf{Từ khóa:} ``sắp xếp topo'', ``thứ tự thực hiện công việc'',
  ``mặc quần áo''.
\item
  \textbf{Kiến thức:} \protect\hyperlink{F-POS-03}{F-POS-03}.
  \textbf{Phương pháp:} \protect\hyperlink{M-21}{M-21}.
\item
  \textbf{Ví dụ:} \protect\hyperlink{POS-05-E01}{POS-05-E01}.
  \textbf{Nguồn:} S3-p54--56.
\end{itemize}

\hypertarget{chux1b0ux1a1ng-iv-ux111ux1ea1i-sux1ed1-bool}{%
\section{Chương IV --- Đại số
Bool}\label{chux1b0ux1a1ng-iv-ux111ux1ea1i-sux1ed1-bool}}

\hypertarget{BOO-01}{%
\subsection{BOO-01 · Hàm Bool từ mô tả; bảng chân trị; dạng nối rời
chính tắc}\label{BOO-01}}

\begin{itemize}
\tightlist
\item
  \textbf{Từ khóa:} ``lập bảng chân trị của \(f\)'', ``bỏ phiếu'', ``đa
  số'', ``dạng nối rời chính tắc'', ``tổng các từ tối tiểu''.
\item
  \textbf{Kiến thức:} \protect\hyperlink{F-BOO-01}{F-BOO-01},
  \protect\hyperlink{F-BOO-02}{F-BOO-02}. \textbf{Phương pháp:}
  \protect\hyperlink{M-23}{M-23}.
\item
  \textbf{Ví dụ:} \protect\hyperlink{BOO-01-E01}{BOO-01-E01},
  \protect\hyperlink{BOO-01-E02}{BOO-01-E02}. \textbf{Nguồn:}
  S4-p10--19.
\end{itemize}

\hypertarget{BOO-02}{%
\subsection{BOO-02 · Biểu đồ Karnaugh: vẽ kar(f), đọc tế bào ↔ đơn
thức}\label{BOO-02}}

\begin{itemize}
\tightlist
\item
  \textbf{Từ khóa:} ``vẽ biểu đồ Karnaugh'', ``tế bào sau là biểu đồ của
  đơn thức nào''.
\item
  \textbf{Kiểm tra trước:} đúng \textbf{khung nhãn của nguồn} (cột
  \(x\bar y, xy, \bar xy, \bar x\bar y\); hàng
  \(z\bar t, zt, \bar zt, \bar z\bar t\)); ô đầu--cuối hàng/cột là kề
  nhau.
\item
  \textbf{Kiến thức:} \protect\hyperlink{F-BOO-03}{F-BOO-03},
  \protect\hyperlink{F-BOO-04}{F-BOO-04}. \textbf{Ví dụ:}
  \protect\hyperlink{BOO-02-E01}{BOO-02-E01}. \textbf{Nguồn:}
  S4-p22--32.
\end{itemize}

\hypertarget{BOO-03}{%
\subsection{BOO-03 · Tìm công thức đa thức tối tiểu bằng
Karnaugh}\label{BOO-03}}

\begin{itemize}
\tightlist
\item
  \textbf{Từ khóa:} ``tìm (tất cả) công thức đa thức tối tiểu'', ``tối
  thiểu hóa'', ``rút gọn hàm Bool'', ``tế bào lớn''.
\item
  \textbf{Kiểm tra trước:} liệt kê \textbf{đủ} tế bào lớn; tế bào bắt
  buộc; ô còn lại có thể có nhiều cách phủ → so sánh độ đơn giản.
\item
  \textbf{Kiến thức:} \protect\hyperlink{F-BOO-03}{F-BOO-03},
  \protect\hyperlink{F-BOO-04}{F-BOO-04}. \textbf{Phương pháp:}
  \protect\hyperlink{M-22}{M-22}.
\item
  \textbf{Ví dụ:} \protect\hyperlink{BOO-03-E01}{BOO-03-E01} \ldots{}
  \protect\hyperlink{BOO-03-E05}{BOO-03-E05}. \textbf{Nguồn:}
  S4-p33--65.
\end{itemize}

\hypertarget{BOO-04}{%
\subsection{BOO-04 · Mạng logic: vẽ mạch, đọc mạch, thiết kế mạch công
tắc}\label{BOO-04}}

\begin{itemize}
\tightlist
\item
  \textbf{Từ khóa:} ``vẽ mạch'', ``cổng AND/OR/NOT/NAND/NOR'', ``viết
  biểu thức \(f\) của sơ đồ'', ``thiết kế mạch điều khiển bởi \ldots{}
  cầu dao''.
\item
  \textbf{Kiến thức:} \protect\hyperlink{F-BOO-05}{F-BOO-05}.
  \textbf{Phương pháp:} \protect\hyperlink{M-23}{M-23}.
\item
  \textbf{Ví dụ:} \protect\hyperlink{BOO-04-E01}{BOO-04-E01} \ldots{}
  \protect\hyperlink{BOO-04-E04}{BOO-04-E04}. \textbf{Nguồn:}
  S4-p66--78.
\end{itemize}

\hypertarget{chux1b0ux1a1ng-v-ux111ux1ed3-thux1ecb}{%
\section{Chương V --- Đồ
thị}\label{chux1b0ux1a1ng-v-ux111ux1ed3-thux1ecb}}

\hypertarget{GRA-01}{%
\subsection{GRA-01 · Phân loại đồ thị; bậc; định lý bắt
tay}\label{GRA-01}}

\begin{itemize}
\tightlist
\item
  \textbf{Từ khóa:} ``bậc của các đỉnh'', ``đơn đồ thị / đa đồ thị / giả
  đồ thị'', ``bậc vào, bậc ra'', ``số cạnh'', ``có tồn tại đồ thị với
  dãy bậc''.
\item
  \textbf{Kiểm tra trước:} khuyên đếm \textbf{2} vào bậc (vô hướng); số
  đỉnh bậc lẻ phải chẵn.
\item
  \textbf{Kiến thức:} \protect\hyperlink{F-GRA-01}{F-GRA-01}. \textbf{Ví
  dụ:} \protect\hyperlink{GRA-01-E01}{GRA-01-E01},
  \protect\hyperlink{GRA-01-E02}{GRA-01-E02}. \textbf{Nguồn:} S5-p3--25.
\end{itemize}

\hypertarget{GRA-02}{%
\subsection{GRA-02 · Ma trận kề}\label{GRA-02}}

\begin{itemize}
\tightlist
\item
  \textbf{Từ khóa:} ``tìm ma trận kề'', ``biểu diễn đồ thị bằng ma
  trận'', ``vẽ đồ thị từ ma trận''.
\item
  \textbf{Kiểm tra trước:} cạnh song song → số \textgreater{} 1; khuyên
  → quy ước (CON-01); vô hướng → ma trận đối xứng.
\item
  \textbf{Kiến thức:} \protect\hyperlink{F-GRA-02}{F-GRA-02}. \textbf{Ví
  dụ:} \protect\hyperlink{GRA-02-E01}{GRA-02-E01}. \textbf{Nguồn:}
  S5-p26--28.
\end{itemize}

\hypertarget{GRA-03}{%
\subsection{GRA-03 · Đẳng cấu đồ thị}\label{GRA-03}}

\begin{itemize}
\tightlist
\item
  \textbf{Từ khóa:} ``hai đồ thị có đẳng cấu không'', ``tìm song ánh''.
\item
  \textbf{Kiểm tra trước:} bất biến (số đỉnh, số cạnh, dãy bậc, \ldots)
  --- khác ⇒ \textbf{không} đẳng cấu; giống \textbf{chưa đủ} ⇒ phải chỉ
  ra song ánh.
\item
  \textbf{Kiến thức:} \protect\hyperlink{F-GRA-03}{F-GRA-03}.
  \textbf{Phương pháp:} \protect\hyperlink{M-24}{M-24}. \textbf{Ví dụ:}
  \protect\hyperlink{GRA-03-E01}{GRA-03-E01}. \textbf{Nguồn:}
  S5-p29--35.
\end{itemize}

\hypertarget{GRA-04}{%
\subsection{GRA-04 · Đường đi, chu trình, liên thông, đỉnh khớp,
cầu}\label{GRA-04}}

\begin{itemize}
\tightlist
\item
  \textbf{Từ khóa:} ``đường đi đơn / sơ cấp'', ``chu trình sơ cấp'',
  ``thành phần liên thông'', ``đỉnh khớp'', ``cầu''.
\item
  \textbf{Kiến thức:} định nghĩa S5-p36--41. \textbf{Ví dụ:}
  \protect\hyperlink{GRA-04-E01}{GRA-04-E01}. \textbf{Nguồn:}
  S5-p36--42.
\end{itemize}

\hypertarget{GRA-05}{%
\subsection{GRA-05 · Đường đi / chu trình Euler; thuật toán
Fleury}\label{GRA-05}}

\begin{itemize}
\tightlist
\item
  \textbf{Từ khóa:} ``đi qua mỗi cạnh đúng một lần'', ``vẽ bằng một
  nét'', ``chu trình Euler'', ``cây cầu Königsberg'', ``Fleury''.
\item
  \textbf{Kiểm tra trước:} \textbf{liên thông} (bỏ qua đỉnh cô lập) +
  đếm đỉnh bậc lẻ: 0 → chu trình; 2 → đường đi (bắt đầu ở đỉnh lẻ).
\item
  \textbf{Kiến thức:} \protect\hyperlink{F-GRA-04}{F-GRA-04}.
  \textbf{Phương pháp:} \protect\hyperlink{M-25}{M-25}. \textbf{Ví dụ:}
  \protect\hyperlink{GRA-05-E01}{GRA-05-E01},
  \protect\hyperlink{GRA-05-E02}{GRA-05-E02}. \textbf{Nguồn:}
  S5-p43--51.
\end{itemize}

\hypertarget{GRA-06}{%
\subsection{GRA-06 · Đường đi ngắn nhất --- Dijkstra}\label{GRA-06}}

\begin{itemize}
\tightlist
\item
  \textbf{Từ khóa:} ``đường đi ngắn nhất'', ``khoảng cách'',
  ``Dijkstra'', ``trọng số dương'', ``ma trận trọng số''.
\item
  \textbf{Kiểm tra trước:} \textbf{mọi trọng số \textgreater{} 0} (nếu
  có cạnh âm → dùng Ford--Bellman); vô hướng hay có hướng.
\item
  \textbf{Kiến thức:} \protect\hyperlink{F-GRA-05}{F-GRA-05}.
  \textbf{Phương pháp:} \protect\hyperlink{M-26}{M-26}. \textbf{Ví dụ:}
  \protect\hyperlink{GRA-06-E01}{GRA-06-E01},
  \protect\hyperlink{GRA-06-E02}{GRA-06-E02},
  \protect\hyperlink{GRA-06-E03}{GRA-06-E03}. \textbf{Nguồn:}
  S5-p52--72.
\end{itemize}

\hypertarget{GRA-07}{%
\subsection{GRA-07 · Đường đi ngắn nhất --- Ford--Bellman; phát hiện
mạch âm}\label{GRA-07}}

\begin{itemize}
\tightlist
\item
  \textbf{Từ khóa:} ``Ford--Bellman'', ``trọng số âm'', ``mạch âm'',
  ``chỉ ra đồ thị có mạch âm''.
\item
  \textbf{Kiểm tra trước:} dừng khi \(L_k=L_{k-1}\) (ổn định) hoặc khi
  \(k=n\) (mạch âm).
\item
  \textbf{Kiến thức:} \protect\hyperlink{F-GRA-06}{F-GRA-06}.
  \textbf{Phương pháp:} \protect\hyperlink{M-27}{M-27}. \textbf{Ví dụ:}
  \protect\hyperlink{GRA-07-E01}{GRA-07-E01},
  \protect\hyperlink{GRA-07-E02}{GRA-07-E02}. \textbf{Nguồn:}
  S5-p73--86.
\end{itemize}

\hypertarget{formulas}{%
\chapter{Sổ tay công thức \& định lý (Formula
Handbook)}\label{formulas}}

Mỗi mục: \textbf{Công thức} · \emph{Ký hiệu} · \emph{Điều kiện} ·
\emph{Lưu ý} · \emph{Ví dụ} · \emph{Nguồn} · \emph{Trạng thái}. Trạng
thái \textbf{SV} = có trong nguồn; \textbf{CC} = đã kiểm tra độc lập
(bảng chân trị/vét cạn/máy tính, xem \texttt{10\_Verification}).

\hypertarget{chux1b0ux1a1ng-i-logic-tux1eadp-hux1ee3p-uxe1nh-xux1ea1-quy-nux1ea1p}{%
\section{Chương I --- Logic, tập hợp, ánh xạ, quy
nạp}\label{chux1b0ux1a1ng-i-logic-tux1eadp-hux1ee3p-uxe1nh-xux1ea1-quy-nux1ea1p}}

\hypertarget{F-LOG-01}{%
\subsection{F-LOG-01 · Năm phép toán mệnh đề}\label{F-LOG-01}}

\begin{longtable}[]{@{}
  >{\centering\arraybackslash}p{(\columnwidth - 12\tabcolsep) * \real{0.1429}}
  >{\centering\arraybackslash}p{(\columnwidth - 12\tabcolsep) * \real{0.1429}}
  >{\centering\arraybackslash}p{(\columnwidth - 12\tabcolsep) * \real{0.1429}}
  >{\centering\arraybackslash}p{(\columnwidth - 12\tabcolsep) * \real{0.1429}}
  >{\centering\arraybackslash}p{(\columnwidth - 12\tabcolsep) * \real{0.1429}}
  >{\centering\arraybackslash}p{(\columnwidth - 12\tabcolsep) * \real{0.1429}}
  >{\centering\arraybackslash}p{(\columnwidth - 12\tabcolsep) * \real{0.1429}}@{}}
\toprule\noalign{}
\begin{minipage}[b]{\linewidth}\centering
\(p\)
\end{minipage} & \begin{minipage}[b]{\linewidth}\centering
\(q\)
\end{minipage} & \begin{minipage}[b]{\linewidth}\centering
\(\neg p\)
\end{minipage} & \begin{minipage}[b]{\linewidth}\centering
\(p\wedge q\)
\end{minipage} & \begin{minipage}[b]{\linewidth}\centering
\(p\vee q\)
\end{minipage} & \begin{minipage}[b]{\linewidth}\centering
\(p\to q\)
\end{minipage} & \begin{minipage}[b]{\linewidth}\centering
\(p\leftrightarrow q\)
\end{minipage} \\
\midrule\noalign{}
\endhead
\bottomrule\noalign{}
\endlastfoot
0 & 0 & 1 & 0 & 0 & 1 & 1 \\
0 & 1 & 1 & 0 & 1 & 1 & 0 \\
1 & 0 & 0 & 0 & 1 & \textbf{0} & 0 \\
1 & 1 & 0 & 1 & 1 & 1 & 1 \\
\end{longtable}

\begin{itemize}
\tightlist
\item
  \emph{Ký hiệu:} \(\neg p\) = \(\bar p\); 1 = Đ/T, 0 = S/F. Đọc
  \(p\to q\): ``nếu \(p\) thì \(q\)'', ``\(p\) là điều kiện \textbf{đủ}
  của \(q\)'', ``\(q\) là điều kiện \textbf{cần} của \(p\)''.
\item
  \emph{Lưu ý:} \(p\to q\) \textbf{chỉ sai khi \(p=1,q=0\)} (tiền đề sai
  ⇒ kéo theo luôn đúng). Dạng mệnh đề \(n\) biến có \(2^n\) dòng.
\item
  \emph{Ví dụ:} \protect\hyperlink{LOG-02-E01}{LOG-02-E01}.
  \emph{Nguồn:} S1-p10--19. \emph{Trạng thái:} SV.
\end{itemize}

\hypertarget{F-LOG-02}{%
\subsection{F-LOG-02 · Tương đương logic, hệ quả logic, quy tắc thay
thế}\label{F-LOG-02}}

\begin{itemize}
\tightlist
\item
  \(E\Leftrightarrow F\) khi và chỉ khi \(E\leftrightarrow F\) là
  \textbf{hằng đúng} (cùng bảng chân trị). \(E\Rightarrow F\) (hệ quả)
  khi \(E\to F\) hằng đúng.
\item
  Quy tắc thay thế: thay biểu thức con bởi biểu thức tương đương thì
  được biểu thức tương đương.
\item
  \emph{Ví dụ:} \(\neg(p\vee q)\Rightarrow\neg p\). \emph{Nguồn:}
  S1-p23--25. \emph{Trạng thái:} SV.
\end{itemize}

\hypertarget{F-LOG-03}{%
\subsection{F-LOG-03 · Các luật logic (11 nhóm)}\label{F-LOG-03}}

\begin{longtable}[]{@{}
  >{\raggedright\arraybackslash}p{(\columnwidth - 4\tabcolsep) * \real{0.0476}}
  >{\raggedright\arraybackslash}p{(\columnwidth - 4\tabcolsep) * \real{0.2381}}
  >{\raggedright\arraybackslash}p{(\columnwidth - 4\tabcolsep) * \real{0.7143}}@{}}
\toprule\noalign{}
\begin{minipage}[b]{\linewidth}\raggedright
\#
\end{minipage} & \begin{minipage}[b]{\linewidth}\raggedright
Tên luật
\end{minipage} & \begin{minipage}[b]{\linewidth}\raggedright
Công thức
\end{minipage} \\
\midrule\noalign{}
\endhead
\bottomrule\noalign{}
\endlastfoot
1 & Phủ định kép & \(\neg\neg p\Leftrightarrow p\) \\
2 & De Morgan & \(\neg(p\wedge q)\Leftrightarrow\neg p\vee\neg q\);
\(\ \neg(p\vee q)\Leftrightarrow\neg p\wedge\neg q\) \\
3 & Giao hoán & \(p\vee q\Leftrightarrow q\vee p\);
\(\ p\wedge q\Leftrightarrow q\wedge p\) \\
4 & Kết hợp & \((p\vee q)\vee r\Leftrightarrow p\vee(q\vee r)\); tương
tự với \(\wedge\) \\
5 & Phân phối &
\(p\vee(q\wedge r)\Leftrightarrow(p\vee q)\wedge(p\vee r)\);
\(\ p\wedge(q\vee r)\Leftrightarrow(p\wedge q)\vee(p\wedge r)\) \\
6 & Lũy đẳng & \(p\vee p\Leftrightarrow p\);
\(\ p\wedge p\Leftrightarrow p\) \\
7 & Trung hòa & \(p\vee 0\Leftrightarrow p\);
\(\ p\wedge 1\Leftrightarrow p\) \\
8 & Phần tử bù & \(p\wedge\neg p\Leftrightarrow 0\);
\(\ p\vee\neg p\Leftrightarrow 1\) \\
9 & Thống trị & \(p\wedge 0\Leftrightarrow 0\);
\(\ p\vee 1\Leftrightarrow 1\) \\
10 & Hấp thụ & \(p\vee(p\wedge q)\Leftrightarrow p\);
\(\ p\wedge(p\vee q)\Leftrightarrow p\) \\
11 & Kéo theo &
\(p\to q\Leftrightarrow\neg p\vee q\Leftrightarrow\neg q\to\neg p\) \\
\end{longtable}

\begin{itemize}
\tightlist
\item
  \emph{Lưu ý:} Mệnh đề \textbf{đảo} \(q\to p\) và \textbf{phản đảo}
  \(\neg q\to\neg p\): chỉ phản đảo tương đương với \(p\to q\).
  \textbf{Phủ định} của \(p\to q\) là \(p\wedge\neg q\) (không phải
  \(\neg p\to\neg q\)). Hệ quả hay dùng (CC):
  \(p\leftrightarrow q\Leftrightarrow(p\to q)\wedge(q\to p)\).
\item
  \emph{Ví dụ:} \protect\hyperlink{LOG-03-E01}{LOG-03-E01}.
  \emph{Nguồn:} S1-p26--29. \emph{Trạng thái:} SV; luật đảo/phủ định kéo
  theo: CC (bảng chân trị).
\end{itemize}

\hypertarget{F-LOG-04}{%
\subsection{F-LOG-04 · Các quy tắc suy diễn}\label{F-LOG-04}}

\begin{longtable}[]{@{}
  >{\raggedright\arraybackslash}p{(\columnwidth - 6\tabcolsep) * \real{0.0417}}
  >{\raggedright\arraybackslash}p{(\columnwidth - 6\tabcolsep) * \real{0.2500}}
  >{\raggedright\arraybackslash}p{(\columnwidth - 6\tabcolsep) * \real{0.2917}}
  >{\raggedright\arraybackslash}p{(\columnwidth - 6\tabcolsep) * \real{0.4167}}@{}}
\toprule\noalign{}
\begin{minipage}[b]{\linewidth}\raggedright
\#
\end{minipage} & \begin{minipage}[b]{\linewidth}\raggedright
Tên (nguồn)
\end{minipage} & \begin{minipage}[b]{\linewidth}\raggedright
Sơ đồ
\end{minipage} & \begin{minipage}[b]{\linewidth}\raggedright
Hằng đúng
\end{minipage} \\
\midrule\noalign{}
\endhead
\bottomrule\noalign{}
\endlastfoot
1 & Khẳng định (Modus Ponens) & \(p\to q,\ p\ \therefore q\) &
\([(p\to q)\wedge p]\to q\) \\
2 & Phủ định (Modus Tollens) & \(p\to q,\ \neg q\ \therefore\neg p\) &
\([(p\to q)\wedge\neg q]\to\neg p\) \\
3 & Tam đoạn luận & \(p\to q,\ q\to r\ \therefore p\to r\) &
\([(p\to q)\wedge(q\to r)]\to(p\to r)\) \\
4 & Tam đoạn luận rời & \(p\vee q,\ \neg q\ \therefore p\) &
\([(p\vee q)\wedge\neg q]\to p\) \\
5 & Nối liền & \(p,\ q\ \therefore p\wedge q\) &
\((p\wedge q)\to(p\wedge q)\) \\
6 & Đơn giản & \(p\wedge q\ \therefore p\) & \((p\wedge q)\to p\) \\
7 & Mâu thuẫn (phản chứng) & \(p_1,\dots,p_n,\neg h\ \therefore 0\) &
\((p_1\wedge\dots\wedge p_n)\to h\Leftrightarrow(p_1\wedge\dots\wedge p_n\wedge\neg h)\to 0\) \\
8 & Chứng minh theo trường hợp &
\(p\to r,\ q\to r\ \therefore(p\vee q)\to r\) &
\([(p\to r)\wedge(q\to r)]\to[(p\vee q)\to r]\) \\
9 & Phản ví dụ & Suy luận sai nếu có \textbf{một} bộ chân trị làm mọi
tiền đề = 1 và kết luận = 0 & --- \\
\end{longtable}

\begin{itemize}
\tightlist
\item
  \emph{Lưu ý:} Không có quy tắc ``\(p\to q,\ q\ \therefore p\)'' (khẳng
  định hệ quả --- \textbf{sai}) hay
  ``\(p\to q,\ \neg p\ \therefore\neg q\)'' (phủ định tiền đề ---
  \textbf{sai}). Xem \protect\hyperlink{T-04}{T-04}.
\item
  \emph{Ví dụ:} \protect\hyperlink{LOG-04-E01}{LOG-04-E01},
  \protect\hyperlink{LOG-05-E01}{LOG-05-E01}. \emph{Nguồn:} S1-p31--48.
  \emph{Trạng thái:} SV; hai suy luận sai: CC.
\end{itemize}

\hypertarget{F-QNT-01}{%
\subsection{F-QNT-01 · Lượng từ và hoán vị lượng từ}\label{F-QNT-01}}

\begin{itemize}
\tightlist
\item
  \(\forall x\in A,\ p(x)\) đúng ⇔ \(p(a)\) đúng với \textbf{mọi}
  \(a\in A\). \(\exists x\in A,\ p(x)\) đúng ⇔ có \textbf{ít nhất một}
  \(a_0\) với \(p(a_0)\) đúng.
\item
  Định lý (S1-p68): (1)
  \(\forall x\forall y\,p\Leftrightarrow\forall y\forall x\,p\); (2)
  \(\exists x\exists y\,p\Leftrightarrow\exists y\exists x\,p\); (3)
  \(\exists x\forall y\,p(x,y)\Rightarrow\forall y\exists x\,p(x,y)\)
  --- \textbf{chiều đảo nói chung sai}.
\item
  Đặc biệt hóa phổ dụng: \(\forall x\in A,\ p(x)\) và \(a\in A\)
  \(\therefore p(a)\).
\item
  \emph{Ví dụ:} \protect\hyperlink{QNT-01-E01}{QNT-01-E01}.
  \emph{Nguồn:} S1-p62--68, p72. \emph{Trạng thái:} SV.
\end{itemize}

\hypertarget{F-QNT-02}{%
\subsection{F-QNT-02 · Phủ định mệnh đề lượng từ}\label{F-QNT-02}}

\[
\overline{\forall x\in A,\ p(x)}\Leftrightarrow\exists x\in A,\ \overline{p(x)},\qquad
\overline{\exists x\in A,\ p(x)}\Leftrightarrow\forall x\in A,\ \overline{p(x)}
\]

\[
\overline{\forall x\,\exists y,\ p(x,y)}\Leftrightarrow\exists x\,\forall y,\ \overline{p(x,y)},\qquad
\overline{\exists x\,\forall y,\ p(x,y)}\Leftrightarrow\forall x\,\exists y,\ \overline{p(x,y)}
\]

\begin{itemize}
\tightlist
\item
  \emph{Quy tắc:} đổi \(\forall\leftrightarrow\exists\), \textbf{giữ
  nguyên miền} (kể cả ``\(\varepsilon>0\)''), phủ định vị từ:
  \(\neg(<)\) là \(\ge\); \(\neg(A\to B)\) là \(A\wedge\neg B\).
\item
  \emph{Ví dụ:} \protect\hyperlink{QNT-02-E01}{QNT-02-E01}.
  \emph{Nguồn:} S1-p69--71. \emph{Trạng thái:} SV.
\end{itemize}

\hypertarget{F-SET-01}{%
\subsection{F-SET-01 · Phép toán và luật tập hợp}\label{F-SET-01}}

\begin{itemize}
\tightlist
\item
  \(x\in A\cup B\Leftrightarrow x\in A\vee x\in B\);
  \(\ x\in A\cap B\Leftrightarrow x\in A\wedge x\in B\);
  \(\ x\in A\setminus B\Leftrightarrow x\in A\wedge x\notin B\).
\item
  Bù: khi \(A\subset B\), \(C_BA=B\setminus A\), ký hiệu \(\bar A\).
\item
  Lũy đẳng, giao hoán, kết hợp, \(\emptyset\cup A=A\),
  \(\emptyset\cap A=\emptyset\); phân phối
  \(A\cup(B\cap C)=(A\cup B)\cap(A\cup C)\),
  \(A\cap(B\cup C)=(A\cap B)\cup(A\cap C)\).
\item
  De Morgan: \(\overline{A\cap B}=\bar A\cup\bar B\),
  \(\ \overline{A\cup B}=\bar A\cap\bar B\).
\item
  \emph{Lưu ý:} mỗi luật tập hợp tương ứng một luật logic
  (\protect\hyperlink{F-LOG-03}{F-LOG-03}) qua định nghĩa phần tử.
\item
  \emph{Ví dụ:} \protect\hyperlink{SET-01-E01}{SET-01-E01}.
  \emph{Nguồn:} S1-p76--79. \emph{Trạng thái:} SV.
\end{itemize}

\hypertarget{F-SET-02}{%
\subsection{F-SET-02 · Tập lũy thừa, tích Đề-các}\label{F-SET-02}}

\begin{itemize}
\tightlist
\item
  \(P(X)\) = tập mọi tập con của \(X\); \(|X|=n\Rightarrow|P(X)|=2^n\)
  \textbf{{[}CC{]}} (mỗi phần tử: chọn/không chọn).
\item
  \(A\times B=\{(x,y)\mid x\in A,\ y\in B\}\);
  \(|A\times B|=|A|\cdot|B|\) \textbf{{[}CC{]}} (nguyên lý nhân);
  \(A\times B\ne B\times A\) nói chung.
\item
  \emph{Nguồn:} S1-p80--82 (slide để dạng câu hỏi ``\(|P(X)|=?\)'',
  ``\(|A\times B|=?\)''). \emph{Trạng thái:} công thức CC.
\end{itemize}

\hypertarget{F-FUN-01}{%
\subsection{F-FUN-01 · Ánh xạ, ảnh, ảnh ngược}\label{F-FUN-01}}

\begin{itemize}
\tightlist
\item
  \(f:X\to Y\) là ánh xạ ⇔
  \(\forall x\in X,\ \exists!\,y\in Y:\ y=f(x)\).
\item
  \(f(A)=\{f(x)\mid x\in A\}\); \(\ f^{-1}(B)=\{x\in X\mid f(x)\in B\}\)
  (\textbf{luôn định nghĩa được}, có thể \(=\emptyset\)).
\item
  \(f=g\) ⇔ cùng \(X,Y\) và \(f(x)=g(x)\ \forall x\in X\).
\item
  \emph{Ví dụ:} \protect\hyperlink{FUN-01-E01}{FUN-01-E01}.
  \emph{Nguồn:} S1-p83--87. \emph{Trạng thái:} SV.
\end{itemize}

\hypertarget{F-FUN-02}{%
\subsection{F-FUN-02 · Đơn ánh, toàn ánh, song ánh}\label{F-FUN-02}}

\begin{longtable}[]{@{}
  >{\raggedright\arraybackslash}p{(\columnwidth - 6\tabcolsep) * \real{0.1200}}
  >{\raggedright\arraybackslash}p{(\columnwidth - 6\tabcolsep) * \real{0.2400}}
  >{\raggedright\arraybackslash}p{(\columnwidth - 6\tabcolsep) * \real{0.4000}}
  >{\raggedright\arraybackslash}p{(\columnwidth - 6\tabcolsep) * \real{0.2400}}@{}}
\toprule\noalign{}
\begin{minipage}[b]{\linewidth}\raggedright
Loại
\end{minipage} & \begin{minipage}[b]{\linewidth}\raggedright
Định nghĩa
\end{minipage} & \begin{minipage}[b]{\linewidth}\raggedright
Tương đương theo phương trình \(f(x)=y\) (\(y\in Y\) là tham số)
\end{minipage} & \begin{minipage}[b]{\linewidth}\raggedright
Phủ định
\end{minipage} \\
\midrule\noalign{}
\endhead
\bottomrule\noalign{}
\endlastfoot
Đơn ánh & \(f(x)=f(x')\Rightarrow x=x'\) & có \textbf{nhiều nhất 1}
nghiệm \(x\in X\) & \(\exists x\ne x'\): \(f(x)=f(x')\) \\
Toàn ánh & \(f(X)=Y\) & có \textbf{ít nhất 1} nghiệm \(x\in X\) &
\(\exists y\): phương trình vô nghiệm \\
Song ánh & đơn ánh + toàn ánh & có \textbf{đúng 1} nghiệm \(x\in X\)
& \\
\end{longtable}

\begin{itemize}
\tightlist
\item
  \emph{Ví dụ:} \protect\hyperlink{FUN-02-E01}{FUN-02-E01}.
  \emph{Nguồn:} S1-p88--93. \emph{Trạng thái:} SV.
\end{itemize}

\hypertarget{F-FUN-03}{%
\subsection{F-FUN-03 · Ánh xạ ngược, tích ánh xạ}\label{F-FUN-03}}

\begin{itemize}
\tightlist
\item
  Chỉ \textbf{song ánh} mới có ánh xạ ngược: \(f^{-1}:Y\to X\),
  \(f^{-1}(y)=x\) với \(f(x)=y\).
\item
  \(g\circ f:X\to Z\), \((g\circ f)(x)=g(f(x))\), cần \(f(X)\subset\)
  miền của \(g\). Nói chung \(g\circ f\ne f\circ g\).
\item
  \emph{Ví dụ:} \protect\hyperlink{FUN-02-E01}{FUN-02-E01},
  \protect\hyperlink{FUN-03-E01}{FUN-03-E01}. \emph{Nguồn:} S1-p94--96
  (ERR-01). \emph{Trạng thái:} SV.
\end{itemize}

\hypertarget{F-IND-01}{%
\subsection{F-IND-01 · Nguyên lý quy nạp toán học}\label{F-IND-01}}

\[
\Big[P(N_0)\ \wedge\ \forall k\ge N_0\ \big(P(k)\to P(k+1)\big)\Big]\ \Longrightarrow\ \forall n\ge N_0,\ P(n)
\]

\begin{itemize}
\tightlist
\item
  \emph{Lưu ý:} \(N_0\) lấy theo đề (\(n\ge1\), \(n>3\) ⇒ \(N_0=4\)).
  \emph{Ví dụ:} \protect\hyperlink{IND-01-E01}{IND-01-E01}.
  \emph{Nguồn:} S1-p97. \emph{Trạng thái:} SV.
\end{itemize}

\hypertarget{chux1b0ux1a1ng-ii-phuxe9p-ux111ux1ebfm-hux1ec7-thux1ee9c-ux111ux1ec7-qui}{%
\section{Chương II --- Phép đếm, hệ thức đệ
qui}\label{chux1b0ux1a1ng-ii-phuxe9p-ux111ux1ebfm-hux1ec7-thux1ee9c-ux111ux1ec7-qui}}

\hypertarget{F-CNT-01}{%
\subsection{F-CNT-01 · Nguyên lý cộng, nguyên lý nhân}\label{F-CNT-01}}

\begin{itemize}
\tightlist
\item
  \textbf{Cộng:} công việc làm bằng \textbf{một trong} các phương pháp
  rời nhau (\(n\) cách hoặc \(m\) cách) ⇒ \(n+m\) cách.
\item
  \textbf{Nhân:} công việc gồm \textbf{các bước liên tiếp} (\(n\) cách
  rồi \(m\) cách) ⇒ \(n\cdot m\) cách.
\item
  \emph{Ví dụ:} \protect\hyperlink{CNT-01-E01}{CNT-01-E01}.
  \emph{Nguồn:} S2-p3--5. \emph{Trạng thái:} SV.
\end{itemize}

\hypertarget{F-CNT-02}{%
\subsection{F-CNT-02 · Nguyên lý chuồng bồ câu
(Dirichlet)}\label{F-CNT-02}}

\(n\) chim trong \(k\) chuồng ⇒ có chuồng chứa \textbf{ít nhất}
\(\left\lceil n/k\right\rceil\) chim (\(\lceil x\rceil\) = số nguyên nhỏ
nhất \(\ge x\)). - \emph{Ví dụ:} \(\lceil 20/7\rceil=3\);
\protect\hyperlink{CNT-02-E01}{CNT-02-E01}. \emph{Nguồn:} S2-p6--7.
\emph{Trạng thái:} SV.

\hypertarget{F-CNT-03}{%
\subsection{F-CNT-03 · Nguyên lý bù trừ}\label{F-CNT-03}}

\[
\lvert A\cup B\rvert=\lvert A\rvert+\lvert B\rvert-\lvert A\cap B\rvert
\]

\[
\begin{aligned}
\lvert A\cup B\cup C\rvert={}&\lvert A\rvert+\lvert B\rvert+\lvert C\rvert-\lvert A\cap B\rvert-\lvert A\cap C\rvert-\lvert B\cap C\rvert\\
&+\lvert A\cap B\cap C\rvert
\end{aligned}
\]

\begin{itemize}
\tightlist
\item
  Tổng quát (4 tập): cộng các \(|A_i|\), trừ các giao đôi, cộng các giao
  ba, trừ giao bốn.
\item
  \emph{Nguồn:} S2-p8--10 (slide để câu hỏi ``\(|A\cup B\cup C|=?\)''
  --- công thức 3 tập là \textbf{CC} bằng đếm vét cạn). \emph{Ví dụ:}
  \protect\hyperlink{CNT-03-E02}{CNT-03-E02}.
\end{itemize}

\hypertarget{F-CNT-04}{%
\subsection{F-CNT-04 · Hoán vị, chỉnh hợp, tổ hợp}\label{F-CNT-04}}

\begin{longtable}[]{@{}
  >{\raggedright\arraybackslash}p{(\columnwidth - 6\tabcolsep) * \real{0.2500}}
  >{\raggedright\arraybackslash}p{(\columnwidth - 6\tabcolsep) * \real{0.2500}}
  >{\raggedright\arraybackslash}p{(\columnwidth - 6\tabcolsep) * \real{0.2500}}
  >{\raggedright\arraybackslash}p{(\columnwidth - 6\tabcolsep) * \real{0.2500}}@{}}
\toprule\noalign{}
\begin{minipage}[b]{\linewidth}\raggedright
Khái niệm
\end{minipage} & \begin{minipage}[b]{\linewidth}\raggedright
Thứ tự?
\end{minipage} & \begin{minipage}[b]{\linewidth}\raggedright
Lặp?
\end{minipage} & \begin{minipage}[b]{\linewidth}\raggedright
Công thức
\end{minipage} \\
\midrule\noalign{}
\endhead
\bottomrule\noalign{}
\endlastfoot
Hoán vị \(n\) phần tử & có & không & \(P_n=n!\), \(0!=1\) \\
Chỉnh hợp chập \(k\) của \(n\) (\(1\le k\le n\)) & có & không &
\(A_n^k=\dfrac{n!}{(n-k)!}\) \\
Tổ hợp chập \(k\) của \(n\) & không & không &
\(C_n^k=\binom nk=\dfrac{n!}{k!(n-k)!}\) \\
\end{longtable}

\begin{itemize}
\tightlist
\item
  Tính chất: \(C_n^{n-k}=C_n^k\); \(\ C_n^k+C_n^{k-1}=C_{n+1}^k\)
  (Pascal).
\item
  Số song ánh \(A\to A\) với \(|A|=n\) là \(n!\).
\item
  \emph{Ví dụ:} \protect\hyperlink{CNT-04-E01}{CNT-04-E01}.
  \emph{Nguồn:} S2-p11--16. \emph{Trạng thái:} SV.
\end{itemize}

\hypertarget{F-CNT-05}{%
\subsection{F-CNT-05 · Hoán vị lặp}\label{F-CNT-05}}

\[
\frac{n!}{n_1!\,n_2!\cdots n_k!},\qquad n_1+n_2+\dots+n_k=n
\]

\begin{itemize}
\tightlist
\item
  \emph{Ký hiệu:} \(n_i\) = số đối tượng giống hệt nhau loại \(i\).
  \emph{Ví dụ:} SUCCESS → \(\frac{7!}{3!1!2!1!}=420\),
  \protect\hyperlink{CNT-05-E01}{CNT-05-E01}. \emph{Nguồn:} S2-p17--18.
  \emph{Trạng thái:} SV, CC.
\end{itemize}

\hypertarget{F-CNT-06}{%
\subsection{F-CNT-06 · Tổ hợp lặp; nghiệm nguyên không
âm}\label{F-CNT-06}}

\[
K_n^k=C_{n+k-1}^{k}
\]

\begin{itemize}
\tightlist
\item
  \(K_n^k\) = số cách chọn \(k\) vật từ \(n\) \textbf{loại} (được chọn
  lại) = số nghiệm nguyên không âm của \(x_1+\dots+x_n=k\) = số cách
  chia \(k\) vật giống nhau vào \(n\) hộp phân biệt.
\item
  \emph{Ràng buộc:} \(x_i\ge a_i\) ⇒ đặt \(x_i'=x_i-a_i\), vế phải giảm
  \(\sum a_i\). \(x_i\le b\) ⇒ (tổng không ràng buộc trên) − (số nghiệm
  với \(x_i\ge b+1\)).
\item
  \emph{Lưu ý:} \(n\) = số \textbf{ẩn/loại/hộp}, \(k\) = vế phải/số vật
  --- đừng đảo. \emph{Ví dụ:}
  \protect\hyperlink{CNT-06-E02}{CNT-06-E02}. \emph{Nguồn:} S2-p19--24.
  \emph{Trạng thái:} SV, CC.
\end{itemize}

\hypertarget{F-REC-01}{%
\subsection{\texorpdfstring{F-REC-01 · Hệ thức đệ qui tuyến tính cấp
\(k\); cấp
1}{F-REC-01 · Hệ thức đệ qui tuyến tính cấp k; cấp 1}}\label{F-REC-01}}

\[
a_0x_n+a_1x_{n-1}+\dots+a_kx_{n-k}=f_n\quad(a_0\ne0)\qquad\text{PTĐT: } a_0\lambda^k+a_1\lambda^{k-1}+\dots+a_k=0
\]

\begin{itemize}
\tightlist
\item
  \(f_n\equiv0\): thuần nhất. Nghiệm xác định duy nhất bởi \(k\) giá trị
  đầu \(x_0,\dots,x_{k-1}\).
\item
  \textbf{Cấp 1:} \(a_0x_n+a_1x_{n-1}=0\) ⇒ \(\lambda_0=-a_1/a_0\),
  \(x_n=C\lambda_0^n\).
\item
  \emph{Ví dụ:} \protect\hyperlink{REC-02-E01}{REC-02-E01}.
  \emph{Nguồn:} S2-p29--32, p41--42 (ERR-02). \emph{Trạng thái:} SV.
\end{itemize}

\hypertarget{F-REC-02}{%
\subsection{F-REC-02 · Thuần nhất cấp 2}\label{F-REC-02}}

PTĐT \(a_0\lambda^2+a_1\lambda+a_2=0\): - Hai nghiệm \textbf{thực phân
biệt} \(\lambda_1\ne\lambda_2\):
\(\ x_n=C_1\lambda_1^n+C_2\lambda_2^n\). - \textbf{Nghiệm kép} thực
\(\lambda_0\): \(\ x_n=(C_1+nC_2)\lambda_0^n\). - \emph{Điều kiện:}
slide chỉ xét nghiệm thực (GAP-05). \emph{Ví dụ:}
\protect\hyperlink{REC-02-E02}{REC-02-E02}. \emph{Nguồn:} S2-p43--46.
\emph{Trạng thái:} SV, CC.

\hypertarget{F-REC-03}{%
\subsection{\texorpdfstring{F-REC-03 · Nghiệm riêng khi
\(f_n=\beta^nP_r(n)\) (Dạng
1)}{F-REC-03 · Nghiệm riêng khi f\_n=\textbackslash beta\^{}nP\_r(n) (Dạng 1)}}\label{F-REC-03}}

\begin{longtable}[]{@{}ll@{}}
\toprule\noalign{}
Trường hợp & Dạng nghiệm riêng \\
\midrule\noalign{}
\endhead
\bottomrule\noalign{}
\endlastfoot
\(\beta\) \textbf{không} là nghiệm PTĐT & \(x_n=\beta^nQ_r(n)\) \\
\(\beta\) là nghiệm \textbf{đơn} & \(x_n=n\,\beta^nQ_r(n)\) \\
\(\beta\) là nghiệm \textbf{kép} & \(x_n=n^2\beta^nQ_r(n)\) \\
\end{longtable}

\begin{itemize}
\tightlist
\item
  \(Q_r(n)=A_rn^r+\dots+A_0\): đa thức tổng quát \textbf{cùng bậc \(r\)}
  với \(P_r(n)\); tìm \(r+1\) hệ số bằng cách thế vào (1) và cho \(n\)
  nhận \(r+1\) giá trị (hoặc đồng nhất hệ số).
\item
  \emph{Lưu ý:} vế phải \textbf{chỉ là đa thức} ⇒ \(\beta=1\); vế phải
  \(2\cdot5^{n+1}=10\cdot5^n\) ⇒ \(\beta=5\), \(r=0\);
  \((2-n)2^{n-2}=\frac{2-n}{4}2^n\) ⇒ \(\beta=2\), \(r=1\).
\item
  \emph{Ví dụ:} \protect\hyperlink{REC-03-E01}{REC-03-E01},
  \protect\hyperlink{REC-03-E02}{REC-03-E02}. \emph{Nguồn:} S2-p48--51.
  \emph{Trạng thái:} SV, CC.
\end{itemize}

\hypertarget{F-REC-04}{%
\subsection{F-REC-04 · Chồng chất nghiệm riêng (Dạng
2)}\label{F-REC-04}}

\(f_n=f_{n1}+\dots+f_{ns}\) (mỗi \(f_{ni}\) thuộc Dạng 1) ⇒
\(x_n=x_{n1}+\dots+x_{ns}\) với \(x_{ni}\) là nghiệm riêng của vế trái
\(=f_{ni}\). - \emph{Ví dụ:} \protect\hyperlink{REC-04-E01}{REC-04-E01}.
\emph{Nguồn:} S2-p52, p64--66 (ERR-09). \emph{Trạng thái:} SV, CC.

\hypertarget{F-REC-05}{%
\subsection{F-REC-05 · Cấu trúc nghiệm hệ không thuần
nhất}\label{F-REC-05}}

\[
\text{Nghiệm TQ của (1)}=\text{Nghiệm TQ của thuần nhất (2)}+\text{Một nghiệm riêng của (1)}
\]

\begin{itemize}
\tightlist
\item
  \textbf{Thứ tự bắt buộc:} tìm hằng số \(C_1,C_2\) từ điều kiện đầu
  \textbf{sau khi} đã cộng nghiệm riêng
  (\protect\hyperlink{T-14}{T-14}).
\item
  \emph{Nguồn:} S2-p47. \emph{Trạng thái:} SV.
\end{itemize}

\hypertarget{chux1b0ux1a1ng-iii-quan-hux1ec7-1}{%
\section{Chương III --- Quan
hệ}\label{chux1b0ux1a1ng-iii-quan-hux1ec7-1}}

\hypertarget{F-REL-01}{%
\subsection{F-REL-01 · Quan hệ, ma trận biểu diễn, bốn tính
chất}\label{F-REL-01}}

\begin{itemize}
\tightlist
\item
  \(R\subseteq A\times B\); \(aRb\Leftrightarrow(a,b)\in R\).
  \(M_R=[m_{ij}]\), \(m_{ij}=1\) nếu \((a_i,b_j)\in R\), ngược lại 0.
\end{itemize}

\begin{longtable}[]{@{}
  >{\raggedright\arraybackslash}p{(\columnwidth - 4\tabcolsep) * \real{0.2273}}
  >{\raggedright\arraybackslash}p{(\columnwidth - 4\tabcolsep) * \real{0.2727}}
  >{\raggedright\arraybackslash}p{(\columnwidth - 4\tabcolsep) * \real{0.5000}}@{}}
\toprule\noalign{}
\begin{minipage}[b]{\linewidth}\raggedright
Tính chất (R trên A)
\end{minipage} & \begin{minipage}[b]{\linewidth}\raggedright
Định nghĩa
\end{minipage} & \begin{minipage}[b]{\linewidth}\raggedright
Dấu hiệu trên \(M_R\)
\end{minipage} \\
\midrule\noalign{}
\endhead
\bottomrule\noalign{}
\endlastfoot
Phản xạ & \(\forall a,\ aRa\) & đường chéo toàn 1 \\
Đối xứng & \(aRb\to bRa\) & \(M_R\) đối xứng: \(m_{ij}=m_{ji}\) \\
Phản xứng & \(aRb\wedge bRa\to a=b\) & \(i\ne j\): không đồng thời
\(m_{ij}=m_{ji}=1\) \\
Bắc cầu & \(aRb\wedge bRc\to aRc\) & \(M_R\odot M_R\le M_R\) (tích Bool)
\textbf{{[}CC{]}} \\
\end{longtable}

\begin{itemize}
\tightlist
\item
  \emph{Lưu ý:} đối xứng và phản xứng \textbf{không} loại trừ nhau (ví
  dụ quan hệ ``='', hoặc \(\emptyset\)). Dấu hiệu bắc cầu bằng tích Bool
  không có trong slide --- dùng để kiểm tra.
\item
  \emph{Ví dụ:} \protect\hyperlink{REL-02-E01}{REL-02-E01}.
  \emph{Nguồn:} S3-p4--18. \emph{Trạng thái:} SV.
\end{itemize}

\hypertarget{F-REL-02}{%
\subsection{F-REL-02 · Quan hệ tương đương, lớp tương đương, phân
hoạch}\label{F-REL-02}}

\begin{itemize}
\tightlist
\item
  Tương đương = phản xạ + đối xứng + bắc cầu.
  \([a]_R=\{b\in A\mid bRa\}\).
\item
  Định lý (viết dạng ⇔, xem ERR-03): \(aRb\Leftrightarrow[a]_R=[b]_R\);
  \(\ [a]_R\ne[b]_R\Leftrightarrow[a]_R\cap[b]_R=\emptyset\). Các lớp
  tương đương tạo \textbf{phân hoạch} của \(A\); ngược lại mỗi phân
  hoạch xác định duy nhất một quan hệ tương đương.
\item
  \emph{Ví dụ:} \protect\hyperlink{REL-03-E02}{REL-03-E02}.
  \emph{Nguồn:} S3-p21--26. \emph{Trạng thái:} SV (dạng ⇔: CC).
\end{itemize}

\hypertarget{F-REL-03}{%
\subsection{\texorpdfstring{F-REL-03 · Đồng dư modulo
\(m\)}{F-REL-03 · Đồng dư modulo m}}\label{F-REL-03}}

\begin{itemize}
\tightlist
\item
  \(a\equiv b\pmod m\Leftrightarrow m\mid(a-b)\) (\(m\) nguyên dương).
  Là quan hệ tương đương trên \(\mathbb Z\); có đúng \(m\) lớp
  \([0]_m,\dots,[m-1]_m\); \(\mathbb Z_m=\{[0]_m,\dots,[m-1]_m\}\);
  \([r]_m=[r+m]_m=\dots\)
\item
  \emph{Ví dụ:} \([1]_8=\{\dots,-15,-7,1,9,17,\dots\}\);
  \protect\hyperlink{REL-03-E01}{REL-03-E01}. \emph{Nguồn:} S3-p22--27.
  \emph{Trạng thái:} SV.
\end{itemize}

\hypertarget{F-POS-01}{%
\subsection{F-POS-01 · Quan hệ thứ tự (poset), thứ tự toàn
phần}\label{F-POS-01}}

\begin{itemize}
\tightlist
\item
  Thứ tự = \textbf{phản xạ + phản xứng + bắc cầu}; \((A,\preceq)\) là
  poset. \(a,b\) so sánh được nếu \(a\preceq b\) hoặc \(b\preceq a\);
  toàn phần nếu mọi cặp so sánh được.
\item
  Ví dụ chuẩn: \((\mathbb Z^+,\mid)\) poset, không toàn phần (5, 7);
  \((\mathbb Z,\mid)\) \textbf{không} poset (\(3\mid-3\), \(-3\mid3\),
  \(3\ne-3\)); \((P(S),\subseteq)\) poset; \((\mathbb Z^+,\le)\) toàn
  phần.
\item
  \emph{Nguồn:} S3-p29--35. \emph{Trạng thái:} SV.
\end{itemize}

\hypertarget{F-POS-02}{%
\subsection{F-POS-02 · Thứ tự từ điển}\label{F-POS-02}}

\begin{itemize}
\tightlist
\item
  Trên \(A\times B\) (hai tập sắp thứ tự toàn phần):
  \((a_1,b_1)\preceq(a_2,b_2)\Leftrightarrow a_1<a_2\ \vee\ (a_1=a_2\wedge b_1\le' b_2)\)
  (ERR-04).
\item
  Trên \(\Sigma^*\): \(s=a_1\dots a_m\preceq t=b_1\dots b_n\) nếu (i)
  \(s\) là \textbf{tiền tố} của \(t\), hoặc (ii) tại vị trí khác nhau
  đầu tiên \(k+1\): \(a_{k+1}<b_{k+1}\).
\item
  \emph{Ví dụ:} \(0110\prec10\); \(0110\prec01100\); discreet \(\prec\)
  discrete; \protect\hyperlink{POS-02-E01}{POS-02-E01}. \emph{Nguồn:}
  S3-p36--41. \emph{Trạng thái:} SV.
\end{itemize}

\hypertarget{F-POS-03}{%
\subsection{F-POS-03 · Hasse, phần tử đặc biệt, chặn, sup/inf,
topo}\label{F-POS-03}}

\begin{itemize}
\tightlist
\item
  \(b\) \textbf{trội trực tiếp} \(a\): \(a\prec b\) và không có \(c\)
  với \(a\prec c\prec b\). Hasse: chỉ vẽ cung trội trực tiếp, hướng từ
  dưới lên, bỏ khuyên và cung bắc cầu.
\item
  \textbf{Tối đại} \(m\): không có \(b\ne m\) với \(m\preceq b\) (không
  có cung đi lên). \textbf{Tối tiểu} tương tự. \textbf{Lớn nhất} \(M\):
  \(a\preceq M\ \forall a\); \textbf{nhỏ nhất} tương tự.
\item
  Poset hữu hạn (khác rỗng) luôn có tối đại, tối tiểu; có thể không duy
  nhất. Lớn nhất (nếu có) là duy nhất và là tối đại duy nhất
  \textbf{{[}CC{]}}.
\item
  Chặn trên của \(A\subseteq S\): \(x\in S\) với
  \(a\preceq x\ \forall a\in A\) (có thể \(x\notin A\)). \(\sup A\) =
  chặn trên nhỏ nhất (\(a\vee b\)), \(\inf A\) = chặn dưới lớn nhất
  (\(a\wedge b\)) --- \textbf{có thể không tồn tại}.
\item
  Sắp xếp topo: lặp lại ``lấy một phần tử tối tiểu, xóa nó'' ⇒ dãy
  \(a_1,a_2,\dots,a_m\) với \(a_i\prec a_j\Rightarrow i<j\).
\item
  \emph{Ví dụ:} \protect\hyperlink{POS-03-E01}{POS-03-E01},
  \protect\hyperlink{POS-04-E01}{POS-04-E01},
  \protect\hyperlink{POS-05-E01}{POS-05-E01}. \emph{Nguồn:} S3-p42--56.
  \emph{Trạng thái:} SV.
\end{itemize}

\hypertarget{chux1b0ux1a1ng-iv-ux111ux1ea1i-sux1ed1-bool-1}{%
\section{Chương IV --- Đại số
Bool}\label{chux1b0ux1a1ng-iv-ux111ux1ea1i-sux1ed1-bool-1}}

\hypertarget{F-BOO-01}{%
\subsection{F-BOO-01 · Tiên đề đại số Bool}\label{F-BOO-01}}

\((A,\wedge,\vee)\), \(A\ne\emptyset\): giao hoán, kết hợp, phân phối
(cả hai chiều), phần tử trung hòa (\(x\wedge1=x\), \(x\vee0=x\)), phần
tử bù (\(x\wedge\bar x=0\), \(x\vee\bar x=1\)). Ví dụ: \(B=\{0,1\}\);
tập dạng mệnh đề \(n\) biến. - \emph{Nguồn:} S4-p5--9. \emph{Trạng
thái:} SV.

\hypertarget{F-BOO-02}{%
\subsection{F-BOO-02 · Hàm Bool, phép toán, dạng nối rời chính
tắc}\label{F-BOO-02}}

\begin{itemize}
\tightlist
\item
  Hàm Bool \(n\) biến \(f:B^n\to B\); bảng chân trị \(2^n\) hàng.
  \(F_n\): tập các hàm Bool \(n\) biến (có \(2^{2^n}\) hàm
  \textbf{{[}CC{]}}).
\item
  \(f\vee g=f+g-fg=\max\{f,g\}\); \(\ f\wedge g=fg=\min\{f,g\}\);
  \(\ \bar f=1-f\).
\item
  \textbf{Từ đơn}: \(x_i\) hoặc \(\bar x_i\). \textbf{Đơn thức}: tích
  khác 0 của các từ đơn. \textbf{Từ tối tiểu}: tích đúng \(n\) từ đơn.
  \textbf{Công thức đa thức}: tổng các đơn thức. \textbf{Dạng nối rời
  chính tắc}: tổng các từ tối tiểu --- mỗi hàng \(f=1\) ↦ một từ tối
  tiểu (biến bằng 1 → \(x\), bằng 0 → \(\bar x\)).
\item
  \emph{Ví dụ:} \protect\hyperlink{BOO-01-E01}{BOO-01-E01}.
  \emph{Nguồn:} S4-p10--19. \emph{Trạng thái:} SV.
\end{itemize}

\hypertarget{F-BOO-03}{%
\subsection{F-BOO-03 · Biểu đồ Karnaugh (khung nhãn của
nguồn)}\label{F-BOO-03}}

\textbf{3 biến} \(x,y,z\) (cột theo \(x,y\); hàng theo \(z\)):

\begin{longtable}[]{@{}lcccc@{}}
\toprule\noalign{}
& \(x\bar y\) & \(xy\) & \(\bar xy\) & \(\bar x\bar y\) \\
\midrule\noalign{}
\endhead
\bottomrule\noalign{}
\endlastfoot
\(z\) & 101 & 111 & 011 & 001 \\
\(\bar z\) & 100 & 110 & 010 & 000 \\
\end{longtable}

\textbf{4 biến} \(x,y,z,t\) (cột theo \(x,y\); hàng theo \(z,t\)):

\begin{longtable}[]{@{}lcccc@{}}
\toprule\noalign{}
& \(x\bar y\) & \(xy\) & \(\bar xy\) & \(\bar x\bar y\) \\
\midrule\noalign{}
\endhead
\bottomrule\noalign{}
\endlastfoot
\(z\bar t\) & 1010 & 1110 & 0110 & 0010 \\
\(zt\) & 1011 & 1111 & 0111 & 0011 \\
\(\bar zt\) & 1001 & 1101 & 0101 & 0001 \\
\(\bar z\bar t\) & 1000 & 1100 & 0100 & 0000 \\
\end{longtable}

\begin{itemize}
\tightlist
\item
  Ô ghi chuỗi \(xyz(t)\). Hai ô kề (kể cả ô đầu--ô cuối cùng hàng/cột)
  khác nhau đúng một biến.
\item
  Định lý: \(\mathrm{kar}(fg)=\mathrm{kar}(f)\cap\mathrm{kar}(g)\);
  \(\ \mathrm{kar}(f\vee g)=\mathrm{kar}(f)\cup\mathrm{kar}(g)\);
  \(\mathrm{kar}(f)\) có đúng một ô ⇔ \(f\) là từ tối tiểu.
\item
  \emph{Lưu ý:} nhiều tài liệu khác dùng khung Gray ``00,01,11,10'' ---
  \textbf{khác khung của nguồn}; trong bài thi theo nguồn hãy dùng khung
  trên.
\item
  \emph{Nguồn:} S4-p22--26 (đọc từ ảnh). \emph{Trạng thái:} SV.
\end{itemize}

\hypertarget{F-BOO-04}{%
\subsection{F-BOO-04 · Tế bào, tế bào lớn, công thức đa thức tối
tiểu}\label{F-BOO-04}}

\begin{itemize}
\tightlist
\item
  \textbf{Tế bào}: hình chữ nhật (nghĩa rộng, kể cả quấn biên) gồm
  \(2^{n-k}\) ô; là kar của \textbf{một} đơn thức gồm \(k\) từ đơn:
  chiếu lên các cạnh, từ đơn nào chứa trọn hình chiếu thì có mặt.
\item
  \textbf{Tế bào lớn} của \(\mathrm{kar}(f)\): tế bào
  \(T\subseteq\mathrm{kar}(f)\) không nằm trong tế bào nào lớn hơn cũng
  \(\subseteq\mathrm{kar}(f)\).
\item
  \textbf{Đơn giản hơn}: \(F=m_1\vee\dots\vee m_k\) đơn giản hơn
  \(G=M_1\vee\dots\vee M_l\) nếu có đơn ánh \(h\) với mỗi \(m_i\) có số
  từ đơn \(\le\) số từ đơn của \(M_{h(i)}\). \textbf{Tối tiểu}: không có
  công thức thực sự đơn giản hơn.
\item
  \emph{Nguồn:} S4-p20--21, p27, p33. \emph{Trạng thái:} SV.
\end{itemize}

\hypertarget{F-BOO-05}{%
\subsection{F-BOO-05 · Cổng logic}\label{F-BOO-05}}

\begin{longtable}[]{@{}lll@{}}
\toprule\noalign{}
Cổng & Biểu thức & Ra = 1 khi \\
\midrule\noalign{}
\endhead
\bottomrule\noalign{}
\endlastfoot
NOT & \(\bar x\) & \(x=0\) \\
AND & \(xy\) (\(x\wedge y\)) & mọi ngõ vào = 1 \\
OR & \(x\vee y\) (\(x+y\)) & có ít nhất một ngõ vào = 1 \\
NAND & \(\overline{xy}\) & không phải mọi ngõ vào = 1 \\
NOR & \(\overline{x\vee y}\) & mọi ngõ vào = 0 \\
\end{longtable}

\begin{itemize}
\tightlist
\item
  \emph{Nguồn:} S4-p66--71. \emph{Trạng thái:} SV.
\end{itemize}

\hypertarget{chux1b0ux1a1ng-v-ux111ux1ed3-thux1ecb-1}{%
\section{Chương V --- Đồ
thị}\label{chux1b0ux1a1ng-v-ux111ux1ed3-thux1ecb-1}}

\hypertarget{F-GRA-01}{%
\subsection{F-GRA-01 · Bậc, định lý bắt tay}\label{F-GRA-01}}

\[
2m=\sum_{v\in V}\deg(v);\qquad\text{có hướng: } m=\sum_{v\in V}\deg^-(v)=\sum_{v\in V}\deg^+(v)
\]

\begin{itemize}
\tightlist
\item
  Số đỉnh bậc lẻ luôn \textbf{chẵn}. Khuyên đếm \textbf{2} vào bậc (vô
  hướng). Có hướng: \(\deg(v)=\deg^-(v)+\deg^+(v)\). Bậc 0: đỉnh cô lập;
  bậc 1: đỉnh treo.
\item
  Phân loại vô hướng: \textbf{đơn đồ thị} (không cạnh song song, không
  khuyên), \textbf{đa đồ thị} (cho phép song song, không khuyên),
  \textbf{giả đồ thị} (cho phép cả hai).
\item
  \emph{Ví dụ:} \protect\hyperlink{GRA-01-E01}{GRA-01-E01}.
  \emph{Nguồn:} S5-p3--25. \emph{Trạng thái:} SV, CC.
\end{itemize}

\hypertarget{F-GRA-02}{%
\subsection{F-GRA-02 · Ma trận kề}\label{F-GRA-02}}

\begin{itemize}
\tightlist
\item
  \(V=\{1,\dots,n\}\), \(A=(a_{ij})\), \(a_{ij}\) = số cạnh (cung) từ
  \(i\) đến \(j\). Vô hướng ⇒ \(A\) đối xứng.
\item
  \textbf{{[}CONFLICT{]} CON-01:} khuyên vô hướng --- ví dụ S5-p28 ghi
  \(a_{ii}=2\) (khớp quy ước bậc: tổng hàng \(i\) = \(\deg(i)\)); định
  nghĩa chữ ở S5-p26 gợi ý 1. Toolkit mặc định theo ví dụ
  (\(a_{ii}=2\)), ghi rõ quy ước trong bài thi.
\item
  Ma trận khoảng cách (trọng số): \(d_{ii}=0\), \(d_{ij}=w(v_iv_j)\) nếu
  có cạnh, \(\infty\) nếu không.
\item
  \emph{Ví dụ:} \protect\hyperlink{GRA-02-E01}{GRA-02-E01}.
  \emph{Nguồn:} S5-p26--28, p53. \emph{Trạng thái:} SV.
\end{itemize}

\hypertarget{F-GRA-03}{%
\subsection{F-GRA-03 · Đẳng cấu và bất biến}\label{F-GRA-03}}

\begin{itemize}
\tightlist
\item
  \(G\cong G'\) ⇔ có song ánh \(f:V\to V'\) với
  \(uv\in E\Leftrightarrow f(u)f(v)\in E'\).
\item
  Bất biến (điều kiện \textbf{cần}): cùng số đỉnh, số cạnh, số đỉnh có
  bậc cho trước; \(\deg v=\deg f(v)\). \textbf{{[}INFERRED{]}} thêm:
  cùng số thành phần liên thông, cùng số chu trình độ dài 3, đỉnh kề của
  ảnh có cùng bậc.
\item
  \emph{Ví dụ:} \protect\hyperlink{GRA-03-E01}{GRA-03-E01}.
  \emph{Nguồn:} S5-p29--35. \emph{Trạng thái:} SV.
\end{itemize}

\hypertarget{F-GRA-04}{%
\subsection{F-GRA-04 · Điều kiện Euler}\label{F-GRA-04}}

\begin{itemize}
\tightlist
\item
  \(G\) vô hướng \textbf{liên thông}: \(G\) có chu trình Euler ⇔
  \textbf{mọi đỉnh bậc chẵn}.
\item
  Liên thông, đúng \textbf{2 đỉnh bậc lẻ} ⇒ có đường đi Euler (bắt đầu ở
  một đỉnh lẻ, kết thúc ở đỉnh lẻ kia). Liên thông, \(2k\) đỉnh bậc lẻ
  (\(k\ge1\)) ⇒ vẽ được bằng \(k\) nét (CON-05).
\item
  \textbf{Fleury:} xuất phát từ một đỉnh (đỉnh lẻ nếu tìm đường đi); qua
  cạnh nào xóa cạnh đó (và đỉnh cô lập); \textbf{không đi qua cầu} trừ
  khi không còn cách nào khác.
\item
  \emph{Ví dụ:} \protect\hyperlink{GRA-05-E01}{GRA-05-E01}.
  \emph{Nguồn:} S5-p48--50. \emph{Trạng thái:} SV, CC.
\end{itemize}

\hypertarget{F-GRA-05}{%
\subsection{F-GRA-05 · Dijkstra}\label{F-GRA-05}}

\begin{itemize}
\tightlist
\item
  \emph{Điều kiện:} \(G\) đơn, liên thông, \textbf{trọng số dương}.
\item
  Bước 1: \(L(u_0)=0\), \(L(v)=\infty\), nhãn \((\infty,-)\). Bước 2:
  với \(v\) chưa cố định kề \(u_i\):
  \(L(v):=\min\{L(v),\ L(u_i)+w(u_iv)\}\); chọn \(v\) có \(L\) nhỏ nhất,
  cố định với nhãn \((L(v),u_i)\). Bước 3: lặp tới khi cố định đủ
  \(n-1\) đỉnh.
\item
  Kết quả: \(0=d(u_0,u_0)<d(u_0,u_1)\le\dots\le d(u_0,u_{n-1})\); truy
  vết theo đỉnh trước.
\item
  \emph{Ví dụ:} \protect\hyperlink{GRA-06-E01}{GRA-06-E01}.
  \emph{Nguồn:} S5-p56--59. \emph{Trạng thái:} SV, CC.
\end{itemize}

\hypertarget{F-GRA-06}{%
\subsection{F-GRA-06 · Ford--Bellman}\label{F-GRA-06}}

\[
L_0(u_0)=0,\quad L_0(v)=\infty\ (v\ne u_0);
\]

\[
L_k(u_0)=0,\qquad L_k(v)=\min\{L_{k-1}(u)+w(uv)\mid u\text{ là đỉnh trước của } v\}.
\]

\begin{itemize}
\tightlist
\item
  Dừng khi \(L_k(v)=L_{k-1}(v)\ \forall v\) (ổn định --- \(L\) là khoảng
  cách). Nếu đến \(k=n\) mà chưa ổn định ⇒ \textbf{có mạch âm}.
\item
  \emph{Lưu ý:} công thức của nguồn không lấy min với \(L_{k-1}(v)\);
  nhiều tài liệu khác có lấy --- với các bài trong nguồn hai cách cho
  cùng kết quả (CC).
\item
  \emph{Ví dụ:} \protect\hyperlink{GRA-07-E01}{GRA-07-E01}.
  \emph{Nguồn:} S5-p73--74. \emph{Trạng thái:} SV, CC.
\end{itemize}

\hypertarget{methods}{%
\chapter{Thư viện phương pháp giải (Method Library)}\label{methods}}

Mỗi phương pháp: \textbf{Mục đích · Dấu hiệu · Điều kiện · Quy trình ·
Trình bày trong bài thi · Không nên dùng khi · Ví dụ · Lỗi thường gặp}.
Tất cả phương pháp đều có căn cứ trong slide (nguồn ghi ở cuối mỗi mục).

\hypertarget{M-01}{%
\subsection{M-01 · Lập bảng chân trị}\label{M-01}}

\begin{itemize}
\tightlist
\item
  \textbf{Mục đích:} xác định chân trị dạng mệnh đề; kiểm tra hằng
  đúng/sai; kiểm tra tương đương (đối chiếu).
\item
  \textbf{Dấu hiệu:} ``lập bảng chân trị'', ``hằng đúng?'', ít biến
  (\(\le4\)).
\item
  \textbf{Quy trình:} (1) đếm \(n\) biến → \(2^n\) dòng, liệt kê theo
  thứ tự nhị phân 00\ldots0 → 11\ldots1; (2) mở cột cho từng
  \textbf{biểu thức con} từ trong ra ngoài; (3) cột cuối = \(E\); (4)
  toàn 1 → hằng đúng, toàn 0 → hằng sai.
\item
  \textbf{Trình bày:} bảng có tiêu đề cột là biểu thức con; kết luận một
  câu.
\item
  \textbf{Không nên dùng:} \(\ge5\) biến (32 dòng) --- chuyển sang
  \protect\hyperlink{M-02}{M-02}/\protect\hyperlink{M-05}{M-05}.
\item
  \textbf{Ví dụ:} \protect\hyperlink{LOG-02-E01}{LOG-02-E01}.
  \textbf{Lỗi:} sai cột \(\to\) (nhớ chỉ sai ở 1→0) ---
  \protect\hyperlink{T-01}{T-01}. \emph{Nguồn:} S1-p19--22.
\end{itemize}

\hypertarget{M-02}{%
\subsection{M-02 · Biến đổi tương đương bằng luật logic}\label{M-02}}

\begin{itemize}
\tightlist
\item
  \textbf{Mục đích:} chứng minh \(E\Leftrightarrow F\) hoặc rút gọn.
\item
  \textbf{Dấu hiệu:} ``chứng minh tương đương'', có \(\to\) cần khử.
\item
  \textbf{Quy trình:} (1) khử \(\to\), \(\leftrightarrow\) bằng luật kéo
  theo; (2) đẩy \(\neg\) vào trong (De Morgan, phủ định kép); (3) phân
  phối / hấp thụ / phần tử bù / thống trị để gom; (4) nếu cần, đưa về
  lại dạng \(\to\).
\item
  \textbf{Trình bày:} mỗi dòng ``\(\Leftrightarrow\ \dots\) (tên
  luật)''; biến đổi từ vế phức tạp sang vế đơn giản.
\item
  \textbf{Không nên dùng:} khi chỉ cần chỉ ra \textbf{không} tương đương
  --- dùng một dòng phản ví dụ của bảng chân trị.
\item
  \textbf{Ví dụ:} \protect\hyperlink{LOG-03-E01}{LOG-03-E01}.
  \textbf{Lỗi:} De Morgan quên đổi \(\wedge\leftrightarrow\vee\)
  \protect\hyperlink{T-02}{T-02}. \emph{Nguồn:} S1-p25--30.
\end{itemize}

\hypertarget{M-03}{%
\subsection{M-03 · Chuỗi quy tắc suy diễn (chứng minh trực
tiếp)}\label{M-03}}

\begin{itemize}
\tightlist
\item
  \textbf{Mục đích:} chứng minh tiền đề ⇒ kết luận.
\item
  \textbf{Quy trình:} (1) đánh số tiền đề; (2) tìm tiền đề ``đơn'' (một
  biến) để kích hoạt Modus Ponens/Tollens/tam đoạn luận rời; (3) mỗi
  dòng mới ghi căn cứ (số dòng + tên quy tắc); (4) kết luận dạng kéo
  theo có thể dùng phản đảo hoặc tam đoạn luận.
\item
  \textbf{Trình bày:} bảng 3 cột: bước -- mệnh đề -- lý do.
\item
  \textbf{Không nên dùng:} khi nghi suy luận sai →
  \protect\hyperlink{M-05}{M-05}.
\item
  \textbf{Ví dụ:} \protect\hyperlink{LOG-04-E01}{LOG-04-E01}.
  \textbf{Lỗi:} ``khẳng định hệ quả'' \protect\hyperlink{T-04}{T-04}.
  \emph{Nguồn:} S1-p31--51.
\end{itemize}

\hypertarget{M-04}{%
\subsection{M-04 · Phản chứng (quy tắc mâu thuẫn)}\label{M-04}}

\begin{itemize}
\tightlist
\item
  \textbf{Mục đích:} chứng minh \((p_1\wedge\dots\wedge p_n)\to h\) bằng
  cách thêm \(\neg h\) và suy ra \(0\).
\item
  \textbf{Dấu hiệu:} kết luận dạng \(a\to b\) hoặc khó suy trực tiếp.
\item
  \textbf{Quy trình:} (1) viết \(\neg h\) (với \(h=a\to b\):
  \(\neg h\Leftrightarrow a\wedge\neg b\)); (2) tách bằng quy tắc đơn
  giản; (3) dùng các tiền đề cho tới khi có \(r\) và \(\neg r\); (4)
  \(r\wedge\neg r\Leftrightarrow0\) ⇒ suy luận đúng.
\item
  \textbf{Ví dụ:} \protect\hyperlink{LOG-04-E02}{LOG-04-E02},
  \protect\hyperlink{LOG-04-E03}{LOG-04-E03}. \textbf{Lỗi:} phủ định
  \(a\to b\) thành \(\neg a\to\neg b\) \protect\hyperlink{T-03}{T-03}.
  \emph{Nguồn:} S1-p44--46, p52--54.
\end{itemize}

\hypertarget{M-05}{%
\subsection{M-05 · Kiểm tra suy luận bằng hệ chân trị / phản ví
dụ}\label{M-05}}

\begin{itemize}
\tightlist
\item
  \textbf{Mục đích:} quyết định suy luận đúng hay sai; nếu sai, đưa phản
  ví dụ.
\item
  \textbf{Quy trình:} (1) đặt kết luận \(=0\) (nếu kết luận \(a\to b\):
  \(a=1,b=0\)); (2) đặt mọi tiền đề \(=1\); (3) suy lần lượt các biến bị
  ép; (4) \textbf{có nghiệm} ⇒ suy luận \textbf{sai} (nghiệm là phản ví
  dụ, phải thử lại đủ tiền đề); \textbf{vô nghiệm} ⇒ suy luận
  \textbf{đúng}.
\item
  \textbf{Ví dụ:} \protect\hyperlink{LOG-05-E01}{LOG-05-E01},
  \protect\hyperlink{LOG-05-E02}{LOG-05-E02}. \textbf{Lỗi:} quên thử lại
  phản ví dụ. \emph{Nguồn:} S1-p48--49, p55--56.
\end{itemize}

\hypertarget{M-06}{%
\subsection{M-06 · Xét chân trị mệnh đề lượng từ}\label{M-06}}

\begin{itemize}
\tightlist
\item
  \textbf{Quy trình:} đọc \textbf{từ trái sang phải}. \(\forall x\)
  đúng: chứng minh với \(x\) tùy ý; sai: chỉ ra \textbf{một} \(x_0\).
  \(\exists x\) đúng: chỉ ra \textbf{một} \(x_0\); sai: chứng minh với
  mọi \(x\). \(\forall x\exists y\): \(y\) được phép \textbf{phụ thuộc}
  \(x\) (ví dụ \(y=-x/2\)). \(\exists x\forall y\): một \(x\) dùng cho
  \textbf{mọi} \(y\).
\item
  \textbf{Ví dụ:} \protect\hyperlink{QNT-01-E01}{QNT-01-E01}.
  \textbf{Lỗi:} đổi thứ tự lượng từ \protect\hyperlink{T-06}{T-06}.
  \emph{Nguồn:} S1-p61--68.
\end{itemize}

\hypertarget{M-07}{%
\subsection{M-07 · Phủ định mệnh đề lượng từ}\label{M-07}}

\begin{itemize}
\tightlist
\item
  \textbf{Quy trình:} (1) quét từ trái: \(\forall\to\exists\),
  \(\exists\to\forall\) (giữ miền); (2) phủ định phần vị từ:
  \(<\to\ge\), \(=\to\ne\), \(\wedge\to\vee\) (De Morgan),
  \(A\to B\ \to\ A\wedge\neg B\).
\item
  \textbf{Ví dụ:} \protect\hyperlink{QNT-02-E01}{QNT-02-E01}.
  \textbf{Lỗi:} phủ định cả miền ``\(\varepsilon>0\)''
  \protect\hyperlink{T-07}{T-07}. \emph{Nguồn:} S1-p69--71.
\end{itemize}

\hypertarget{M-08}{%
\subsection{M-08 · Chứng minh đẳng thức tập hợp}\label{M-08}}

\begin{itemize}
\tightlist
\item
  \textbf{Cách 1 (phần tử):} \(x\in VT\Leftrightarrow\dots\) (dịch sang
  logic, dùng luật logic) \(\Leftrightarrow x\in VP\).
\item
  \textbf{Cách 2 (luật tập hợp):} biến đổi bằng
  \protect\hyperlink{F-SET-01}{F-SET-01}, dùng
  \(A\setminus B=A\cap\bar B\) \textbf{{[}CC{]}}.
\item
  \textbf{Kiểm tra nhanh:} sơ đồ Ven / thử với tập cụ thể nhỏ.
\item
  \textbf{Ví dụ:} \protect\hyperlink{SET-01-E01}{SET-01-E01}.
  \emph{Nguồn:} S1-p76--79.
\end{itemize}

\hypertarget{M-09}{%
\subsection{\texorpdfstring{M-09 · Xét đơn/toàn/song ánh qua phương
trình
\(f(x)=y\)}{M-09 · Xét đơn/toàn/song ánh qua phương trình f(x)=y}}\label{M-09}}

\begin{itemize}
\tightlist
\item
  \textbf{Quy trình:} (1) kiểm tra \(f\) là ánh xạ (xác định với mọi
  \(x\in X\), giá trị thuộc \(Y\)); (2) cố định \(y\in Y\), giải
  \(f(x)=y\) \textbf{trong \(X\)}; (3) đếm nghiệm: \(\le1\) mọi \(y\) →
  đơn; \(\ge1\) mọi \(y\) → toàn; đúng 1 → song ánh, \(f^{-1}(y)\) =
  nghiệm đó; (4) để bác bỏ: một cặp \(x\ne x'\) cùng ảnh, hoặc một \(y\)
  vô nghiệm.
\item
  \textbf{Ví dụ:} \protect\hyperlink{FUN-02-E01}{FUN-02-E01},
  \protect\hyperlink{FUN-02-E02}{FUN-02-E02}. \textbf{Lỗi:} giải nghiệm
  ngoài \(X\) \protect\hyperlink{T-08}{T-08}. \emph{Nguồn:} S1-p88--94.
\end{itemize}

\hypertarget{M-10}{%
\subsection{M-10 · Quy nạp toán học}\label{M-10}}

\begin{itemize}
\tightlist
\item
  \textbf{Quy trình:} (1) Đặt \(P(n)\) = ``\ldots{}''. (2) \textbf{Cơ
  sở:} kiểm tra \(P(N_0)\). (3) \textbf{Quy nạp:} giả sử \(P(k)\) đúng
  (\(k\ge N_0\)), chứng minh \(P(k+1)\): viết vế của \(P(k+1)\),
  \textbf{tách ra phần bằng vế của \(P(k)\)}, thay bằng giả thiết quy
  nạp. (4) Kết luận theo nguyên lý quy nạp.
\item
  \textbf{Không nên dùng:} mệnh đề không phụ thuộc số tự nhiên.
\item
  \textbf{Ví dụ:} \protect\hyperlink{IND-01-E01}{IND-01-E01},
  \protect\hyperlink{IND-01-E03}{IND-01-E03}. \textbf{Lỗi:} sai \(N_0\),
  không dùng giả thiết \protect\hyperlink{T-09}{T-09}. \emph{Nguồn:}
  S1-p97--100.
\end{itemize}

\hypertarget{M-11}{%
\subsection{M-11 · Đếm bằng chia trường hợp + nguyên lý
nhân}\label{M-11}}

\begin{itemize}
\tightlist
\item
  \textbf{Quy trình:} (1) mô tả đối tượng bằng các ``vị trí''; (2) vị
  trí \textbf{bị ràng buộc nhiều nhất} chọn trước (chữ số cuối khi
  ``chẵn/chia hết cho 5'', chữ số đầu \(\ne0\)); (3) nếu ràng buộc phụ
  thuộc nhau (chữ số 0) → chia trường hợp rời nhau; (4) nhân trong mỗi
  trường hợp, cộng các trường hợp.
\item
  \textbf{Ví dụ:} \protect\hyperlink{CNT-01-E01}{CNT-01-E01}.
  \textbf{Lỗi:} quên chữ số đầu \(\ne0\) \protect\hyperlink{T-10}{T-10}.
  \emph{Nguồn:} S2-p3--5.
\end{itemize}

\hypertarget{M-12}{%
\subsection{M-12 · Nguyên lý chuồng bồ câu}\label{M-12}}

\begin{itemize}
\tightlist
\item
  \textbf{Quy trình:} (1) xác định ``chim'' (đối tượng) và ``chuồng''
  (nhóm) sao cho hai chim cùng chuồng ⇒ điều cần chứng minh; (2) đếm: số
  chim \textgreater{} số chuồng (hoặc dùng \(\lceil n/k\rceil\)); (3)
  kết luận.
\item
  \textbf{Ví dụ:} \protect\hyperlink{CNT-02-E01}{CNT-02-E01}.
  \emph{Nguồn:} S2-p6--7.
\end{itemize}

\hypertarget{M-13}{%
\subsection{M-13 · Bù trừ / đếm phần bù}\label{M-13}}

\begin{itemize}
\tightlist
\item
  \textbf{Quy trình:} (1) đặt tên tập \(A_i\) = các đối tượng có tính
  chất \(i\); (2) ``hoặc / ít nhất một'' → \(|\bigcup A_i|\) theo
  \protect\hyperlink{F-CNT-03}{F-CNT-03}; (3) ``không có tính chất nào''
  = Tổng \(-|\bigcup A_i|\); (4) ràng buộc ``\(\le b\)'' = tổng − số
  trường hợp ``\(\ge b+1\)''.
\item
  \textbf{Ví dụ:} \protect\hyperlink{CNT-03-E02}{CNT-03-E02},
  \protect\hyperlink{CNT-06-E02}{CNT-06-E02}. \emph{Nguồn:} S2-p8--10,
  p22--24.
\end{itemize}

\hypertarget{M-14}{%
\subsection{M-14 · Chọn đúng công thức tổ hợp (cây quyết
định)}\label{M-14}}

\begin{longtable}[]{@{}
  >{\raggedright\arraybackslash}p{(\columnwidth - 4\tabcolsep) * \real{0.4000}}
  >{\raggedright\arraybackslash}p{(\columnwidth - 4\tabcolsep) * \real{0.2400}}
  >{\raggedright\arraybackslash}p{(\columnwidth - 4\tabcolsep) * \real{0.3600}}@{}}
\toprule\noalign{}
\begin{minipage}[b]{\linewidth}\raggedright
Thứ tự quan trọng?
\end{minipage} & \begin{minipage}[b]{\linewidth}\raggedright
Được lặp?
\end{minipage} & \begin{minipage}[b]{\linewidth}\raggedright
Dùng
\end{minipage} \\
\midrule\noalign{}
\endhead
\bottomrule\noalign{}
\endlastfoot
Có, lấy hết \(n\) & Không & \(P_n=n!\) \\
Có, lấy \(k\) & Không & \(A_n^k\) \\
Có, lấy \(k\) & Có & \(n^k\) (nguyên lý nhân) \\
Không & Không & \(C_n^k\) \\
Không & Có & \(K_n^k=C_{n+k-1}^k\) \\
Sắp xếp \(n\) vật có nhóm giống nhau & --- &
\(\frac{n!}{n_1!\cdots n_k!}\) \\
\end{longtable}

\begin{itemize}
\tightlist
\item
  \textbf{Mẹo thử thứ tự:} đổi chỗ hai phần tử được chọn --- nếu ra kết
  quả \textbf{khác} thì có thứ tự. \textbf{Ví dụ:}
  \protect\hyperlink{CNT-04-E01}{CNT-04-E01}. \textbf{Lỗi:}
  \protect\hyperlink{T-11}{T-11}. \emph{Nguồn:} S2-p11--21.
\end{itemize}

\hypertarget{M-15}{%
\subsection{M-15 · Đổi biến cho nghiệm nguyên có ràng buộc}\label{M-15}}

\begin{itemize}
\tightlist
\item
  \textbf{Quy trình:} (1) đưa mọi ràng buộc về dạng \(x_i\ge a_i\) hoặc
  \(x_i\le b_i\) (nguyên: \(x>4\Leftrightarrow x\ge5\)); (2) đặt
  \(x_i'=x_i-a_i\) cho ràng buộc dưới; (3) ràng buộc trên: \(p=q-r\) với
  \(q\) = không xét cận trên, \(r\) = thay cận trên bằng
  \(x_i\ge b_i+1\); (4) mỗi số đếm dùng \(K_n^k=C_{n+k-1}^k\).
\item
  \textbf{Ví dụ:} \protect\hyperlink{CNT-06-E02}{CNT-06-E02},
  \protect\hyperlink{CNT-06-E03}{CNT-06-E03}. \textbf{Lỗi:}
  \protect\hyperlink{T-12}{T-12}. \emph{Nguồn:} S2-p22--24.
\end{itemize}

\hypertarget{M-16}{%
\subsection{M-16 · Lập hệ thức đệ qui}\label{M-16}}

\begin{itemize}
\tightlist
\item
  \textbf{Quy trình:} (1) đặt \(x_n\) = đại lượng cần đếm với tham số
  \(n\); (2) tính trực tiếp vài giá trị đầu (\(x_1, x_2\)); (3) với
  \(n\) lớn, chia theo \textbf{bước đầu tiên} (hoặc cuối cùng) thành các
  trường hợp \textbf{rời nhau}, mỗi trường hợp quy về bài toán cỡ nhỏ
  hơn; (4) cộng (nguyên lý cộng) ⇒ hệ thức + điều kiện đầu; (5) kiểm tra
  bằng liệt kê \(n\) nhỏ.
\item
  \textbf{Ví dụ:} \protect\hyperlink{REC-01-E01}{REC-01-E01},
  \protect\hyperlink{REC-01-E02}{REC-01-E02}. \emph{Nguồn:} S2-p34--39.
\end{itemize}

\hypertarget{M-17}{%
\subsection{M-17 · Giải hệ thức đệ qui tuyến tính (quy trình đầy
đủ)}\label{M-17}}

\begin{enumerate}
\def\labelenumi{\arabic{enumi}.}
\tightlist
\item
  \textbf{Chuẩn hóa:} viết vế trái theo \(x_{n}, x_{n-1}, x_{n-2}\)
  (hoặc \(x_{n+2}, x_{n+1}, x_n\)) --- giữ nguyên hệ số \(a_0,a_1,a_2\);
  xác định vế phải \(f_n\).
\item
  \textbf{PTĐT} \(a_0\lambda^2+a_1\lambda+a_2=0\); giải; viết nghiệm TQ
  thuần nhất \(\bar x_n\) (\protect\hyperlink{F-REC-02}{F-REC-02}).
\item
  Nếu \(f_n\ne0\): viết \(f_n=\beta^nP_r(n)\) (hoặc tổng các dạng). Với
  \textbf{mỗi} \(\beta\): xác định bội \(m\in\{0,1,2\}\) của \(\beta\)
  trong PTĐT → dạng \(n^m\beta^nQ_r(n)\)
  (\protect\hyperlink{F-REC-03}{F-REC-03}).
\item
  Thế vào hệ thức, cho \(n\) nhận \(r+1\) giá trị (hoặc đồng nhất hệ số)
  → hệ phương trình → hệ số.
\item
  Nghiệm TQ \(=\bar x_n+x_n^*\); thế điều kiện đầu → \(C_1,C_2\).
\item
  \textbf{Kiểm tra:} thế \(n=0,1\) (điều kiện đầu) và tính \(x_2\) hai
  cách (từ hệ thức và từ công thức).
\end{enumerate}

\begin{itemize}
\tightlist
\item
  \textbf{Không nên dùng:} PTĐT có nghiệm phức, hệ thức cấp \(\ge3\)
  (ngoài phạm vi, GAP-05) --- \protect\hyperlink{T-31}{T-31}.
\item
  \textbf{Ví dụ:} \protect\hyperlink{REC-02-E02}{REC-02-E02},
  \protect\hyperlink{REC-03-E02}{REC-03-E02},
  \protect\hyperlink{REC-04-E01}{REC-04-E01}. \textbf{Lỗi:}
  \protect\hyperlink{T-13}{T-13}, \protect\hyperlink{T-14}{T-14},
  \protect\hyperlink{T-15}{T-15}. \emph{Nguồn:} S2-p41--66.
\end{itemize}

\hypertarget{M-18}{%
\subsection{M-18 · Kiểm tra tính chất quan hệ}\label{M-18}}

\begin{itemize}
\tightlist
\item
  \textbf{Quan hệ hữu hạn (liệt kê/ma trận):} phản xạ --- xem đường
  chéo; đối xứng --- mỗi \((a,b)\) có \((b,a)\); phản xứng --- không có
  cặp \((a,b),(b,a)\) với \(a\ne b\); bắc cầu --- với \textbf{mỗi}
  \((a,b)\) và mỗi \((b,c)\) kiểm tra \((a,c)\).
\item
  \textbf{Quan hệ cho bằng quy tắc trên tập vô hạn:} ``có'' → chứng minh
  tổng quát bằng định nghĩa; ``không'' → \textbf{một} phản ví dụ số cụ
  thể.
\item
  \textbf{Trình bày:} bảng 4 tính chất × (Có/Không + lý do/phản ví dụ).
  Tương đương = PX + ĐX + BC; Thứ tự = PX + PXứng + BC.
\item
  \textbf{Ví dụ:} \protect\hyperlink{REL-02-E01}{REL-02-E01},
  \protect\hyperlink{REL-02-E02}{REL-02-E02},
  \protect\hyperlink{POS-01-E01}{POS-01-E01}. \textbf{Lỗi:}
  \protect\hyperlink{T-16}{T-16}, \protect\hyperlink{T-17}{T-17}.
  \emph{Nguồn:} S3-p7--18.
\end{itemize}

\hypertarget{M-19}{%
\subsection{M-19 · Tìm lớp tương đương}\label{M-19}}

\begin{itemize}
\tightlist
\item
  \textbf{Quy trình:} (1) chứng minh \(R\) tương đương; (2) lấy phần tử
  \(a\) chưa thuộc lớp nào, tính \([a]=\{b\mid bRa\}\); (3) lặp tới khi
  phủ hết \(A\); (4) kiểm tra: các lớp rời nhau, hợp bằng \(A\) (phân
  hoạch).
\item
  \textbf{Ví dụ:} \protect\hyperlink{REL-03-E01}{REL-03-E01},
  \protect\hyperlink{REL-03-E02}{REL-03-E02}. \emph{Nguồn:} S3-p23--27.
\end{itemize}

\hypertarget{M-20}{%
\subsection{M-20 · Vẽ Hasse và đọc phần tử đặc biệt, chặn}\label{M-20}}

\begin{itemize}
\tightlist
\item
  \textbf{Quy trình:} (1) liệt kê các cặp \(a\prec b\); (2) loại cặp bắc
  cầu (có \(c\) ở giữa) → còn cặp trội trực tiếp; (3) xếp tầng: phần tử
  tối tiểu ở dưới; vẽ đoạn từ \(a\) lên \(b\); (4) đọc: tối đại = đỉnh
  không có cạnh đi lên; tối tiểu = không có cạnh đi xuống; lớn nhất =
  phần tử mà \textbf{mọi} phần tử đi lên được tới; (5) chặn trên của
  \(A\) = các phần tử đi lên được từ \textbf{mọi} phần tử của \(A\); sup
  = chặn trên nhỏ nhất (phải so sánh được và nhỏ hơn mọi chặn trên
  khác).
\item
  \textbf{Ví dụ:} \protect\hyperlink{POS-03-E01}{POS-03-E01},
  \protect\hyperlink{POS-04-E01}{POS-04-E01}. \textbf{Lỗi:}
  \protect\hyperlink{T-18}{T-18}, \protect\hyperlink{T-19}{T-19}.
  \emph{Nguồn:} S3-p42--53.
\end{itemize}

\hypertarget{M-21}{%
\subsection{M-21 · Sắp xếp topo}\label{M-21}}

\begin{itemize}
\tightlist
\item
  \textbf{Quy trình:} lặp: chọn một phần tử \textbf{tối tiểu} của tập
  còn lại, ghi ra, xóa khỏi poset. Kết quả không duy nhất. \textbf{Kiểm
  tra:} mọi cặp \(a\prec b\) có \(a\) đứng trước \(b\).
\item
  \textbf{Ví dụ:} \protect\hyperlink{POS-05-E01}{POS-05-E01}.
  \emph{Nguồn:} S3-p54--56.
\end{itemize}

\hypertarget{M-22}{%
\subsection{M-22 · Thuật toán Karnaugh 5 bước}\label{M-22}}

\begin{enumerate}
\def\labelenumi{\arabic{enumi}.}
\tightlist
\item
  Vẽ \(\mathrm{kar}(f)\): mỗi đơn thức của \(f\) → tô tế bào tương ứng
  (hoặc tô từng hàng \(f=1\) của bảng chân trị).
\item
  Liệt kê \textbf{tất cả} tế bào lớn (thử lần lượt: 8 ô → 4 ô → 2 ô → 1
  ô, kể cả quấn biên và 4 góc).
\item
  Tế bào lớn \textbf{bắt buộc}: chứa một ô chỉ nằm trong nó.
\item
  Nếu chưa phủ hết: với mỗi ô chưa phủ, xét các tế bào lớn chứa nó →
  liệt kê mọi phủ; loại phủ không tối tiểu.
\item
  Viết công thức đa thức tương ứng; loại công thức có công thức khác
  thực sự đơn giản hơn.
\end{enumerate}

\begin{itemize}
\tightlist
\item
  \textbf{Trình bày:} vẽ bản đồ, đánh số ô; liệt kê tế bào lớn kèm đơn
  thức; nêu ô quyết định khi chọn bắt buộc.
\item
  \textbf{Ví dụ:} \protect\hyperlink{BOO-03-E01}{BOO-03-E01},
  \protect\hyperlink{BOO-03-E02}{BOO-03-E02}. \textbf{Lỗi:}
  \protect\hyperlink{T-20}{T-20}, \protect\hyperlink{T-21}{T-21}.
  \emph{Nguồn:} S4-p38--65.
\end{itemize}

\hypertarget{M-23}{%
\subsection{M-23 · Từ bảng chân trị / mô tả đến công thức và
mạch}\label{M-23}}

\begin{itemize}
\tightlist
\item
  \textbf{Quy trình:} (1) mã hóa biến (1 = bật/tán thành) và \(f\) (1 =
  đèn sáng/thông qua); (2) lập bảng chân trị; (3) dạng nối rời chính
  tắc: mỗi hàng \(f=1\) → từ tối tiểu; (4) (nên) rút gọn bằng Karnaugh;
  (5) vẽ mạch: mỗi đơn thức → cổng AND (biến có gạch → qua NOT), các đơn
  thức → cổng OR.
\item
  \textbf{Đọc mạch:} đi từ ngõ vào tới ngõ ra, ghi biểu thức ở đầu ra
  của từng cổng.
\item
  \textbf{Ví dụ:} \protect\hyperlink{BOO-04-E01}{BOO-04-E01},
  \protect\hyperlink{BOO-04-E02}{BOO-04-E02}. \emph{Nguồn:} S4-p12--19,
  p66--78.
\end{itemize}

\hypertarget{M-24}{%
\subsection{M-24 · Kiểm tra đẳng cấu}\label{M-24}}

\begin{itemize}
\tightlist
\item
  \textbf{Quy trình:} (1) so bất biến: số đỉnh, số cạnh, dãy bậc (sắp
  giảm); khác ⇒ kết luận \textbf{không đẳng cấu}; (2) nếu giống: ghép
  các đỉnh \textbf{bậc duy nhất} trước, rồi theo quan hệ kề với các đỉnh
  đã ghép; (3) kiểm tra \textbf{mọi cạnh} \(uv\) ↦ \(f(u)f(v)\) là cạnh
  (và số cạnh bằng nhau ⇒ đủ); (4) nếu mọi cách ghép đều thất bại ⇒
  không đẳng cấu (nêu bất biến phân biệt, ví dụ hai đỉnh bậc 3 kề nhau ở
  \(G\) nhưng không kề ở \(G'\)).
\item
  \textbf{Ví dụ:} \protect\hyperlink{GRA-03-E01}{GRA-03-E01}.
  \textbf{Lỗi:} \protect\hyperlink{T-24}{T-24}. \emph{Nguồn:}
  S5-p29--35.
\end{itemize}

\hypertarget{M-25}{%
\subsection{M-25 · Euler: kiểm tra điều kiện + Fleury}\label{M-25}}

\begin{itemize}
\tightlist
\item
  \textbf{Quy trình:} (1) kiểm tra liên thông; (2) tính bậc mọi đỉnh;
  (3) 0 đỉnh lẻ → chu trình Euler, 2 đỉnh lẻ → đường đi Euler (xuất phát
  ở đỉnh lẻ), \textgreater2 → không có; (4) Fleury: đi một cạnh, xóa
  cạnh; ưu tiên cạnh \textbf{không phải cầu} của đồ thị còn lại; (5)
  kiểm tra: dãy dùng đủ \(m\) cạnh, mỗi cạnh một lần.
\item
  \textbf{Ví dụ:} \protect\hyperlink{GRA-05-E01}{GRA-05-E01},
  \protect\hyperlink{GRA-05-E02}{GRA-05-E02}. \textbf{Lỗi:}
  \protect\hyperlink{T-25}{T-25}. \emph{Nguồn:} S5-p48--51.
\end{itemize}

\hypertarget{M-26}{%
\subsection{M-26 · Dijkstra dạng bảng}\label{M-26}}

\begin{itemize}
\tightlist
\item
  \textbf{Quy trình:} (1) kiểm tra trọng số dương; (2) bảng: cột = đỉnh,
  dòng = bước; dòng 0: nguồn \(0^*\), còn lại \((\infty,-)\); (3) mỗi
  bước: cập nhật các đỉnh \textbf{chưa cố định kề với đỉnh vừa cố định}:
  \(L(v)=\min\{L(v),L(u_i)+w\}\), nhãn \((L,u_i)\) nếu giảm; (4) đánh
  dấu * đỉnh có \(L\) nhỏ nhất; đỉnh đã cố định ghi ``--''; (5) truy vết
  ngược theo đỉnh trước để có đường đi.
\item
  \textbf{Trình bày:} bảng + ``Đường đi ngắn nhất \(u_0\to v\): \ldots{}
  , độ dài \(d\)''.
\item
  \textbf{Không nên dùng:} có trọng số âm →
  \protect\hyperlink{M-27}{M-27}.
\item
  \textbf{Ví dụ:} \protect\hyperlink{GRA-06-E01}{GRA-06-E01}.
  \textbf{Lỗi:} \protect\hyperlink{T-26}{T-26},
  \protect\hyperlink{T-27}{T-27}. \emph{Nguồn:} S5-p56--72.
\end{itemize}

\hypertarget{M-27}{%
\subsection{M-27 · Ford--Bellman dạng bảng}\label{M-27}}

\begin{itemize}
\tightlist
\item
  \textbf{Quy trình:} (1) dòng \(k=0\): nguồn 0, còn lại \((\infty,-)\);
  (2) dòng \(k\): với mỗi \(v\ne u_0\), \(L_k(v)=\min\) trên \textbf{mọi
  cung vào} \(uv\) của \(L_{k-1}(u)+w(uv)\) (dùng giá trị \textbf{dòng
  trước}); ghi \((L_k(v),u)\); (3) dòng \(k\) trùng dòng \(k-1\) ⇒ dừng,
  đọc khoảng cách; (4) tới \(k=n\) chưa ổn định ⇒ có mạch âm --- chỉ ra
  một mạch âm bằng truy vết nhãn.
\item
  \textbf{Ví dụ:} \protect\hyperlink{GRA-07-E01}{GRA-07-E01},
  \protect\hyperlink{GRA-07-E02}{GRA-07-E02}. \textbf{Lỗi:}
  \protect\hyperlink{T-28}{T-28}. \emph{Nguồn:} S5-p73--86.
\end{itemize}

\hypertarget{catalog}{%
\chapter{Danh mục ví dụ (Example Catalog)}\label{catalog}}

Tổng số: \textbf{83 ví dụ}; trong đó 17 ví dụ \textbf{{[}AI-GENERATED
PRACTICE{]}}, còn lại đề lấy từ slide. Cột ``Trang'' là số trang trong
bản PDF.

\begin{longtable}[]{@{}
  >{\raggedright\arraybackslash}p{(\columnwidth - 6\tabcolsep) * \real{0.1818}}
  >{\raggedright\arraybackslash}p{(\columnwidth - 6\tabcolsep) * \real{0.4545}}
  >{\raggedright\arraybackslash}p{(\columnwidth - 6\tabcolsep) * \real{0.2955}}
  >{\raggedright\arraybackslash}p{(\columnwidth - 6\tabcolsep) * \real{0.0682}}@{}}
\toprule\noalign{}
\begin{minipage}[b]{\linewidth}\raggedright
Mã
\end{minipage} & \begin{minipage}[b]{\linewidth}\raggedright
Tên
\end{minipage} & \begin{minipage}[b]{\linewidth}\raggedright
Độ khó · Biến thể · Nguồn đề
\end{minipage} & \begin{minipage}[b]{\linewidth}\raggedright
Trang
\end{minipage} \\
\midrule\noalign{}
\endhead
\bottomrule\noalign{}
\endlastfoot
\protect\hyperlink{LOG-01-E01}{LOG-01-E01} & Nhận biết mệnh đề; sơ cấp
hay phức hợp & D1 · V0 · Đề {[}SV{]} S1-p8, S1-p9 &
\pageref{LOG-01-E01} \\
\protect\hyperlink{LOG-02-E01}{LOG-02-E01} & Lập bảng chân trị & D1 · V0
· Đề {[}SV{]} S1-p22 (Bài tập) & \pageref{LOG-02-E01} \\
\protect\hyperlink{LOG-02-E02}{LOG-02-E02} & Hằng đúng, hằng sai --- bẫy
& D1 · V4 · Đề {[}SV{]} S1-p21 & \pageref{LOG-02-E02} \\
\protect\hyperlink{LOG-03-E01}{LOG-03-E01} & Chứng minh tương đương bằng
luật logic & D2 · V0 · Đề {[}SV{]} S1-p29 & \pageref{LOG-03-E01} \\
\protect\hyperlink{LOG-03-E02}{LOG-03-E02} & Rút gọn dạng mệnh đề; hệ
quả logic & D1 · V1 · {[}AI-GENERATED PRACTICE{]} &
\pageref{LOG-03-E02} \\
\protect\hyperlink{LOG-04-E01}{LOG-04-E01} & Chứng minh suy luận bằng
chuỗi quy tắc & D2 · V0 · Đề {[}SV{]} S1-p50 & \pageref{LOG-04-E01} \\
\protect\hyperlink{LOG-04-E02}{LOG-04-E02} & Chứng minh suy luận bằng
phản chứng & D2 · V3 · Đề {[}SV{]} S1-p52 & \pageref{LOG-04-E02} \\
\protect\hyperlink{LOG-04-E03}{LOG-04-E03} & Phản chứng hoặc tam đoạn
luận nối tiếp & D2 · V1 · Đề {[}SV{]} S1-p46 & \pageref{LOG-04-E03} \\
\protect\hyperlink{LOG-05-E01}{LOG-05-E01} & Suy luận lời văn --- tìm
phản ví dụ & D3 · V2 · Đề {[}SV{]} S1-p49 & \pageref{LOG-05-E01} \\
\protect\hyperlink{LOG-05-E02}{LOG-05-E02} & Suy luận sai --- và kết
luận đúng là gì & D2 · V4 · Đề {[}SV{]} S1-p55 & \pageref{LOG-05-E02} \\
\protect\hyperlink{QNT-01-E01}{QNT-01-E01} & Chân trị mệnh đề lượng từ
một và hai biến & D1--D2 · V0 · Đề {[}SV{]} S1-p61, p63, p65--67 &
\pageref{QNT-01-E01} \\
\protect\hyperlink{QNT-01-E02}{QNT-01-E02} & Thứ tự lượng từ và miền xác
định --- bẫy & D2 · V4 · {[}AI-GENERATED PRACTICE{]} &
\pageref{QNT-01-E02} \\
\protect\hyperlink{QNT-02-E01}{QNT-02-E01} & Phủ định mệnh đề lượng từ &
D2 · V0 · Đề {[}SV{]} S1-p71 (+ một câu AI) & \pageref{QNT-02-E01} \\
\protect\hyperlink{SET-01-E01}{SET-01-E01} & Chứng minh đẳng thức tập
hợp & D2 · V0 · {[}AI-GENERATED PRACTICE{]} & \pageref{SET-01-E01} \\
\protect\hyperlink{SET-02-E01}{SET-02-E01} & Tập các tập con, tích
Đề-các & D1 · V0 · Đề {[}SV{]} S1-p80--81 (câu hỏi mở) &
\pageref{SET-02-E01} \\
\protect\hyperlink{FUN-01-E01}{FUN-01-E01} & Ảnh và ảnh ngược & D1 · V0
· Đề {[}SV{]} S1-p87 & \pageref{FUN-01-E01} \\
\protect\hyperlink{FUN-02-E01}{FUN-02-E01} & Đơn, toàn, song ánh; ánh xạ
ngược & D2 · V0 · Đề {[}SV{]} S1-p88, p90, p92, p94 &
\pageref{FUN-02-E01} \\
\protect\hyperlink{FUN-02-E02}{FUN-02-E02} & Hàm phân thức: đơn ánh
nhưng không toàn ánh --- bẫy tập đích & D3 · V4 · {[}AI-GENERATED
PRACTICE{]} & \pageref{FUN-02-E02} \\
\protect\hyperlink{FUN-03-E01}{FUN-03-E01} & Tích ánh xạ, kể cả hàm từng
khúc & D2 · V0 · Đề {[}SV{]} S1-p96 & \pageref{FUN-03-E01} \\
\protect\hyperlink{IND-01-E01}{IND-01-E01} & Quy nạp: tổng số lẻ & D1 ·
V0 · Đề {[}SV{]} S1-p97--99 & \pageref{IND-01-E01} \\
\protect\hyperlink{IND-01-E02}{IND-01-E02} & Quy nạp: tổng
\(1+2+\dots+n\) & D1 · V1 · Đề {[}SV{]} S1-p100 &
\pageref{IND-01-E02} \\
\protect\hyperlink{IND-01-E03}{IND-01-E03} & Quy nạp bất đẳng thức, bước
cơ sở \(n_0=4\) & D2 · V4 · Đề {[}SV{]} S1-p100 &
\pageref{IND-01-E03} \\
\protect\hyperlink{IND-01-E04}{IND-01-E04} & Chia hết: quy nạp và chứng
minh theo trường hợp & D2 · V3 · Đề {[}SV{]} S1-p47 &
\pageref{IND-01-E04} \\
\protect\hyperlink{CNT-01-E01}{CNT-01-E01} & Đếm số chẵn có 3 chữ số
khác nhau (chia trường hợp chữ số 0) & D2 · V0 · Đề {[}SV{]} S2-p5 &
\pageref{CNT-01-E01} \\
\protect\hyperlink{CNT-01-E02}{CNT-01-E02} & Số 4 chữ số khác nhau chia
hết cho 5 & D2 · V1 · {[}AI-GENERATED PRACTICE{]} &
\pageref{CNT-01-E02} \\
\protect\hyperlink{CNT-02-E01}{CNT-02-E01} & Nguyên lý chuồng bồ câu &
D2 · V0 · Đề {[}SV{]} S2-p6--7 & \pageref{CNT-02-E01} \\
\protect\hyperlink{CNT-03-E01}{CNT-03-E01} & Bù trừ hai tập & D1 · V0 ·
Đề {[}SV{]} S2-p10 & \pageref{CNT-03-E01} \\
\protect\hyperlink{CNT-03-E02}{CNT-03-E02} & Bù trừ ba tập: chia hết cho
2, 3 hoặc 5 & D2 · V3 · {[}AI-GENERATED PRACTICE{]} &
\pageref{CNT-03-E02} \\
\protect\hyperlink{CNT-04-E01}{CNT-04-E01} & Hoán vị, chỉnh hợp hay tổ
hợp? & D1 · V4 · Đề {[}SV{]} S2-p12, p14, p16 & \pageref{CNT-04-E01} \\
\protect\hyperlink{CNT-05-E01}{CNT-05-E01} & Hoán vị lặp --- sắp xếp chữ
cái & D1 · V0 · Đề {[}SV{]} S2-p18 & \pageref{CNT-05-E01} \\
\protect\hyperlink{CNT-06-E01}{CNT-06-E01} & Tổ hợp lặp cơ bản & D1 · V0
· Đề {[}SV{]} S2-p20 & \pageref{CNT-06-E01} \\
\protect\hyperlink{CNT-06-E02}{CNT-06-E02} & Nghiệm nguyên không âm với
ràng buộc trên và dưới & D3 · V0 · Đề {[}SV{]} S2-p22 &
\pageref{CNT-06-E02} \\
\protect\hyperlink{CNT-06-E03}{CNT-06-E03} & Chia vật giống nhau có cận
trên và cận dưới & D2 · V2 · Đề {[}SV{]} S2-p25 (slide không có đáp án)
& \pageref{CNT-06-E03} \\
\protect\hyperlink{REC-01-E01}{REC-01-E01} & Lập hệ thức: cầu thang bước
1 hoặc 2 bậc & D2 · V0 · Đề {[}SV{]} S2-p34 & \pageref{REC-01-E01} \\
\protect\hyperlink{REC-01-E02}{REC-01-E02} & Lập hệ thức: tháp Hà Nội &
D2 · V0 · Đề {[}SV{]} S2-p27 & \pageref{REC-01-E02} \\
\protect\hyperlink{REC-01-E03}{REC-01-E03} & Lập hệ thức: chuỗi bit
không có hai số 0 liên tiếp & D2 · V1 · {[}AI-GENERATED PRACTICE{]} &
\pageref{REC-01-E03} \\
\protect\hyperlink{REC-02-E01}{REC-02-E01} & Thuần nhất cấp 1 & D1 · V0
· Đề {[}SV{]} S2-p42 & \pageref{REC-02-E01} \\
\protect\hyperlink{REC-02-E02}{REC-02-E02} & Thuần nhất cấp 2: nghiệm
phân biệt và nghiệm kép & D2 · V0 · Đề {[}SV{]} S2-p44 &
\pageref{REC-02-E02} \\
\protect\hyperlink{REC-02-E03}{REC-02-E03} & Giải hệ thức cầu thang
(nghiệm vô tỉ) & D2 · V3 · Đề: nối tiếp
\protect\hyperlink{REC-01-E01}{REC-01-E01} {[}SV{]}; lời giải
{[}AI-GENERATED PRACTICE{]}, {[}CC{]} & \pageref{REC-02-E03} \\
\protect\hyperlink{REC-03-E01}{REC-03-E01} & Vế phải đa thức,
\(\beta=1\) là nghiệm đơn & D2 · V0 · Đề {[}SV{]} S2-p53 a) &
\pageref{REC-03-E01} \\
\protect\hyperlink{REC-03-E02}{REC-03-E02} & \(\beta\) là nghiệm kép, có
điều kiện đầu & D3 · V4 · Đề {[}SV{]} S2-p53 b), S2-p57 &
\pageref{REC-03-E02} \\
\protect\hyperlink{REC-03-E03}{REC-03-E03} & \(\beta\) không là nghiệm,
đa thức bậc 2 & D3 · V1 · Đề {[}SV{]} S2-p53 c), S2-p60 &
\pageref{REC-03-E03} \\
\protect\hyperlink{REC-03-E04}{REC-03-E04} & Bài tập S2-p67 (1):
\(\beta=1\) nghiệm đơn & D2 · V1 · Đề {[}SV{]} S2-p67 (không có lời
giải) & \pageref{REC-03-E04} \\
\protect\hyperlink{REC-03-E05}{REC-03-E05} & Bài tập S2-p67 (2): nghiệm
âm, \(\beta=3\) không là nghiệm & D2 · V1 · Đề {[}SV{]} S2-p67 &
\pageref{REC-03-E05} \\
\protect\hyperlink{REC-03-E06}{REC-03-E06} & Bài tập S2-p67 (3): vế phải
hằng, \(\beta=1\) nghiệm kép & D2 · V4 · Đề {[}SV{]} S2-p67 &
\pageref{REC-03-E06} \\
\protect\hyperlink{REC-03-E07}{REC-03-E07} & Bài tập S2-p67 (4): đa thức
bậc 2, \(\beta=1\) không là nghiệm & D2 · V1 · Đề {[}SV{]} S2-p67 &
\pageref{REC-03-E07} \\
\protect\hyperlink{REC-03-E08}{REC-03-E08} & Bài tập S2-p67 (5): viết
lại vế phải \(2\cdot5^{n+1}\) & D2 · V2 · Đề {[}SV{]} S2-p67 &
\pageref{REC-03-E08} \\
\protect\hyperlink{REC-03-E09}{REC-03-E09} & Bài tập S2-p67 (6):
\(\beta=8\) nghiệm kép & D2 · V0 · Đề {[}SV{]} S2-p67 &
\pageref{REC-03-E09} \\
\protect\hyperlink{REC-04-E01}{REC-04-E01} & Vế phải tổng ba dạng (chồng
chất) & D3 · V3 · Đề {[}SV{]} S2-p53 d), S2-p64 &
\pageref{REC-04-E01} \\
\protect\hyperlink{REL-01-E01}{REL-01-E01} & Quan hệ ↔ ma trận biểu diễn
& D1 · V0 · Đề {[}SV{]} S3-p13--15 & \pageref{REL-01-E01} \\
\protect\hyperlink{REL-02-E01}{REL-02-E01} & Bảng bốn tính chất của quan
hệ hữu hạn & D1 · V0 · Đề {[}SV{]} S3-p7, p9, p11 &
\pageref{REL-02-E01} \\
\protect\hyperlink{REL-02-E02}{REL-02-E02} & Quan hệ trên \(\mathbb Z\):
\(aRb\Leftrightarrow ab\ge0\) --- bẫy bắc cầu & D2 · V4 ·
{[}AI-GENERATED PRACTICE{]} & \pageref{REL-02-E02} \\
\protect\hyperlink{REL-03-E01}{REL-03-E01} & Đồng dư modulo \(m\) và lớp
tương đương & D1 · V0 · Đề {[}SV{]} S3-p22, p24, p27 &
\pageref{REL-03-E01} \\
\protect\hyperlink{REL-03-E02}{REL-03-E02} & Từ phân hoạch sang quan hệ
tương đương & D2 · V2 · {[}AI-GENERATED PRACTICE{]} dựa trên S3-p26 &
\pageref{REL-03-E02} \\
\protect\hyperlink{POS-01-E01}{POS-01-E01} & Có phải poset? Thứ tự toàn
phần? & D2 · V0 · Đề {[}SV{]} S3-p31--35 & \pageref{POS-01-E01} \\
\protect\hyperlink{POS-02-E01}{POS-02-E01} & Thứ tự từ điển & D1 · V0 ·
Đề {[}SV{]} S3-p36, p40--41 & \pageref{POS-02-E01} \\
\protect\hyperlink{POS-03-E01}{POS-03-E01} & Hasse, tối đại, tối tiểu,
lớn nhất, nhỏ nhất & D2 · V0 · Đề {[}SV{]} S3-p48 &
\pageref{POS-03-E01} \\
\protect\hyperlink{POS-04-E01}{POS-04-E01} & Chặn trên, chặn dưới, sup,
inf & D2 · V4 · {[}AI-GENERATED PRACTICE{]} (định nghĩa S3-p50--53) &
\pageref{POS-04-E01} \\
\protect\hyperlink{POS-05-E01}{POS-05-E01} & Sắp xếp topo & D1 · V0 ·
{[}AI-GENERATED PRACTICE{]} (thuật toán S3-p54--56) &
\pageref{POS-05-E01} \\
\protect\hyperlink{BOO-01-E01}{BOO-01-E01} & Hàm bỏ phiếu đa số: bảng
chân trị → dạng nối rời chính tắc → tối tiểu & D1 · V0 · Đề {[}SV{]}
S4-p12--13 & \pageref{BOO-01-E01} \\
\protect\hyperlink{BOO-01-E02}{BOO-01-E02} & Dạng nối rời chính tắc từ
công thức & D1 · V1 · {[}AI-GENERATED PRACTICE{]} &
\pageref{BOO-01-E02} \\
\protect\hyperlink{BOO-02-E01}{BOO-02-E01} & Tế bào ↔ đơn thức & D1 · V0
· Đề {[}SV{]} S4-p28--32, p34--37 & \pageref{BOO-02-E01} \\
\protect\hyperlink{BOO-03-E01}{BOO-03-E01} & Karnaugh: một phủ tối tiểu
duy nhất & D2 · V0 · Đề {[}SV{]} S4-p41 & \pageref{BOO-03-E01} \\
\protect\hyperlink{BOO-03-E02}{BOO-03-E02} & Karnaugh: hai công thức tối
tiểu & D3 · V0 · Đề {[}SV{]} S4-p53 & \pageref{BOO-03-E02} \\
\protect\hyperlink{BOO-03-E03}{BOO-03-E03} & Karnaugh: hai phủ nhưng chỉ
một tối tiểu --- bẫy & D3 · V4 · Đề {[}SV{]} S4-p62 &
\pageref{BOO-03-E03} \\
\protect\hyperlink{BOO-03-E04}{BOO-03-E04} & Karnaugh 3 biến & D1 · V1 ·
{[}AI-GENERATED PRACTICE{]} & \pageref{BOO-03-E04} \\
\protect\hyperlink{BOO-03-E05}{BOO-03-E05} & Karnaugh 4 biến có tế bào
bốn góc và hai lựa chọn & D2 · V3 · {[}AI-GENERATED PRACTICE{]} &
\pageref{BOO-03-E05} \\
\protect\hyperlink{BOO-04-E01}{BOO-04-E01} & Thiết kế mạch 2 công tắc &
D1 · V0 · Đề {[}SV{]} S4-p75--76 & \pageref{BOO-04-E01} \\
\protect\hyperlink{BOO-04-E02}{BOO-04-E02} & Thiết kế mạch 3 công tắc
(đèn cầu thang) & D2 · V1 · Đề {[}SV{]} S4-p77--78 &
\pageref{BOO-04-E02} \\
\protect\hyperlink{BOO-04-E03}{BOO-04-E03} & Đọc mạch: kết quả là hằng 0
--- bẫy & D1 · V4 · Đề {[}SV{]} S4-p74 & \pageref{BOO-04-E03} \\
\protect\hyperlink{BOO-04-E04}{BOO-04-E04} & Vẽ mạch từ công thức --- và
rút gọn trước khi vẽ & D2 · V3 · Đề {[}SV{]} S4-p72 &
\pageref{BOO-04-E04} \\
\protect\hyperlink{GRA-01-E01}{GRA-01-E01} & Bậc đỉnh: vô hướng có
khuyên; có hướng & D1 · V0 · Đề {[}SV{]} S5-p20, p21, p24 &
\pageref{GRA-01-E01} \\
\protect\hyperlink{GRA-01-E02}{GRA-01-E02} & Định lý bắt tay: tồn tại đồ
thị? số cạnh? --- bẫy & D2 · V4 · {[}AI-GENERATED PRACTICE{]} &
\pageref{GRA-01-E02} \\
\protect\hyperlink{GRA-02-E01}{GRA-02-E01} & Ma trận kề & D1 · V0 · Đề
{[}SV{]} S5-p27, p28 & \pageref{GRA-02-E01} \\
\protect\hyperlink{GRA-03-E01}{GRA-03-E01} & Đẳng cấu: bác bỏ bằng bất
biến, khẳng định bằng song ánh & D2 · V0 · Đề (a) {[}SV{]} S5-p32; (b)
{[}AI-GENERATED PRACTICE{]} & \pageref{GRA-03-E01} \\
\protect\hyperlink{GRA-04-E01}{GRA-04-E01} & Đường đi đơn/sơ cấp, chu
trình sơ cấp, đỉnh khớp, cầu & D1 · V0 · Đề {[}SV{]} S5-p42 (+ đỉnh
khớp, cầu theo định nghĩa S5-p38) & \pageref{GRA-04-E01} \\
\protect\hyperlink{GRA-05-E01}{GRA-05-E01} & Chu trình Euler bằng Fleury
& D2 · V0 · Đề {[}SV{]} S5-p51 & \pageref{GRA-05-E01} \\
\protect\hyperlink{GRA-05-E02}{GRA-05-E02} & Bài toán 7 cây cầu
Königsberg & D1 · V2 · Đề {[}SV{]} S5-p44--47 & \pageref{GRA-05-E02} \\
\protect\hyperlink{GRA-06-E01}{GRA-06-E01} & Dijkstra trên đồ thị vô
hướng & D2 · V0 · Đề {[}SV{]} S5-p60 & \pageref{GRA-06-E01} \\
\protect\hyperlink{GRA-06-E02}{GRA-06-E02} & Dijkstra trên đồ thị có
hướng cho bằng ma trận trọng số & D2 · V2 · Đề {[}SV{]} S5-p66--68 &
\pageref{GRA-06-E02} \\
\protect\hyperlink{GRA-06-E03}{GRA-06-E03} & Dijkstra: đường đi ngắn
nhất từ \(a\) đến \(z\) & D2 · V1 · Đề {[}SV{]} S5-p71 &
\pageref{GRA-06-E03} \\
\protect\hyperlink{GRA-07-E01}{GRA-07-E01} & Ford--Bellman phát hiện
mạch âm & D3 · V4 · Đề {[}SV{]} S5-p75 & \pageref{GRA-07-E01} \\
\protect\hyperlink{GRA-07-E02}{GRA-07-E02} & Ford--Bellman hội tụ & D2 ·
V0 · Đề {[}SV{]} S5-p84 & \pageref{GRA-07-E02} \\
\end{longtable}

\hypertarget{ex-ch1}{%
\chapter{Ví dụ Chương I --- Logic, lượng từ, tập hợp, ánh xạ, quy
nạp}\label{ex-ch1}}

\begin{quote}
Nhãn nguồn: \textbf{{[}SV{]}} = đề lấy từ slide (vị trí ghi kèm);
\textbf{{[}AI-GENERATED PRACTICE{]}} = đề do toolkit tạo để lấp khoảng
trống kỹ năng. Nhãn lời giải: \textbf{{[}CC{]}} = đã kiểm tra độc lập
(bảng chân trị/vét cạn bằng \texttt{10\_Verification}).
\end{quote}

\hypertarget{LOG-01-E01}{%
\subsection{LOG-01-E01 · Nhận biết mệnh đề; sơ cấp hay phức
hợp}\label{LOG-01-E01}}

\textbf{D1 · V0 · Đề {[}SV{]} S1-p8, S1-p9 · Lời giải
{[}SV{]}+{[}INFERRED{]}}

\begin{itemize}
\tightlist
\item
  \textbf{Đề.} (a) Câu nào là mệnh đề: (1) Paris là thành phố của Mỹ;
  (2) \(n\) là số tự nhiên; (3) Con nhà ai mà xinh thế!; (4) 3 là số
  nguyên tố; (5) Toán rời rạc là môn bắt buộc của ngành Tin học; (6) Bạn
  có khỏe không?; (7) \(x^2+1\) luôn dương. (b) Mệnh đề nào sơ cấp, phức
  hợp: (i) 2 không là số nguyên tố; (ii) 2 là số nguyên tố; (iii) Nếu
  \(3>4\) thì trời mưa; (iv) An đang xem phim hay An đang học bài; (v)
  Hôm nay trời đẹp và \(1+1=3\).
\item
  \textbf{Nhận diện.} ``có phải mệnh đề'', ``sơ cấp/phức hợp'' →
  \protect\hyperlink{LOG-01}{LOG-01}.
\item
  \textbf{Kiến thức.} Mệnh đề = khẳng định có chân trị xác định (đúng
  hoặc sai); câu hỏi, cảm thán, mệnh lệnh không là mệnh đề (S1-p6). Phức
  hợp = xây từ mệnh đề khác bằng liên từ hoặc ``không'' (S1-p9).
\item
  \textbf{Phương pháp.} Với mỗi câu: (1) là câu khẳng định? (2) có biến
  tự do không? (3) chân trị xác định không?
\item
  \textbf{Lời giải.} (a) (1) mệnh đề, chân trị 0. (2) \textbf{không} ---
  chứa biến \(n\) chưa xác định, đây là vị từ (S1-p59). (3) không ---
  câu cảm thán. (4) mệnh đề, chân trị 1. (5) mệnh đề (chân trị phụ thuộc
  chương trình đào tạo, nhưng xác định) \textbf{{[}INFERRED{]}}. (6)
  không --- câu hỏi. (7) mệnh đề \textbf{nếu hiểu ``luôn'' là ``với mọi
  \(x\in\mathbb R\)''}: \(\forall x\in\mathbb R,\ x^2+1>0\), chân trị 1
  \textbf{{[}INFERRED{]}}.

  \begin{enumerate}
  \def\labelenumi{(\alph{enumi})}
  \setcounter{enumi}{1}
  \item
    \begin{enumerate}
    \def\labelenumii{(\roman{enumii})}
    \tightlist
    \item
      phức hợp (có ``không''); (ii) sơ cấp; (iii) phức hợp (kéo theo)
      --- chân trị \textbf{1} vì tiền đề \(3>4\) sai; (iv) phức hợp
      (tuyển); (v) phức hợp (hội) --- chân trị 0 vì \(1+1=3\) sai.
    \end{enumerate}
  \end{enumerate}
\item
  \textbf{Kết luận.} Mệnh đề: (1), (4), (5), (7). Sơ cấp: (ii); phức
  hợp: (i), (iii), (iv), (v).
\item
  \textbf{Kiểm tra.} Đối chiếu bảng chân trị
  \protect\hyperlink{F-LOG-01}{F-LOG-01}: \(0\to q\) luôn 1;
  \(q\wedge0=0\).
\item
  \textbf{Biến thể (V2).} ``\(x+1=3\)'' --- không là mệnh đề;
  ``\(\exists x\in\mathbb R,\ x+1=3\)'' --- là mệnh đề đúng (lượng từ
  hóa biến tự do).
\item
  \textbf{Lỗi thường gặp.} Coi câu chứa biến là mệnh đề; nghĩ ``Nếu
  \(3>4\) thì trời mưa'' sai vì ``vô lý'' ---
  \protect\hyperlink{T-01}{T-01}.
\item
  \textbf{Liên kết.} \protect\hyperlink{F-LOG-01}{F-LOG-01},
  \protect\hyperlink{QNT-01}{QNT-01}.
\end{itemize}

\hypertarget{LOG-02-E01}{%
\subsection{LOG-02-E01 · Lập bảng chân trị}\label{LOG-02-E01}}

\textbf{D1 · V0 · Đề {[}SV{]} S1-p22 (Bài tập) · Lời giải {[}CC{]}}

\begin{itemize}
\tightlist
\item
  \textbf{Đề.} Lập bảng chân trị của \(E(p,q)=\neg(p\wedge q)\wedge p\)
  và \(F(p,q,r)=p\wedge(q\vee r)\leftrightarrow\neg q\).
\item
  \textbf{Nhận diện.} ``lập bảng chân trị'' →
  \protect\hyperlink{LOG-02}{LOG-02}.
\item
  \textbf{Kiến thức.} \protect\hyperlink{F-LOG-01}{F-LOG-01}; \(n\) biến
  → \(2^n\) dòng. Quy ước ưu tiên: \(\leftrightarrow\) yếu nhất, nên
  \(F=\big(p\wedge(q\vee r)\big)\leftrightarrow\neg q\)
  \textbf{{[}INFERRED{]}} (slide không ghi ngoặc).
\item
  \textbf{Phương pháp.} \protect\hyperlink{M-01}{M-01}.
\item
  \textbf{Lời giải.}
\end{itemize}

\begin{longtable}[]{@{}ccccc@{}}
\toprule\noalign{}
\(p\) & \(q\) & \(p\wedge q\) & \(\neg(p\wedge q)\) & \(E\) \\
\midrule\noalign{}
\endhead
\bottomrule\noalign{}
\endlastfoot
0 & 0 & 0 & 1 & 0 \\
0 & 1 & 0 & 1 & 0 \\
1 & 0 & 0 & 1 & 1 \\
1 & 1 & 1 & 0 & 0 \\
\end{longtable}

\begin{longtable}[]{@{}ccccccc@{}}
\toprule\noalign{}
\(p\) & \(q\) & \(r\) & \(q\vee r\) & \(p\wedge(q\vee r)\) & \(\neg q\)
& \(F\) \\
\midrule\noalign{}
\endhead
\bottomrule\noalign{}
\endlastfoot
0 & 0 & 0 & 0 & 0 & 1 & 0 \\
0 & 0 & 1 & 1 & 0 & 1 & 0 \\
0 & 1 & 0 & 1 & 0 & 0 & 1 \\
0 & 1 & 1 & 1 & 0 & 0 & 1 \\
1 & 0 & 0 & 0 & 0 & 1 & 0 \\
1 & 0 & 1 & 1 & 1 & 1 & 1 \\
1 & 1 & 0 & 1 & 1 & 0 & 0 \\
1 & 1 & 1 & 1 & 1 & 0 & 0 \\
\end{longtable}

\begin{itemize}
\tightlist
\item
  \textbf{Kết luận.} \(E\) đúng duy nhất khi \(p=1,q=0\) (nên
  \(E\Leftrightarrow p\wedge\neg q\)). \(F\) đúng tại
  \((p,q,r)=(0,1,0),(0,1,1),(1,0,1)\).
\item
  \textbf{Kiểm tra.} Biến đổi:
  \(E\Leftrightarrow(\neg p\vee\neg q)\wedge p\Leftrightarrow(\neg p\wedge p)\vee(\neg q\wedge p)\Leftrightarrow p\wedge\neg q\)
  --- khớp cột \(E\). Bảng \(F\) được kiểm tra bằng chương trình.
\item
  \textbf{Biến thể (V2).} Nếu đề hiểu
  \(F'=p\wedge\big((q\vee r)\leftrightarrow\neg q\big)\) thì \(F'\) đúng
  tại \((1,0,1)\) duy nhất --- ngoặc thay đổi kết quả, nên ghi rõ quy
  ước.
\item
  \textbf{Lỗi thường gặp.} Liệt kê thiếu dòng; tính \(\leftrightarrow\)
  như \(\to\).
\item
  \textbf{Liên kết.} \protect\hyperlink{M-01}{M-01},
  \protect\hyperlink{LOG-03}{LOG-03}.
\end{itemize}

\hypertarget{LOG-02-E02}{%
\subsection{LOG-02-E02 · Hằng đúng, hằng sai --- bẫy}\label{LOG-02-E02}}

\textbf{D1 · V4 · Đề {[}SV{]} S1-p21 · Lời giải {[}CC{]}}

\begin{itemize}
\tightlist
\item
  \textbf{Đề.} Xét \(B(p,q)=p\to(p\vee q)\) và
  \(C(p,q)=p\wedge\neg(p\to q)\). Dạng nào là hằng đúng, hằng sai?
\item
  \textbf{Nhận diện.} ``hằng đúng / hằng sai'' →
  \protect\hyperlink{LOG-02}{LOG-02}.
\item
  \textbf{Kiến thức.} Hằng đúng: luôn 1; hằng sai (mâu thuẫn): luôn 0
  (S1-p21).
\item
  \textbf{Phương pháp.} \protect\hyperlink{M-02}{M-02}, đối chiếu
  \protect\hyperlink{M-01}{M-01}.
\item
  \textbf{Lời giải.}
  \(B\Leftrightarrow\neg p\vee p\vee q\Leftrightarrow1\vee q\Leftrightarrow1\)
  ⇒ \textbf{hằng đúng}.
  \(C\Leftrightarrow p\wedge\neg(\neg p\vee q)\Leftrightarrow p\wedge(p\wedge\neg q)\Leftrightarrow p\wedge\neg q\):
  bằng 1 khi \(p=1,q=0\) và bằng 0 khi \(p=0\) ⇒ \textbf{không} hằng
  đúng, \textbf{không} hằng sai.
\item
  \textbf{Kết luận.} \(B\) hằng đúng; \(C\Leftrightarrow p\wedge\neg q\)
  (khả đúng).
\item
  \textbf{Kiểm tra.} Bảng 4 dòng: \(B=(1,1,1,1)\); \(C=(0,0,1,0)\).
\item
  \textbf{Ghi chú nguồn.} Slide đặt \(C\) ngay sau định nghĩa hằng sai
  nhưng \textbf{không} nói \(C\) hằng sai; như đã viết, \(C\) không hằng
  sai (ERR-13). Dạng hằng sai gần giống là
  \(p\wedge\neg(q\to p)\Leftrightarrow p\wedge q\wedge\neg p\Leftrightarrow0\).
\item
  \textbf{Biến thể (V1).} \((p\wedge q)\to p\) --- hằng đúng (quy tắc
  đơn giản); \((p\to q)\wedge(p\wedge\neg q)\) --- hằng sai.
\item
  \textbf{Lỗi thường gặp.} Kết luận ``hằng sai'' chỉ vì thấy dấu phủ
  định; thử một dòng rồi kết luận cho mọi dòng ---
  \protect\hyperlink{T-05}{T-05}.
\item
  \textbf{Liên kết.} \protect\hyperlink{F-LOG-03}{F-LOG-03}.
\end{itemize}

\hypertarget{LOG-03-E01}{%
\subsection{LOG-03-E01 · Chứng minh tương đương bằng luật
logic}\label{LOG-03-E01}}

\textbf{D2 · V0 · Đề {[}SV{]} S1-p29 · Lời giải {[}SV{]} S1-p30,
{[}CC{]}}

\begin{itemize}
\item
  \textbf{Đề.} Chứng minh
  \((\neg p\to r)\wedge(q\to r)\Leftrightarrow(p\to q)\to r\).
\item
  \textbf{Nhận diện.} ``chứng minh ⇔'' →
  \protect\hyperlink{LOG-03}{LOG-03}.
\item
  \textbf{Kiến thức.} Luật kéo theo, phủ định kép, phân phối, De Morgan
  (\protect\hyperlink{F-LOG-03}{F-LOG-03}).
\item
  \textbf{Phương pháp.} \protect\hyperlink{M-02}{M-02}: khử \(\to\) →
  gom bằng phân phối → đưa về \(\to\).
\item
  \textbf{Lời giải.} \[
  \begin{aligned}
  (\neg p\to r)\wedge(q\to r)
  &\Leftrightarrow(p\vee r)\wedge(\neg q\vee r) &&\text{(kéo theo, phủ định kép)}\\
  &\Leftrightarrow(p\wedge\neg q)\vee r &&\text{(phân phối)}\\
  &\Leftrightarrow\neg(\neg p\vee q)\vee r &&\text{(De Morgan, phủ định kép)}\\
  &\Leftrightarrow\neg(p\to q)\vee r &&\text{(kéo theo)}\\
  &\Leftrightarrow(p\to q)\to r &&\text{(kéo theo)}
  \end{aligned}
  \]
\item
  \textbf{Kết luận.} Hai dạng tương đương logic.
\item
  \textbf{Kiểm tra.} Bảng chân trị 8 dòng: hai vế cùng bằng 0 đúng tại
  các bộ \((p,q,r)\) = \((0,0,0)\), \((0,1,0)\), \((1,1,0)\) và bằng 1 ở
  5 dòng còn lại {[}CC{]}.
\item
  \textbf{Biến thể (V1).} Chứng minh
  \((p\to r)\wedge(q\to r)\Leftrightarrow(p\vee q)\to r\) (cơ sở quy tắc
  chứng minh theo trường hợp):
  \((\neg p\vee r)\wedge(\neg q\vee r)\Leftrightarrow(\neg p\wedge\neg q)\vee r\Leftrightarrow\neg(p\vee q)\vee r\Leftrightarrow(p\vee q)\to r\).
\item
  \textbf{Lỗi thường gặp.} Phân phối sai chiều:
  \((p\vee r)\wedge(\neg q\vee r)\) là ``\(r\) chung'' nên gom thành
  \((p\wedge\neg q)\vee r\), không phải \((p\vee\neg q)\wedge r\) ---
  \protect\hyperlink{T-02}{T-02}.
\item
  \textbf{Liên kết.} \protect\hyperlink{M-02}{M-02},
  \protect\hyperlink{F-LOG-04}{F-LOG-04} (quy tắc 8).
\end{itemize}

\hypertarget{LOG-03-E02}{%
\subsection{LOG-03-E02 · Rút gọn dạng mệnh đề; hệ quả
logic}\label{LOG-03-E02}}

\textbf{D1 · V1 · {[}AI-GENERATED PRACTICE{]} · Lời giải {[}CC{]}}

\begin{itemize}
\tightlist
\item
  \textbf{Đề.} Rút gọn \(E=\neg\big(p\vee(\neg p\wedge q)\big)\) và
  chứng tỏ \(E\Rightarrow\neg p\).
\item
  \textbf{Nhận diện.} ``rút gọn'', ``hệ quả logic'' →
  \protect\hyperlink{LOG-03}{LOG-03}.
\item
  \textbf{Kiến thức.} De Morgan, phân phối, phần tử bù, trung hòa;
  \(E\Rightarrow F\) ⇔ \(E\to F\) hằng đúng
  (\protect\hyperlink{F-LOG-02}{F-LOG-02}).
\item
  \textbf{Phương pháp.} \protect\hyperlink{M-02}{M-02}.
\item
  \textbf{Lời giải.} \[
  \begin{aligned}
  E&\Leftrightarrow\neg p\wedge\neg(\neg p\wedge q) &&\text{(De Morgan)}\\
  &\Leftrightarrow\neg p\wedge(p\vee\neg q) &&\text{(De Morgan, phủ định kép)}\\
  &\Leftrightarrow(\neg p\wedge p)\vee(\neg p\wedge\neg q) &&\text{(phân phối)}\\
  &\Leftrightarrow0\vee(\neg p\wedge\neg q)\Leftrightarrow\neg p\wedge\neg q &&\text{(phần tử bù, trung hòa)}
  \end{aligned}
  \]
\end{itemize}

Khi đó \(E\to\neg p\Leftrightarrow(\neg p\wedge\neg q)\to\neg p\) là
hằng đúng (quy tắc đơn giản) ⇒ \(E\Rightarrow\neg p\).

\begin{itemize}
\tightlist
\item
  \textbf{Kết luận.} \(E\Leftrightarrow\neg p\wedge\neg q\);
  \(E\Rightarrow\neg p\) (nhưng \(\neg p\not\Rightarrow E\):
  \(p=0,q=1\)).
\item
  \textbf{Kiểm tra.} Bảng: \(E=(1,0,0,0)\) tại \(pq=00,01,10,11\) ---
  trùng \(\neg p\wedge\neg q\).
\item
  \textbf{Biến thể (V1).} Rút gọn
  \((p\wedge q)\vee(p\wedge\neg q)\Leftrightarrow p\wedge(q\vee\neg q)\Leftrightarrow p\).
\item
  \textbf{Lỗi thường gặp.} Nhầm ``\(\Rightarrow\)'' với
  ``\(\Leftrightarrow\)''.
\item
  \textbf{Liên kết.} \protect\hyperlink{F-LOG-03}{F-LOG-03}.
\end{itemize}

\hypertarget{LOG-04-E01}{%
\subsection{LOG-04-E01 · Chứng minh suy luận bằng chuỗi quy
tắc}\label{LOG-04-E01}}

\textbf{D2 · V0 · Đề {[}SV{]} S1-p50 · Lời giải {[}SV{]} S1-p51 (đọc từ
ảnh), {[}CC{]}}

\begin{itemize}
\tightlist
\item
  \textbf{Đề.} Chứng minh: \(p\to(q\to r)\), \(\ p\vee s\),
  \(\ t\to q\), \(\ \neg s\ \ \therefore\ \neg r\to\neg t\).
\item
  \textbf{Nhận diện.} Sơ đồ tiền đề/kết luận →
  \protect\hyperlink{LOG-04}{LOG-04}; có tiền đề đơn \(\neg s\) → khởi
  động bằng tam đoạn luận rời.
\item
  \textbf{Kiến thức.} \protect\hyperlink{F-LOG-04}{F-LOG-04} (1, 3, 4),
  luật phản đảo.
\item
  \textbf{Phương pháp.} \protect\hyperlink{M-03}{M-03}.
\item
  \textbf{Lời giải.}
\end{itemize}

\begin{longtable}[]{@{}lll@{}}
\toprule\noalign{}
Bước & Mệnh đề & Lý do \\
\midrule\noalign{}
\endhead
\bottomrule\noalign{}
\endlastfoot
1 & \(p\vee s\) & tiền đề \\
2 & \(\neg s\) & tiền đề \\
3 & \(p\) & 1, 2: tam đoạn luận rời \\
4 & \(p\to(q\to r)\) & tiền đề \\
5 & \(q\to r\) & 3, 4: Modus Ponens \\
6 & \(t\to q\) & tiền đề \\
7 & \(t\to r\) & 6, 5: tam đoạn luận \\
8 & \(\neg r\to\neg t\) & 7: luật kéo theo (phản đảo) \\
\end{longtable}

\begin{itemize}
\tightlist
\item
  \textbf{Kết luận.} Suy luận đúng.
\item
  \textbf{Kiểm tra.} Vét cạn 32 bộ chân trị: không bộ nào làm 4 tiền đề
  đúng mà kết luận sai {[}CC{]}.
\item
  \textbf{Biến thể (V1).} Thay tiền đề \(\neg s\) bằng \(s\): suy luận
  \textbf{sai} (phản ví dụ \(p=0,q=1,r=0,s=1,t=1\): mọi tiền đề đúng,
  \(\neg r\to\neg t=1\to0=0\)).
\item
  \textbf{Lỗi thường gặp.} Dùng \(t\to q\) và \(q\) để suy ra \(t\)
  (khẳng định hệ quả) --- \protect\hyperlink{T-04}{T-04}.
\item
  \textbf{Liên kết.} \protect\hyperlink{M-03}{M-03},
  \protect\hyperlink{LOG-05}{LOG-05}.
\end{itemize}

\hypertarget{LOG-04-E02}{%
\subsection{LOG-04-E02 · Chứng minh suy luận bằng phản
chứng}\label{LOG-04-E02}}

\textbf{D2 · V3 · Đề {[}SV{]} S1-p52 · Lời giải {[}SV{]} S1-p53--54
(ảnh), {[}CC{]}}

\begin{itemize}
\tightlist
\item
  \textbf{Đề.} Kiểm tra suy luận: \(p\to q\), \(\ \neg r\vee s\),
  \(\ p\vee r\ \ \therefore\ \neg q\to s\).
\item
  \textbf{Nhận diện.} Kết luận dạng kéo theo → phản chứng
  (\protect\hyperlink{LOG-04}{LOG-04}).
\item
  \textbf{Kiến thức.} Quy tắc mâu thuẫn;
  \(\neg(a\to b)\Leftrightarrow a\wedge\neg b\).
\item
  \textbf{Phương pháp.} \protect\hyperlink{M-04}{M-04}.
\item
  \textbf{Lời giải.}
\end{itemize}

\begin{longtable}[]{@{}lll@{}}
\toprule\noalign{}
Bước & Mệnh đề & Lý do \\
\midrule\noalign{}
\endhead
\bottomrule\noalign{}
\endlastfoot
1 & \(\neg(\neg q\to s)\) & giả thiết phản chứng \\
2 & \(\neg q\wedge\neg s\) & 1: kéo theo + De Morgan \\
3 & \(\neg q\) và \(\neg s\) & 2: quy tắc đơn giản \\
4 & \(p\to q\) & tiền đề \\
5 & \(\neg p\) & 3, 4: quy tắc phủ định (Modus Tollens) \\
6 & \(\neg r\vee s\) & tiền đề \\
7 & \(\neg r\) & 3 (\(\neg s\)), 6: tam đoạn luận rời \\
8 & \(\neg p\wedge\neg r\) & 5, 7: nối liền \\
9 & \(\neg(p\vee r)\) & 8: De Morgan \\
10 & \(p\vee r\) & tiền đề \\
11 & \(0\) & 9, 10: phần tử bù \\
\end{longtable}

\begin{itemize}
\tightlist
\item
  \textbf{Kết luận.} Suy luận đúng (quy tắc phản chứng).
\item
  \textbf{Kiểm tra.} Vét cạn 16 bộ {[}CC{]}. Cách trực tiếp:
  \(\neg q\to\neg p\) (phản đảo 4); \(\neg p\to r\) (từ 10); \(r\to s\)
  (từ 6) ⇒ tam đoạn luận: \(\neg q\to s\).
\item
  \textbf{Biến thể (V1).} Bỏ tiền đề \(p\vee r\): suy luận sai (phản ví
  dụ \(p=q=r=s=0\)).
\item
  \textbf{Lỗi thường gặp.} Viết giả thiết phản chứng là \(q\to\neg s\)
  --- \protect\hyperlink{T-03}{T-03}.
\item
  \textbf{Liên kết.} \protect\hyperlink{M-04}{M-04}.
\end{itemize}

\hypertarget{LOG-04-E03}{%
\subsection{LOG-04-E03 · Phản chứng hoặc tam đoạn luận nối
tiếp}\label{LOG-04-E03}}

\textbf{D2 · V1 · Đề {[}SV{]} S1-p46 · Lời giải {[}CC{]}}

\begin{itemize}
\tightlist
\item
  \textbf{Đề.} Chứng minh: \(p\to r\), \(\ \neg p\to q\),
  \(\ q\to s\ \ \therefore\ \neg r\to s\).
\item
  \textbf{Nhận diện.} Chuỗi kéo theo →
  \protect\hyperlink{LOG-04}{LOG-04}.
\item
  \textbf{Kiến thức.} \protect\hyperlink{F-LOG-04}{F-LOG-04}; phản đảo.
\item
  \textbf{Phương pháp.} \protect\hyperlink{M-04}{M-04} (slide yêu cầu
  phản chứng); cách 2 \protect\hyperlink{M-03}{M-03}.
\item
  \textbf{Lời giải (phản chứng).} Giả sử
  \(\neg(\neg r\to s)\Leftrightarrow\neg r\wedge\neg s\). Từ \(\neg r\)
  và \(p\to r\): \(\neg p\) (Modus Tollens). Từ \(\neg p\) và
  \(\neg p\to q\): \(q\) (Modus Ponens). Từ \(q\) và \(q\to s\): \(s\)
  (Modus Ponens). Có \(s\wedge\neg s\Leftrightarrow0\) --- mâu thuẫn.
\item
  \textbf{Kết luận.} Suy luận đúng.
\item
  \textbf{Kiểm tra (cách 2).} \(\neg r\to\neg p\) (phản đảo),
  \(\neg p\to q\), \(q\to s\) ⇒ tam đoạn luận hai lần: \(\neg r\to s\).
  Vét cạn {[}CC{]}.
\item
  \textbf{Biến thể (V2).} Lời văn: ``Nếu trời mưa thì đường ướt; nếu
  trời không mưa thì tôi đi bộ; nếu tôi đi bộ thì tôi mệt. Vậy nếu đường
  không ướt thì tôi mệt.'' --- cùng cấu trúc ⇒ đúng.
\item
  \textbf{Lỗi thường gặp.} Suy \(\neg p\to\neg r\) từ \(p\to r\) (đảo
  sai) --- \protect\hyperlink{T-03}{T-03}.
\item
  \textbf{Liên kết.} \protect\hyperlink{LOG-04-E02}{LOG-04-E02}.
\end{itemize}

\hypertarget{LOG-05-E01}{%
\subsection{LOG-05-E01 · Suy luận lời văn --- tìm phản ví
dụ}\label{LOG-05-E01}}

\textbf{D3 · V2 · Đề {[}SV{]} S1-p49 · Lời giải {[}SV{]} S1-p49,
{[}CC{]}}

\begin{itemize}
\tightlist
\item
  \textbf{Đề.} ``Ông Minh nói rằng nếu không được tăng lương thì ông ta
  sẽ nghỉ việc. Mặt khác, nếu ông ấy nghỉ việc và vợ ông ấy bị mất việc
  thì phải bán xe. Biết rằng nếu vợ ông Minh hay đi làm trễ thì trước
  sau gì cũng sẽ bị mất việc và cuối cùng ông Minh đã được tăng lương.
  Suy ra nếu ông Minh không bán xe thì vợ ông ta đã không đi làm trễ.''
  Suy luận có đúng không?
\item
  \textbf{Nhận diện.} ``Suy luận sau có đúng không'' + lời văn →
  \protect\hyperlink{LOG-05}{LOG-05}.
\item
  \textbf{Kiến thức.} Phản ví dụ (\protect\hyperlink{F-LOG-04}{F-LOG-04}
  quy tắc 9).
\item
  \textbf{Phương pháp.} Đặt biến → \protect\hyperlink{M-05}{M-05}.
\item
  \textbf{Lời giải.} \(p\): ông Minh được tăng lương; \(q\): ông Minh
  nghỉ việc; \(r\): vợ mất việc; \(s\): phải bán xe; \(t\): vợ hay đi
  làm trễ. Suy luận: \(\neg p\to q\), \(\ q\wedge r\to s\),
  \(\ t\to r\), \(\ p\ \ \therefore\ \neg s\to\neg t\). Đặt kết luận
  \(=0\): \(\neg s=1,\ \neg t=0\) ⇒ \(s=0,\ t=1\). Tiền đề \(p=1\).
  \(t\to r=1\) với \(t=1\) ⇒ \(r=1\). \(q\wedge r\to s=1\) với
  \(r=1,s=0\) ⇒ \(q=0\). Thử lại \(\neg p\to q\): \(0\to0=1\) ✓.
\item
  \textbf{Kết luận.} Có bộ \(p=1,q=0,r=1,s=0,t=1\) làm mọi tiền đề đúng,
  kết luận sai ⇒ \textbf{suy luận sai}. (Tình huống: được tăng lương,
  không nghỉ việc, vợ đi trễ và mất việc, không bán xe.)
\item
  \textbf{Kiểm tra.} Vét cạn 32 bộ: phản ví dụ trên là duy nhất có
  \(s=0,t=1\) {[}CC{]}.
\item
  \textbf{Biến thể (V3).} Đổi kết luận thành ``Nếu ông Minh nghỉ việc và
  vợ hay đi làm trễ thì phải bán xe'' (\(q\wedge t\to s\)):
  \textbf{đúng} --- từ \(t\) có \(r\) (\(t\to r\)), có \(q\wedge r\),
  suy ra \(s\).
\item
  \textbf{Lỗi thường gặp.} Đặt biến cho câu phủ định (đặt ``không tăng
  lương'' là \(p\)) làm rối dấu; thử ngẫu nhiên vài bộ rồi kết luận
  ``đúng''.
\item
  \textbf{Liên kết.} \protect\hyperlink{M-05}{M-05}.
\end{itemize}

\hypertarget{LOG-05-E02}{%
\subsection{LOG-05-E02 · Suy luận sai --- và kết luận đúng là
gì}\label{LOG-05-E02}}

\textbf{D2 · V4 · Đề {[}SV{]} S1-p55 · Lời giải {[}SV{]} S1-p56 (ảnh),
{[}CC{]}}

\begin{itemize}
\tightlist
\item
  \textbf{Đề.} Kiểm tra suy luận: \(p\), \(\ p\to r\),
  \(\ p\to(q\vee\neg r)\), \(\ \neg q\vee\neg s\ \ \therefore\ s\).
\item
  \textbf{Nhận diện.} Kết luận đơn biến; nghi sai →
  \protect\hyperlink{LOG-05}{LOG-05}.
\item
  \textbf{Kiến thức.} ``Hệ (tiền đề = 1, kết luận = 0) vô nghiệm ⇔ suy
  luận đúng'' (S1-p56).
\item
  \textbf{Phương pháp.} \protect\hyperlink{M-05}{M-05}.
\item
  \textbf{Lời giải.} Hệ: (1) \(p=1\); (2) \(p\to r=1\); (3)
  \(p\to(q\vee\neg r)=1\); (4) \(\neg q\vee\neg s=1\); (5) \(s=0\). Từ
  (5): (4) đúng tự động. Từ (1),(2): \(r=1\). Từ (1),(3):
  \(q\vee\neg r=1\) ⇒ \(q=1\). Thử lại: \(p=q=r=1,s=0\) thỏa cả 5 phương
  trình.
\item
  \textbf{Kết luận.} Hệ có nghiệm ⇒ \textbf{suy luận sai} (phản ví dụ
  \(p=q=r=1\), \(s=0\)).
\item
  \textbf{Kiểm tra.} Vét cạn 16 bộ: đây là phản ví dụ duy nhất {[}CC{]}.
\item
  \textbf{Biến thể (V4).} Với cùng tiền đề, kết luận \textbf{\(\neg s\)}
  là đúng: \(p\) ⇒ \(r\) ⇒ (vì \(p\to(q\vee\neg r)\)) \(q\) ⇒ (tam đoạn
  luận rời với \(\neg q\vee\neg s\)) \(\neg s\).
\item
  \textbf{Lỗi thường gặp.} Dừng khi thấy (4) ``tự đúng'' mà không giải
  hết các biến.
\item
  \textbf{Liên kết.} \protect\hyperlink{M-05}{M-05},
  \protect\hyperlink{LOG-04}{LOG-04}.
\end{itemize}

\hypertarget{QNT-01-E01}{%
\subsection{QNT-01-E01 · Chân trị mệnh đề lượng từ một và hai
biến}\label{QNT-01-E01}}

\textbf{D1--D2 · V0 · Đề {[}SV{]} S1-p61, p63, p65--67 · Lời giải
{[}SV{]}+{[}CC{]}}

\begin{itemize}
\tightlist
\item
  \textbf{Đề.} (a) Với \(x\in\mathbb R\): (1)
  \(\forall x,\ x^2+3x+1\le0\); (2) \(\exists x,\ x^2+3x+1\le0\); (3)
  \(\forall x,\ x^2+1\ge2x\); (4) \(\exists x,\ x^2+1<0\). (b) Với
  \(x,y\in\mathbb R\), xét \(p(x,y)\): ``\(x+2y<1\)'':
  \(\forall x\forall y\); \(\forall x\exists y\);
  \(\exists x\forall y\); \(\exists x\exists y\).
\item
  \textbf{Nhận diện.} ``đúng hay sai'' + \(\forall,\exists\) →
  \protect\hyperlink{QNT-01}{QNT-01}.
\item
  \textbf{Kiến thức.} \protect\hyperlink{F-QNT-01}{F-QNT-01}.
\item
  \textbf{Phương pháp.} \protect\hyperlink{M-06}{M-06}.
\item
  \textbf{Lời giải.} (a) (1) \textbf{Sai}: \(x=0\) cho \(1\le0\) sai.
  (2) \textbf{Đúng}: \(x=-1\): \(1-3+1=-1\le0\). (3) \textbf{Đúng}:
  \(x^2+1-2x=(x-1)^2\ge0\ \forall x\). (4) \textbf{Sai}:
  \(x^2+1\ge1>0\ \forall x\).

  \begin{enumerate}
  \def\labelenumi{(\alph{enumi})}
  \setcounter{enumi}{1}
  \tightlist
  \item
    \(\forall x\forall y\): \textbf{sai}, \(x_0=0,y_0=1\): \(2\ge1\).
    \(\forall x\exists y\): \textbf{đúng}, với \(x=a\) chọn
    \(y_a=-a/2\): \(a+2y_a=0<1\). \(\exists x\forall y\): \textbf{sai},
    với mọi \(a\), chọn \(y=-a/2+2\): \(a+2y=4\ge1\).
    \(\exists x\exists y\): \textbf{đúng}, \(x_0=y_0=0\).
  \end{enumerate}
\item
  \textbf{Kết luận.} (a) S, Đ, Đ, S. (b) S, Đ, S, Đ.
\item
  \textbf{Kiểm tra.} Định lý hoán vị (3):
  \(\exists x\forall y\Rightarrow\forall y\exists x\) --- ở đây vế trái
  sai nên không mâu thuẫn; \(\forall x\exists y\) đúng nhưng
  \(\exists x\forall y\) sai minh họa \textbf{chiều đảo không đúng}.
\item
  \textbf{Biến thể (V1).} S1-p61: \(p(x)\): ``\(x^2-2x+1=0\)'' chỉ đúng
  tại \(x=1\) ⇒ \(\exists x\,p(x)\) đúng, \(\forall x\,p(x)\) sai;
  \(p(x)\): ``\(x^2-2x+3=0\)'' (\(\Delta'<0\)) ⇒ \(\exists x\,p(x)\)
  sai.
\item
  \textbf{Lỗi thường gặp.} Ở \(\exists x\forall y\) chọn \(y\) trước rồi
  chọn \(x\) theo \(y\) --- \protect\hyperlink{T-06}{T-06}.
\item
  \textbf{Liên kết.} \protect\hyperlink{QNT-02}{QNT-02}.
\end{itemize}

\hypertarget{QNT-01-E02}{%
\subsection{QNT-01-E02 · Thứ tự lượng từ và miền xác định ---
bẫy}\label{QNT-01-E02}}

\textbf{D2 · V4 · {[}AI-GENERATED PRACTICE{]} · Lời giải {[}CC{]}}

\begin{itemize}
\tightlist
\item
  \textbf{Đề.} Xét \(p(x,y)\): ``\(x+y=0\)''. Xác định chân trị: (a)
  \(\forall x\in\mathbb Z\ \exists y\in\mathbb Z,\ p(x,y)\); (b)
  \(\exists y\in\mathbb Z\ \forall x\in\mathbb Z,\ p(x,y)\); (c)
  \(\forall x\in\mathbb N\ \exists y\in\mathbb N,\ p(x,y)\).
\item
  \textbf{Nhận diện.} Hai lượng từ khác loại + đổi miền →
  \protect\hyperlink{QNT-01}{QNT-01}.
\item
  \textbf{Kiến thức.} \protect\hyperlink{F-QNT-01}{F-QNT-01};
  \(\mathbb N=\{0,1,2,\dots\}\) (S1-p58).
\item
  \textbf{Phương pháp.} \protect\hyperlink{M-06}{M-06}.
\item
  \textbf{Lời giải.} (a) Đúng: với mỗi \(x\), chọn \(y=-x\in\mathbb Z\).
  (b) Sai: giả sử có \(y_0\) với \(x+y_0=0\ \forall x\); lấy \(x=0\) ⇒
  \(y_0=0\); lấy \(x=1\) ⇒ \(1=0\), vô lý. (c) Sai: \(x=1\) cần
  \(y=-1\notin\mathbb N\).
\item
  \textbf{Kết luận.} Đ, S, S.
\item
  \textbf{Kiểm tra.} (b) phù hợp định lý: \(\exists y\forall x\) sai
  trong khi \(\forall x\exists y\) đúng --- chiều
  \(\forall\exists\Rightarrow\exists\forall\) không đúng.
\item
  \textbf{Biến thể (V1).}
  \(\exists x\in\mathbb N\ \forall y\in\mathbb N,\ x\le y\): đúng
  (\(x=0\)); trên \(\mathbb Z\): sai.
\item
  \textbf{Lỗi thường gặp.} Bỏ qua miền (\(\mathbb N\) vs \(\mathbb Z\))
  --- \protect\hyperlink{T-06}{T-06}.
\item
  \textbf{Liên kết.} \protect\hyperlink{QNT-01-E01}{QNT-01-E01}.
\end{itemize}

\hypertarget{QNT-02-E01}{%
\subsection{QNT-02-E01 · Phủ định mệnh đề lượng từ}\label{QNT-02-E01}}

\textbf{D2 · V0 · Đề {[}SV{]} S1-p71 (+ một câu AI) · Lời giải {[}SV{]}}

\begin{itemize}
\tightlist
\item
  \textbf{Đề.} Phủ định: (a) \(\forall x\in A,\ 2x+1\le0\); (b)
  \(\forall\varepsilon>0,\ \exists\delta>0,\ \forall x\in\mathbb R,\ |x-a|<\delta\to|f(x)-f(a)|<\varepsilon\);
  (c) {[}AI{]}
  \(\forall x\in\mathbb R,\ \exists y\in\mathbb R,\ x+2y<1\).
\item
  \textbf{Nhận diện.} ``phủ định'' → \protect\hyperlink{QNT-02}{QNT-02}.
\item
  \textbf{Kiến thức.} \protect\hyperlink{F-QNT-02}{F-QNT-02};
  \(\neg(A\to B)\Leftrightarrow A\wedge\neg B\).
\item
  \textbf{Phương pháp.} \protect\hyperlink{M-07}{M-07}.
\item
  \textbf{Lời giải.} (a) \(\exists x\in A,\ 2x+1>0\). (b)
  \(\exists\varepsilon>0,\ \forall\delta>0,\ \exists x\in\mathbb R,\ |x-a|<\delta\ \wedge\ |f(x)-f(a)|\ge\varepsilon\).
  (c) \(\exists x\in\mathbb R,\ \forall y\in\mathbb R,\ x+2y\ge1\).
\item
  \textbf{Kết luận.} Như trên.
\item
  \textbf{Kiểm tra.} (c): mệnh đề gốc đúng
  (\protect\hyperlink{QNT-01-E01}{QNT-01-E01}) nên phủ định phải sai ---
  đúng vậy (với mọi \(x\), chọn \(y=-x/2\) cho \(x+2y=0<1\)).
\item
  \textbf{Biến thể (V2).} Phủ định ``Mọi sinh viên đều có ít nhất một
  môn đạt điểm A'' ⇒ ``Có sinh viên mà mọi môn đều không đạt điểm A''.
\item
  \textbf{Lỗi thường gặp.} Viết \(\varepsilon\le0\) (phủ định cả miền);
  viết \(|x-a|\ge\delta\to\dots\) --- \protect\hyperlink{T-07}{T-07}.
\item
  \textbf{Liên kết.} \protect\hyperlink{M-07}{M-07}.
\end{itemize}

\hypertarget{SET-01-E01}{%
\subsection{SET-01-E01 · Chứng minh đẳng thức tập
hợp}\label{SET-01-E01}}

\textbf{D2 · V0 · {[}AI-GENERATED PRACTICE{]} · Lời giải {[}CC{]}}

\begin{itemize}
\item
  \textbf{Đề.} Chứng minh
  \(A\setminus(B\cup C)=(A\setminus B)\cap(A\setminus C)\).
\item
  \textbf{Nhận diện.} Đẳng thức tập hợp →
  \protect\hyperlink{SET-01}{SET-01}.
\item
  \textbf{Kiến thức.} Định nghĩa \(\cup,\cap,\setminus\)
  (\protect\hyperlink{F-SET-01}{F-SET-01}); De Morgan logic.
\item
  \textbf{Phương pháp.} \protect\hyperlink{M-08}{M-08} cách 1.
\item
  \textbf{Lời giải.} \[
  \begin{aligned}
  x\in A\setminus(B\cup C)&\Leftrightarrow x\in A\wedge\neg(x\in B\vee x\in C)\\
  &\Leftrightarrow x\in A\wedge x\notin B\wedge x\notin C &&\text{(De Morgan)}\\
  &\Leftrightarrow(x\in A\wedge x\notin B)\wedge(x\in A\wedge x\notin C) &&\text{(lũy đẳng, giao hoán, kết hợp)}\\
  &\Leftrightarrow x\in(A\setminus B)\cap(A\setminus C).
  \end{aligned}
  \]
\item
  \textbf{Kết luận.} Hai tập bằng nhau.
\item
  \textbf{Kiểm tra.} Cách 2:
  \(A\setminus(B\cup C)=A\cap\overline{B\cup C}=A\cap\bar B\cap\bar C=(A\cap\bar B)\cap(A\cap\bar C)\).
  Thử \(A=\{1,2,3,4\},B=\{1\},C=\{2\}\): hai vế \(=\{3,4\}\); kiểm tra
  vét cạn trên mọi bộ tập con của \(\{1,2,3\}\) {[}CC{]}.
\item
  \textbf{Biến thể (V4).}
  \(A\setminus(B\cap C)=(A\setminus B)\cap(A\setminus C)\)?
  \textbf{Sai}: \(A=\{1,2\},B=\{1\},C=\{2\}\): vế trái \(\{1,2\}\), vế
  phải \(\emptyset\). Đúng là
  \(A\setminus(B\cap C)=(A\setminus B)\cup(A\setminus C)\).
\item
  \textbf{Lỗi thường gặp.} Chứng minh bằng một ví dụ cụ thể (chỉ dùng
  được để \textbf{bác bỏ}).
\item
  \textbf{Liên kết.} \protect\hyperlink{F-SET-01}{F-SET-01}.
\end{itemize}

\hypertarget{SET-02-E01}{%
\subsection{SET-02-E01 · Tập các tập con, tích
Đề-các}\label{SET-02-E01}}

\textbf{D1 · V0 · Đề {[}SV{]} S1-p80--81 (câu hỏi mở) · Lời giải
{[}CC{]}}

\begin{itemize}
\tightlist
\item
  \textbf{Đề.} (a) \(Y=\{1,2,3\}\), tìm \(P(Y)\). (b) \(|X|=n\) thì
  \(|P(X)|=?\) (c) \(A=\{a,b\}\), \(B=\{1,2,3\}\): tìm \(A\times B\),
  \(|A\times B|\); so sánh với \(B\times A\).
\item
  \textbf{Nhận diện.} \(P(\cdot)\), \(\times\) →
  \protect\hyperlink{SET-02}{SET-02}.
\item
  \textbf{Kiến thức.} \protect\hyperlink{F-SET-02}{F-SET-02}.
\item
  \textbf{Phương pháp.} Liệt kê tập con theo số phần tử 0, 1, 2, 3;
  nguyên lý nhân (\protect\hyperlink{F-CNT-01}{F-CNT-01}) cho
  \(|A\times B|\).
\item
  \textbf{Lời giải.} (a) \(P(Y)\) gồm 8 tập con: \(\emptyset\);
  \(\{1\}\), \(\{2\}\), \(\{3\}\); \(\{1,2\}\), \(\{1,3\}\),
  \(\{2,3\}\); \(\{1,2,3\}\).

  \begin{enumerate}
  \def\labelenumi{(\alph{enumi})}
  \setcounter{enumi}{1}
  \tightlist
  \item
    Mỗi phần tử có 2 lựa chọn (thuộc/không thuộc tập con) ⇒ \(2^n\).
  \item
    \(A\times B\) gồm \((a,1)\), \((a,2)\), \((a,3)\), \((b,1)\),
    \((b,2)\), \((b,3)\); \(|A\times B|=2\cdot3=6\);
    \((1,a)\in B\times A\) nhưng \((1,a)\notin A\times B\) ⇒
    \(A\times B\ne B\times A\).
  \end{enumerate}
\item
  \textbf{Kết luận.} \(|P(Y)|=8\); \(|P(X)|=2^n\);
  \(|A\times B|=|A||B|=6\).
\item
  \textbf{Kiểm tra.} \(\sum_k C_3^k=1+3+3+1=8=2^3\).
\item
  \textbf{Biến thể (V3).} \(|P(A\times B)|=2^{6}=64\);
  \(|P(\emptyset)|=1\) (chỉ có \(\emptyset\)).
\item
  \textbf{Lỗi thường gặp.} Quên \(\emptyset\) hoặc chính \(Y\) trong
  \(P(Y)\).
\item
  \textbf{Liên kết.} \protect\hyperlink{CNT-04}{CNT-04}.
\end{itemize}

\hypertarget{FUN-01-E01}{%
\subsection{FUN-01-E01 · Ảnh và ảnh ngược}\label{FUN-01-E01}}

\textbf{D1 · V0 · Đề {[}SV{]} S1-p87 · Lời giải {[}SV{]}, {[}CC{]}}

\begin{itemize}
\tightlist
\item
  \textbf{Đề.} \(f:\mathbb R\to\mathbb R\), \(f(x)=x^2+1\). Tìm
  \(f([1,3])\), \(f([-2,-1])\), \(f([-1,3])\), \(f((1,5))\),
  \(f^{-1}(1)\), \(f^{-1}(2)\), \(f^{-1}(-5)\), \(f^{-1}([2,5])\).
\item
  \textbf{Nhận diện.} \(f(A)\), \(f^{-1}(B)\) →
  \protect\hyperlink{FUN-01}{FUN-01}.
\item
  \textbf{Kiến thức.} \protect\hyperlink{F-FUN-01}{F-FUN-01}; \(f\) giảm
  trên \((-\infty,0]\), tăng trên \([0,+\infty)\), \(\min f=f(0)=1\).
\item
  \textbf{Phương pháp.} Ảnh: dùng tính đơn điệu trên từng khúc; ảnh
  ngược: giải bất phương trình \(f(x)\in B\).
\item
  \textbf{Lời giải.} \(f([1,3])=[2,10]\) (tăng).
  \(f([-2,-1])=[f(-1),f(-2)]=[2,5]\) (giảm). \(f([-1,3])\): chứa \(0\) ⇒
  min \(=1\), max \(=\max(f(-1),f(3))=10\) ⇒ \([1,10]\).
  \(f((1,5))=(2,26)\). \(f^{-1}(1)\): \(x^2=0\) ⇒ \(\{0\}\).
  \(f^{-1}(2)\): \(x^2=1\) ⇒ \(\{-1,1\}\). \(f^{-1}(-5)\): \(x^2=-6\) vô
  nghiệm ⇒ \(\emptyset\). \(f^{-1}([2,5])\): \(1\le x^2\le4\) ⇒
  \([-2,-1]\cup[1,2]\).
\item
  \textbf{Kết luận.} Như trên.
\item
  \textbf{Kiểm tra.} Lưới số \(x\in[-4,4]\) bước 0.01 xác nhận
  \(f^{-1}([2,5])\) và \(f([-1,3])\) {[}CC{]}.
\item
  \textbf{Biến thể (V1).} \(f([-3,2])=[1,10]\); \(f^{-1}([0,1])=\{0\}\);
  \(f^{-1}((5,10])=[-3,-2)\cup(2,3]\).
\item
  \textbf{Lỗi thường gặp.} Tính \(f([-1,3])=[f(-1),f(3)]=[2,10]\) (quên
  điểm cực tiểu \(x=0\)) --- \protect\hyperlink{T-08}{T-08}.
\item
  \textbf{Liên kết.} \protect\hyperlink{FUN-02-E01}{FUN-02-E01}.
\end{itemize}

\hypertarget{FUN-02-E01}{%
\subsection{FUN-02-E01 · Đơn, toàn, song ánh; ánh xạ
ngược}\label{FUN-02-E01}}

\textbf{D2 · V0 · Đề {[}SV{]} S1-p88, p90, p92, p94 · Lời giải {[}SV{]},
{[}CC{]}}

\begin{itemize}
\tightlist
\item
  \textbf{Đề.} Xét: (a) \(f:\mathbb N\to\mathbb R\), \(f(x)=x^2+1\); (b)
  \(g:\mathbb R\to\mathbb R\), \(g(x)=x^2+1\); (c)
  \(h:\mathbb R\to\mathbb R\), \(h(x)=x^3+1\); (d)
  \(k:\mathbb R\to\mathbb R\), \(k(x)=2x+1\) --- tìm \(k^{-1}\).
\item
  \textbf{Nhận diện.} ``đơn/toàn/song ánh'', ``ánh xạ ngược'' →
  \protect\hyperlink{FUN-02}{FUN-02}.
\item
  \textbf{Kiến thức.} \protect\hyperlink{F-FUN-02}{F-FUN-02},
  \protect\hyperlink{F-FUN-03}{F-FUN-03}.
\item
  \textbf{Phương pháp.} \protect\hyperlink{M-09}{M-09}: giải \(f(x)=y\)
  trong miền.
\item
  \textbf{Lời giải.} (a)
  \(x^2+1=x'^2+1\Rightarrow x^2=x'^2\Rightarrow x=x'\) (vì \(x,x'\ge0\))
  ⇒ \textbf{đơn ánh}; \(y=0\): \(x^2=-1\) vô nghiệm ⇒ không toàn ánh.
  (b) \(g(1)=g(-1)=2\) ⇒ \textbf{không đơn ánh}; \(y=0\) vô nghiệm ⇒
  không toàn ánh. (c) \(x^3+1=y\Leftrightarrow x=\sqrt[3]{y-1}\), nghiệm
  duy nhất với mọi \(y\in\mathbb R\) ⇒ \textbf{song ánh}. (d)
  \(2x+1=y\Leftrightarrow x=\frac{y-1}{2}\), duy nhất ⇒ song ánh,
  \(k^{-1}(y)=\frac{y-1}{2}\) (hay \(k^{-1}(x)=\frac{x-1}{2}\)).
\item
  \textbf{Kết luận.} (a) đơn ánh, không toàn ánh; (b) không đơn, không
  toàn; (c), (d) song ánh; \(k^{-1}(y)=\frac{y-1}{2}\).
\item
  \textbf{Kiểm tra.} \(k(k^{-1}(y))=2\cdot\frac{y-1}{2}+1=y\);
  \(k^{-1}(k(x))=\frac{2x+1-1}{2}=x\).
\item
  \textbf{Ghi chú nguồn.} Slide viết ``\(f^{-1}(x)=(y-1)/2\)'' --- lẫn
  biến (ERR-01).
\item
  \textbf{Biến thể (V4).} \(g:[0,+\infty)\to[1,+\infty)\),
  \(g(x)=x^2+1\) là song ánh, \(g^{-1}(y)=\sqrt{y-1}\) --- cùng công
  thức, đổi miền/đích thì kết quả đổi.
\item
  \textbf{Lỗi thường gặp.} Kết luận theo công thức mà bỏ qua miền/đích
  --- \protect\hyperlink{T-08}{T-08}.
\item
  \textbf{Liên kết.} \protect\hyperlink{M-09}{M-09}.
\end{itemize}

\hypertarget{FUN-02-E02}{%
\subsection{FUN-02-E02 · Hàm phân thức: đơn ánh nhưng không toàn ánh ---
bẫy tập đích}\label{FUN-02-E02}}

\textbf{D3 · V4 · {[}AI-GENERATED PRACTICE{]} · Lời giải {[}CC{]}}

\begin{itemize}
\tightlist
\item
  \textbf{Đề.} \(f:\mathbb R\setminus\{1\}\to\mathbb R\),
  \(f(x)=\dfrac{2x+1}{x-1}\). \(f\) có đơn ánh, toàn ánh không? Sửa tập
  đích để \(f\) là song ánh và tìm \(f^{-1}\).
\item
  \textbf{Nhận diện.} Miền khuyết điểm, hàm phân thức →
  \protect\hyperlink{FUN-02}{FUN-02}.
\item
  \textbf{Kiến thức.} \protect\hyperlink{F-FUN-02}{F-FUN-02}.
\item
  \textbf{Phương pháp.} \protect\hyperlink{M-09}{M-09}.
\item
  \textbf{Lời giải.} Với \(y\in\mathbb R\):
  \(\frac{2x+1}{x-1}=y\Leftrightarrow2x+1=yx-y\) (vì \(x\ne1\))
  \(\Leftrightarrow x(y-2)=y+1\). • \(y\ne2\): \(x=\frac{y+1}{y-2}\), và
  \(x\ne1\) (vì \(\frac{y+1}{y-2}=1\Leftrightarrow y+1=y-2\) vô lý) ⇒
  đúng một nghiệm. • \(y=2\): \(0\cdot x=3\) vô nghiệm. Vậy mỗi \(y\) có
  \(\le1\) nghiệm ⇒ \textbf{đơn ánh}; \(y=2\) không có tạo ảnh ⇒
  \textbf{không toàn ánh}.
\item
  \textbf{Kết luận.} Đơn ánh, không toàn ánh.
  \(f:\mathbb R\setminus\{1\}\to\mathbb R\setminus\{2\}\) là song ánh,
  \(f^{-1}(y)=\dfrac{y+1}{y-2}\).
\item
  \textbf{Kiểm tra.}
  \(f(f^{-1}(y))=\dfrac{2\frac{y+1}{y-2}+1}{\frac{y+1}{y-2}-1}\); tử
  \(=\frac{3y}{y-2}\), mẫu \(=\frac{3}{y-2}\) ⇒ \(f(f^{-1}(y))=y\) ✓.
  Thử số: \(f(4)=3\), \(f^{-1}(3)=4\) ✓.
\item
  \textbf{Biến thể (V1).} \(f:\mathbb R\to\mathbb R\),
  \(f(x)=\frac{2x}{x^2+1}\): không đơn ánh
  (\(f(2)=f(\tfrac12)=\tfrac45\)), không toàn ánh (ảnh là \([-1,1]\)).
\item
  \textbf{Lỗi thường gặp.} Quên kiểm tra nghiệm \(x\) tìm được có thuộc
  miền (\(x\ne1\)).
\item
  \textbf{Liên kết.} \protect\hyperlink{FUN-02-E01}{FUN-02-E01}.
\end{itemize}

\hypertarget{FUN-03-E01}{%
\subsection{FUN-03-E01 · Tích ánh xạ, kể cả hàm từng
khúc}\label{FUN-03-E01}}

\textbf{D2 · V0 · Đề {[}SV{]} S1-p96 · Lời giải {[}CC{]}}

\begin{itemize}
\tightlist
\item
  \textbf{Đề.} Tìm \(g\circ f\) và \(f\circ g\)
  (\(\mathbb R\to\mathbb R\)): (a) \(f(x)=x^2+1\), \(g(x)=x+1\); (b)
  \(f(x)=\begin{cases}x^2&x>0\\x+1&x\le0\end{cases}\), \(g(x)=2x+1\).
\item
  \textbf{Nhận diện.} ``tìm \(g\circ f\), \(f\circ g\)'' →
  \protect\hyperlink{FUN-03}{FUN-03}.
\item
  \textbf{Kiến thức.} \protect\hyperlink{F-FUN-03}{F-FUN-03}:
  \((g\circ f)(x)=g(f(x))\).
\item
  \textbf{Phương pháp.} Thay biểu thức hàm trong vào biến của hàm ngoài;
  với hàm ngoài từng khúc, đặt điều kiện lên \textbf{giá trị hàm trong}.
\item
  \textbf{Lời giải.} (a) \(g\circ f(x)=(x^2+1)+1=x^2+2\);
  \(\ f\circ g(x)=(x+1)^2+1=x^2+2x+2\).

  \begin{enumerate}
  \def\labelenumi{(\alph{enumi})}
  \setcounter{enumi}{1}
  \tightlist
  \item
    \(g\circ f(x)=2f(x)+1=\begin{cases}2x^2+1&x>0\\2x+3&x\le0\end{cases}\).
    \(\ f\circ g(x)=f(2x+1)\): điều kiện
    \(2x+1>0\Leftrightarrow x>-\tfrac12\) ⇒
    \(f\circ g(x)=\begin{cases}(2x+1)^2&x>-\tfrac12\\2x+2&x\le-\tfrac12\end{cases}\).
  \end{enumerate}
\item
  \textbf{Kết luận.} Như trên; \(g\circ f\ne f\circ g\).
\item
  \textbf{Kiểm tra.} Thử \(x=-0.25\): \(g(-0.25)=0.5>0\) ⇒
  \(f(0.5)=0.25\); công thức: \((2\cdot(-0.25)+1)^2=0.25\) ✓. \(x=-1\):
  \(g(-1)=-1\), \(f(-1)=0\); công thức \(2(-1)+2=0\) ✓. Kiểm tra trên
  lưới {[}CC{]}.
\item
  \textbf{Biến thể (V1).} \(f(x)=2x-3\), \(g(x)=x^2\):
  \(g\circ f=(2x-3)^2\), \(f\circ g=2x^2-3\).
\item
  \textbf{Lỗi thường gặp.} Ở \(f\circ g\) giữ điều kiện ``\(x>0\)'' thay
  vì ``\(g(x)>0\)''.
\item
  \textbf{Liên kết.} \protect\hyperlink{F-FUN-03}{F-FUN-03}.
\end{itemize}

\hypertarget{IND-01-E01}{%
\subsection{IND-01-E01 · Quy nạp: tổng số lẻ}\label{IND-01-E01}}

\textbf{D1 · V0 · Đề {[}SV{]} S1-p97--99 · Lời giải {[}SV{]}, {[}CC{]}}

\begin{itemize}
\tightlist
\item
  \textbf{Đề.} Chứng minh \(1+3+5+\dots+(2n-1)=n^2\) với mọi \(n\ge1\).
\item
  \textbf{Nhận diện.} ``chứng minh với mọi \(n\)'' + tổng →
  \protect\hyperlink{IND-01}{IND-01}.
\item
  \textbf{Kiến thức.} \protect\hyperlink{F-IND-01}{F-IND-01}.
  \textbf{Phương pháp.} \protect\hyperlink{M-10}{M-10}.
\item
  \textbf{Lời giải.} \(P(n)\): ``\(1+3+\dots+(2n-1)=n^2\)''. Cơ sở:
  \(P(1)\): \(1=1^2\) ✓. Quy nạp: giả sử \(P(k)\):
  \(1+3+\dots+(2k-1)=k^2\). Khi đó
  \(1+3+\dots+(2k-1)+(2k+1)=k^2+2k+1=(k+1)^2\), tức \(P(k+1)\) đúng.
\item
  \textbf{Kết luận.} Theo nguyên lý quy nạp, \(P(n)\) đúng
  \(\forall n\ge1\).
\item
  \textbf{Kiểm tra.} \(n=4\): \(1+3+5+7=16=4^2\); máy kiểm \(n\le200\)
  {[}CC{]}.
\item
  \textbf{Biến thể (V1).} \(2+4+\dots+2n=n(n+1)\).
\item
  \textbf{Lỗi thường gặp.} Viết số hạng thứ \(k+1\) là \(2k-1\) thay vì
  \(2(k+1)-1=2k+1\) --- \protect\hyperlink{T-09}{T-09}.
\item
  \textbf{Liên kết.} \protect\hyperlink{IND-01-E02}{IND-01-E02}.
\end{itemize}

\hypertarget{IND-01-E02}{%
\subsection{\texorpdfstring{IND-01-E02 · Quy nạp: tổng
\(1+2+\dots+n\)}{IND-01-E02 · Quy nạp: tổng 1+2+\textbackslash dots+n}}\label{IND-01-E02}}

\textbf{D1 · V1 · Đề {[}SV{]} S1-p100 · Lời giải {[}CC{]}}

\begin{itemize}
\tightlist
\item
  \textbf{Đề.} Chứng minh \(1+2+\dots+(n-1)+n=\dfrac{n(n+1)}{2}\),
  \(\forall n\ge1\).
\item
  \textbf{Nhận diện / Kiến thức / Phương pháp.}
  \protect\hyperlink{IND-01}{IND-01},
  \protect\hyperlink{F-IND-01}{F-IND-01},
  \protect\hyperlink{M-10}{M-10}.
\item
  \textbf{Lời giải.} Cơ sở: \(n=1\): \(1=\frac{1\cdot2}{2}\) ✓. Giả sử
  \(1+\dots+k=\frac{k(k+1)}2\). Khi đó
  \(1+\dots+k+(k+1)=\frac{k(k+1)}2+(k+1)=\frac{(k+1)(k+2)}{2}\) --- đúng
  dạng \(P(k+1)\).
\item
  \textbf{Kết luận.} Đúng \(\forall n\ge1\).
\item
  \textbf{Kiểm tra.} \(n=10\): \(55=\frac{10\cdot11}{2}\) ✓.
\item
  \textbf{Biến thể (V3).} \(1^2+2^2+\dots+n^2=\frac{n(n+1)(2n+1)}{6}\);
  bước quy nạp:
  \(\frac{k(k+1)(2k+1)}6+(k+1)^2=\frac{(k+1)(k+2)(2k+3)}{6}\).
\item
  \textbf{Lỗi thường gặp.} Chỉ kiểm tra vài giá trị rồi kết luận.
\item
  \textbf{Liên kết.} \protect\hyperlink{IND-01-E01}{IND-01-E01}.
\end{itemize}

\hypertarget{IND-01-E03}{%
\subsection{\texorpdfstring{IND-01-E03 · Quy nạp bất đẳng thức, bước cơ
sở
\(n_0=4\)}{IND-01-E03 · Quy nạp bất đẳng thức, bước cơ sở n\_0=4}}\label{IND-01-E03}}

\textbf{D2 · V4 · Đề {[}SV{]} S1-p100 · Lời giải {[}CC{]}}

\begin{itemize}
\tightlist
\item
  \textbf{Đề.} Chứng minh \(\forall n>3,\ 2^n<n!\).
\item
  \textbf{Nhận diện.} ``\(n>3\)'' ⇒ \(N_0=4\) →
  \protect\hyperlink{IND-01}{IND-01}.
\item
  \textbf{Kiến thức.} \protect\hyperlink{F-IND-01}{F-IND-01}.
  \textbf{Phương pháp.} \protect\hyperlink{M-10}{M-10}.
\item
  \textbf{Lời giải.} Cơ sở \(n=4\): \(2^4=16<24=4!\) ✓. Giả sử
  \(2^k<k!\) với \(k\ge4\). Khi đó
  \(2^{k+1}=2\cdot2^k<2\cdot k!<(k+1)\cdot k!=(k+1)!\) (vì
  \(k+1\ge5>2\)).
\item
  \textbf{Kết luận.} \(2^n<n!\) với mọi \(n\ge4\).
\item
  \textbf{Kiểm tra.} \(n=3\): \(8<6\) \textbf{sai} --- đúng là phải bắt
  đầu từ 4; máy kiểm \(4\le n\le60\) {[}CC{]}.
\item
  \textbf{Biến thể (V1).} \(\forall n\ge5,\ 2^n>n^2\): cơ sở \(32>25\);
  bước: \(2^{k+1}>2k^2\ge(k+1)^2\) vì \(k^2-2k-1>0\) khi \(k\ge3\).
\item
  \textbf{Lỗi thường gặp.} Lấy cơ sở \(n=1\) (\(2<1\) sai) rồi bế tắc;
  dùng bất đẳng thức ``\(2<k+1\)'' mà không nêu điều kiện \(k\ge4\) ---
  \protect\hyperlink{T-09}{T-09}.
\item
  \textbf{Liên kết.} \protect\hyperlink{M-10}{M-10}.
\end{itemize}

\hypertarget{IND-01-E04}{%
\subsection{IND-01-E04 · Chia hết: quy nạp và chứng minh theo trường
hợp}\label{IND-01-E04}}

\textbf{D2 · V3 · Đề {[}SV{]} S1-p47 · Lời giải {[}CC{]}}

\begin{itemize}
\tightlist
\item
  \textbf{Đề.} Chứng minh \((n^3-4n)\ \vdots\ 3\) với mọi
  \(n\in\mathbb N\).
\item
  \textbf{Nhận diện.} Slide đặt bài này ở \textbf{quy tắc chứng minh
  theo trường hợp}; cũng giải được bằng quy nạp →
  \protect\hyperlink{IND-01}{IND-01},
  \protect\hyperlink{F-LOG-04}{F-LOG-04} (8).
\item
  \textbf{Kiến thức.} \protect\hyperlink{F-IND-01}{F-IND-01}; chứng minh
  theo trường hợp.
\item
  \textbf{Phương pháp.} Cách 1 theo trường hợp \(n\bmod3\); cách 2
  \protect\hyperlink{M-10}{M-10}.
\item
  \textbf{Lời giải.} Cách 1: \(n^3-4n=n(n^2-4)=n(n-2)(n+2)\). Nếu
  \(n\equiv0\pmod3\): \(3\mid n\). Nếu \(n\equiv1\): \(n+2\equiv0\). Nếu
  \(n\equiv2\): \(n-2\equiv0\). Mọi trường hợp đều có một thừa số chia
  hết cho 3. Cách 2: \(n=0\): \(0\ \vdots\ 3\). Giả sử \(3\mid k^3-4k\).
  \((k+1)^3-4(k+1)=(k^3-4k)+3(k^2+k-1)\) chia hết cho 3.
\item
  \textbf{Kết luận.} \(3\mid n^3-4n\) với mọi \(n\in\mathbb N\).
\item
  \textbf{Kiểm tra.} \(n=5\): \(125-20=105=3\cdot35\) ✓; máy kiểm
  \(n\le300\) {[}CC{]}.
\item
  \textbf{Biến thể (V1).} \(6\mid n^3-n\): \(n^3-n=(n-1)n(n+1)\) là tích
  3 số liên tiếp (chia hết cho 2 và 3).
\item
  \textbf{Lỗi thường gặp.} Quên trường hợp \(n\equiv2\).
\item
  \textbf{Liên kết.} \protect\hyperlink{IND-01-E01}{IND-01-E01}.
\end{itemize}

\hypertarget{ex-ch2a}{%
\chapter{Ví dụ Chương II (a) --- Phép đếm}\label{ex-ch2a}}

\hypertarget{CNT-01-E01}{%
\subsection{CNT-01-E01 · Đếm số chẵn có 3 chữ số khác nhau (chia trường
hợp chữ số 0)}\label{CNT-01-E01}}

\textbf{D2 · V0 · Đề {[}SV{]} S2-p5 · Lời giải {[}SV{]}, {[}CC{]}}

\begin{itemize}
\tightlist
\item
  \textbf{Đề.} Cho \(X=\{0,1,2,3,4,5\}\). Có bao nhiêu số tự nhiên có 3
  chữ số khác nhau lấy từ \(X\) và chia hết cho 2?
\item
  \textbf{Nhận diện.} ``bao nhiêu số \ldots{} chữ số khác nhau \ldots{}
  chia hết cho 2'', tập có số 0 → \protect\hyperlink{CNT-01}{CNT-01}.
\item
  \textbf{Kiến thức.} Nguyên lý cộng, nhân
  (\protect\hyperlink{F-CNT-01}{F-CNT-01}).
\item
  \textbf{Phương pháp.} \protect\hyperlink{M-11}{M-11}: chọn chữ số cuối
  \(c\) trước; tách trường hợp \(c=0\) / \(c\ne0\) vì 0 ảnh hưởng chữ số
  đầu.
\item
  \textbf{Lời giải.} Số \(\overline{abc}\), \(c\in\{0,2,4\}\),
  \(a\ne0\). TH1 \(c=0\): \(c\): 1 cách; \(a\in X\setminus\{0\}\): 5
  cách; \(b\in X\setminus\{a,0\}\): 4 cách ⇒ \(1\cdot5\cdot4=20\). TH2
  \(c\in\{2,4\}\): 2 cách; \(a\in X\setminus\{c,0\}\): 4 cách;
  \(b\in X\setminus\{a,c\}\): 4 cách ⇒ \(2\cdot4\cdot4=32\).
\item
  \textbf{Kết luận.} \(20+32=52\) số.
\item
  \textbf{Kiểm tra.} Cách bù: số 3 chữ số khác nhau:
  \(5\cdot5\cdot4=100\); số lẻ (\(c\in\{1,3,5\}\)):
  \(3\cdot4\cdot4=48\); \(100-48=52\) ✓. Liệt kê bằng máy {[}CC{]}.
\item
  \textbf{Biến thể (V1).} Cùng \(X\), số 3 chữ số khác nhau \textbf{chia
  hết cho 5}: \(c=0\): 20; \(c=5\): \(a\in X\setminus\{0,5\}\) 4 cách,
  \(b\) 4 cách: 16 ⇒ \textbf{36}.
\item
  \textbf{Lỗi thường gặp.} Tính \(3\cdot5\cdot4\) (không tách trường hợp
  0) --- \protect\hyperlink{T-10}{T-10}.
\item
  \textbf{Liên kết.} \protect\hyperlink{CNT-01-E02}{CNT-01-E02}.
\end{itemize}

\hypertarget{CNT-01-E02}{%
\subsection{CNT-01-E02 · Số 4 chữ số khác nhau chia hết cho
5}\label{CNT-01-E02}}

\textbf{D2 · V1 · {[}AI-GENERATED PRACTICE{]} · Lời giải {[}CC{]}}

\begin{itemize}
\tightlist
\item
  \textbf{Đề.} Từ \(X=\{0,1,2,3,4,5\}\) lập được bao nhiêu số có 4 chữ
  số khác nhau chia hết cho 5?
\item
  \textbf{Nhận diện / Kiến thức / Phương pháp.}
  \protect\hyperlink{CNT-01}{CNT-01},
  \protect\hyperlink{F-CNT-01}{F-CNT-01},
  \protect\hyperlink{M-11}{M-11}.
\item
  \textbf{Lời giải.} \(\overline{abcd}\), \(d\in\{0,5\}\). \(d=0\):
  \(a\): 5, \(b\): 4, \(c\): 3 ⇒ 60. \(d=5\):
  \(a\in X\setminus\{0,5\}\): 4, \(b\): 4, \(c\): 3 ⇒ 48.
\item
  \textbf{Kết luận.} \(60+48=108\) số.
\item
  \textbf{Kiểm tra.} Liệt kê bằng máy {[}CC{]}.
\item
  \textbf{Biến thể (V2).} ``Có bao nhiêu số \textbf{lẻ} có 4 chữ số khác
  nhau'': \(d\in\{1,3,5\}\): 3; \(a\): 4 (≠0, ≠d); \(b\): 4; \(c\): 3 ⇒
  \(144\).
\item
  \textbf{Lỗi thường gặp.} Chọn \(a\) trước \(d\) rồi không biết \(d\)
  còn mấy cách.
\item
  \textbf{Liên kết.} \protect\hyperlink{CNT-01-E01}{CNT-01-E01}.
\end{itemize}

\hypertarget{CNT-02-E01}{%
\subsection{CNT-02-E01 · Nguyên lý chuồng bồ câu}\label{CNT-02-E01}}

\textbf{D2 · V0 · Đề {[}SV{]} S2-p6--7 · Lời giải {[}SV{]}, {[}CC{]}}

\begin{itemize}
\tightlist
\item
  \textbf{Đề.} (a) 20 bồ câu trong 7 chuồng: chứng tỏ có chuồng chứa
  \(\ge3\) con. (b) Lấy \(A\subset X=\{1,\dots,9\}\), \(|A|=6\). Chứng
  minh trong \(A\) có hai phần tử có tổng bằng 10.
\item
  \textbf{Nhận diện.} ``chứng minh tồn tại / luôn có'' →
  \protect\hyperlink{CNT-02}{CNT-02}.
\item
  \textbf{Kiến thức.} \protect\hyperlink{F-CNT-02}{F-CNT-02}.
\item
  \textbf{Phương pháp.} \protect\hyperlink{M-12}{M-12}.
\item
  \textbf{Lời giải.} (a) \(\lceil20/7\rceil=3\). (b) Lập 5 ``chuồng'':
  \(\{1,9\},\{2,8\},\{3,7\},\{4,6\},\{5\}\) (hai số cùng chuồng có tổng
  10, trừ chuồng \(\{5\}\) chỉ có 1 phần tử). 6 phần tử của \(A\) vào 5
  chuồng ⇒ có chuồng chứa 2 phần tử; chuồng đó không thể là \(\{5\}\) ⇒
  đó là một cặp tổng 10.
\item
  \textbf{Kết luận.} Điều phải chứng minh.
\item
  \textbf{Kiểm tra.} Vét cạn cả \(C_9^6=84\) tập con {[}CC{]}. Con số 6
  là tối ưu: \(A=\{1,2,3,4,5\}\) (5 phần tử) không có cặp tổng 10.
\item
  \textbf{Biến thể (V2).} Trong 367 người luôn có 2 người cùng ngày sinh
  (tối đa 366 ngày): \(\lceil367/366\rceil=2\).
\item
  \textbf{Lỗi thường gặp.} Chọn chuồng sao cho hai chim cùng chuồng
  không suy ra được điều cần chứng minh.
\item
  \textbf{Liên kết.} \protect\hyperlink{M-12}{M-12}.
\end{itemize}

\hypertarget{CNT-03-E01}{%
\subsection{CNT-03-E01 · Bù trừ hai tập}\label{CNT-03-E01}}

\textbf{D1 · V0 · Đề {[}SV{]} S2-p10 · Lời giải {[}SV{]}}

\begin{itemize}
\tightlist
\item
  \textbf{Đề.} Lớp ngoại ngữ: 24 HS học tiếng Pháp, 26 học tiếng Anh, 15
  học cả hai. Mọi HS học ít nhất một môn. Lớp có bao nhiêu người?
\item
  \textbf{Nhận diện.} ``học A hoặc B, cả hai'' →
  \protect\hyperlink{CNT-03}{CNT-03}.
\item
  \textbf{Kiến thức / Phương pháp.}
  \protect\hyperlink{F-CNT-03}{F-CNT-03},
  \protect\hyperlink{M-13}{M-13}.
\item
  \textbf{Lời giải.} \(A\): học Pháp, \(B\): học Anh.
  \(|A\cup B|=24+26-15=35\).
\item
  \textbf{Kết luận.} 35 học sinh.
\item
  \textbf{Kiểm tra.} Chỉ Pháp \(24-15=9\); chỉ Anh \(26-15=11\); cả hai
  15: \(9+11+15=35\) ✓.
\item
  \textbf{Biến thể (V2).} Nếu lớp có 40 HS: số HS không học môn nào
  \(=40-35=5\).
\item
  \textbf{Lỗi thường gặp.} Cộng \(24+26=50\) (đếm hai lần nhóm học cả
  hai).
\item
  \textbf{Liên kết.} \protect\hyperlink{CNT-03-E02}{CNT-03-E02}.
\end{itemize}

\hypertarget{CNT-03-E02}{%
\subsection{CNT-03-E02 · Bù trừ ba tập: chia hết cho 2, 3 hoặc
5}\label{CNT-03-E02}}

\textbf{D2 · V3 · {[}AI-GENERATED PRACTICE{]} · Lời giải {[}CC{]}}

\begin{itemize}
\tightlist
\item
  \textbf{Đề.} Có bao nhiêu số nguyên trong \(\{1,\dots,1000\}\) chia
  hết cho ít nhất một trong các số 2, 3, 5?
\item
  \textbf{Nhận diện.} ``ít nhất một trong'' ba tính chất →
  \protect\hyperlink{CNT-03}{CNT-03}.
\item
  \textbf{Kiến thức.} \protect\hyperlink{F-CNT-03}{F-CNT-03} (3 tập); số
  bội của \(d\) trong \(1..N\) là \(\lfloor N/d\rfloor\); chia hết cho 2
  và 3 ⇔ chia hết cho 6.
\item
  \textbf{Phương pháp.} \protect\hyperlink{M-13}{M-13}.
\item
  \textbf{Lời giải.} \(|A_2|=500\), \(|A_3|=333\), \(|A_5|=200\);
  \(|A_2\cap A_3|=|A_6|=166\), \(|A_2\cap A_5|=|A_{10}|=100\),
  \(|A_3\cap A_5|=|A_{15}|=66\); \(|A_2\cap A_3\cap A_5|=|A_{30}|=33\).
  \(|A_2\cup A_3\cup A_5|=500+333+200-166-100-66+33=734\).
\item
  \textbf{Kết luận.} 734 số.
\item
  \textbf{Kiểm tra.} Đếm bằng máy {[}CC{]}. Số không chia hết cho 2, 3,
  5: \(1000-734=266\).
\item
  \textbf{Biến thể (V2).} Số trong \(1..100\) không chia hết cho 2 và
  không chia hết cho 3: \(100-(50+33-16)=33\).
\item
  \textbf{Lỗi thường gặp.} Quên cộng lại giao ba tập; dùng
  \(\lceil\ \rceil\) thay cho \(\lfloor\ \rfloor\).
\item
  \textbf{Liên kết.} \protect\hyperlink{M-13}{M-13}.
\end{itemize}

\hypertarget{CNT-04-E01}{%
\subsection{CNT-04-E01 · Hoán vị, chỉnh hợp hay tổ
hợp?}\label{CNT-04-E01}}

\textbf{D1 · V4 · Đề {[}SV{]} S2-p12, p14, p16 · Lời giải
{[}SV{]}+{[}CC{]}}

\begin{itemize}
\tightlist
\item
  \textbf{Đề.} (a) Có bao nhiêu số gồm 5 chữ số khác nhau tạo từ
  \(\{1,2,3,4,5\}\)? (b) Có bao nhiêu số tự nhiên gồm 3 chữ số tạo thành
  từ \(1,2,3,4,5,6\)? (c) Lớp 30 HS, có bao nhiêu cách chọn 10 bạn? (d)
  \(|A|=n\): có bao nhiêu song ánh \(A\to A\)?
\item
  \textbf{Nhận diện.} Thứ tự? lặp? → \protect\hyperlink{CNT-04}{CNT-04}.
\item
  \textbf{Kiến thức.} \protect\hyperlink{F-CNT-04}{F-CNT-04}.
  \textbf{Phương pháp.} \protect\hyperlink{M-14}{M-14}.
\item
  \textbf{Lời giải.} (a) Có thứ tự, không lặp, lấy hết: \(P_5=5!=120\).
  (b) Slide cho đáp số \(A_6^3=\frac{6!}{3!}=120\) --- \textbf{ngầm hiểu
  các chữ số khác nhau}. Nếu đề cho phép lặp chữ số: \(6^3=216\). (c)
  Không thứ tự: \(C_{30}^{10}=30\,045\,015\). (d) Mỗi song ánh là một
  hoán vị của \(A\): \(n!\).
\item
  \textbf{Kết luận.} 120; 120 (khác nhau) hoặc 216 (cho lặp); 30 045
  015; \(n!\).
\item
  \textbf{Kiểm tra.} \(C_{30}^{10}=C_{30}^{20}\); tính bằng máy
  {[}CC{]}.
\item
  \textbf{Biến thể (V2).} Chọn 1 lớp trưởng, 1 lớp phó, 1 thủ quỹ từ 30
  HS: có thứ tự (vai trò khác nhau) ⇒ \(A_{30}^3=24\,360\).
\item
  \textbf{Lỗi thường gặp.} Đề (b) không ghi ``khác nhau'' --- phải nêu
  giả thiết rõ --- \protect\hyperlink{T-11}{T-11}.
\item
  \textbf{Liên kết.} \protect\hyperlink{CNT-05}{CNT-05},
  \protect\hyperlink{CNT-06}{CNT-06}.
\end{itemize}

\hypertarget{CNT-05-E01}{%
\subsection{CNT-05-E01 · Hoán vị lặp --- sắp xếp chữ
cái}\label{CNT-05-E01}}

\textbf{D1 · V0 · Đề {[}SV{]} S2-p18 · Lời giải {[}SV{]}, {[}CC{]}}

\begin{itemize}
\tightlist
\item
  \textbf{Đề.} Có bao nhiêu chuỗi ký tự khác nhau tạo bằng cách sắp xếp
  các chữ cái của từ SUCCESS?
\item
  \textbf{Nhận diện.} Sắp xếp có chữ lặp →
  \protect\hyperlink{CNT-05}{CNT-05}.
\item
  \textbf{Kiến thức.} \protect\hyperlink{F-CNT-05}{F-CNT-05}.
\item
  \textbf{Phương pháp.} \protect\hyperlink{M-14}{M-14}: đếm số lần lặp
  của từng chữ rồi áp công thức hoán vị lặp.
\item
  \textbf{Lời giải.} SUCCESS: S×3, U×1, C×2, E×1; \(n=7\). Số chuỗi
  \(=\frac{7!}{3!\,1!\,2!\,1!}=\frac{5040}{12}=420\).
\item
  \textbf{Kết luận.} 420 chuỗi.
\item
  \textbf{Kiểm tra.} Đếm tập hoán vị khác nhau bằng máy {[}CC{]}.
\item
  \textbf{Biến thể (V1).} MISSISSIPPI: M1, I4, S4, P2:
  \(\frac{11!}{1!4!4!2!}=34\,650\).
\item
  \textbf{Lỗi thường gặp.} Tính \(7!\) (coi các chữ S là khác nhau).
\item
  \textbf{Liên kết.} \protect\hyperlink{M-14}{M-14}.
\end{itemize}

\hypertarget{CNT-06-E01}{%
\subsection{CNT-06-E01 · Tổ hợp lặp cơ bản}\label{CNT-06-E01}}

\textbf{D1 · V0 · Đề {[}SV{]} S2-p20 · Lời giải {[}SV{]}}

\begin{itemize}
\tightlist
\item
  \textbf{Đề.} Có 3 loại nón A, B, C. An mua 2 cái. Có bao nhiêu cách
  chọn?
\item
  \textbf{Nhận diện.} Chọn từ các \textbf{loại}, được lặp, không thứ tự
  → \protect\hyperlink{CNT-06}{CNT-06}.
\item
  \textbf{Kiến thức.} \protect\hyperlink{F-CNT-06}{F-CNT-06}:
  \(K_n^k=C_{n+k-1}^k\) với \(n=3\) loại, \(k=2\) vật.
\item
  \textbf{Phương pháp.} \protect\hyperlink{M-14}{M-14}: không thứ tự +
  có lặp ⇒ tổ hợp lặp.
\item
  \textbf{Lời giải.} \(K_3^2=C_4^2=6\): AA, AB, AC, BB, BC, CC.
\item
  \textbf{Kết luận.} 6 cách.
\item
  \textbf{Kiểm tra.} Liệt kê đủ 6 ✓.
\item
  \textbf{Biến thể (V1).} Mua 5 cái từ 3 loại: \(K_3^5=C_7^5=21\).
\item
  \textbf{Lỗi thường gặp.} Đảo \(n,k\): \(C_{n+k-1}^{k}\) với \(n\) = số
  \textbf{loại} --- \protect\hyperlink{T-12}{T-12}.
\item
  \textbf{Liên kết.} \protect\hyperlink{CNT-06-E02}{CNT-06-E02}.
\end{itemize}

\hypertarget{CNT-06-E02}{%
\subsection{CNT-06-E02 · Nghiệm nguyên không âm với ràng buộc trên và
dưới}\label{CNT-06-E02}}

\textbf{D3 · V0 · Đề {[}SV{]} S2-p22 · Lời giải {[}SV{]} S2-p23--24,
{[}CC{]}}

\begin{itemize}
\tightlist
\item
  \textbf{Đề.} Tìm số nghiệm nguyên không âm của \(x_1+x_2+x_3+x_4=20\)
  (1) thỏa \(x_1\le3\); \(x_2\ge2\); \(x_3>4\) (∗).
\item
  \textbf{Nhận diện.} ``số nghiệm nguyên không âm'' + ``≤, ≥,
  \textgreater{}'' → \protect\hyperlink{CNT-06}{CNT-06}.
\item
  \textbf{Kiến thức.} \protect\hyperlink{F-CNT-06}{F-CNT-06}.
  \textbf{Phương pháp.} \protect\hyperlink{M-15}{M-15},
  \protect\hyperlink{M-13}{M-13}.
\item
  \textbf{Lời giải.} Viết (∗) thành \(x_1\le3,\ x_2\ge2,\ x_3\ge5\). Xét
  (∗∗) \(x_2\ge2,x_3\ge5\) và (∗∗∗) \(x_1\ge4,x_2\ge2,x_3\ge5\). Gọi
  \(p,q,r\) là số nghiệm thỏa (∗), (∗∗), (∗∗∗): \(p=q-r\). \(q\): đặt
  \(x_2'=x_2-2\), \(x_3'=x_3-5\) ⇒ \(x_1+x_2'+x_3'+x_4=13\) ⇒
  \(q=K_4^{13}=C_{16}^{13}=560\). \(r\): thêm \(x_1'=x_1-4\) ⇒ tổng
  \(=9\) ⇒ \(r=K_4^9=C_{12}^9=220\).
\item
  \textbf{Kết luận.} \(p=560-220=340\) nghiệm.
\item
  \textbf{Kiểm tra.} Đếm vét cạn bằng máy: 340 {[}CC{]}.
\item
  \textbf{Biến thể (V3).} Thêm \(x_4\le6\): trừ thêm trường hợp
  \(x_4\ge7\) và cộng lại giao --- bù trừ 2 tập:
  \(p'=C_{16}^{3}-C_{12}^{3}-C_{9}^{3}+C_{5}^{3}=560-220-84+10=266\)
  \textbf{{[}CC{]}}.
\item
  \textbf{Lỗi thường gặp.} Đổi \(x_3>4\) thành \(x_3\ge4\); lấy
  \(C_{16}^{4}\) (sai chỉ số trên) --- \protect\hyperlink{T-12}{T-12}.
\item
  \textbf{Liên kết.} \protect\hyperlink{CNT-06-E03}{CNT-06-E03}.
\end{itemize}

\hypertarget{CNT-06-E03}{%
\subsection{CNT-06-E03 · Chia vật giống nhau có cận trên và cận
dưới}\label{CNT-06-E03}}

\textbf{D2 · V2 · Đề {[}SV{]} S2-p25 (slide không có đáp án) · Lời giải
{[}CC{]}}

\begin{itemize}
\tightlist
\item
  \textbf{Đề.} Có 12 trái táo chia cho 3 bạn A, B, C. A lấy ít nhất 4
  trái, B và C lấy ít nhất 2 trái, C không lấy quá 5 trái. Có bao nhiêu
  cách chia?
\item
  \textbf{Nhận diện.} Chia vật \textbf{giống nhau} cho người khác nhau,
  có ràng buộc → \protect\hyperlink{CNT-06}{CNT-06} (lời văn của bài
  nghiệm nguyên).
\item
  \textbf{Kiến thức.} \protect\hyperlink{F-CNT-06}{F-CNT-06}.
  \textbf{Phương pháp.} \protect\hyperlink{M-15}{M-15}.
\item
  \textbf{Lời giải.} \(a+b+c=12\), \(a\ge4\), \(b\ge2\), \(2\le c\le5\).
  Đặt \(a'=a-4\), \(b'=b-2\), \(c'=c-2\): \(a'+b'+c'=4\),
  \(0\le c'\le3\). Không cận trên: \(K_3^4=C_6^4=15\). Trường hợp
  \(c'\ge4\): \(c''=c'-4\), \(a'+b'+c''=0\): \(K_3^0=1\).
\item
  \textbf{Kết luận.} \(15-1=14\) cách.
\item
  \textbf{Kiểm tra.} Liệt kê: \(c=2\): \(a+b=10\), \(a\ge4,b\ge2\):
  \(a=4..8\) (5 cách); \(c=3\): 4; \(c=4\): 3; \(c=5\): 2 ⇒
  \(5+4+3+2=14\) ✓ {[}CC{]}.
\item
  \textbf{Biến thể (V1).} Bỏ ràng buộc ``C không quá 5'': 15 cách.
\item
  \textbf{Lỗi thường gặp.} Trừ cận trên trước khi đổi biến cận dưới
  (dùng \(c\ge6\) trong phương trình gốc nhưng quên các cận dưới khác).
\item
  \textbf{Liên kết.} \protect\hyperlink{CNT-06-E02}{CNT-06-E02}.
\end{itemize}

\hypertarget{ex-ch2b}{%
\chapter{Ví dụ Chương II (b) --- Hệ thức đệ qui}\label{ex-ch2b}}

\begin{quote}
Mọi nghiệm trong phần này đã được thế lại vào hệ thức với \(n\) từ 2 đến
25 và vào điều kiện đầu bằng
\texttt{10\_Verification/verify\_examples.py} \textbf{{[}CC{]}}.
\end{quote}

\hypertarget{REC-01-E01}{%
\subsection{REC-01-E01 · Lập hệ thức: cầu thang bước 1 hoặc 2
bậc}\label{REC-01-E01}}

\textbf{D2 · V0 · Đề {[}SV{]} S2-p34 · Lời giải {[}SV{]} S2-p34--36,
{[}CC{]}}

\begin{itemize}
\tightlist
\item
  \textbf{Đề.} Cầu thang có \(n\) bậc, mỗi bước đi 1 hoặc 2 bậc. Gọi
  \(x_n\) là số cách đi hết cầu thang. Tìm hệ thức đệ qui cho \(x_n\).
\item
  \textbf{Nhận diện.} ``tìm một hệ thức đệ qui cho \(x_n\)'' →
  \protect\hyperlink{REC-01}{REC-01}.
\item
  \textbf{Kiến thức.} Nguyên lý cộng
  (\protect\hyperlink{F-CNT-01}{F-CNT-01}).
\item
  \textbf{Phương pháp.} \protect\hyperlink{M-16}{M-16}: chia theo bước
  đầu tiên.
\item
  \textbf{Lời giải.} \(x_1=1\); \(x_2=2\) (1+1 hoặc 2). Với \(n>2\): TH1
  bước đầu 1 bậc → còn \(n-1\) bậc: \(x_{n-1}\) cách; TH2 bước đầu 2 bậc
  → còn \(n-2\) bậc: \(x_{n-2}\) cách. Hai trường hợp rời nhau ⇒
  \(x_n=x_{n-1}+x_{n-2}\).
\item
  \textbf{Kết luận.}
  \(\begin{cases}x_n-x_{n-1}-x_{n-2}=0\\ x_1=1,\ x_2=2\end{cases}\)
  (thuần nhất cấp 2). Giải: \protect\hyperlink{REC-02-E03}{REC-02-E03}.
\item
  \textbf{Kiểm tra.} \(n=3\): liệt kê 111, 12, 21 → 3 \(=x_2+x_1\) ✓;
  \(n=4\): 5 ✓.
\item
  \textbf{Biến thể (V1).} Mỗi bước 1, 2 hoặc 3 bậc:
  \(x_n=x_{n-1}+x_{n-2}+x_{n-3}\), \(x_1=1,x_2=2,x_3=4\) (cấp 3 --- lập
  được, \textbf{giải ngoài phạm vi slide}).
\item
  \textbf{Lỗi thường gặp.} Thiếu điều kiện đầu \(x_2\); trường hợp chồng
  nhau.
\item
  \textbf{Liên kết.} \protect\hyperlink{M-16}{M-16}.
\end{itemize}

\hypertarget{REC-01-E02}{%
\subsection{REC-01-E02 · Lập hệ thức: tháp Hà Nội}\label{REC-01-E02}}

\textbf{D2 · V0 · Đề {[}SV{]} S2-p27 · Lời giải {[}SV{]} S2-p28,
p38--39, {[}CC{]}}

\begin{itemize}
\tightlist
\item
  \textbf{Đề.} 3 cọc A, B, C và \(n\) đĩa khác cỡ ở cọc A (đĩa nhỏ trên
  đĩa lớn). Chuyển hết sang cọc C, mỗi lần một đĩa, không đặt đĩa lớn
  lên đĩa nhỏ. Gọi \(x_n\) là số lần chuyển. Tìm hệ thức đệ qui và
  \(x_n\).
\item
  \textbf{Nhận diện.} \protect\hyperlink{REC-01}{REC-01}.
\item
  \textbf{Kiến thức / Phương pháp.} \protect\hyperlink{M-16}{M-16},
  \protect\hyperlink{F-REC-03}{F-REC-03},
  \protect\hyperlink{F-IND-01}{F-IND-01}.
\item
  \textbf{Lời giải.} \(x_1=1\). Với \(n>1\): chuyển \(n-1\) đĩa trên
  sang B qua C (\(x_{n-1}\) lần), chuyển đĩa lớn nhất sang C (1 lần),
  chuyển \(n-1\) đĩa từ B sang C (\(x_{n-1}\) lần) ⇒ \(x_n=2x_{n-1}+1\),
  \(x_1=1\). Giải: PTĐT \(\lambda-2=0\), \(\lambda=2\); vế phải
  \(1=1^n\cdot1\), \(\beta=1\) không là nghiệm ⇒ nghiệm riêng \(A\):
  \(A=2A+1\Rightarrow A=-1\). Nghiệm TQ \(x_n=C2^n-1\);
  \(x_1=2C-1=1\Rightarrow C=1\).
\item
  \textbf{Kết luận.} \(x_n=2x_{n-1}+1\), \(x_1=1\); \(x_n=2^n-1\). Với
  64 đĩa, 1 giây/lần: \(2^{64}-1\approx1{,}8\times10^{19}\) giây.
\item
  \textbf{Kiểm tra (quy nạp --- S2-p39).} \(n=1\): \(2^1-1=1\) ✓. Giả sử
  \(x_k=2^k-1\): \(x_{k+1}=2(2^k-1)+1=2^{k+1}-1\) ✓.
\item
  \textbf{Biến thể (V2).} Gửi tiết kiệm: mỗi năm lãi 5\% rồi gửi thêm 1
  (triệu): \(x_n=1{,}05x_{n-1}+1\), \(x_0=a\) ⇒ nghiệm riêng
  \(A=1{,}05A+1\Rightarrow A=-20\); \(x_n=(a+20)\,1{,}05^n-20\).
\item
  \textbf{Lỗi thường gặp.} Viết \(x_n=2x_{n-1}\) (quên lần chuyển đĩa
  lớn).
\item
  \textbf{Liên kết.} \protect\hyperlink{REC-02}{REC-02}.
\end{itemize}

\hypertarget{REC-01-E03}{%
\subsection{REC-01-E03 · Lập hệ thức: chuỗi bit không có hai số 0 liên
tiếp}\label{REC-01-E03}}

\textbf{D2 · V1 · {[}AI-GENERATED PRACTICE{]} · Lời giải {[}CC{]}}

\begin{itemize}
\tightlist
\item
  \textbf{Đề.} Gọi \(x_n\) là số chuỗi bit độ dài \(n\) không chứa hai
  bit 0 liên tiếp. Tìm hệ thức đệ qui và điều kiện đầu.
\item
  \textbf{Nhận diện / Phương pháp.} \protect\hyperlink{REC-01}{REC-01},
  \protect\hyperlink{M-16}{M-16}: chia theo bit \textbf{cuối}.
\item
  \textbf{Kiến thức.} Nguyên lý cộng
  (\protect\hyperlink{F-CNT-01}{F-CNT-01}); hai trường hợp bit cuối rời
  nhau.
\item
  \textbf{Lời giải.} \(x_1=2\) (0, 1); \(x_2=3\) (01, 10, 11).
  \(n\ge3\): bit cuối là 1 → \(n-1\) bit đầu là chuỗi hợp lệ bất kỳ:
  \(x_{n-1}\); bit cuối là 0 → bit kế cuối buộc là 1, \(n-2\) bit đầu
  hợp lệ: \(x_{n-2}\).
\item
  \textbf{Kết luận.} \(x_n=x_{n-1}+x_{n-2}\), \(x_1=2\), \(x_2=3\).
\item
  \textbf{Kiểm tra.} \(n=3\): 010, 011, 101, 110, 111 → 5 \(=3+2\) ✓;
  vét cạn \(n\le15\) {[}CC{]}.
\item
  \textbf{Biến thể (V2).} ``không chứa hai bit 1 liên tiếp'' --- đổi vai
  0/1, cùng đáp số.
\item
  \textbf{Lỗi thường gặp.} Chia theo bit cuối nhưng quên ràng buộc ở bit
  kế cuối.
\item
  \textbf{Liên kết.} \protect\hyperlink{REC-01-E01}{REC-01-E01} (cùng hệ
  thức, khác điều kiện đầu).
\end{itemize}

\hypertarget{REC-02-E01}{%
\subsection{REC-02-E01 · Thuần nhất cấp 1}\label{REC-02-E01}}

\textbf{D1 · V0 · Đề {[}SV{]} S2-p42 · Lời giải {[}SV{]}}

\begin{itemize}
\tightlist
\item
  \textbf{Đề.} Giải \(2x_n-3x_{n-1}=0\), \(x_0=1\).
\item
  \textbf{Nhận diện.} Vế phải 0, cấp 1 →
  \protect\hyperlink{REC-02}{REC-02}.
\item
  \textbf{Kiến thức / Phương pháp.}
  \protect\hyperlink{F-REC-01}{F-REC-01},
  \protect\hyperlink{M-17}{M-17}.
\item
  \textbf{Lời giải.} PTĐT \(2\lambda-3=0\Rightarrow\lambda_0=\frac32\).
  Nghiệm TQ \(x_n=C\left(\frac32\right)^n\); \(x_0=C=1\).
\item
  \textbf{Kết luận.} \(x_n=\left(\frac32\right)^n\).
\item
  \textbf{Kiểm tra.} \(x_1=\frac32\); từ hệ thức
  \(2x_1=3x_0\Rightarrow x_1=\frac32\) ✓.
\item
  \textbf{Biến thể (V1).} \(x_n=-3x_{n-1}\), \(x_0=2\) ⇒
  \(x_n=2(-3)^n\).
\item
  \textbf{Lỗi thường gặp.} Viết \(\lambda=-\frac32\) (sai dấu khi chuyển
  vế).
\item
  \textbf{Liên kết.} \protect\hyperlink{REC-02-E02}{REC-02-E02}.
\end{itemize}

\hypertarget{REC-02-E02}{%
\subsection{REC-02-E02 · Thuần nhất cấp 2: nghiệm phân biệt và nghiệm
kép}\label{REC-02-E02}}

\textbf{D2 · V0 · Đề {[}SV{]} S2-p44 · Lời giải {[}SV{]} S2-p45--46 +
toolkit tìm hằng số, {[}CC{]}}

\begin{itemize}
\tightlist
\item
  \textbf{Đề.} Giải (a) \(2x_n-3x_{n-1}+x_{n-2}=0\); (b)
  \(4x_{n+1}-12x_n+9x_{n-1}=0\), \(x_0=2\), \(x_1=4\).
\item
  \textbf{Nhận diện.} Thuần nhất cấp 2 →
  \protect\hyperlink{REC-02}{REC-02}.
\item
  \textbf{Kiến thức.} \protect\hyperlink{F-REC-02}{F-REC-02}.
  \textbf{Phương pháp.} \protect\hyperlink{M-17}{M-17}.
\item
  \textbf{Lời giải.} (a) PTĐT
  \(2\lambda^2-3\lambda+1=0\Rightarrow\lambda_1=1,\ \lambda_2=\frac12\)
  ⇒ \(x_n=C_1+C_2\left(\frac12\right)^n\).

  \begin{enumerate}
  \def\labelenumi{(\alph{enumi})}
  \setcounter{enumi}{1}
  \tightlist
  \item
    PTĐT \(4\lambda^2-12\lambda+9=0\Leftrightarrow(2\lambda-3)^2=0\) ⇒
    nghiệm kép \(\lambda_0=\frac32\) ⇒
    \(x_n=(C_1+nC_2)\left(\frac32\right)^n\). \(x_0=C_1=2\);
    \(x_1=(2+C_2)\frac32=4\Rightarrow C_2=\frac23\). Vậy
    \(x_n=\left(2+\frac{2n}{3}\right)\left(\frac32\right)^n=(3+n)\left(\frac32\right)^{n-1}\).
  \end{enumerate}
\item
  \textbf{Kết luận.} (a) \(x_n=C_1+C_2\left(\frac12\right)^n\); (b)
  \(x_n=(3+n)\left(\frac32\right)^{n-1}\) (khớp đáp số S2-p44).
\item
  \textbf{Kiểm tra.} (b) \(x_2\) từ hệ thức:
  \(4x_2=12\cdot4-9\cdot2=30\Rightarrow x_2=7{,}5\); công thức:
  \(5\cdot\frac32=7{,}5\) ✓.
\item
  \textbf{Biến thể (V1).} \(x_n=5x_{n-1}-6x_{n-2}\), \(x_0=1,x_1=0\):
  \(\lambda=2,3\); \(C_1+C_2=1\), \(2C_1+3C_2=0\) ⇒
  \(x_n=3\cdot2^n-2\cdot3^n\).
\item
  \textbf{Lỗi thường gặp.} Với nghiệm kép viết
  \(C_1\lambda_0^n+C_2\lambda_0^n\) (thực chất chỉ một hằng số) ---
  \protect\hyperlink{T-15}{T-15}.
\item
  \textbf{Liên kết.} \protect\hyperlink{F-REC-02}{F-REC-02}.
\end{itemize}

\hypertarget{REC-02-E03}{%
\subsection{REC-02-E03 · Giải hệ thức cầu thang (nghiệm vô
tỉ)}\label{REC-02-E03}}

\textbf{D2 · V3 · Đề: nối tiếp
\protect\hyperlink{REC-01-E01}{REC-01-E01} {[}SV{]}; lời giải
{[}AI-GENERATED PRACTICE{]}, {[}CC{]}}

\begin{itemize}
\item
  \textbf{Đề.} Giải \(x_n=x_{n-1}+x_{n-2}\), \(x_1=1\), \(x_2=2\).
\item
  \textbf{Nhận diện / Phương pháp.} \protect\hyperlink{REC-02}{REC-02},
  \protect\hyperlink{M-17}{M-17}.
\item
  \textbf{Kiến thức.} \protect\hyperlink{F-REC-02}{F-REC-02} --- PTĐT có
  hai nghiệm thực phân biệt (vô tỉ).
\item
  \textbf{Lời giải.} PTĐT
  \(\lambda^2-\lambda-1=0\Rightarrow\lambda_{1,2}=\frac{1\pm\sqrt5}{2}\)
  (thực phân biệt).
  \(x_n=C_1\left(\frac{1+\sqrt5}{2}\right)^n+C_2\left(\frac{1-\sqrt5}{2}\right)^n\).
  Từ \(x_1=1,x_2=2\) (dùng \(\lambda^2=\lambda+1\)):
  \(C_1=\frac{5+\sqrt5}{10}\), \(C_2=\frac{5-\sqrt5}{10}\).
\item
  \textbf{Kết luận.} \[
  \begin{aligned}
  x_n&=\frac{5+\sqrt5}{10}\left(\frac{1+\sqrt5}{2}\right)^n+\frac{5-\sqrt5}{10}\left(\frac{1-\sqrt5}{2}\right)^n\\
  &=\frac{1}{\sqrt5}\left[\left(\frac{1+\sqrt5}{2}\right)^{n+1}-\left(\frac{1-\sqrt5}{2}\right)^{n+1}\right].
  \end{aligned}
  \]
\item
  \textbf{Kiểm tra.} \(n=1\):
  \(\frac{(5+\sqrt5)(1+\sqrt5)+(5-\sqrt5)(1-\sqrt5)}{20}=\frac{20}{20}=1\)
  ✓; so khớp dãy 1, 2, 3, 5, 8, \ldots{} tới \(n=25\) {[}CC{]}.
\item
  \textbf{Biến thể (V1).} Chuỗi bit
  \protect\hyperlink{REC-01-E03}{REC-01-E03} (\(x_1=2,x_2=3\)): cùng
  \(\lambda\), \(C_1=\frac{5+3\sqrt5}{10}\),
  \(C_2=\frac{5-3\sqrt5}{10}\).
\item
  \textbf{Lỗi thường gặp.} Áp điều kiện đầu như thể có \(x_0\) (đề cho
  \(x_1,x_2\)).
\item
  \textbf{Liên kết.} \protect\hyperlink{REC-01-E01}{REC-01-E01}.
\end{itemize}

\hypertarget{REC-03-E01}{%
\subsection{\texorpdfstring{REC-03-E01 · Vế phải đa thức, \(\beta=1\) là
nghiệm
đơn}{REC-03-E01 · Vế phải đa thức, \textbackslash beta=1 là nghiệm đơn}}\label{REC-03-E01}}

\textbf{D2 · V0 · Đề {[}SV{]} S2-p53 a) · Lời giải {[}SV{]} S2-p54--56,
{[}CC{]}}

\begin{itemize}
\tightlist
\item
  \textbf{Đề.} Giải \(2x_n-3x_{n-1}+x_{n-2}=4n+1\) (1).
\item
  \textbf{Nhận diện.} Vế phải đa thức bậc 1 ⇒ \(\beta=1\) →
  \protect\hyperlink{REC-03}{REC-03}.
\item
  \textbf{Kiến thức.} \protect\hyperlink{F-REC-03}{F-REC-03},
  \protect\hyperlink{F-REC-05}{F-REC-05}. \textbf{Phương pháp.}
  \protect\hyperlink{M-17}{M-17}.
\item
  \textbf{Lời giải.} Thuần nhất:
  \(2\lambda^2-3\lambda+1=0\Rightarrow\lambda=1,\frac12\) ⇒
  \(\bar x_n=C_1+C_2\left(\frac12\right)^n\). \(f_n=4n+1=1^n P_1(n)\),
  \(\beta=1\) là \textbf{nghiệm đơn} ⇒ \(x_n^*=n(an+b)\). Thế vào (1):
  \(2n(an+b)-3(n-1)[a(n-1)+b]+(n-2)[a(n-2)+b]=4n+1\). Cho \(n=0\):
  \(a+b=1\); \(n=1\): \(3a+b=5\) ⇒ \(a=2\), \(b=-1\) ⇒
  \(x_n^*=n(2n-1)\).
\item
  \textbf{Kết luận.} \(x_n=C_1+C_2\left(\frac12\right)^n+n(2n-1)\).
\item
  \textbf{Kiểm tra.} Thế \(x_n^*\):
  \(2n(2n-1)-3(n-1)(2n-3)+(n-2)(2n-5)=4n+1\) (khai triển: hệ số \(n^2\):
  \(4-6+2=0\); \(n\): \(-2+15-9=4\); hằng: \(-9+10=1\)) ✓.
\item
  \textbf{Biến thể (V1).} Thêm \(x_0=0,x_1=1\): \(C_1+C_2=0\),
  \(C_1+\frac{C_2}{2}+1=1\) ⇒ \(C_1=C_2=0\) ⇒ \(x_n=n(2n-1)\).
\item
  \textbf{Lỗi thường gặp.} Chọn \(x_n^*=an+b\) (quên nhân \(n\) vì
  \(\beta=1\) là nghiệm) --- hệ vô nghiệm ---
  \protect\hyperlink{T-13}{T-13}, \protect\hyperlink{T-15}{T-15}.
\item
  \textbf{Liên kết.} \protect\hyperlink{REC-03-E04}{REC-03-E04}.
\end{itemize}

\hypertarget{REC-03-E02}{%
\subsection{\texorpdfstring{REC-03-E02 · \(\beta\) là nghiệm kép, có
điều kiện
đầu}{REC-03-E02 · \textbackslash beta là nghiệm kép, có điều kiện đầu}}\label{REC-03-E02}}

\textbf{D3 · V4 · Đề {[}SV{]} S2-p53 b), S2-p57 · Lời giải {[}SV{]}
S2-p57--59 (ảnh), {[}CC{]}}

\begin{itemize}
\tightlist
\item
  \textbf{Đề.} Giải \(x_{n+1}-6x_n+9x_{n-1}=(18n+12)3^n\), \(x_0=2\),
  \(x_1=0\).
\item
  \textbf{Nhận diện.} \(f_n=3^n\cdot(18n+12)\), \(\beta=3\); PTĐT có
  nghiệm kép 3 → \protect\hyperlink{REC-03}{REC-03}.
\item
  \textbf{Kiến thức.} \protect\hyperlink{F-REC-03}{F-REC-03} (TH3).
  \textbf{Phương pháp.} \protect\hyperlink{M-17}{M-17}.
\item
  \textbf{Lời giải.} PTĐT \(\lambda^2-6\lambda+9=0\Rightarrow\lambda=3\)
  (kép) ⇒ \(\bar x_n=(C_1+nC_2)3^n\). \(\beta=3\) nghiệm kép ⇒
  \(x_n^*=n^2(an+b)3^n\). Thế:
  \((n+1)^2[a(n+1)+b]3^{n+1}-6n^2[an+b]3^n+9(n-1)^2[a(n-1)+b]3^{n-1}=(18n+12)3^n\).
  \(n=0\): \(6b=12\); \(n=1\): \(54a+18b=90\) ⇒ \(b=2\), \(a=1\) ⇒
  \(x_n^*=n^2(n+2)3^n\). Nghiệm TQ \(x_n=(C_1+nC_2)3^n+n^2(n+2)3^n\).
  \(x_0=C_1=2\); \(x_1=3(2+C_2)+9=0\Rightarrow C_2=-5\).
\item
  \textbf{Kết luận.} \(x_n=(n^3+2n^2-5n+2)\,3^n\).
\item
  \textbf{Kiểm tra.} \(x_2\) từ hệ thức (\(n=1\)):
  \(x_2=6\cdot0-9\cdot2+30\cdot3=72\); công thức:
  \((8+8-10+2)\cdot9=72\) ✓.
\item
  \textbf{Biến thể (V4).} Vế phải chỉ là \(3^n\): \(x_n^*=An^23^n\), thế
  ⇒ \(6A\cdot3^n=3^n\) ⇒ \(A=\frac16\), \(x_n^*=\frac{n^2 3^n}{6}\).
\item
  \textbf{Lỗi thường gặp.} Chọn \(x_n^*=(an+b)3^n\) hoặc \(n(an+b)3^n\)
  --- trùng nghiệm thuần nhất, hệ vô nghiệm ---
  \protect\hyperlink{T-15}{T-15}.
\item
  \textbf{Liên kết.} \protect\hyperlink{REC-03-E09}{REC-03-E09}.
\end{itemize}

\hypertarget{REC-03-E03}{%
\subsection{\texorpdfstring{REC-03-E03 · \(\beta\) không là nghiệm, đa
thức bậc
2}{REC-03-E03 · \textbackslash beta không là nghiệm, đa thức bậc 2}}\label{REC-03-E03}}

\textbf{D3 · V1 · Đề {[}SV{]} S2-p53 c), S2-p60 · Lời giải {[}SV{]}
S2-p60--63, {[}CC{]}}

\begin{itemize}
\tightlist
\item
  \textbf{Đề.} Giải \(4x_{n+1}-12x_n+9x_{n-1}=(2n^2+29n+56)2^{n-1}\),
  \(x_0=1\), \(x_1=-2\).
\item
  \textbf{Nhận diện.} \(\beta=2\), \(r=2\); PTĐT nghiệm kép \(\frac32\)
  → \protect\hyperlink{REC-03}{REC-03}.
\item
  \textbf{Kiến thức / Phương pháp.}
  \protect\hyperlink{F-REC-03}{F-REC-03} (TH1),
  \protect\hyperlink{M-17}{M-17}.
\item
  \textbf{Lời giải.} \(\bar x_n=(C_1+nC_2)\left(\frac32\right)^n\).
  \(\beta=2\) không là nghiệm ⇒ \(x_n^*=(an^2+bn+c)2^n\). Thế, chia hai
  vế cho \(2^{n-1}\):
  \(16[a(n+1)^2+b(n+1)+c]-24[an^2+bn+c]+9[a(n-1)^2+b(n-1)+c]=2n^2+29n+56\).
  Đồng nhất hệ số: \(n^2\): \(a=2\); \(n\): \(14a+b=29\Rightarrow b=1\);
  hằng: \(25a+7b+c=56\Rightarrow c=-1\). (Slide cho \(n=-1,0,1\) ---
  cùng kết quả.) ⇒ \(x_n^*=(2n^2+n-1)2^n\).
  \(x_n=(C_1+nC_2)\left(\frac32\right)^n+(2n^2+n-1)2^n\). \(x_0\):
  \(C_1-1=1\Rightarrow C_1=2\); \(x_1\):
  \(\frac32(2+C_2)+4=-2\Rightarrow C_2=-6\).
\item
  \textbf{Kết luận.} \(x_n=(2-6n)\left(\frac32\right)^n+(2n^2+n-1)2^n\).
\item
  \textbf{Kiểm tra.} \(x_2\): hệ thức (\(n=1\)):
  \(4x_2=12(-2)-9\cdot1+87\cdot1=54\Rightarrow x_2=13{,}5\); công thức:
  \((2-12)\cdot\frac94+9\cdot4=-22{,}5+36=13{,}5\) ✓.
\item
  \textbf{Biến thể (V4).} Vế phải \(\left(\frac32\right)^n\):
  \(\beta=\frac32\) nghiệm kép ⇒ \(x_n^*=An^2\left(\frac32\right)^n\),
  \(A=\frac1{12}\).
\item
  \textbf{Lỗi thường gặp.} Quên chia/nhân lũy thừa \(2^{n\pm1}\) khi
  thế; tính \(C_1,C_2\) từ nghiệm thuần nhất ---
  \protect\hyperlink{T-14}{T-14}.
\item
  \textbf{Liên kết.} \protect\hyperlink{REC-03-E02}{REC-03-E02}.
\end{itemize}

\hypertarget{REC-03-E04}{%
\subsection{\texorpdfstring{REC-03-E04 · Bài tập S2-p67 (1): \(\beta=1\)
nghiệm
đơn}{REC-03-E04 · Bài tập S2-p67 (1): \textbackslash beta=1 nghiệm đơn}}\label{REC-03-E04}}

\textbf{D2 · V1 · Đề {[}SV{]} S2-p67 (không có lời giải) · Lời giải
{[}CC{]}}

\begin{itemize}
\tightlist
\item
  \textbf{Đề.} \(x_n+4x_{n-1}-5x_{n-2}=12n+8\), \(x_0=0\), \(x_1=-5\).
\item
  \textbf{Nhận diện / Phương pháp.} \protect\hyperlink{REC-03}{REC-03},
  \protect\hyperlink{M-17}{M-17}.
\item
  \textbf{Kiến thức.} \protect\hyperlink{F-REC-03}{F-REC-03} (TH2:
  \(\beta=1\) nghiệm đơn), \protect\hyperlink{F-REC-05}{F-REC-05}.
\item
  \textbf{Lời giải.} PTĐT
  \(\lambda^2+4\lambda-5=0\Rightarrow\lambda=1,\ -5\) ⇒
  \(\bar x_n=C_1+C_2(-5)^n\). \(\beta=1\) nghiệm đơn ⇒
  \(x_n^*=n(an+b)\). Thế và khai triển: hệ số \(n^2\): \(0\); \(n\):
  \(12a=12\); hằng: \(-16a+6b=8\) ⇒ \(a=1,b=4\) ⇒ \(x_n^*=n^2+4n\).
  \(x_0\): \(C_1+C_2=0\); \(x_1\): \(C_1-5C_2+5=-5\) ⇒ \(C_2=\frac53\),
  \(C_1=-\frac53\).
\item
  \textbf{Kết luận.} \(x_n=-\dfrac53+\dfrac53(-5)^n+n^2+4n\).
\item
  \textbf{Kiểm tra.} \(x_2=-4x_1+5x_0+32=52\); công thức:
  \(-\frac53+\frac{125}{3}+12=52\) ✓.
\item
  \textbf{Biến thể (V1).} Vế phải hằng 6: \(x_n^*=An\),
  \(6A=6\Rightarrow x_n^*=n\).
\item
  \textbf{Lỗi thường gặp.} \protect\hyperlink{T-13}{T-13} (coi vế phải
  đa thức là ``không có \(\beta\)'').
\item
  \textbf{Liên kết.} \protect\hyperlink{REC-03-E01}{REC-03-E01}.
\end{itemize}

\hypertarget{REC-03-E05}{%
\subsection{\texorpdfstring{REC-03-E05 · Bài tập S2-p67 (2): nghiệm âm,
\(\beta=3\) không là
nghiệm}{REC-03-E05 · Bài tập S2-p67 (2): nghiệm âm, \textbackslash beta=3 không là nghiệm}}\label{REC-03-E05}}

\textbf{D2 · V1 · Đề {[}SV{]} S2-p67 · Lời giải {[}CC{]}}

\begin{itemize}
\tightlist
\item
  \textbf{Đề.} \(2x_{n+2}+5x_{n+1}+2x_n=(35n+51)3^n\), \(x_0=3\),
  \(x_1=0\).
\item
  \textbf{Nhận diện.} Vế phải \((35n+51)3^n\): \(\beta=3\), \(r=1\); hệ
  thức viết theo \(x_{n+2},x_{n+1},x_n\) →
  \protect\hyperlink{REC-03}{REC-03}.
\item
  \textbf{Kiến thức.} \protect\hyperlink{F-REC-02}{F-REC-02},
  \protect\hyperlink{F-REC-03}{F-REC-03} (TH1).
\item
  \textbf{Phương pháp.} \protect\hyperlink{M-17}{M-17}.
\item
  \textbf{Lời giải.} PTĐT
  \(2\lambda^2+5\lambda+2=0\Rightarrow\lambda=-2,\ -\frac12\) ⇒
  \(\bar x_n=C_1(-2)^n+C_2\left(-\frac12\right)^n\). \(\beta=3\) không
  là nghiệm ⇒ \(x_n^*=(an+b)3^n\). Thế, chia \(3^n\):
  \(18(an+2a+b)+15(an+a+b)+2(an+b)=35an+51a+35b\) ⇒ \(35a=35\),
  \(51a+35b=51\) ⇒ \(a=1,b=0\) ⇒ \(x_n^*=n3^n\). \(x_0\): \(C_1+C_2=3\);
  \(x_1\): \(-2C_1-\frac{C_2}{2}+3=0\) ⇒ \(C_1=1\), \(C_2=2\).
\item
  \textbf{Kết luận.} \(x_n=(-2)^n+2\left(-\dfrac12\right)^n+n\,3^n\).
\item
  \textbf{Kiểm tra.} \(n=0\) trong hệ thức:
  \(2x_2+0+6=51\Rightarrow x_2=22{,}5\); công thức:
  \(4+\frac12+18=22{,}5\) ✓.
\item
  \textbf{Biến thể (V4).} Vế phải \(9(-2)^n\): \(\beta=-2\) là nghiệm
  đơn ⇒ \(x_n^*=An(-2)^n\), \(6A=9\Rightarrow A=\frac32\).
\item
  \textbf{Lỗi thường gặp.} Sai dấu nghiệm PTĐT; quên rằng hệ thức viết
  theo \(x_{n+2}\) nên thế \(n+2,n+1,n\).
\item
  \textbf{Liên kết.} \protect\hyperlink{REC-03-E03}{REC-03-E03}.
\end{itemize}

\hypertarget{REC-03-E06}{%
\subsection{\texorpdfstring{REC-03-E06 · Bài tập S2-p67 (3): vế phải
hằng, \(\beta=1\) nghiệm
kép}{REC-03-E06 · Bài tập S2-p67 (3): vế phải hằng, \textbackslash beta=1 nghiệm kép}}\label{REC-03-E06}}

\textbf{D2 · V4 · Đề {[}SV{]} S2-p67 · Lời giải {[}CC{]}}

\begin{itemize}
\tightlist
\item
  \textbf{Đề.} \(x_{n+2}-2x_{n+1}+x_n=2\), \(x_0=1\), \(x_1=0\).
\item
  \textbf{Nhận diện.} Vế phải hằng số = đa thức bậc 0 ⇒ \(\beta=1\);
  PTĐT có nghiệm kép 1 → \protect\hyperlink{REC-03}{REC-03}.
\item
  \textbf{Kiến thức.} \protect\hyperlink{F-REC-03}{F-REC-03} (TH3).
\item
  \textbf{Phương pháp.} \protect\hyperlink{M-17}{M-17}.
\item
  \textbf{Lời giải.} PTĐT \((\lambda-1)^2=0\): \(\lambda=1\) kép ⇒
  \(\bar x_n=C_1+nC_2\). Vế phải \(2=1^n\cdot2\): \(\beta=1\)
  \textbf{nghiệm kép} ⇒ \(x_n^*=An^2\). Thế:
  \(A[(n+2)^2-2(n+1)^2+n^2]=2A=2\Rightarrow A=1\). \(x_n=C_1+nC_2+n^2\);
  \(x_0\): \(C_1=1\); \(x_1\): \(1+C_2+1=0\Rightarrow C_2=-2\).
\item
  \textbf{Kết luận.} \(x_n=n^2-2n+1=(n-1)^2\).
\item
  \textbf{Kiểm tra.} \(x_2=2x_1-x_0+2=1=(2-1)^2\) ✓.
\item
  \textbf{Biến thể (V1).} Vế phải \(2^n\): \(\beta=2\) không là nghiệm ⇒
  \(x_n^*=A2^n\), \(A(4-4+1)=1\Rightarrow x_n^*=2^n\).
\item
  \textbf{Lỗi thường gặp.} Thử \(x_n^*=A\) hoặc \(An\) (đều nằm trong
  nghiệm thuần nhất) --- \protect\hyperlink{T-15}{T-15}.
\item
  \textbf{Liên kết.} \protect\hyperlink{REC-03-E09}{REC-03-E09}.
\end{itemize}

\hypertarget{REC-03-E07}{%
\subsection{\texorpdfstring{REC-03-E07 · Bài tập S2-p67 (4): đa thức bậc
2, \(\beta=1\) không là
nghiệm}{REC-03-E07 · Bài tập S2-p67 (4): đa thức bậc 2, \textbackslash beta=1 không là nghiệm}}\label{REC-03-E07}}

\textbf{D2 · V1 · Đề {[}SV{]} S2-p67 · Lời giải {[}CC{]}}

\begin{itemize}
\tightlist
\item
  \textbf{Đề.} \(2x_n-5x_{n-1}+2x_{n-2}=-n^2-2n+3\), \(x_0=1\),
  \(x_1=3\).
\item
  \textbf{Nhận diện.} Vế phải đa thức bậc 2 ⇒ \(\beta=1\), \(r=2\) →
  \protect\hyperlink{REC-03}{REC-03}.
\item
  \textbf{Kiến thức.} \protect\hyperlink{F-REC-03}{F-REC-03} (TH1).
\item
  \textbf{Phương pháp.} \protect\hyperlink{M-17}{M-17}, đồng nhất hệ số.
\item
  \textbf{Lời giải.} PTĐT
  \(2\lambda^2-5\lambda+2=0\Rightarrow\lambda=2,\ \frac12\). \(\beta=1\)
  không là nghiệm ⇒ \(x_n^*=an^2+bn+c\). Đồng nhất: \(n^2\): \(-a=-1\);
  \(n\): \(2a-b=-2\); hằng: \(3a+b-c=3\) ⇒ \(a=1,b=4,c=4\) ⇒
  \(x_n^*=(n+2)^2\). \(x_n=C_12^n+C_2\left(\frac12\right)^n+(n+2)^2\);
  \(x_0\): \(C_1+C_2=-3\); \(x_1\): \(2C_1+\frac{C_2}2=-6\) ⇒
  \(C_1=-3\), \(C_2=0\).
\item
  \textbf{Kết luận.} \(x_n=(n+2)^2-3\cdot2^n\).
\item
  \textbf{Kiểm tra.}
  \(2x_2=5\cdot3-2\cdot1+(-4-4+3)=8\Rightarrow x_2=4\); công thức
  \(16-12=4\) ✓.
\item
  \textbf{Biến thể (V4).} Vế phải \(2^n\): \(\beta=2\) nghiệm đơn ⇒
  \(x_n^*=An2^n\), \(A=\frac23\).
\item
  \textbf{Lỗi thường gặp.} Đồng nhất sai hệ số khi khai triển
  \((n-1)^2,(n-2)^2\).
\item
  \textbf{Liên kết.} \protect\hyperlink{REC-03-E04}{REC-03-E04}.
\end{itemize}

\hypertarget{REC-03-E08}{%
\subsection{\texorpdfstring{REC-03-E08 · Bài tập S2-p67 (5): viết lại vế
phải
\(2\cdot5^{n+1}\)}{REC-03-E08 · Bài tập S2-p67 (5): viết lại vế phải 2\textbackslash cdot5\^{}\{n+1\}}}\label{REC-03-E08}}

\textbf{D2 · V2 · Đề {[}SV{]} S2-p67 · Lời giải {[}CC{]}}

\begin{itemize}
\tightlist
\item
  \textbf{Đề.} \(x_{n+2}-8x_{n+1}+15x_n=2\cdot5^{n+1}\), \(x_0=-1\),
  \(x_1=-2\).
\item
  \textbf{Nhận diện.} \(2\cdot5^{n+1}=10\cdot5^n\) ⇒ \(\beta=5\),
  \(r=0\).
\item
  \textbf{Kiến thức.} \protect\hyperlink{F-REC-03}{F-REC-03} (TH2).
\item
  \textbf{Phương pháp.} \protect\hyperlink{M-17}{M-17}.
\item
  \textbf{Lời giải.} PTĐT
  \(\lambda^2-8\lambda+15=0\Rightarrow\lambda=3,5\). \(\beta=5\) nghiệm
  đơn ⇒ \(x_n^*=An5^n\). Thế:
  \(A5^n[25(n+2)-40(n+1)+15n]=10A\cdot5^n=10\cdot5^n\Rightarrow A=1\).
  \(x_n=C_13^n+C_25^n+n5^n\); \(x_0\): \(C_1+C_2=-1\); \(x_1\):
  \(3C_1+5C_2+5=-2\) ⇒ \(C_1=1\), \(C_2=-2\).
\item
  \textbf{Kết luận.} \(x_n=3^n+(n-2)5^n\).
\item
  \textbf{Kiểm tra.} \(x_2=8(-2)-15(-1)+10=9\); công thức \(9+0=9\) ✓.
\item
  \textbf{Biến thể (V1).} Vế phải \(3^n\): \(\beta=3\) nghiệm đơn ⇒
  \(x_n^*=An3^n\), \(A=-\frac16\).
\item
  \textbf{Lỗi thường gặp.} Lấy \(\beta=5\) nhưng giữ hệ số 2 thay vì 10
  --- \protect\hyperlink{T-13}{T-13}.
\item
  \textbf{Liên kết.} \protect\hyperlink{F-REC-03}{F-REC-03}.
\end{itemize}

\hypertarget{REC-03-E09}{%
\subsection{\texorpdfstring{REC-03-E09 · Bài tập S2-p67 (6): \(\beta=8\)
nghiệm
kép}{REC-03-E09 · Bài tập S2-p67 (6): \textbackslash beta=8 nghiệm kép}}\label{REC-03-E09}}

\textbf{D2 · V0 · Đề {[}SV{]} S2-p67 · Lời giải {[}CC{]}}

\begin{itemize}
\tightlist
\item
  \textbf{Đề.} \(x_{n+2}-16x_{n+1}+64x_n=128\cdot8^n\), \(x_0=2\),
  \(x_1=32\).
\item
  \textbf{Nhận diện.} \(\beta=8\), PTĐT \((\lambda-8)^2=0\) →
  \protect\hyperlink{REC-03}{REC-03}.
\item
  \textbf{Kiến thức.} \protect\hyperlink{F-REC-03}{F-REC-03} (TH3).
\item
  \textbf{Phương pháp.} \protect\hyperlink{M-17}{M-17}.
\item
  \textbf{Lời giải.} PTĐT \((\lambda-8)^2=0\): \(\lambda=8\) kép.
  \(\beta=8\) nghiệm kép ⇒ \(x_n^*=An^28^n\). Thế:
  \(64A\cdot8^n[(n+2)^2-2(n+1)^2+n^2]=128A\cdot8^n\Rightarrow A=1\).
  \(x_n=(C_1+nC_2)8^n+n^28^n\); \(x_0\): \(C_1=2\); \(x_1\):
  \(8(2+C_2)+8=32\Rightarrow C_2=1\).
\item
  \textbf{Kết luận.} \(x_n=(n^2+n+2)\,8^n\).
\item
  \textbf{Kiểm tra.} \(x_2=16\cdot32-64\cdot2+128=512\); công thức
  \(8\cdot64=512\) ✓.
\item
  \textbf{Biến thể (V1).} Vế phải hằng 64: \(\beta=1\) không là nghiệm ⇒
  \(49A=64\), \(x_n^*=\frac{64}{49}\).
\item
  \textbf{Lỗi thường gặp.} \protect\hyperlink{T-15}{T-15}.
\item
  \textbf{Liên kết.} \protect\hyperlink{REC-03-E02}{REC-03-E02}.
\end{itemize}

\hypertarget{REC-04-E01}{%
\subsection{REC-04-E01 · Vế phải tổng ba dạng (chồng
chất)}\label{REC-04-E01}}

\textbf{D3 · V3 · Đề {[}SV{]} S2-p53 d), S2-p64 · Lời giải {[}SV{]}
S2-p64--66, {[}CC{]}}

\begin{itemize}
\tightlist
\item
  \textbf{Đề.} Giải \(x_n-4x_{n-1}+3x_{n-2}=20+(2-n)2^{n-2}+3\cdot4^n\)
  (1).
\item
  \textbf{Nhận diện.} Ba dạng với \(\beta=1,2,4\) →
  \protect\hyperlink{REC-04}{REC-04}.
\item
  \textbf{Kiến thức.} \protect\hyperlink{F-REC-04}{F-REC-04}.
  \textbf{Phương pháp.} \protect\hyperlink{M-17}{M-17}.
\item
  \textbf{Lời giải.} PTĐT
  \(\lambda^2-4\lambda+3=0\Rightarrow\lambda=1,3\) ⇒
  \(\bar x_n=C_1+C_23^n\). \((1')\) vế phải 20: \(\beta=1\) nghiệm đơn ⇒
  \(An\): \(A[n-4(n-1)+3(n-2)]=-2A=20\Rightarrow A=-10\). \((1'')\) vế
  phải \((2-n)2^{n-2}\): \(\beta=2\) không là nghiệm ⇒ \((an+b)2^n\);
  chia \(2^{n-2}\): \(-an+(2a-b)=2-n\Rightarrow a=1,b=0\) ⇒ \(n2^n\).
  \((1''')\) vế phải \(3\cdot4^n\): \(\beta=4\) không là nghiệm ⇒
  \(A4^n\): \(A4^{n-2}(16-16+3)=3\cdot4^n\Rightarrow A=16\) ⇒
  \(16\cdot4^n=4^{n+2}\).
\item
  \textbf{Kết luận.} \(x_n=C_1+C_2\,3^n-10n+n\,2^n+4^{n+2}\).
\item
  \textbf{Kiểm tra.} Thế nghiệm riêng tổng vào (1) với \(n=2..25\)
  {[}CC{]}.
\item
  \textbf{Biến thể (V3).} \(x_n-4x_{n-1}+3x_{n-2}=3^n+1\): cả
  \(\beta=3\) và \(\beta=1\) đều là nghiệm đơn ⇒
  \(x_n^*=\frac32n3^n-\frac n2\).
\item
  \textbf{Lỗi thường gặp.} Gộp \((2-n)2^{n-2}\) và \(3\cdot4^n\) chung
  một dạng; xét bội của \(\beta\) chung cho cả vế phải ---
  \protect\hyperlink{T-15}{T-15}.
\item
  \textbf{Ghi chú nguồn.} S2-p65 gọi đây là ``Trường hợp 4'' --- thực
  chất là Dạng 2 (ERR-09).
\item
  \textbf{Liên kết.} \protect\hyperlink{F-REC-04}{F-REC-04}.
\end{itemize}

\hypertarget{ex-ch3}{%
\chapter{Ví dụ Chương III --- Quan hệ, quan hệ thứ tự}\label{ex-ch3}}

\hypertarget{REL-01-E01}{%
\subsection{REL-01-E01 · Quan hệ ↔ ma trận biểu diễn}\label{REL-01-E01}}

\textbf{D1 · V0 · Đề {[}SV{]} S3-p13--15 · Lời giải {[}SV{]}}

\begin{itemize}
\tightlist
\item
  \textbf{Đề.} (a) \(R\) từ \(A=\{1,2,3,4\}\) đến \(B=\{u,v,w\}\):
  \(R=\{(1,u),(1,v),(2,w),(3,w),(4,u)\}\). Tìm \(M_R\). (b) \(R\) từ
  \(A=\{1,2,3\}\) đến \(B=\{1,2\}\): \(aRb\Leftrightarrow a>b\). Tìm
  \(M_R\). (c) \(R\) từ \(\{a_1,a_2,a_3\}\) đến \(\{b_1,\dots,b_5\}\) có
  \(M_R=\begin{pmatrix}0&1&0&0&0\\1&0&1&1&0\\1&0&1&0&1\end{pmatrix}\).
  Liệt kê \(R\).
\item
  \textbf{Nhận diện.} ``ma trận biểu diễn'' →
  \protect\hyperlink{REL-01}{REL-01}.
\item
  \textbf{Kiến thức.} \protect\hyperlink{F-REL-01}{F-REL-01}: hàng =
  phần tử của \(A\), cột = phần tử của \(B\).
\item
  \textbf{Phương pháp.} Duyệt từng cặp \((a_i,b_j)\in R\) và ghi 1 vào ô
  (hàng \(i\), cột \(j\)); chiều ngược lại đọc từng ô bằng 1.
\item
  \textbf{Lời giải.} (a)
  \(M_R=\begin{pmatrix}1&1&0\\0&0&1\\0&0&1\\1&0&0\end{pmatrix}\) (hàng
  1,2,3,4; cột \(u,v,w\)). (b) \(R=\{(2,1),(3,1),(3,2)\}\),
  \(M_R=\begin{pmatrix}0&0\\1&0\\1&1\end{pmatrix}\). (c)
  \(R=\{(a_1,b_2),(a_2,b_1),(a_2,b_3),(a_2,b_4),(a_3,b_1),(a_3,b_3),(a_3,b_5)\}\).
\item
  \textbf{Kết luận.} Như trên; \(M_R\) cấp \(|A|\times|B|\).
\item
  \textbf{Kiểm tra.} Số số 1 trong ma trận = số cặp: (a) 5, (b) 3, (c) 7
  ✓.
\item
  \textbf{Biến thể (V1).} Quan hệ ``\(a\) là ước của \(b\)'' trên
  \(\{1,2,3,4\}\) (S3-p6):
  \(M_R=\begin{pmatrix}1&1&1&1\\0&1&0&1\\0&0&1&0\\0&0&0&1\end{pmatrix}\)
  (8 cặp).
\item
  \textbf{Lỗi thường gặp.} Đảo hàng/cột (viết \(M_R\) cấp
  \(|B|\times|A|\)).
\item
  \textbf{Liên kết.} \protect\hyperlink{REL-02}{REL-02}.
\end{itemize}

\hypertarget{REL-02-E01}{%
\subsection{REL-02-E01 · Bảng bốn tính chất của quan hệ hữu
hạn}\label{REL-02-E01}}

\textbf{D1 · V0 · Đề {[}SV{]} S3-p7, p9, p11 · Lời giải {[}CC{]}}

\begin{itemize}
\tightlist
\item
  \textbf{Đề.} Trên \(A=\{1,2,3,4\}\) xét tính phản xạ, đối xứng, phản
  xứng, bắc cầu của:
\end{itemize}

\[
\begin{aligned}
R_1&=\{(1,1),(1,2),(2,1),(2,2),(3,4),(4,1),(4,4)\},\\
R_2&=\{(1,1),(1,2),(1,4),(2,2),(3,3),(4,1),(4,4)\},\\
R_3&=\{(1,1),(1,2),(2,1)\},\\
R_4&=\{(1,1),(1,2),(2,1),(2,2),(1,3),(2,3)\}.
\end{aligned}
\]

\begin{itemize}
\tightlist
\item
  \textbf{Nhận diện.} \protect\hyperlink{REL-02}{REL-02}. \textbf{Kiến
  thức.} \protect\hyperlink{F-REL-01}{F-REL-01}. \textbf{Phương pháp.}
  \protect\hyperlink{M-18}{M-18}.
\item
  \textbf{Lời giải.}
\end{itemize}

\begin{longtable}[]{@{}
  >{\raggedright\arraybackslash}p{(\columnwidth - 8\tabcolsep) * \real{0.2000}}
  >{\raggedright\arraybackslash}p{(\columnwidth - 8\tabcolsep) * \real{0.2000}}
  >{\raggedright\arraybackslash}p{(\columnwidth - 8\tabcolsep) * \real{0.2000}}
  >{\raggedright\arraybackslash}p{(\columnwidth - 8\tabcolsep) * \real{0.2000}}
  >{\raggedright\arraybackslash}p{(\columnwidth - 8\tabcolsep) * \real{0.2000}}@{}}
\toprule\noalign{}
\begin{minipage}[b]{\linewidth}\raggedright
\end{minipage} & \begin{minipage}[b]{\linewidth}\raggedright
Phản xạ
\end{minipage} & \begin{minipage}[b]{\linewidth}\raggedright
Đối xứng
\end{minipage} & \begin{minipage}[b]{\linewidth}\raggedright
Phản xứng
\end{minipage} & \begin{minipage}[b]{\linewidth}\raggedright
Bắc cầu
\end{minipage} \\
\midrule\noalign{}
\endhead
\bottomrule\noalign{}
\endlastfoot
\(R_1\) & Không: \((3,3)\notin R_1\) & Không: \((3,4)\), thiếu \((4,3)\)
& Không: \((1,2),(2,1)\) & Không: \((3,4),(4,1)\) mà
\((3,1)\notin R_1\) \\
\(R_2\) & \textbf{Có} & Không: \((1,2)\), thiếu \((2,1)\) & Không:
\((1,4),(4,1)\) & Không: \((4,1),(1,2)\) mà \((4,2)\notin R_2\) \\
\(R_3\) & Không: \((2,2)\notin R_3\) & \textbf{Có} & Không:
\((1,2),(2,1)\) & Không: \((2,1),(1,2)\) mà \((2,2)\notin R_3\) \\
\(R_4\) & Không: \((3,3)\notin R_4\) & Không: \((1,3)\) & Không:
\((1,2),(2,1)\) & \textbf{Có} \\
\end{longtable}

\begin{itemize}
\tightlist
\item
  \textbf{Kết luận.} Như bảng (khớp các khẳng định của slide: \(R_1\)
  không phản xạ, \(R_2\) phản xạ, \(R_3\) đối xứng, \(R_4\) bắc cầu).
\item
  \textbf{Kiểm tra.} Chương trình kiểm tra 4 tính chất theo định nghĩa
  {[}CC{]}. Bắc cầu \(R_4\): mọi cặp nối \((a,b),(b,c)\):
  \((1,1)(1,x)\), \((1,2)(2,x)\), \((2,1)(1,x)\), \((2,2)(2,x)\) đều cho
  \((1,x),(2,x)\in R_4\) ✓.
\item
  \textbf{Biến thể (V3).} Trên ma trận {[}AI{]}:
  \(M=\begin{pmatrix}1&1&0\\0&1&1\\0&0&1\end{pmatrix}\) (trên
  \(\{u,v,w\}\)): phản xạ (chéo toàn 1), phản xứng, \textbf{không} bắc
  cầu (\(uRv\), \(vRw\), \(u\not R w\)).
\item
  \textbf{Lỗi thường gặp.} Kết luận ``đối xứng'' khi chỉ thấy một cặp
  đối xứng; quên kiểm tra bắc cầu với cặp có \(a=c\) (như
  \((2,1),(1,2)\Rightarrow(2,2)\)) --- \protect\hyperlink{T-16}{T-16}.
\item
  \textbf{Liên kết.} \protect\hyperlink{M-18}{M-18}.
\end{itemize}

\hypertarget{REL-02-E02}{%
\subsection{\texorpdfstring{REL-02-E02 · Quan hệ trên \(\mathbb Z\):
\(aRb\Leftrightarrow ab\ge0\) --- bẫy bắc
cầu}{REL-02-E02 · Quan hệ trên \textbackslash mathbb Z: aRb\textbackslash Leftrightarrow ab\textbackslash ge0 --- bẫy bắc cầu}}\label{REL-02-E02}}

\textbf{D2 · V4 · {[}AI-GENERATED PRACTICE{]} · Lời giải {[}CC{]}}

\begin{itemize}
\tightlist
\item
  \textbf{Đề.} Trên \(\mathbb Z\): \(aRb\Leftrightarrow ab\ge0\). \(R\)
  có là quan hệ tương đương không?
\item
  \textbf{Nhận diện.} Quan hệ cho bằng quy tắc trên tập vô hạn →
  \protect\hyperlink{REL-02}{REL-02},
  \protect\hyperlink{REL-03}{REL-03}.
\item
  \textbf{Kiến thức / Phương pháp.}
  \protect\hyperlink{F-REL-01}{F-REL-01},
  \protect\hyperlink{M-18}{M-18}.
\item
  \textbf{Lời giải.} Phản xạ: \(a\cdot a=a^2\ge0\) ✓. Đối xứng:
  \(ab=ba\) ✓. Bắc cầu: \textbf{sai} --- \((-1)\cdot0\ge0\) và
  \(0\cdot1\ge0\) nhưng \((-1)\cdot1<0\). Phản xứng: sai ---
  \(1\cdot2\ge0\), \(2\cdot1\ge0\), \(1\ne2\).
\item
  \textbf{Kết luận.} Không phải quan hệ tương đương (không bắc cầu),
  cũng không phải quan hệ thứ tự.
\item
  \textbf{Kiểm tra.} Kiểm tra vét cạn trên \(\{-3,\dots,3\}\) tìm thấy
  đúng loại phản ví dụ chứa 0 {[}CC{]}.
\item
  \textbf{Biến thể (V4).} Trên \(\mathbb Z\setminus\{0\}\),
  \(aRb\Leftrightarrow ab>0\): \textbf{là} quan hệ tương đương (bỏ 0 thì
  bắc cầu đúng: cùng dấu), có 2 lớp: số dương, số âm.
\item
  \textbf{Lỗi thường gặp.} Thử vài số dương rồi kết luận bắc cầu ---
  \protect\hyperlink{T-16}{T-16}.
\item
  \textbf{Liên kết.} \protect\hyperlink{REL-03-E01}{REL-03-E01}.
\end{itemize}

\hypertarget{REL-03-E01}{%
\subsection{\texorpdfstring{REL-03-E01 · Đồng dư modulo \(m\) và lớp
tương
đương}{REL-03-E01 · Đồng dư modulo m và lớp tương đương}}\label{REL-03-E01}}

\textbf{D1 · V0 · Đề {[}SV{]} S3-p22, p24, p27 · Lời giải {[}SV{]}}

\begin{itemize}
\tightlist
\item
  \textbf{Đề.} (a) Chứng minh quan hệ đồng dư modulo \(m\)
  (\(aRb\Leftrightarrow m\mid a-b\)) là quan hệ tương đương trên
  \(\mathbb Z\). (b) Tìm các lớp tương đương modulo 8 chứa 0 và chứa 1.
  (c) Có bao nhiêu lớp đồng dư modulo \(m\)?
\item
  \textbf{Nhận diện.} \protect\hyperlink{REL-03}{REL-03}. \textbf{Kiến
  thức.} \protect\hyperlink{F-REL-02}{F-REL-02},
  \protect\hyperlink{F-REL-03}{F-REL-03}. \textbf{Phương pháp.}
  \protect\hyperlink{M-18}{M-18}, \protect\hyperlink{M-19}{M-19}.
\item
  \textbf{Lời giải.} (a) Phản xạ: \(a-a=0=0\cdot m\). Đối xứng:
  \(a-b=km\Rightarrow b-a=(-k)m\). Bắc cầu: \(a-b=km\), \(b-c=jm\) ⇒
  \(a-c=(k+j)m\). (b)
  \([0]_8=\{a: 8\mid a\}=\{\dots,-16,-8,0,8,16,\dots\}\);
  \([1]_8=\{a: a\text{ chia 8 dư 1}\}=\{\dots,-15,-7,1,9,17,\dots\}\).
  (c) Đúng \(m\) lớp \([0]_m,\dots,[m-1]_m\) (theo số dư), tạo phân
  hoạch của \(\mathbb Z\); \(\mathbb Z_m=\{[0]_m,\dots,[m-1]_m\}\).
\item
  \textbf{Kết luận.} Như trên.
\item
  \textbf{Kiểm tra.} \(-7-1=-8\) chia hết cho 8 ⇒ \(-7\in[1]_8\) ✓;
  \([0]_8\cap[1]_8=\emptyset\) (số dư khác nhau).
\item
  \textbf{Biến thể (V1).} Lớp modulo 3:
  \([0]_3=\{\dots,-3,0,3,\dots\}\), \([1]_3=\{\dots,-2,1,4,\dots\}\),
  \([2]_3=\{\dots,-1,2,5,\dots\}\); \([5]_3=[2]_3\).
\item
  \textbf{Lỗi thường gặp.} Ghi \(-7\in[7]_8\) ``vì giống 7'' --- dùng số
  dư không âm: \(-7=8\cdot(-1)+1\).
\item
  \textbf{Liên kết.} \protect\hyperlink{M-19}{M-19}.
\end{itemize}

\hypertarget{REL-03-E02}{%
\subsection{REL-03-E02 · Từ phân hoạch sang quan hệ tương
đương}\label{REL-03-E02}}

\textbf{D2 · V2 · {[}AI-GENERATED PRACTICE{]} dựa trên S3-p26 · Lời giải
{[}CC{]}}

\begin{itemize}
\tightlist
\item
  \textbf{Đề.} Trên \(A=\{1,2,3,4,5,6\}\) cho phân hoạch
  \(\{\{1,2\},\{3\},\{4,5,6\}\}\). Viết quan hệ tương đương \(R\) tương
  ứng; \(|R|=?\); tìm \([5]_R\).
\item
  \textbf{Nhận diện.} ``phân hoạch'' ↔ quan hệ tương đương →
  \protect\hyperlink{REL-03}{REL-03}.
\item
  \textbf{Kiến thức.} \protect\hyperlink{F-REL-02}{F-REL-02}:
  \(aRb\Leftrightarrow a,b\) cùng thuộc một khối.
\item
  \textbf{Phương pháp.} \protect\hyperlink{M-19}{M-19} theo chiều ngược:
  mỗi khối \(B\) của phân hoạch góp các cặp \(B\times B\).
\item
  \textbf{Lời giải.} \(R=\{1,2\}^2\cup\{3\}^2\cup\{4,5,6\}^2\), tức \[
  \begin{aligned}
  R=\{&(1,1),(1,2),(2,1),(2,2),(3,3),\\
  &(4,4),(4,5),(4,6),(5,4),(5,5),(5,6),(6,4),(6,5),(6,6)\}.
  \end{aligned}
  \] \(|R|=2^2+1^2+3^2=14\). \([5]_R=\{4,5,6\}\).
\item
  \textbf{Kết luận.} \(|R|=14\), \([5]_R=\{4,5,6\}\).
\item
  \textbf{Kiểm tra.} \(R\) phản xạ, đối xứng, bắc cầu (máy kiểm) và các
  lớp đúng bằng các khối {[}CC{]}.
\item
  \textbf{Biến thể (V1).} S3-p21: trên \(\mathbb R\),
  \(aRb\Leftrightarrow a-b\in\mathbb Z\) là tương đương;
  \([x]=\{x+k\mid k\in\mathbb Z\}\), ví dụ
  \([0{,}5]=\{\dots,-0{,}5;\ 0{,}5;\ 1{,}5;\dots\}\).
\item
  \textbf{Lỗi thường gặp.} Quên các cặp \((a,a)\).
\item
  \textbf{Liên kết.} \protect\hyperlink{M-19}{M-19}.
\end{itemize}

\hypertarget{POS-01-E01}{%
\subsection{POS-01-E01 · Có phải poset? Thứ tự toàn
phần?}\label{POS-01-E01}}

\textbf{D2 · V0 · Đề {[}SV{]} S3-p31--35 · Lời giải {[}SV{]}}

\begin{itemize}
\tightlist
\item
  \textbf{Đề.} Xét: (a) \((\mathbb Z^+,\mid)\); (b)
  \((\mathbb Z,\mid)\); (c) \((P(S),\subseteq)\); (d)
  \((\mathbb Z^+,\le)\). Cái nào là poset? Cái nào là thứ tự toàn phần?
\item
  \textbf{Nhận diện.} \protect\hyperlink{POS-01}{POS-01}. \textbf{Kiến
  thức.} \protect\hyperlink{F-POS-01}{F-POS-01}. \textbf{Phương pháp.}
  \protect\hyperlink{M-18}{M-18}.
\item
  \textbf{Lời giải.} (a) Phản xạ: \(x=1\cdot x\). Bắc cầu: \(b=ka\),
  \(c=jb\) ⇒ \(c=(jk)a\). Phản xứng: \(b=ka\), \(a=jb\) ⇒ \(a=jka\) ⇒
  \(jk=1\) ⇒ \(j=k=1\) (vì \(j,k\) nguyên dương) ⇒ \(a=b\). ⇒ poset;
  \textbf{không toàn phần} (5 và 7 không so sánh được). (b) Không phản
  xứng: \(3\mid-3\), \(-3\mid3\), \(3\ne-3\) ⇒ \textbf{không} poset. (c)
  \(A\subseteq A\); \(A\subseteq B\subseteq C\Rightarrow A\subseteq C\);
  \(A\subseteq B,B\subseteq A\Rightarrow A=B\) ⇒ poset; không toàn phần
  khi \(|S|\ge2\) (\(\{a\},\{b\}\)). (d) poset toàn phần.
\item
  \textbf{Kết luận.} (a), (c), (d) là poset; chỉ (d) toàn phần.
\item
  \textbf{Kiểm tra.} (b): lỗi xảy ra vì \(k=j=-1\) cũng cho \(jk=1\)
  trên \(\mathbb Z\).
\item
  \textbf{Biến thể (V4).} \((\{1,2,4,8\},\mid)\) là thứ tự toàn phần
  (một xích); \((\{2,3,6\},\mid)\) không toàn phần.
\item
  \textbf{Lỗi thường gặp.} Chứng minh phản xứng trên \(\mathbb Z\) giống
  trên \(\mathbb Z^+\) --- \protect\hyperlink{T-17}{T-17}.
\item
  \textbf{Liên kết.} \protect\hyperlink{POS-03}{POS-03}.
\end{itemize}

\hypertarget{POS-02-E01}{%
\subsection{POS-02-E01 · Thứ tự từ điển}\label{POS-02-E01}}

\textbf{D1 · V0 · Đề {[}SV{]} S3-p36, p40--41 · Lời giải {[}SV{]}}

\begin{itemize}
\tightlist
\item
  \textbf{Đề.} So sánh theo thứ tự từ điển: (a) discreet / discrete; (b)
  discreet / discreetness; (c) \(0110\) / \(10\); (d) \(0110\) /
  \(01100\). (e) Với thứ tự ``\(a_i\le b_i\ \forall i\)'' trên chuỗi bit
  cùng độ dài, \(0110\) và \(1000\) có so sánh được không?
\item
  \textbf{Nhận diện.} \protect\hyperlink{POS-02}{POS-02}. \textbf{Kiến
  thức.} \protect\hyperlink{F-POS-02}{F-POS-02}.
\item
  \textbf{Phương pháp.} So từng ký tự từ trái sang; dừng ở vị trí khác
  nhau đầu tiên hoặc khi chuỗi ngắn hơn hết ký tự (tiền tố).
\item
  \textbf{Lời giải.} (a) Giống tới ``discre'', vị trí 7: e \textless{} t
  ⇒ discreet \(\prec\) discrete. (b) discreet là tiền tố ⇒ discreet
  \(\prec\) discreetness. (c) Vị trí 1: \(0<1\) ⇒ \(0110\prec10\) (độ
  dài không quan trọng). (d) \(0110\) là tiền tố ⇒ \(0110\prec01100\).
  (e) Không: vị trí 1 có \(0<1\) nhưng vị trí 2 có \(1>0\).
\item
  \textbf{Kết luận.} Như trên.
\item
  \textbf{Kiểm tra.} So sánh chuỗi bằng máy (thứ tự từ điển chuẩn trùng
  định nghĩa S3-p39) ✓.
\item
  \textbf{Biến thể (V1).} Sắp tăng: b, ab, a, aab ⇒ a \(\prec\) aab
  \(\prec\) ab \(\prec\) b.
\item
  \textbf{Lỗi thường gặp.} So độ dài trước (đó là thứ tự khác) ---
  \protect\hyperlink{T-29}{T-29}.
\item
  \textbf{Liên kết.} \protect\hyperlink{F-POS-02}{F-POS-02}.
\end{itemize}

\hypertarget{POS-03-E01}{%
\subsection{POS-03-E01 · Hasse, tối đại, tối tiểu, lớn nhất, nhỏ
nhất}\label{POS-03-E01}}

\textbf{D2 · V0 · Đề {[}SV{]} S3-p48 · Lời giải {[}SV{]}, {[}CC{]}}

\begin{itemize}
\tightlist
\item
  \textbf{Đề.} Vẽ biểu đồ Hasse và tìm phần tử tối đại, tối tiểu, lớn
  nhất, nhỏ nhất của poset \(S=\{2,4,5,10,12,20,25\}\) với quan hệ chia
  hết \(\mid\).
\item
  \textbf{Nhận diện.} \protect\hyperlink{POS-03}{POS-03}. \textbf{Kiến
  thức.} \protect\hyperlink{F-POS-03}{F-POS-03}. \textbf{Phương pháp.}
  \protect\hyperlink{M-20}{M-20}.
\item
  \textbf{Lời giải.} Cặp \(a\mid b\) (\(a\ne b\)): 2\textbar4,
  2\textbar10, 2\textbar12, 2\textbar20, 4\textbar12, 4\textbar20,
  5\textbar10, 5\textbar20, 5\textbar25, 10\textbar20. Bỏ cặp bắc cầu:
  2\textbar12 (qua 4), 2\textbar20 (qua 4 hoặc 10), 5\textbar20 (qua
  10). \textbf{Cung Hasse:} 2--4, 2--10, 4--12, 4--20, 5--10, 5--25,
  10--20. Tầng dưới: 2, 5; tầng giữa: 4, 10, 25; tầng trên: 12, 20. Tối
  đại: không có cung đi lên: 12, 20, 25. Tối tiểu: không có cung đi
  xuống: 2, 5. Lớn nhất / nhỏ nhất: \textbf{không có} (có hơn một tối
  đại / tối tiểu).
\item
  \textbf{Kết luận.} Tối đại 12, 20, 25; tối tiểu 2, 5; không có lớn
  nhất, nhỏ nhất.
\item
  \textbf{Kiểm tra.} Chương trình liệt kê tối đại/tối tiểu theo định
  nghĩa {[}CC{]}. 25 vừa thuộc tầng giữa vừa tối đại --- tối đại không
  nhất thiết ở tầng cao nhất.
\item
  \textbf{Biến thể (V1).} \((\{1,2,3,4,6,12\},\mid)\): cung 1--2, 1--3,
  2--4, 2--6, 3--6, 4--12, 6--12; nhỏ nhất 1, lớn nhất 12.
\item
  \textbf{Lỗi thường gặp.} Vẽ cả cung 2--12, 2--20 (bắc cầu); cho rằng
  tối đại là lớn nhất --- \protect\hyperlink{T-18}{T-18}.
\item
  \textbf{Liên kết.} \protect\hyperlink{POS-04-E01}{POS-04-E01}.
\end{itemize}

\hypertarget{POS-04-E01}{%
\subsection{POS-04-E01 · Chặn trên, chặn dưới, sup,
inf}\label{POS-04-E01}}

\textbf{D2 · V4 · {[}AI-GENERATED PRACTICE{]} (định nghĩa S3-p50--53) ·
Lời giải {[}CC{]}}

\begin{itemize}
\tightlist
\item
  \textbf{Đề.} Poset \((S,\mid)\), \(S=\{2,3,4,6,12,18\}\). Tìm chặn
  trên, chặn dưới, sup, inf (nếu có) của \(A=\{2,3\}\), \(B=\{12,18\}\),
  \(C=\{4,6\}\).
\item
  \textbf{Nhận diện.} \protect\hyperlink{POS-04}{POS-04}. \textbf{Kiến
  thức.} \protect\hyperlink{F-POS-03}{F-POS-03}. \textbf{Phương pháp.}
  \protect\hyperlink{M-20}{M-20}.
\item
  \textbf{Lời giải.} Hasse: 2--4, 2--6, 3--6, 4--12, 6--12, 6--18.
  \(A=\{2,3\}\): chặn trên = phần tử chia hết cho cả 2 và 3 trong \(S\):
  \(6,12,18\); nhỏ nhất trong chúng: 6 (\(6\mid12\), \(6\mid18\)) ⇒
  \(\sup A=6\). Chặn dưới: ước chung của 2, 3 trong \(S\): \textbf{không
  có} (1 \(\notin S\)) ⇒ \(\inf A\) không tồn tại. \(B=\{12,18\}\): chặn
  trên: không có (\(36\notin S\)) ⇒ không có sup. Chặn dưới: \(2,3,6\);
  lớn nhất: 6 (\(2\mid6\), \(3\mid6\)) ⇒ \(\inf B=6\). \(C=\{4,6\}\):
  chặn trên: 12 ⇒ \(\sup C=12\); chặn dưới: 2 ⇒ \(\inf C=2\).
\item
  \textbf{Kết luận.} \(\sup A=6\), \(\inf A\) không có; \(\sup B\) không
  có, \(\inf B=6\); \(\sup C=12\), \(\inf C=2\).
\item
  \textbf{Kiểm tra.} Chương trình vét cạn theo định nghĩa {[}CC{]}. Chặn
  trên 6 của \(A\) \textbf{không thuộc} \(A\) --- hợp lệ (S3-p50).
\item
  \textbf{Biến thể (V2).} Trong \((P(\{a,b,c\}),\subseteq)\):
  \(\{a\}\vee\{b\}=\{a,b\}\), \(\{a,b\}\wedge\{b,c\}=\{b\}\) (sup = hợp,
  inf = giao).
\item
  \textbf{Lỗi thường gặp.} Chọn sup trong số các chặn trên \textbf{không
  so sánh được} (12 và 18 đều là chặn trên của \(A\) nhưng không nhỏ
  nhất); cho rằng chặn phải thuộc \(A\) ---
  \protect\hyperlink{T-19}{T-19}.
\item
  \textbf{Liên kết.} \protect\hyperlink{POS-03-E01}{POS-03-E01}.
\end{itemize}

\hypertarget{POS-05-E01}{%
\subsection{POS-05-E01 · Sắp xếp topo}\label{POS-05-E01}}

\textbf{D1 · V0 · {[}AI-GENERATED PRACTICE{]} (thuật toán S3-p54--56) ·
Lời giải {[}CC{]}}

\begin{itemize}
\tightlist
\item
  \textbf{Đề.} Công việc \(a,b,c,d,e,f\) có quan hệ ``phải làm trước''
  (trội trực tiếp): \(a\prec b\), \(a\prec c\), \(e\prec c\),
  \(b\prec d\), \(c\prec d\), \(d\prec f\). Tìm một sắp xếp topo.
\item
  \textbf{Nhận diện.} \protect\hyperlink{POS-05}{POS-05}. \textbf{Phương
  pháp.} \protect\hyperlink{M-21}{M-21}.
\item
  \textbf{Kiến thức.} \protect\hyperlink{F-POS-03}{F-POS-03}: mỗi poset
  hữu hạn có phần tử tối tiểu; bỏ nó đi vẫn là poset (S3-p54).
\item
  \textbf{Lời giải.} Tối tiểu ban đầu: \(a\), \(e\). Lấy \(a\) → còn tối
  tiểu \(b\), \(e\). Lấy \(e\) → tối tiểu \(b\), \(c\). Lấy \(b\) →
  \(c\). Lấy \(c\) → \(d\). Lấy \(d\) → \(f\).
\item
  \textbf{Kết luận.} Một sắp xếp topo: \(a, e, b, c, d, f\) (không duy
  nhất; ví dụ \(e, a, c, b, d, f\) cũng đúng).
\item
  \textbf{Kiểm tra.} Mọi cặp \(x\prec y\) (kể cả cặp bắc cầu như
  \(a\prec f\)) có \(x\) đứng trước \(y\) --- máy kiểm {[}CC{]}.
\item
  \textbf{Biến thể (V4).} Thêm \(f\prec a\): tạo vòng
  \(a\prec b\prec d\prec f\prec a\) ⇒ không còn là thứ tự (vi phạm phản
  xứng) ⇒ \textbf{không} có sắp xếp topo (thuật toán kẹt: không còn phần
  tử tối tiểu).
\item
  \textbf{Lỗi thường gặp.} Lấy phần tử tối đại trước; quên rằng có nhiều
  đáp án đúng.
\item
  \textbf{Liên kết.} \protect\hyperlink{M-21}{M-21}.
\end{itemize}

\hypertarget{ex-ch4}{%
\chapter{Ví dụ Chương IV --- Hàm Bool, Karnaugh, mạch
logic}\label{ex-ch4}}

\begin{quote}
Bản đồ Karnaugh dùng \textbf{đúng khung của nguồn}
(\protect\hyperlink{F-BOO-03}{F-BOO-03}): cột
\(x\bar y\,|\,xy\,|\,\bar xy\,|\,\bar x\bar y\); hàng
\(z\bar t\,|\,zt\,|\,\bar zt\,|\,\bar z\bar t\) (3 biến: hàng
\(z\,|\,\bar z\)). Ô có 1 là ô thuộc \(\mathrm{kar}(f)\). Mọi kết quả
tối tiểu đã được kiểm tra bằng vét cạn
\texttt{10\_Verification/boolmin.py} \textbf{{[}CC{]}}.
\end{quote}

\hypertarget{BOO-01-E01}{%
\subsection{BOO-01-E01 · Hàm bỏ phiếu đa số: bảng chân trị → dạng nối
rời chính tắc → tối tiểu}\label{BOO-01-E01}}

\textbf{D1 · V0 · Đề {[}SV{]} S4-p12--13 · Lời giải {[}CC{]}}

\begin{itemize}
\tightlist
\item
  \textbf{Đề.} Quyết định \(f\) thông qua (=1) khi đa số trong 3 phiếu
  \(x,y,z\) tán thành (=1). Lập bảng chân trị, viết dạng nối rời chính
  tắc và công thức đa thức tối tiểu.
\item
  \textbf{Nhận diện.} Mô tả lời → hàm Bool →
  \protect\hyperlink{BOO-01}{BOO-01},
  \protect\hyperlink{BOO-03}{BOO-03}.
\item
  \textbf{Kiến thức.} \protect\hyperlink{F-BOO-02}{F-BOO-02},
  \protect\hyperlink{F-BOO-03}{F-BOO-03}. \textbf{Phương pháp.}
  \protect\hyperlink{M-23}{M-23}, \protect\hyperlink{M-22}{M-22}.
\item
  \textbf{Lời giải.} \(f=1\) khi có ít nhất 2 số 1: hàng
  \(011,101,110,111\) (các hàng khác \(f=0\); khớp bảng S4-p13). Dạng
  nối rời chính tắc: \(f=\bar xyz\vee x\bar yz\vee xy\bar z\vee xyz\).
  Karnaugh: hàng \(z\): \(x\bar y\)=1, \(xy\)=1, \(\bar xy\)=1,
  \(\bar x\bar y\)=0; hàng \(\bar z\): 0, 1, 0, 0. Tế bào lớn: \(xz\) (ô
  \(x\bar yz\), \(xyz\)), \(yz\) (\(xyz\), \(\bar xyz\)), \(xy\)
  (\(xyz\), \(xy\bar z\)) --- cả ba bắt buộc.
\item
  \textbf{Kết luận.} \(f=xy\vee yz\vee xz\).
\item
  \textbf{Kiểm tra.} So bảng chân trị của \(xy\vee yz\vee xz\) với
  ``tổng \(\ge2\)'' trên 8 hàng {[}CC{]}.
\item
  \textbf{Biến thể (V1).} ``Thông qua khi \textbf{cả 3} tán thành'':
  \(f=xyz\); ``khi có ít nhất 1'': \(f=x\vee y\vee z\).
\item
  \textbf{Lỗi thường gặp.} Viết từ tối tiểu cho hàng \(f=0\).
\item
  \textbf{Liên kết.} \protect\hyperlink{BOO-04}{BOO-04}.
\end{itemize}

\hypertarget{BOO-01-E02}{%
\subsection{BOO-01-E02 · Dạng nối rời chính tắc từ công
thức}\label{BOO-01-E02}}

\textbf{D1 · V1 · {[}AI-GENERATED PRACTICE{]} · Lời giải {[}CC{]}}

\begin{itemize}
\tightlist
\item
  \textbf{Đề.} Viết dạng nối rời chính tắc của \(f(x,y,z)=x\vee\bar yz\)
  và của \(\bar f\).
\item
  \textbf{Nhận diện.} \protect\hyperlink{BOO-01}{BOO-01}. \textbf{Kiến
  thức.} \protect\hyperlink{F-BOO-02}{F-BOO-02}.
\item
  \textbf{Phương pháp.} Bổ sung biến thiếu:
  \(x=x(y\vee\bar y)(z\vee\bar z)\); hoặc đọc hàng \(f=1\).
\item
  \textbf{Lời giải.} \(f=1\) khi \(x=1\) (4 hàng \(100,101,110,111\))
  hoặc \((x,y,z)=(0,0,1)\). ⇒
  \(f=x\bar y\bar z\vee x\bar yz\vee xy\bar z\vee xyz\vee\bar x\bar yz\).
  Các hàng còn lại \(000,010,011\) ⇒
  \(\bar f=\bar x\bar y\bar z\vee\bar xy\bar z\vee\bar xyz\)
  \(\big(=\bar x(y\vee\bar z)\big)\).
\item
  \textbf{Kết luận.} \(f\) có 5 từ tối tiểu; \(\bar f\) có 3; tổng
  \(5+3=2^3\).
\item
  \textbf{Kiểm tra.}
  \(\bar f=\overline{x\vee\bar yz}=\bar x\wedge(y\vee\bar z)\) (De
  Morgan) --- khớp {[}CC{]}.
\item
  \textbf{Biến thể (V2).} Số hàm Bool 3 biến: \(2^{2^3}=256\).
\item
  \textbf{Lỗi thường gặp.} Bỏ sót từ tối tiểu khi bổ sung biến.
\item
  \textbf{Liên kết.} \protect\hyperlink{F-BOO-02}{F-BOO-02}.
\end{itemize}

\hypertarget{BOO-02-E01}{%
\subsection{BOO-02-E01 · Tế bào ↔ đơn thức}\label{BOO-02-E01}}

\textbf{D1 · V0 · Đề {[}SV{]} S4-p28--32, p34--37 · Lời giải {[}SV{]}}

\begin{itemize}
\tightlist
\item
  \textbf{Đề.} (a) Vẽ kar của các đơn thức \(x\bar yz\bar t\),
  \(\bar yz\bar t\), \(\bar y\bar t\), \(\bar t\) (4 biến). (b) Tế bào
  gồm 4 ô: hàng \(z\bar t\) và hàng \(\bar z\bar t\), cột \(xy\) và
  \(\bar xy\) là kar của đơn thức nào?
\item
  \textbf{Nhận diện.} \protect\hyperlink{BOO-02}{BOO-02}. \textbf{Kiến
  thức.} \protect\hyperlink{F-BOO-03}{F-BOO-03},
  \protect\hyperlink{F-BOO-04}{F-BOO-04}.
\item
  \textbf{Phương pháp.} Vẽ: \(\mathrm{kar}\) của tích = giao các dải của
  từng từ đơn. Đọc: chiếu tế bào lên hai cạnh.
\item
  \textbf{Lời giải.} (a) \(x\bar yz\bar t\): 1 ô (hàng \(z\bar t\), cột
  \(x\bar y\)). \(\bar yz\bar t\): 2 ô (hàng \(z\bar t\), cột
  \(x\bar y\) và \(\bar x\bar y\) --- kề nhau qua biên).
  \(\bar y\bar t\): 4 ô ở \textbf{4 góc}. \(\bar t\): 8 ô (hàng
  \(z\bar t\) và \(\bar z\bar t\)).

  \begin{enumerate}
  \def\labelenumi{(\alph{enumi})}
  \setcounter{enumi}{1}
  \tightlist
  \item
    Chiếu lên cạnh trên: cột \(xy,\bar xy\) nằm trọn trong \(y\) (không
    trọn trong \(x\) hay \(\bar x\)); chiếu lên cạnh phải: hàng
    \(z\bar t,\bar z\bar t\) nằm trọn trong \(\bar t\) ⇒ đơn thức
    \(y\bar t\).
  \end{enumerate}
\item
  \textbf{Kết luận.} Số ô \(=2^{4-k}\) với \(k\) = số từ đơn: 1, 2, 4,
  8; (b) \(y\bar t\) (khớp S4-p32).
\item
  \textbf{Kiểm tra.}
  \(\mathrm{kar}(\bar y\bar t)=\mathrm{kar}(\bar y)\cap\mathrm{kar}(\bar t)\)
  (định lý S4-p26) = 2 cột biên ∩ 2 hàng biên = 4 góc ✓.
\item
  \textbf{Biến thể (V1).} S4-p37: 2 ô hàng \(\bar zt\), cột
  \(xy,\bar xy\) ⇒ \(y\bar zt\).
\item
  \textbf{Lỗi thường gặp.} Không coi cột đầu--cuối (hay 4 góc) là kề
  nhau --- \protect\hyperlink{T-20}{T-20}.
\item
  \textbf{Liên kết.} \protect\hyperlink{BOO-03}{BOO-03}.
\end{itemize}

\hypertarget{BOO-03-E01}{%
\subsection{BOO-03-E01 · Karnaugh: một phủ tối tiểu duy
nhất}\label{BOO-03-E01}}

\textbf{D2 · V0 · Đề {[}SV{]} S4-p41 · Lời giải {[}SV{]} S4-p42--52
(ảnh), {[}CC{]}}

\begin{itemize}
\tightlist
\item
  \textbf{Đề.} Tìm tất cả công thức đa thức tối tiểu của
  \(f(x,y,z,t)=xyzt\vee x\bar y\vee x\bar z\vee yz\vee xy\bar z\vee xy\bar t\).
\item
  \textbf{Nhận diện.} ``tìm tất cả công thức đa thức tối tiểu'' →
  \protect\hyperlink{BOO-03}{BOO-03}.
\item
  \textbf{Kiến thức.} \protect\hyperlink{F-BOO-04}{F-BOO-04}.
  \textbf{Phương pháp.} \protect\hyperlink{M-22}{M-22}.
\item
  \textbf{Lời giải.} B1 --- kar(f) (10 ô):
\end{itemize}

\begin{longtable}[]{@{}lcccc@{}}
\toprule\noalign{}
& \(x\bar y\) & \(xy\) & \(\bar xy\) & \(\bar x\bar y\) \\
\midrule\noalign{}
\endhead
\bottomrule\noalign{}
\endlastfoot
\(z\bar t\) & 1 & 1 & 1 & 0 \\
\(zt\) & 1 & 1 & 1 & 0 \\
\(\bar zt\) & 1 & 1 & 0 & 0 \\
\(\bar z\bar t\) & 1 & 1 & 0 & 0 \\
\end{longtable}

B2 --- tế bào lớn: \(x\) (2 cột đầu, 8 ô), \(yz\) (cột \(xy,\bar xy\) ×
hàng \(z\bar t,zt\)). B3 --- bắt buộc: ô \(x\bar y\bar z\bar t\) chỉ
thuộc \(x\); ô \(\bar xyz\bar t\) chỉ thuộc \(yz\) ⇒ chọn cả hai. B4 ---
đã phủ hết 10 ô ⇒ phủ tối tiểu duy nhất.

\begin{itemize}
\tightlist
\item
  \textbf{Kết luận.} \(f=x\vee yz\) (duy nhất).
\item
  \textbf{Kiểm tra.} Vét cạn: tế bào lớn \(\{x, yz\}\), phủ tối tiểu duy
  nhất {[}CC{]}. Đại số:
  \(x\bar y\vee x\bar z\vee xy\bar t\vee xyzt\vee xy\bar z=x\) (mọi ô có
  \(x=1\) đều được phủ), cộng \(yz\).
\item
  \textbf{Ghi chú nguồn.} Dấu bù trong biểu thức đọc từ ảnh slide; bản
  đồ Bước 1 của slide (10 ô) xác nhận \(\mathrm{kar}(f)\) như trên.
\item
  \textbf{Biến thể (V1).} Bỏ hạng tử \(yz\): \(f'=x\) (tế bào lớn duy
  nhất là 8 ô cột \(x\)).
\item
  \textbf{Lỗi thường gặp.} Dừng ở các tế bào 4 ô (ví dụ \(x\bar y\),
  \(xy\)) mà không gộp thành 8 ô \(x\) ---
  \protect\hyperlink{T-21}{T-21}.
\item
  \textbf{Liên kết.} \protect\hyperlink{BOO-03-E02}{BOO-03-E02}.
\end{itemize}

\hypertarget{BOO-03-E02}{%
\subsection{BOO-03-E02 · Karnaugh: hai công thức tối
tiểu}\label{BOO-03-E02}}

\textbf{D3 · V0 · Đề {[}SV{]} S4-p53 · Lời giải {[}SV{]} S4-p53--61
(ảnh), {[}CC{]}}

\begin{itemize}
\tightlist
\item
  \textbf{Đề.} Tìm các công thức đa thức tối tiểu của
  \(f=\bar yzt\vee\bar y\bar z\bar t\vee y\bar z\bar t\vee xyzt\vee\bar xz\bar t\).
\item
  \textbf{Nhận diện.} \protect\hyperlink{BOO-03}{BOO-03}. \textbf{Phương
  pháp.} \protect\hyperlink{M-22}{M-22}.
\item
  \textbf{Kiến thức.} \protect\hyperlink{F-BOO-03}{F-BOO-03},
  \protect\hyperlink{F-BOO-04}{F-BOO-04}.
\item
  \textbf{Lời giải.} B1 --- kar(f), đánh số ô như slide:
\end{itemize}

\begin{longtable}[]{@{}lcccc@{}}
\toprule\noalign{}
& \(x\bar y\) & \(xy\) & \(\bar xy\) & \(\bar x\bar y\) \\
\midrule\noalign{}
\endhead
\bottomrule\noalign{}
\endlastfoot
\(z\bar t\) & & & ① & ② \\
\(zt\) & ③ & ④ & & ⑤ \\
\(\bar zt\) & & & & \\
\(\bar z\bar t\) & ⑥ & ⑦ & ⑧ & ⑨ \\
\end{longtable}

B2 --- tế bào lớn: \(\bar z\bar t\) \{⑥⑦⑧⑨\}, \(\bar x\bar t\) \{①②⑧⑨\}
(hàng \(z\bar t\) và \(\bar z\bar t\) kề nhau qua biên), \(xzt\) \{③④\},
\(\bar yzt\) \{③⑤\} (cột đầu--cuối kề nhau), \(\bar x\bar yz\) \{②⑤\}.
B3 --- bắt buộc: ô ⑥ chỉ thuộc \(\bar z\bar t\); ô ① chỉ thuộc
\(\bar x\bar t\); ô ④ chỉ thuộc \(xzt\). Ba tế bào này phủ ①②③④⑥⑦⑧⑨, còn
ô ⑤. B4 --- ô ⑤ thuộc 2 tế bào lớn \(\bar x\bar yz\) và \(\bar yzt\) ⇒ 2
phủ, mỗi phủ 4 đơn thức, số từ đơn bằng nhau (2+2+3+3).

\begin{itemize}
\tightlist
\item
  \textbf{Kết luận.} Hai công thức tối tiểu:
  \(f=\bar z\bar t\vee\bar x\bar t\vee xzt\vee\bar x\bar yz\) và
  \(f=\bar z\bar t\vee\bar x\bar t\vee xzt\vee\bar yzt\).
\item
  \textbf{Kiểm tra.} Vét cạn liệt kê đúng 5 tế bào lớn, 3 bắt buộc, 2
  phủ tối tiểu {[}CC{]}.
\item
  \textbf{Biến thể (V2).} Thêm ô \(\bar x\bar y\bar zt\) vào kar(f): cột
  \(\bar x\bar y\) đủ 4 ô nên \(\bar x\bar y\) thành tế bào lớn
  \textbf{bắt buộc} (chứa ô mới) và phủ luôn ô ⑤ ⇒ công thức tối tiểu
  duy nhất \(\bar z\bar t\vee\bar x\bar t\vee xzt\vee\bar x\bar y\)
  \textbf{{[}CC{]}}.
\item
  \textbf{Lỗi thường gặp.} Bỏ sót tế bào quấn biên \(\bar yzt\) hoặc
  \(\bar x\bar t\) --- \protect\hyperlink{T-20}{T-20}.
\item
  \textbf{Liên kết.} \protect\hyperlink{BOO-03-E03}{BOO-03-E03}.
\end{itemize}

\hypertarget{BOO-03-E03}{%
\subsection{BOO-03-E03 · Karnaugh: hai phủ nhưng chỉ một tối tiểu ---
bẫy}\label{BOO-03-E03}}

\textbf{D3 · V4 · Đề {[}SV{]} S4-p62 · Lời giải {[}SV{]} S4-p63--65,
{[}CC{]}}

\begin{itemize}
\tightlist
\item
  \textbf{Đề.} Xác định các công thức đa thức tối tiểu của
  \(f=xz(\bar y\vee\bar t)\vee\bar x\bar z\bar t\vee z(yt\vee\bar x\bar y)\).
\item
  \textbf{Nhận diện.} \protect\hyperlink{BOO-03}{BOO-03}. \textbf{Phương
  pháp.} \protect\hyperlink{M-22}{M-22} (chú ý Bước 5).
\item
  \textbf{Kiến thức.} \protect\hyperlink{F-BOO-04}{F-BOO-04} --- định
  nghĩa ``đơn giản hơn'', ``tối tiểu''.
\item
  \textbf{Lời giải.} Khai triển:
  \(xz\bar y\vee xz\bar t\vee\bar x\bar z\bar t\vee yzt\vee\bar x\bar yz\).
  Kar(f):
\end{itemize}

\begin{longtable}[]{@{}lcccc@{}}
\toprule\noalign{}
& \(x\bar y\) & \(xy\) & \(\bar xy\) & \(\bar x\bar y\) \\
\midrule\noalign{}
\endhead
\bottomrule\noalign{}
\endlastfoot
\(z\bar t\) & 1 & 1 & 0 & 1 \\
\(zt\) & 1 & 1 & 1 & 1 \\
\(\bar zt\) & 0 & 0 & 0 & 0 \\
\(\bar z\bar t\) & 0 & 0 & 1 & 1 \\
\end{longtable}

Tế bào lớn: \(xz\), \(\bar yz\), \(zt\), \(\bar x\bar z\bar t\),
\(\bar x\bar y\bar t\). Bắt buộc: \(xz\) (ô \(xyz\bar t\)), \(zt\) (ô
\(\bar xyzt\)), \(\bar x\bar z\bar t\) (ô \(\bar xy\bar z\bar t\)). Còn
ô (hàng 1, cột 4) \(=\bar x\bar yz\bar t\): thuộc \(\bar yz\) hoặc
\(\bar x\bar y\bar t\) ⇒ 2 phủ:
\(F_1=xz\vee zt\vee\bar x\bar z\bar t\vee\bar x\bar y\bar t\),
\(\ F_2=xz\vee zt\vee\bar x\bar z\bar t\vee\bar yz\). Bước 5: \(F_2\)
thực sự đơn giản hơn \(F_1\) (\(\bar yz\) có 2 từ đơn \(<3\)) ⇒ loại
\(F_1\).

\begin{itemize}
\tightlist
\item
  \textbf{Kết luận.} Công thức tối tiểu duy nhất:
  \(f=xz\vee zt\vee\bar x\bar z\bar t\vee\bar yz\).
\item
  \textbf{Kiểm tra.} Vét cạn: \(F_1,F_2\) đều biểu diễn \(f\), chỉ
  \(F_2\) tối tiểu {[}CC{]}.
\item
  \textbf{Biến thể (V3).} Bỏ ô \(\bar x\bar yz\bar t\) khỏi kar(f) thì
  chỉ còn 3 tế bào bắt buộc và \(f=xz\vee zt\vee\bar x\bar z\bar t\).
\item
  \textbf{Lỗi thường gặp.} Liệt kê cả hai phủ là ``tối tiểu'' mà không
  so sánh độ đơn giản --- \protect\hyperlink{T-21}{T-21}.
\item
  \textbf{Liên kết.} \protect\hyperlink{F-BOO-04}{F-BOO-04}.
\end{itemize}

\hypertarget{BOO-03-E04}{%
\subsection{BOO-03-E04 · Karnaugh 3 biến}\label{BOO-03-E04}}

\textbf{D1 · V1 · {[}AI-GENERATED PRACTICE{]} · Lời giải {[}CC{]}}

\begin{itemize}
\tightlist
\item
  \textbf{Đề.} \(f(x,y,z)=1\) đúng tại các bộ
  \(xyz\in\{001,010,011,101,111\}\). Tìm công thức đa thức tối tiểu.
\item
  \textbf{Nhận diện.} Cho tập bộ giá trị làm \(f=1\), hỏi công thức tối
  tiểu → \protect\hyperlink{BOO-03}{BOO-03}.
\item
  \textbf{Kiến thức.} \protect\hyperlink{F-BOO-03}{F-BOO-03} (bản đồ 3
  biến), \protect\hyperlink{F-BOO-04}{F-BOO-04}.
\item
  \textbf{Phương pháp.} \protect\hyperlink{M-22}{M-22}.
\item
  \textbf{Lời giải.} Kar(f):
\end{itemize}

\begin{longtable}[]{@{}lcccc@{}}
\toprule\noalign{}
& \(x\bar y\) & \(xy\) & \(\bar xy\) & \(\bar x\bar y\) \\
\midrule\noalign{}
\endhead
\bottomrule\noalign{}
\endlastfoot
\(z\) & 1 & 1 & 1 & 1 \\
\(\bar z\) & 0 & 0 & 1 & 0 \\
\end{longtable}

Tế bào lớn: \(z\) (cả hàng \(z\), 4 ô), \(\bar xy\) (cột \(\bar xy\), 2
ô). Cả hai bắt buộc (ô \(x\bar yz\) chỉ thuộc \(z\); ô \(\bar xy\bar z\)
chỉ thuộc \(\bar xy\)).

\begin{itemize}
\tightlist
\item
  \textbf{Kết luận.} \(f=z\vee\bar xy\).
\item
  \textbf{Kiểm tra.} \(z\vee\bar xy\) bằng 1 đúng tại 5 bộ đã cho
  {[}CC{]}.
\item
  \textbf{Biến thể (V2).} \(\bar f\) (3 ô \(000,100,110\)):
  \(\bar f=x\bar z\vee\bar y\bar z\).
\item
  \textbf{Lỗi thường gặp.} Gom \(\bar xy\bar z\) với \(\bar x\bar yz\)
  (không kề nhau --- khác 2 biến).
\item
  \textbf{Liên kết.} \protect\hyperlink{M-22}{M-22}.
\end{itemize}

\hypertarget{BOO-03-E05}{%
\subsection{BOO-03-E05 · Karnaugh 4 biến có tế bào bốn góc và hai lựa
chọn}\label{BOO-03-E05}}

\textbf{D2 · V3 · {[}AI-GENERATED PRACTICE{]} · Lời giải {[}CC{]}}

\begin{itemize}
\tightlist
\item
  \textbf{Đề.} \(f(x,y,z,t)=1\) tại
  \(xyzt\in\{0000,0010,1000,1010,0101,1101,0111,1111,0001\}\). Tìm các
  công thức đa thức tối tiểu.
\item
  \textbf{Nhận diện.} Như trên, 4 biến →
  \protect\hyperlink{BOO-03}{BOO-03}.
\item
  \textbf{Kiến thức.} \protect\hyperlink{F-BOO-03}{F-BOO-03},
  \protect\hyperlink{F-BOO-04}{F-BOO-04}.
\item
  \textbf{Phương pháp.} \protect\hyperlink{M-22}{M-22}.
\item
  \textbf{Lời giải.} Kar(f):
\end{itemize}

\begin{longtable}[]{@{}lcccc@{}}
\toprule\noalign{}
& \(x\bar y\) & \(xy\) & \(\bar xy\) & \(\bar x\bar y\) \\
\midrule\noalign{}
\endhead
\bottomrule\noalign{}
\endlastfoot
\(z\bar t\) & 1 & 0 & 0 & 1 \\
\(zt\) & 0 & 1 & 1 & 0 \\
\(\bar zt\) & 0 & 1 & 1 & 1 \\
\(\bar z\bar t\) & 1 & 0 & 0 & 1 \\
\end{longtable}

Tế bào lớn: \(\bar y\bar t\) (4 góc), \(yt\) (khối giữa 2×2),
\(\bar x\bar zt\) (hàng \(\bar zt\), cột \(\bar xy,\bar x\bar y\)),
\(\bar x\bar y\bar z\) (cột cuối, hàng \(\bar zt,\bar z\bar t\)). Bắt
buộc: \(\bar y\bar t\), \(yt\). Còn ô \(\bar x\bar y\bar zt\) ⇒ chọn
\(\bar x\bar zt\) hoặc \(\bar x\bar y\bar z\) (cùng 3 từ đơn).

\begin{itemize}
\tightlist
\item
  \textbf{Kết luận.} \(f=\bar y\bar t\vee yt\vee\bar x\bar zt\) hoặc
  \(f=\bar y\bar t\vee yt\vee\bar x\bar y\bar z\).
\item
  \textbf{Kiểm tra.} Vét cạn cho đúng 2 phủ tối tiểu {[}CC{]}.
\item
  \textbf{Biến thể (V4).} Thêm ô \(1001\) (\(x\bar y\bar zt\)): xuất
  hiện tế bào lớn \(\bar y\bar z\) (hàng \(\bar zt,\bar z\bar t\) × cột
  \(x\bar y,\bar x\bar y\)) và \(\bar zt\) (cả hàng \(\bar zt\)) ⇒ hai
  công thức tối tiểu \(\bar y\bar t\vee yt\vee\bar zt\) và
  \(\bar y\bar t\vee yt\vee\bar y\bar z\) \textbf{{[}CC{]}}.
\item
  \textbf{Lỗi thường gặp.} Không nhận ra 4 góc là một tế bào ---
  \protect\hyperlink{T-20}{T-20}.
\item
  \textbf{Liên kết.} \protect\hyperlink{BOO-03-E02}{BOO-03-E02}.
\end{itemize}

\hypertarget{BOO-04-E01}{%
\subsection{BOO-04-E01 · Thiết kế mạch 2 công tắc}\label{BOO-04-E01}}

\textbf{D1 · V0 · Đề {[}SV{]} S4-p75--76 · Lời giải {[}SV{]}}

\begin{itemize}
\tightlist
\item
  \textbf{Đề.} Đèn điều khiển bởi 2 cầu dao \(x,y\) (1 = bật). Đèn sáng
  (\(F=1\)) khi cả hai cùng bật hoặc cùng tắt. Thiết kế mạch.
\item
  \textbf{Nhận diện.} \protect\hyperlink{BOO-04}{BOO-04}. \textbf{Phương
  pháp.} \protect\hyperlink{M-23}{M-23}.
\item
  \textbf{Kiến thức.} \protect\hyperlink{F-BOO-02}{F-BOO-02} (dạng nối
  rời chính tắc), \protect\hyperlink{F-BOO-05}{F-BOO-05} (cổng).
\item
  \textbf{Lời giải.} Bảng: \(F(1,1)=1\), \(F(1,0)=0\), \(F(0,1)=0\),
  \(F(0,0)=1\) ⇒ \(F=xy\vee\bar x\bar y\) (không rút gọn được: hai ô
  không kề). Mạch: AND(\(x,y\)); NOT \(x\), NOT \(y\) →
  AND(\(\bar x,\bar y\)); hai đầu ra → OR.
\item
  \textbf{Kết luận.} \(F=xy\vee\bar x\bar y\) (2 cổng AND, 2 NOT, 1 OR).
\item
  \textbf{Kiểm tra.} Thế 4 bộ giá trị ✓.
\item
  \textbf{Biến thể (V2).} Đèn sáng khi \textbf{đúng một} công tắc bật:
  \(F=x\bar y\vee\bar xy\) (phần bù).
\item
  \textbf{Lỗi thường gặp.} Nhầm \(\bar x\bar y\) với \(\overline{xy}\)
  (NAND) --- \protect\hyperlink{T-22}{T-22}.
\item
  \textbf{Liên kết.} \protect\hyperlink{BOO-04-E02}{BOO-04-E02}.
\end{itemize}

\hypertarget{BOO-04-E02}{%
\subsection{BOO-04-E02 · Thiết kế mạch 3 công tắc (đèn cầu
thang)}\label{BOO-04-E02}}

\textbf{D2 · V1 · Đề {[}SV{]} S4-p77--78 · Lời giải {[}SV{]}, {[}CC{]}}

\begin{itemize}
\tightlist
\item
  \textbf{Đề.} Ba cầu dao \(x,y,z\). \(F(x,y,z)=1\) khi có 1 hoặc 3 cầu
  dao bật. Lập bảng, công thức, mạch; có rút gọn được không?
\item
  \textbf{Nhận diện.} \protect\hyperlink{BOO-04}{BOO-04}. \textbf{Phương
  pháp.} \protect\hyperlink{M-23}{M-23}, \protect\hyperlink{M-22}{M-22}.
\item
  \textbf{Kiến thức.} \protect\hyperlink{F-BOO-02}{F-BOO-02},
  \protect\hyperlink{F-BOO-04}{F-BOO-04},
  \protect\hyperlink{F-BOO-05}{F-BOO-05}.
\item
  \textbf{Lời giải.} \(F=1\) tại \(111,100,010,001\) ⇒
  \(F=xyz\vee x\bar y\bar z\vee\bar xy\bar z\vee\bar x\bar yz\). Trên
  kar(F) 4 ô này đôi một \textbf{không kề} (bàn cờ) ⇒ mỗi ô là một tế
  bào lớn ⇒ công thức trên đã tối tiểu. Mạch: 4 cổng AND 3 ngõ vào (dùng
  NOT cho biến có gạch) → 1 cổng OR 4 ngõ vào.
\item
  \textbf{Kết luận.}
  \(F=xyz\vee x\bar y\bar z\vee\bar xy\bar z\vee\bar x\bar yz\), không
  rút gọn được dạng đa thức.
\item
  \textbf{Kiểm tra.} \(F=1\Leftrightarrow x+y+z\) lẻ (8 hàng) {[}CC{]};
  vét cạn xác nhận phủ tối tiểu là 4 từ tối tiểu.
\item
  \textbf{Ghi chú nguồn.} Slide ghi ``\(F(x,y)\)'' --- đúng phải là
  \(F(x,y,z)\) (ERR-07).
\item
  \textbf{Biến thể (V1).} 3 cầu dao, đèn sáng khi \textbf{đúng 2} cầu
  dao bật: \(F=xy\bar z\vee x\bar yz\vee\bar xyz\).
\item
  \textbf{Lỗi thường gặp.} Cố gộp ô chéo nhau trên bản đồ.
\item
  \textbf{Liên kết.} \protect\hyperlink{BOO-04-E01}{BOO-04-E01}.
\end{itemize}

\hypertarget{BOO-04-E03}{%
\subsection{BOO-04-E03 · Đọc mạch: kết quả là hằng 0 ---
bẫy}\label{BOO-04-E03}}

\textbf{D1 · V4 · Đề {[}SV{]} S4-p74 · Lời giải {[}SV{]}+{[}CC{]}}

\begin{itemize}
\tightlist
\item
  \textbf{Đề.} Sơ đồ: cổng OR với ngõ vào \(x,y,z\); ba cổng NOT với ngõ
  vào \(x,y,z\) nối vào cổng AND; đầu ra OR và đầu ra AND nối vào một
  cổng AND cuối. Viết biểu thức \(f\) và nhận xét.
\item
  \textbf{Nhận diện.} ``Viết biểu thức \(f\) của sơ đồ'' →
  \protect\hyperlink{BOO-04}{BOO-04}.
\item
  \textbf{Phương pháp.} \protect\hyperlink{M-23}{M-23} (đọc mạch từ trái
  sang phải).
\item
  \textbf{Kiến thức.} \protect\hyperlink{F-BOO-05}{F-BOO-05}; De Morgan
  (\protect\hyperlink{F-LOG-03}{F-LOG-03}).
\item
  \textbf{Lời giải.} Đầu ra OR: \(x\vee y\vee z\). Đầu ra AND giữa:
  \(\bar x\bar y\bar z\). Đầu ra cuối:
  \(f=(x\vee y\vee z)\,\bar x\bar y\bar z\). Theo De Morgan
  \(\bar x\bar y\bar z=\overline{x\vee y\vee z}\) ⇒ \(f=u\cdot\bar u=0\)
  với \(u=x\vee y\vee z\).
\item
  \textbf{Kết luận.} \(f\equiv0\) --- mạch luôn cho ra 0 (slide chỉ yêu
  cầu viết biểu thức; nhận xét hằng 0 là phân tích bổ sung, ERR-08).
\item
  \textbf{Kiểm tra.} 8 bộ giá trị đều cho 0 {[}CC{]}.
\item
  \textbf{Biến thể (V2).} Thay cổng AND cuối bằng OR:
  \(f=u\vee\bar u\equiv1\).
\item
  \textbf{Lỗi thường gặp.} Dừng ở biểu thức mà không rút gọn; viết
  \(\overline{xyz}\) thay cho \(\bar x\bar y\bar z\) ---
  \protect\hyperlink{T-22}{T-22}.
\item
  \textbf{Liên kết.} \protect\hyperlink{F-BOO-05}{F-BOO-05}.
\end{itemize}

\hypertarget{BOO-04-E04}{%
\subsection{BOO-04-E04 · Vẽ mạch từ công thức --- và rút gọn trước khi
vẽ}\label{BOO-04-E04}}

\textbf{D2 · V3 · Đề {[}SV{]} S4-p72 · Lời giải {[}SV{]} (mạch), rút gọn
{[}CC{]}}

\begin{itemize}
\tightlist
\item
  \textbf{Đề.} Vẽ mạch cho
  \(f=xz\vee\bar yz\vee xt\vee\bar y\bar t\vee\bar x\bar y\bar z\). Có
  mạch nào ít cổng hơn không?
\item
  \textbf{Nhận diện.} \protect\hyperlink{BOO-04}{BOO-04} +
  \protect\hyperlink{BOO-03}{BOO-03}.
\item
  \textbf{Phương pháp.} \protect\hyperlink{M-23}{M-23},
  \protect\hyperlink{M-22}{M-22}.
\item
  \textbf{Kiến thức.} \protect\hyperlink{F-BOO-04}{F-BOO-04},
  \protect\hyperlink{F-BOO-05}{F-BOO-05}.
\item
  \textbf{Lời giải.} Mạch nguồn: 5 cổng AND (\(xz\); \(\bar y\) qua NOT
  với \(z\); \(xt\); \(\bar y,\bar t\) qua NOT; \(\bar x,\bar y,\bar z\)
  qua NOT) → 1 cổng OR 5 ngõ vào. Rút gọn: mọi ô có \(y=0\) đều thuộc
  kar(f) (\(\bar yz\) phủ \(z=1\); \(\bar y\bar t\) phủ
  \(\bar z\bar t\); ô \(\bar x\bar y\bar zt\) do \(\bar x\bar y\bar z\);
  ô \(x\bar y\bar zt\) do \(xt\)) ⇒ \(\bar y\) là tế bào lớn 8 ô. Tế bào
  lớn: \(\bar y\), \(xz\), \(xt\) --- đều bắt buộc.
\item
  \textbf{Kết luận.} \(f=\bar y\vee xz\vee xt\) --- chỉ cần 1 NOT, 2
  AND, 1 OR.
\item
  \textbf{Kiểm tra.} Vét cạn: hai công thức cùng 11 ô; phủ tối tiểu duy
  nhất \(\bar y\vee xt\vee xz\) {[}CC{]}.
\item
  \textbf{Biến thể (V1).} Có thể viết thêm \(f=\bar y\vee x(z\vee t)\)
  (không còn là công thức đa thức nhưng ít cổng hơn: 1 NOT, 1 OR, 1 AND,
  1 OR).
\item
  \textbf{Lỗi thường gặp.} Vẽ mạch ngay từ công thức chưa rút gọn khi đề
  yêu cầu ``mạch đơn giản''.
\item
  \textbf{Liên kết.} \protect\hyperlink{BOO-03-E01}{BOO-03-E01}.
\end{itemize}

\hypertarget{ex-ch5}{%
\chapter{Ví dụ Chương V --- Đồ thị}\label{ex-ch5}}

\begin{quote}
Hình đồ thị trong slide chỉ có dạng ảnh; toolkit ghi lại \textbf{danh
sách cạnh} đọc từ ảnh để có thể dựng lại. Bảng Dijkstra/Ford--Bellman
được sinh bằng \texttt{10\_Verification/graph\_tables.py} theo đúng quy
trình slide và trùng khớp từng ô với bảng của nguồn \textbf{{[}CC{]}}.
\end{quote}

\hypertarget{GRA-01-E01}{%
\subsection{GRA-01-E01 · Bậc đỉnh: vô hướng có khuyên; có
hướng}\label{GRA-01-E01}}

\textbf{D1 · V0 · Đề {[}SV{]} S5-p20, p21, p24 · Lời giải {[}SV{]},
{[}CC{]}}

\begin{itemize}
\tightlist
\item
  \textbf{Đề.} (a) Đồ thị vô hướng cạnh \(ab, ac, bc, bd, cd\) và một
  khuyên tại \(b\): tìm bậc, phân loại đồ thị. (b) Đồ thị vô hướng cạnh
  \(ab, ac, bc, be, bf, cf\), đỉnh \(d\) cô lập: tìm bậc. (c) Đồ thị có
  hướng cung \(b\to a\), \(a\to c\), \(c\to b\), \(b\to e\), \(b\to f\),
  \(c\to f\) và đỉnh \(d\): tìm bậc vào, bậc ra.
\item
  \textbf{Nhận diện.} ``bậc của các đỉnh'' →
  \protect\hyperlink{GRA-01}{GRA-01}. \textbf{Kiến thức.}
  \protect\hyperlink{F-GRA-01}{F-GRA-01}.
\item
  \textbf{Phương pháp.} Đếm cạnh kề từng đỉnh (khuyên +2; có hướng: đếm
  cung vào/ra); kiểm tra tổng bậc \(=2m\).
\item
  \textbf{Lời giải.} (a) \(\deg a=2\), \(\deg b=3+2=5\) (khuyên đếm 2),
  \(\deg c=3\), \(\deg d=2\); có khuyên ⇒ \textbf{giả đồ thị}. (b)
  \(\deg a=2,\ \deg b=4,\ \deg c=3,\ \deg d=0\) (cô lập), \(\deg e=1\)
  (treo), \(\deg f=2\). (c) \(a\): \(\deg^-=1,\deg^+=1\); \(b\): 1, 3;
  \(c\): 1, 2; \(d\): 0, 0; \(e\): 1, 0; \(f\): 2, 0.
\item
  \textbf{Kết luận.} Như trên (khớp S5-p20, p24).
\item
  \textbf{Kiểm tra.} (a) tổng bậc \(12=2\cdot6\) (5 cạnh + 1 khuyên).
  (b) tổng \(12=2\cdot6\); đỉnh bậc lẻ: \(c,e\) --- 2 đỉnh (chẵn) ✓. (c)
  \(\sum\deg^-=\sum\deg^+=6=m\) ✓.
\item
  \textbf{Biến thể (V1).} Ma trận kề của (c): xem
  \protect\hyperlink{GRA-02-E01}{GRA-02-E01} biến thể.
\item
  \textbf{Lỗi thường gặp.} Đếm khuyên là 1 ---
  \protect\hyperlink{T-23}{T-23}.
\item
  \textbf{Liên kết.} \protect\hyperlink{F-GRA-01}{F-GRA-01}.
\end{itemize}

\hypertarget{GRA-01-E02}{%
\subsection{GRA-01-E02 · Định lý bắt tay: tồn tại đồ thị? số cạnh? ---
bẫy}\label{GRA-01-E02}}

\textbf{D2 · V4 · {[}AI-GENERATED PRACTICE{]} · Lời giải {[}CC{]}}

\begin{itemize}
\tightlist
\item
  \textbf{Đề.} (a) Có đồ thị nào có 5 đỉnh đều bậc 3? (b) Đơn đồ thị 10
  đỉnh, mỗi đỉnh bậc 3 có bao nhiêu cạnh? (c) Có đơn đồ thị 4 đỉnh với
  bậc \(3,3,3,1\) không?
\item
  \textbf{Nhận diện.} ``có tồn tại'', ``số cạnh'' →
  \protect\hyperlink{GRA-01}{GRA-01}. \textbf{Kiến thức.}
  \protect\hyperlink{F-GRA-01}{F-GRA-01}.
\item
  \textbf{Phương pháp.} Định lý bắt tay để bác bỏ; với đơn đồ thị \(n\)
  đỉnh thêm điều kiện \(\deg v\le n-1\) và lập luận kề.
\item
  \textbf{Lời giải.} (a) Tổng bậc \(15\) lẻ, mâu thuẫn \(2m=\sum\deg\) ⇒
  \textbf{không}. (b) \(2m=30\Rightarrow m=15\). (c) Tổng \(10\) chẵn
  nhưng \textbf{vẫn không}: trong đơn đồ thị 4 đỉnh, đỉnh bậc 3 kề cả 3
  đỉnh còn lại; có 3 đỉnh như vậy ⇒ đỉnh thứ tư kề cả 3, bậc
  \(\ge3\ne1\).
\item
  \textbf{Kết luận.} (a) không; (b) 15; (c) không.
\item
  \textbf{Kiểm tra.} Vét cạn mọi đơn đồ thị 4 đỉnh (64 đồ thị): không có
  dãy bậc \((3,3,3,1)\) {[}CC{]}.
\item
  \textbf{Biến thể (V1).} Dãy \(3,3,3,3\): có (\(K_4\), 6 cạnh).
\item
  \textbf{Lỗi thường gặp.} Coi ``tổng bậc chẵn'' là điều kiện đủ.
\item
  \textbf{Liên kết.} \protect\hyperlink{GRA-01-E01}{GRA-01-E01}.
\end{itemize}

\hypertarget{GRA-02-E01}{%
\subsection{GRA-02-E01 · Ma trận kề}\label{GRA-02-E01}}

\textbf{D1 · V0 · Đề {[}SV{]} S5-p27, p28 · Lời giải {[}SV{]}}

\begin{itemize}
\tightlist
\item
  \textbf{Đề.} Tìm ma trận kề của: (a) đồ thị cạnh
  \(ab, ac, bc, bd, cd\); (b) đa/giả đồ thị cạnh \(ab\) (2 cạnh song
  song), \(ac\), \(bc\), \(be\), \(bf\), \(cf\), khuyên tại \(e\); \(d\)
  cô lập.
\item
  \textbf{Nhận diện.} \protect\hyperlink{GRA-02}{GRA-02}. \textbf{Kiến
  thức.} \protect\hyperlink{F-GRA-02}{F-GRA-02}.
\item
  \textbf{Phương pháp.} Cố định thứ tự đỉnh; \(a_{ij}\) = số cạnh giữa
  \(i\) và \(j\); vô hướng thì điền đối xứng.
\item
  \textbf{Lời giải.} (a) Thứ tự \(a,b,c,d\):
  \(\begin{pmatrix}0&1&1&0\\1&0&1&1\\1&1&0&1\\0&1&1&0\end{pmatrix}\).
  (b) Thứ tự \(a,\dots,f\):
  \(\begin{pmatrix}0&2&1&0&0&0\\2&0&1&0&1&1\\1&1&0&0&0&1\\0&0&0&0&0&0\\0&1&0&0&2&0\\0&1&1&0&0&0\end{pmatrix}\).
\item
  \textbf{Kết luận.} Như trên (khớp S5-p27, p28).
\item
  \textbf{Kiểm tra.} Ma trận đối xứng (vô hướng). Với quy ước khuyên =
  2: tổng hàng = bậc (\(\deg b=5\), \(\deg e=3\)).
\item
  \textbf{Ghi chú nguồn --- {[}CONFLICT{]} CON-01.} Định nghĩa S5-p26
  (``số cạnh từ \(i\) đến \(j\)'') cho \(a_{ee}=1\), ví dụ S5-p28 ghi 2.
  Khi làm bài: \textbf{ghi rõ quy ước} đang dùng.
\item
  \textbf{Biến thể (V2).} Đồ thị có hướng
  \protect\hyperlink{GRA-01-E01}{GRA-01-E01}(c), thứ tự \(a..f\): hàng
  \(a\): \((0,0,1,0,0,0)\); \(b\): \((1,0,0,0,1,1)\); \(c\):
  \((0,1,0,0,0,1)\); \(d,e,f\): hàng 0 --- \textbf{không} đối xứng; tổng
  hàng = bậc ra, tổng cột = bậc vào.
\item
  \textbf{Lỗi thường gặp.} Ghi 1 cho cạnh song song; không nêu quy ước
  khuyên --- \protect\hyperlink{T-23}{T-23},
  \protect\hyperlink{T-30}{T-30}.
\item
  \textbf{Liên kết.} \protect\hyperlink{F-GRA-02}{F-GRA-02}.
\end{itemize}

\hypertarget{GRA-03-E01}{%
\subsection{GRA-03-E01 · Đẳng cấu: bác bỏ bằng bất biến, khẳng định bằng
song ánh}\label{GRA-03-E01}}

\textbf{D2 · V0 · Đề (a) {[}SV{]} S5-p32; (b) {[}AI-GENERATED
PRACTICE{]} · Lời giải {[}CC{]}}

\begin{itemize}
\tightlist
\item
  \textbf{Đề.} (a) \(G_1\): cạnh \(ab,bc,cd,de,ea,ce\); \(G_2\): cạnh
  \(ab,bc,ac,ae,ad,cd\). Đẳng cấu không? (b) \(G_1\) như (a) và \(H\):
  đỉnh \(1..5\), cạnh \(12,23,34,45,51,24\). Đẳng cấu không?
\item
  \textbf{Nhận diện.} \protect\hyperlink{GRA-03}{GRA-03}. \textbf{Kiến
  thức.} \protect\hyperlink{F-GRA-03}{F-GRA-03}. \textbf{Phương pháp.}
  \protect\hyperlink{M-24}{M-24}.
\item
  \textbf{Lời giải.} (a) Cùng 5 đỉnh, 6 cạnh; dãy bậc \(G_1\):
  \((3,3,2,2,2)\); \(G_2\): \((4,3,2,2,1)\) --- \(G_2\) có đỉnh bậc 1
  (\(e\)), \(G_1\) không ⇒ \textbf{không đẳng cấu}.

  \begin{enumerate}
  \def\labelenumi{(\alph{enumi})}
  \setcounter{enumi}{1}
  \tightlist
  \item
    Dãy bậc đều \((3,3,2,2,2)\). Ở \(G_1\) hai đỉnh bậc 3 là \(c,e\) (kề
    nhau), đỉnh kề chung của chúng là \(d\); ở \(H\): \(2,4\) (kề nhau),
    kề chung \(3\). Ghép \(c\mapsto2\), \(e\mapsto4\), \(d\mapsto3\);
    \(b\) (kề \(c\)) \(\mapsto1\) (kề 2); \(a\mapsto5\). Kiểm tra 6
    cạnh: \(ab\mapsto51\), \(bc\mapsto12\), \(cd\mapsto23\),
    \(de\mapsto34\), \(ea\mapsto45\), \(ce\mapsto24\) --- đều là cạnh
    của \(H\) ⇒ \textbf{đẳng cấu}.
  \end{enumerate}
\item
  \textbf{Kết luận.} (a) không; (b) có,
  \(f=(a,b,c,d,e)\mapsto(5,1,2,3,4)\).
\item
  \textbf{Kiểm tra.} Thuật toán kiểm tra đẳng cấu của \texttt{networkx}
  xác nhận cả hai kết luận {[}CC{]}.
\item
  \textbf{Biến thể (V4).} Cùng dãy bậc \textbf{không} đủ: \(C_6\) (chu
  trình 6 đỉnh) và hai tam giác rời nhau đều có 6 đỉnh bậc 2, 6 cạnh
  nhưng một liên thông, một không ⇒ không đẳng cấu.
\item
  \textbf{Lỗi thường gặp.} Kết luận đẳng cấu chỉ vì cùng số đỉnh, cạnh,
  dãy bậc --- \protect\hyperlink{T-24}{T-24}.
\item
  \textbf{Liên kết.} \protect\hyperlink{M-24}{M-24}.
\end{itemize}

\hypertarget{GRA-04-E01}{%
\subsection{GRA-04-E01 · Đường đi đơn/sơ cấp, chu trình sơ cấp, đỉnh
khớp, cầu}\label{GRA-04-E01}}

\textbf{D1 · V0 · Đề {[}SV{]} S5-p42 (+ đỉnh khớp, cầu theo định nghĩa
S5-p38) · Lời giải {[}SV{]}+{[}CC{]}}

\begin{itemize}
\tightlist
\item
  \textbf{Đề.} Đồ thị: \(e_1=ab\), \(e_2=bc\), \(e_3=cd\), \(e_4=db\),
  \(e_5=bf\), \(e_6=be\), \(e_7=fe\), \(e_8=eg\). (a)
  \((a,e_1,b,e_2,c,e_3,d,e_4,b)\) là đường đi gì, độ dài? (b) Liệt kê
  các chu trình sơ cấp. (c) Tìm đỉnh khớp và cầu.
\item
  \textbf{Nhận diện.} \protect\hyperlink{GRA-04}{GRA-04}. \textbf{Kiến
  thức.} Định nghĩa S5-p38--41.
\item
  \textbf{Phương pháp.} Áp định nghĩa; với khớp/cầu: xóa thử đỉnh/cạnh
  rồi kiểm tra liên thông.
\item
  \textbf{Lời giải.} (a) Đường đi từ \(a\) tới \(b\) độ dài 4; không lặp
  cạnh ⇒ \textbf{đơn}; lặp đỉnh \(b\) ⇒ \textbf{không sơ cấp}. Viết gọn
  \((a,b,c,d,b)\) (đơn đồ thị). (b) \((b,c,d,b)\), \((b,f,e,b)\). (c)
  Xóa \(b\): \(a\) tách ra ⇒ \(b\) là đỉnh khớp; xóa \(e\): \(g\) tách
  ra ⇒ \(e\) là đỉnh khớp. Cầu: \(e_1=ab\), \(e_8=eg\).
\item
  \textbf{Kết luận.} (a) đơn, không sơ cấp, độ dài 4; (b) 2 chu trình sơ
  cấp; (c) đỉnh khớp \(b,e\); cầu \(ab, eg\).
\item
  \textbf{Kiểm tra.} \texttt{networkx}: articulation points \(\{b,e\}\),
  bridges \(\{ab, eg\}\) {[}CC{]}.
\item
  \textbf{Biến thể (V1).} Cạnh \(bc\) có phải cầu? Không --- nằm trên
  chu trình \((b,c,d,b)\) (cạnh là cầu ⇔ không thuộc chu trình nào
  \textbf{{[}INFERRED{]}}).
\item
  \textbf{Lỗi thường gặp.} Nhầm ``đơn'' (không lặp \textbf{cạnh}) với
  ``sơ cấp'' (không lặp \textbf{đỉnh}).
\item
  \textbf{Liên kết.} \protect\hyperlink{GRA-05}{GRA-05} (Fleury dùng
  khái niệm cầu).
\end{itemize}

\hypertarget{GRA-05-E01}{%
\subsection{GRA-05-E01 · Chu trình Euler bằng Fleury}\label{GRA-05-E01}}

\textbf{D2 · V0 · Đề {[}SV{]} S5-p51 · Lời giải {[}SV{]}, {[}CC{]}}

\begin{itemize}
\tightlist
\item
  \textbf{Đề.} Đồ thị cạnh
  \(ab, bc, cd, ag, bh, bg, hg, gf, fe, cf, ce, df\). Tìm chu trình
  Euler.
\item
  \textbf{Nhận diện.} \protect\hyperlink{GRA-05}{GRA-05}. \textbf{Kiến
  thức.} \protect\hyperlink{F-GRA-04}{F-GRA-04}. \textbf{Phương pháp.}
  \protect\hyperlink{M-25}{M-25}.
\item
  \textbf{Lời giải.} Liên thông. Bậc: \(a2, b4, c4, d2, e2, f4, g4, h2\)
  --- \textbf{tất cả chẵn} ⇒ có chu trình Euler. Fleury từ \(a\):
  \(a\to b\to c\to f\to d\to c\to e\to f\to g\to h\to b\to g\to a\) (mỗi
  bước tránh cầu của đồ thị còn lại; ví dụ khi tới \(f\) lần đầu, cạnh
  \(fd\) chưa là cầu vì \(d\) còn nối với \(c\)).
\item
  \textbf{Kết luận.} Chu trình Euler: \(abcfdcefghbga\) (12 cạnh).
\item
  \textbf{Kiểm tra.} Dãy dùng đúng 12 cạnh, mỗi cạnh một lần, về lại
  \(a\) {[}CC{]}.
\item
  \textbf{Biến thể (V2).} Xóa cạnh \(ab\): \(a\) và \(b\) thành bậc lẻ ⇒
  không còn chu trình Euler nhưng có \textbf{đường đi} Euler từ \(a\)
  tới \(b\), ví dụ \(a\,g\,b\,c\,d\,f\,c\,e\,f\,g\,h\,b\) (11 cạnh)
  {[}CC{]}.
\item
  \textbf{Lỗi thường gặp.} Không kiểm tra liên thông; đi qua cầu sớm làm
  kẹt --- \protect\hyperlink{T-25}{T-25}.
\item
  \textbf{Liên kết.} \protect\hyperlink{GRA-05-E02}{GRA-05-E02}.
\end{itemize}

\hypertarget{GRA-05-E02}{%
\subsection{GRA-05-E02 · Bài toán 7 cây cầu
Königsberg}\label{GRA-05-E02}}

\textbf{D1 · V2 · Đề {[}SV{]} S5-p44--47 · Lời giải {[}INFERRED{]} (bậc
đọc từ hình), {[}CC{]}}

\begin{itemize}
\tightlist
\item
  \textbf{Đề.} 4 vùng đất nối bởi 7 cây cầu. Có thể đi qua mỗi cầu đúng
  một lần không?
\item
  \textbf{Nhận diện.} ``đi qua \ldots{} đúng một lần'' → Euler →
  \protect\hyperlink{GRA-05}{GRA-05}.
\item
  \textbf{Kiến thức.} \protect\hyperlink{F-GRA-04}{F-GRA-04}.
\item
  \textbf{Phương pháp.} \protect\hyperlink{M-25}{M-25}: mô hình đa đồ
  thị, đếm đỉnh bậc lẻ.
\item
  \textbf{Lời giải.} Mô hình đa đồ thị: đỉnh = vùng đất, cạnh = cầu.
  Bậc: đảo giữa 5, ba vùng còn lại đều 3 ⇒ \textbf{4 đỉnh bậc lẻ} ⇒
  không có đường đi Euler (cần 0 hoặc 2 đỉnh lẻ).
\item
  \textbf{Kết luận.} Không thể. Cần ít nhất \(k=2\) nét để vẽ đồ thị.
\item
  \textbf{Kiểm tra.} Tổng bậc \(5+3+3+3=14=2\cdot7\) ✓.
\item
  \textbf{Biến thể (V3).} Xây thêm một cầu nối hai vùng bậc 3 ⇒ còn 2
  đỉnh lẻ ⇒ có đường đi Euler (không có chu trình).
\item
  \textbf{Lỗi thường gặp.} Dùng đơn đồ thị (gộp các cầu song song) làm
  sai bậc.
\item
  \textbf{Liên kết.} \protect\hyperlink{GRA-05-E01}{GRA-05-E01}.
\end{itemize}

\hypertarget{GRA-06-E01}{%
\subsection{GRA-06-E01 · Dijkstra trên đồ thị vô
hướng}\label{GRA-06-E01}}

\textbf{D2 · V0 · Đề {[}SV{]} S5-p60 · Lời giải {[}SV{]} S5-p61--65,
{[}CC{]}}

\begin{itemize}
\tightlist
\item
  \textbf{Đề.} Tìm đường đi ngắn nhất từ \(u\) đến các đỉnh còn lại.
  Cạnh (trọng số): ur 4, uy 1, rs 7, rt 3, ry 2, st 3, sx 1, tx 1, tz 4,
  xw 3, yz 3, zw 5.
\item
  \textbf{Nhận diện.} Trọng số dương →
  \protect\hyperlink{GRA-06}{GRA-06}. \textbf{Kiến thức.}
  \protect\hyperlink{F-GRA-05}{F-GRA-05}. \textbf{Phương pháp.}
  \protect\hyperlink{M-26}{M-26}.
\item
  \textbf{Lời giải.}
\end{itemize}

\begin{longtable}[]{@{}lllllllll@{}}
\toprule\noalign{}
Bước & u & r & s & t & x & y & z & w \\
\midrule\noalign{}
\endhead
\bottomrule\noalign{}
\endlastfoot
0 & 0* & (∞,-) & (∞,-) & (∞,-) & (∞,-) & (∞,-) & (∞,-) & (∞,-) \\
1 & -- & (4,u) & (∞,-) & (∞,-) & (∞,-) & (1,u)* & (∞,-) & (∞,-) \\
2 & -- & (3,y)* & (∞,-) & (∞,-) & (∞,-) & -- & (4,y) & (∞,-) \\
3 & -- & -- & (10,r) & (6,r) & (∞,-) & -- & (4,y)* & (∞,-) \\
4 & -- & -- & (10,r) & (6,r)* & (∞,-) & -- & -- & (9,z) \\
5 & -- & -- & (9,t) & -- & (7,t)* & -- & -- & (9,z) \\
6 & -- & -- & (8,x)* & -- & -- & -- & -- & (9,z) \\
7 & -- & -- & -- & -- & -- & -- & -- & (9,z)* \\
\end{longtable}

\begin{itemize}
\tightlist
\item
  \textbf{Kết luận.} \(d(u,\cdot)\): \(y=1\) (\(u y\)); \(r=3\)
  (\(u y r\)); \(z=4\) (\(u y z\)); \(t=6\) (\(u y r t\)); \(x=7\)
  (\(u y r t x\)); \(s=8\) (\(u y r t x s\)); \(w=9\) (\(u y z w\)).
\item
  \textbf{Kiểm tra.} Trùng bảng S5-p64 và cây đường đi S5-p65;
  \texttt{networkx} cho cùng khoảng cách {[}CC{]}. Bước 3: \(r\) (3) và
  \(z\) (4) --- chọn \(r\) trước vì nhỏ hơn.
\item
  \textbf{Biến thể (V4).} Dijkstra \textbf{sai} khi có trọng số âm: đồ
  thị có hướng \(a\to b\) (2), \(a\to c\) (3), \(c\to b\) (−2): Dijkstra
  cố định \(b=2\) ở bước 1, nhưng thực tế \(d(a,b)=1\) (qua \(c\)) ⇒
  phải dùng Ford--Bellman.
\item
  \textbf{Lỗi thường gặp.} Cập nhật cả đỉnh đã cố định; chọn đỉnh nhỏ
  nhất chỉ trong số đỉnh vừa cập nhật ---
  \protect\hyperlink{T-26}{T-26}, \protect\hyperlink{T-27}{T-27}.
\item
  \textbf{Liên kết.} \protect\hyperlink{GRA-07}{GRA-07}.
\end{itemize}

\hypertarget{GRA-06-E02}{%
\subsection{GRA-06-E02 · Dijkstra trên đồ thị có hướng cho bằng ma trận
trọng số}\label{GRA-06-E02}}

\textbf{D2 · V2 · Đề {[}SV{]} S5-p66--68 · Lời giải {[}SV{]} S5-p69,
{[}CC{]}}

\begin{itemize}
\tightlist
\item
  \textbf{Đề.} \(V=\{v_1,\dots,v_7\}\), ma trận trọng số (hàng = đỉnh
  đầu): \[
  D=\begin{pmatrix}0&5&31&40&\infty&\infty&\infty\\ \infty&0&27&\infty&73&\infty&\infty\\ \infty&26&0&8&49&25&38\\ \infty&\infty&\infty&0&\infty&16&\infty\\ \infty&70&\infty&\infty&0&\infty&9\\ \infty&\infty&\infty&\infty&23&0&12\\ \infty&\infty&\infty&\infty&10&\infty&0\end{pmatrix}
  \]
\end{itemize}

Tìm đường đi ngắn nhất từ \(v_1\) đến các đỉnh khác.

\begin{itemize}
\tightlist
\item
  \textbf{Nhận diện.} Ma trận không đối xứng ⇒ \textbf{có hướng}; trọng
  số dương → \protect\hyperlink{GRA-06}{GRA-06}.
\item
  \textbf{Kiến thức.} \protect\hyperlink{F-GRA-05}{F-GRA-05}; ma trận
  trọng số (\protect\hyperlink{F-GRA-02}{F-GRA-02}).
\item
  \textbf{Phương pháp.} \protect\hyperlink{M-26}{M-26}: đỉnh kề của
  \(u_i\) = các cột \(j\) có \(d_{ij}<\infty\) ở hàng \(u_i\).
\item
  \textbf{Lời giải.}
\end{itemize}

\begin{longtable}[]{@{}llllllll@{}}
\toprule\noalign{}
Bước & v1 & v2 & v3 & v4 & v5 & v6 & v7 \\
\midrule\noalign{}
\endhead
\bottomrule\noalign{}
\endlastfoot
0 & 0* & (∞,-) & (∞,-) & (∞,-) & (∞,-) & (∞,-) & (∞,-) \\
1 & -- & (5,v1)* & (31,v1) & (40,v1) & (∞,-) & (∞,-) & (∞,-) \\
2 & -- & -- & (31,v1)* & (40,v1) & (78,v2) & (∞,-) & (∞,-) \\
3 & -- & -- & -- & (39,v3)* & (78,v2) & (56,v3) & (69,v3) \\
4 & -- & -- & -- & -- & (78,v2) & (55,v4)* & (69,v3) \\
5 & -- & -- & -- & -- & (78,v2) & -- & (67,v6)* \\
6 & -- & -- & -- & -- & (77,v7)* & -- & -- \\
\end{longtable}

\begin{itemize}
\tightlist
\item
  \textbf{Kết luận.} \(v_2\): 5; \(v_3\): 31; \(v_4\): 39
  (\(v_1v_3v_4\)); \(v_6\): 55 (\(v_1v_3v_4v_6\)); \(v_7\): 67
  (\(\dots v_6v_7\)); \(v_5\): 77 (\(v_1v_3v_4v_6v_7v_5\)).
\item
  \textbf{Kiểm tra.} \(31+8+16+12+10=77\); trùng bảng S5-p69 {[}CC{]}. Ở
  bước 2, \(v_3\) qua \(v_2\) là \(5+27=32>31\) nên giữ \((31,v_1)\).
\item
  \textbf{Biến thể (V1).} Bỏ cung \(v_7\to v_5\): \(d(v_1,v_5)=78\) (qua
  \(v_2\) hoặc qua \(v_6\): \(55+23\)) {[}CC{]}.
\item
  \textbf{Lỗi thường gặp.} Đọc ma trận theo cột (đảo hướng cung) ---
  \protect\hyperlink{T-30}{T-30}.
\item
  \textbf{Liên kết.} \protect\hyperlink{GRA-06-E01}{GRA-06-E01}.
\end{itemize}

\hypertarget{GRA-06-E03}{%
\subsection{\texorpdfstring{GRA-06-E03 · Dijkstra: đường đi ngắn nhất từ
\(a\) đến
\(z\)}{GRA-06-E03 · Dijkstra: đường đi ngắn nhất từ a đến z}}\label{GRA-06-E03}}

\textbf{D2 · V1 · Đề {[}SV{]} S5-p71 · Lời giải {[}SV{]} S5-p72,
{[}CC{]}}

\begin{itemize}
\tightlist
\item
  \textbf{Đề.} Đồ thị vô hướng, cạnh (trọng số): ab 4, ac 3, bc 2, bd 5,
  cd 3, ce 6, de 1, df 5, eg 5, fg 2, fz 7, gz 4. Tìm đường đi ngắn nhất
  \(a\to z\) và độ dài.
\item
  \textbf{Nhận diện.} Vô hướng, trọng số dương, hỏi đường đi \(a\to z\)
  → \protect\hyperlink{GRA-06}{GRA-06}.
\item
  \textbf{Kiến thức.} \protect\hyperlink{F-GRA-05}{F-GRA-05}.
\item
  \textbf{Phương pháp.} \protect\hyperlink{M-26}{M-26}.
\item
  \textbf{Lời giải.}
\end{itemize}

\begin{longtable}[]{@{}lllllllll@{}}
\toprule\noalign{}
Bước & a & b & c & d & e & f & g & z \\
\midrule\noalign{}
\endhead
\bottomrule\noalign{}
\endlastfoot
0 & 0* & (∞,-) & (∞,-) & (∞,-) & (∞,-) & (∞,-) & (∞,-) & (∞,-) \\
1 & -- & (4,a) & (3,a)* & (∞,-) & (∞,-) & (∞,-) & (∞,-) & (∞,-) \\
2 & -- & (4,a)* & -- & (6,c) & (9,c) & (∞,-) & (∞,-) & (∞,-) \\
3 & -- & -- & -- & (6,c)* & (9,c) & (∞,-) & (∞,-) & (∞,-) \\
4 & -- & -- & -- & -- & (7,d)* & (11,d) & (∞,-) & (∞,-) \\
5 & -- & -- & -- & -- & -- & (11,d)* & (12,e) & (∞,-) \\
6 & -- & -- & -- & -- & -- & -- & (12,e)* & (18,f) \\
7 & -- & -- & -- & -- & -- & -- & -- & (16,g)* \\
\end{longtable}

\begin{itemize}
\tightlist
\item
  \textbf{Kết luận.} Đường đi \(a\to c\to d\to e\to g\to z\), độ dài
  \(3+3+1+5+4=16\).
\item
  \textbf{Kiểm tra.} Truy vết:
  \(z\leftarrow g\leftarrow e\leftarrow d\leftarrow c\leftarrow a\);
  trùng S5-p72 {[}CC{]}.
\item
  \textbf{Biến thể (V1).} Đường đi ngắn nhất \(a\to f\): \(a\,c\,d\,f\),
  độ dài 11.
\item
  \textbf{Lỗi thường gặp.} Dừng khi lần đầu thấy nhãn của \(z\) (18) ---
  chỉ dừng khi \(z\) được \textbf{cố định} (*).
\item
  \textbf{Liên kết.} \protect\hyperlink{M-26}{M-26}.
\end{itemize}

\hypertarget{GRA-07-E01}{%
\subsection{GRA-07-E01 · Ford--Bellman phát hiện mạch
âm}\label{GRA-07-E01}}

\textbf{D3 · V4 · Đề {[}SV{]} S5-p75 · Lời giải {[}SV{]} S5-p76--83,
{[}CC{]}}

\begin{itemize}
\tightlist
\item
  \textbf{Đề.} Đồ thị có hướng 6 đỉnh, cung: \(1\to2\) (7), \(1\to4\)
  (8), \(2\to3\) (4), \(2\to5\) (2), \(2\to6\) (1), \(4\to2\) (2),
  \(4\to5\) (2), \(5\to6\) (3), \(6\to3\) (2), \(6\to4\) (−6). Tìm đường
  đi ngắn nhất từ 1 hoặc chỉ ra mạch âm.
\item
  \textbf{Nhận diện.} Có trọng số âm →
  \protect\hyperlink{GRA-07}{GRA-07}. \textbf{Kiến thức.}
  \protect\hyperlink{F-GRA-06}{F-GRA-06}. \textbf{Phương pháp.}
  \protect\hyperlink{M-27}{M-27}.
\item
  \textbf{Lời giải.}
\end{itemize}

\begin{longtable}[]{@{}lllllll@{}}
\toprule\noalign{}
k & 1 & 2 & 3 & 4 & 5 & 6 \\
\midrule\noalign{}
\endhead
\bottomrule\noalign{}
\endlastfoot
0 & 0 & (∞,-) & (∞,-) & (∞,-) & (∞,-) & (∞,-) \\
1 & 0 & (7,1) & (∞,-) & (8,1) & (∞,-) & (∞,-) \\
2 & 0 & (7,1) & (11,2) & (8,1) & (9,2) & (8,2) \\
3 & 0 & (7,1) & (10,6) & (2,6) & (9,2) & (8,2) \\
4 & 0 & (4,4) & (10,6) & (2,6) & (4,4) & (8,2) \\
5 & 0 & (4,4) & (8,2) & (2,6) & (4,4) & (5,2) \\
6 & 0 & (4,4) & (7,6) & (−1,6) & (4,4) & (5,2) \\
\end{longtable}

\begin{itemize}
\tightlist
\item
  \textbf{Kết luận.} \(k=n=6\) mà \(L_6\ne L_5\) ⇒ \textbf{có mạch âm},
  ví dụ \(4\to2\to6\to4\) có độ dài \(2+1-6=-3\).
\item
  \textbf{Kiểm tra.} Trùng S5-p76--83; \texttt{networkx} xác nhận có chu
  trình âm {[}CC{]}.
\item
  \textbf{Biến thể (V4).} Đổi \(w(6\to4)=-3\): mạch \(4\to2\to6\to4\) có
  độ dài \(0\) (không âm) ⇒ ổn định tại \(k=5\) với \(L=(0,7,10,5,7,8)\)
  {[}CC{]}.
\item
  \textbf{Lỗi thường gặp.} Dùng giá trị \textbf{cùng dòng} vừa cập nhật
  thay cho dòng \(k-1\) (kết quả trung gian khác bảng) ---
  \protect\hyperlink{T-28}{T-28}.
\item
  \textbf{Liên kết.} \protect\hyperlink{GRA-07-E02}{GRA-07-E02}.
\end{itemize}

\hypertarget{GRA-07-E02}{%
\subsection{GRA-07-E02 · Ford--Bellman hội tụ}\label{GRA-07-E02}}

\textbf{D2 · V0 · Đề {[}SV{]} S5-p84 · Lời giải {[}SV{]} S5-p85--86,
{[}CC{]}}

\begin{itemize}
\tightlist
\item
  \textbf{Đề.} Như \protect\hyperlink{GRA-07-E01}{GRA-07-E01} nhưng
  \(w(6\to4)=-2\).
\item
  \textbf{Phương pháp.} \protect\hyperlink{M-27}{M-27}.
\item
  \textbf{Nhận diện.} Có cung trọng số âm →
  \protect\hyperlink{GRA-07}{GRA-07}.
\item
  \textbf{Kiến thức.} \protect\hyperlink{F-GRA-06}{F-GRA-06}.
\item
  \textbf{Lời giải.}
\end{itemize}

\begin{longtable}[]{@{}lllllll@{}}
\toprule\noalign{}
k & 1 & 2 & 3 & 4 & 5 & 6 \\
\midrule\noalign{}
\endhead
\bottomrule\noalign{}
\endlastfoot
0 & 0 & (∞,-) & (∞,-) & (∞,-) & (∞,-) & (∞,-) \\
1 & 0 & (7,1) & (∞,-) & (8,1) & (∞,-) & (∞,-) \\
2 & 0 & (7,1) & (11,2) & (8,1) & (9,2) & (8,2) \\
3 & 0 & (7,1) & (10,6) & (6,6) & (9,2) & (8,2) \\
4 & 0 & (7,1) & (10,6) & (6,6) & (8,4) & (8,2) \\
5 & 0 & (7,1) & (10,6) & (6,6) & (8,4) & (8,2) \\
\end{longtable}

\begin{itemize}
\tightlist
\item
  \textbf{Kết luận.} Ổn định tại \(k=5\). \(d(1,2)=7\) (\(1\,2\));
  \(d(1,6)=8\) (\(1\,2\,6\)); \(d(1,3)=10\) (\(1\,2\,6\,3\));
  \(d(1,4)=6\) (\(1\,2\,6\,4\)); \(d(1,5)=8\) (\(1\,2\,6\,4\,5\)). Cây
  đường đi như S5-p86.
\item
  \textbf{Kiểm tra.} Mạch \(4\to2\to6\to4\) dài \(+1>0\);
  \texttt{networkx} Bellman--Ford cho cùng khoảng cách {[}CC{]}.
\item
  \textbf{Biến thể (V2).} Hỏi riêng \(d(1,5)\): không dừng ở \((9,2)\)
  của \(k=2\) --- giá trị cuối là 8.
\item
  \textbf{Lỗi thường gặp.} Dừng ở \(k\) đầu tiên một cột không đổi (phải
  là \textbf{cả dòng} không đổi).
\item
  \textbf{Liên kết.} \protect\hyperlink{F-GRA-06}{F-GRA-06}.
\end{itemize}

\hypertarget{traps}{%
\chapter{Danh mục bẫy \& lỗi thường gặp (Common Traps)}\label{traps}}

Mỗi mục: \textbf{Tình huống · Vì sao sai · Cách sửa · Đối chiếu · Liên
quan · Căn cứ}. Căn cứ là định nghĩa/định lý trong nguồn (ghi vị trí)
hoặc phản ví dụ đã kiểm tra.

\hypertarget{T-01}{%
\subsection{T-01 · Kéo theo có tiền đề sai; câu chứa biến}\label{T-01}}

\begin{itemize}
\tightlist
\item
  \textbf{Tình huống:} đánh giá ``Nếu \(3>4\) thì trời mưa'' là sai; coi
  ``\(n\) là số tự nhiên'' là mệnh đề.
\item
  \textbf{Vì sao sai:} \(p\to q\) chỉ sai khi \(p=1,q=0\); câu chứa biến
  tự do là vị từ.
\item
  \textbf{Cách sửa:} tra bảng \protect\hyperlink{F-LOG-01}{F-LOG-01};
  kiểm tra biến tự do trước. \textbf{Liên quan:}
  \protect\hyperlink{LOG-01}{LOG-01},
  \protect\hyperlink{LOG-02}{LOG-02}. \textbf{Căn cứ:} S1-p14, p59.
\end{itemize}

\hypertarget{T-02}{%
\subsection{T-02 · De Morgan và phân phối sai chiều}\label{T-02}}

\begin{itemize}
\tightlist
\item
  \textbf{Tình huống:} \(\neg(p\wedge q)\to\neg p\wedge\neg q\); gom
  \((p\vee r)\wedge(\neg q\vee r)\) thành \((p\vee\neg q)\wedge r\).
\item
  \textbf{Vì sao sai:} De Morgan đổi \(\wedge\leftrightarrow\vee\); phần
  tử chung \(r\) phải đứng ngoài với phép \textbf{còn lại}.
\item
  \textbf{Cách sửa:} viết đúng
  \(\neg(p\wedge q)\Leftrightarrow\neg p\vee\neg q\);
  \((a\vee r)\wedge(b\vee r)\Leftrightarrow(a\wedge b)\vee r\).
  \textbf{Đối chiếu:} \(p=1,q=0\): \(\neg(p\wedge q)=1\) nhưng
  \(\neg p\wedge\neg q=0\). \textbf{Liên quan:}
  \protect\hyperlink{LOG-03}{LOG-03},
  \protect\hyperlink{SET-01}{SET-01}. \textbf{Căn cứ:} S1-p26--27.
\end{itemize}

\hypertarget{T-03}{%
\subsection{T-03 · Đảo, phản đảo, phủ định của kéo theo}\label{T-03}}

\begin{itemize}
\tightlist
\item
  \textbf{Tình huống:} từ \(p\to r\) suy ra \(\neg p\to\neg r\); phủ
  định \(a\to b\) thành \(\neg a\to\neg b\) hoặc \(a\to\neg b\).
\item
  \textbf{Vì sao sai:} chỉ phản đảo \(\neg r\to\neg p\) tương đương
  \(p\to r\); \(\neg(a\to b)\Leftrightarrow a\wedge\neg b\).
\item
  \textbf{Cách sửa:} dùng \protect\hyperlink{F-LOG-03}{F-LOG-03} luật
  11. \textbf{Đối chiếu:} \(p=0,r=1\): \(p\to r=1\) nhưng
  \(\neg p\to\neg r=0\). \textbf{Liên quan:}
  \protect\hyperlink{LOG-04}{LOG-04},
  \protect\hyperlink{QNT-02}{QNT-02}. \textbf{Căn cứ:} S1-p29.
\end{itemize}

\hypertarget{T-04}{%
\subsection{T-04 · Khẳng định hệ quả / phủ định tiền đề}\label{T-04}}

\begin{itemize}
\tightlist
\item
  \textbf{Tình huống:} \(p\to q,\ q\ \therefore p\) hoặc
  \(p\to q,\ \neg p\ \therefore\neg q\).
\item
  \textbf{Vì sao sai:} không phải hằng đúng: \(p=0,q=1\) làm tiền đề
  đúng, kết luận sai.
\item
  \textbf{Cách sửa:} chỉ dùng 9 quy tắc
  \protect\hyperlink{F-LOG-04}{F-LOG-04}. \textbf{Liên quan:}
  \protect\hyperlink{LOG-04}{LOG-04},
  \protect\hyperlink{LOG-05}{LOG-05}. \textbf{Căn cứ:} S1-p32--48 + phản
  ví dụ {[}CC{]}.
\end{itemize}

\hypertarget{T-05}{%
\subsection{T-05 · ``Có dấu phủ định nên là hằng sai''}\label{T-05}}

\begin{itemize}
\tightlist
\item
  \textbf{Tình huống:} kết luận \(p\wedge\neg(p\to q)\) là hằng sai.
\item
  \textbf{Vì sao sai:}
  \(p\wedge\neg(p\to q)\Leftrightarrow p\wedge\neg q\) đúng tại
  \(p=1,q=0\).
\item
  \textbf{Cách sửa:} rút gọn hoặc lập đủ bảng. \textbf{Liên quan:}
  \protect\hyperlink{LOG-02-E02}{LOG-02-E02}. \textbf{Căn cứ:} S1-p21
  (ERR-13).
\end{itemize}

\hypertarget{T-06}{%
\subsection{T-06 · Đổi thứ tự lượng từ; bỏ qua miền}\label{T-06}}

\begin{itemize}
\tightlist
\item
  \textbf{Tình huống:} coi \(\forall x\exists y\) như
  \(\exists y\forall x\); quên \(\mathbb N\) không có số âm.
\item
  \textbf{Vì sao sai:} chỉ có
  \(\exists x\forall y\Rightarrow\forall y\exists x\) (S1-p68); miền
  thay đổi chân trị.
\item
  \textbf{Cách sửa:} đọc từ trái sang; ở \(\forall x\exists y\), \(y\)
  được phụ thuộc \(x\). \textbf{Đối chiếu:} ``\(x+y=0\)'' trên
  \(\mathbb Z\) --- \protect\hyperlink{QNT-01-E02}{QNT-01-E02}.
  \textbf{Liên quan:} \protect\hyperlink{QNT-01}{QNT-01}.
\end{itemize}

\hypertarget{T-07}{%
\subsection{T-07 · Phủ định lượng từ: phủ định cả miền}\label{T-07}}

\begin{itemize}
\tightlist
\item
  \textbf{Tình huống:} phủ định \(\forall\varepsilon>0\) thành
  \(\exists\varepsilon\le0\).
\item
  \textbf{Vì sao sai:} miền không đổi; chỉ đổi lượng từ và phủ định vị
  từ.
\item
  \textbf{Cách sửa:} \protect\hyperlink{M-07}{M-07}. \textbf{Liên quan:}
  \protect\hyperlink{QNT-02-E01}{QNT-02-E01}. \textbf{Căn cứ:}
  S1-p69--71.
\end{itemize}

\hypertarget{T-08}{%
\subsection{T-08 · Ánh xạ: bỏ qua miền, tập đích, điểm cực
trị}\label{T-08}}

\begin{itemize}
\tightlist
\item
  \textbf{Tình huống:} ``\(x^2+1\) đơn ánh'' mà không xét miền;
  \(f([-1,3])=[f(-1),f(3)]\); giải \(f(x)=y\) được nghiệm ngoài miền.
\item
  \textbf{Vì sao sai:} tính chất ánh xạ phụ thuộc \(X,Y\); ảnh của đoạn
  chứa điểm cực trị không phải đoạn nối hai đầu mút.
\item
  \textbf{Cách sửa:} \protect\hyperlink{M-09}{M-09}; xét tính đơn điệu
  từng khúc. \textbf{Liên quan:} \protect\hyperlink{FUN-01}{FUN-01},
  \protect\hyperlink{FUN-02}{FUN-02}. \textbf{Căn cứ:} S1-p87--90.
\end{itemize}

\hypertarget{T-09}{%
\subsection{\texorpdfstring{T-09 · Quy nạp: sai bước cơ sở, sai số hạng
thứ \(k+1\), không dùng giả
thiết}{T-09 · Quy nạp: sai bước cơ sở, sai số hạng thứ k+1, không dùng giả thiết}}\label{T-09}}

\begin{itemize}
\tightlist
\item
  \textbf{Tình huống:} bắt đầu \(n=1\) cho mệnh đề ``\(n>3\)''; số hạng
  cuối của \(P(k+1)\) viết sai.
\item
  \textbf{Cách sửa:} \(N_0\) theo đề; viết rõ \(P(k+1)\) trước khi biến
  đổi; chỉ ra chỗ dùng giả thiết quy nạp. \textbf{Liên quan:}
  \protect\hyperlink{IND-01}{IND-01}. \textbf{Căn cứ:} S1-p97--100.
\end{itemize}

\hypertarget{T-10}{%
\subsection{T-10 · Chữ số đầu bằng 0; không tách trường hợp chữ số
0}\label{T-10}}

\begin{itemize}
\tightlist
\item
  \textbf{Tình huống:} đếm số chẵn 3 chữ số khác nhau từ \(\{0..5\}\)
  bằng \(3\cdot5\cdot4\).
\item
  \textbf{Vì sao sai:} khi chữ số cuối là 0, chữ số đầu có 5 lựa chọn;
  khi khác 0, chỉ còn 4.
\item
  \textbf{Cách sửa:} \protect\hyperlink{M-11}{M-11}. \textbf{Liên quan:}
  \protect\hyperlink{CNT-01-E01}{CNT-01-E01}. \textbf{Căn cứ:} S2-p5.
\end{itemize}

\hypertarget{T-11}{%
\subsection{T-11 · Nhầm có thứ tự / không thứ tự; có lặp / không
lặp}\label{T-11}}

\begin{itemize}
\tightlist
\item
  \textbf{Tình huống:} dùng \(C_n^k\) cho bài lập số; dùng \(A_n^k\) khi
  đề không nói ``khác nhau''.
\item
  \textbf{Cách sửa:} bảng \protect\hyperlink{M-14}{M-14}; ghi rõ giả
  thiết khi đề mơ hồ. \textbf{Liên quan:}
  \protect\hyperlink{CNT-04-E01}{CNT-04-E01}. \textbf{Căn cứ:}
  S2-p13--16.
\end{itemize}

\hypertarget{T-12}{%
\subsection{\texorpdfstring{T-12 · Tổ hợp lặp: đảo \(n,k\); ``\(>\)''
thành ``\(\ge\)''
sai}{T-12 · Tổ hợp lặp: đảo n,k; ``\textgreater'' thành ``\textbackslash ge'' sai}}\label{T-12}}

\begin{itemize}
\tightlist
\item
  \textbf{Tình huống:} \(K_n^k\) tính thành \(C_{n+k-1}^{n}\) với
  \(n,k\) đảo; \(x>4\) đổi thành \(x\ge4\).
\item
  \textbf{Cách sửa:} \(n\) = số loại/ẩn/hộp, \(k\) = vế phải; với số
  nguyên \(x>4\Leftrightarrow x\ge5\). \textbf{Liên quan:}
  \protect\hyperlink{CNT-06}{CNT-06}. \textbf{Căn cứ:} S2-p19--24.
\end{itemize}

\hypertarget{T-13}{%
\subsection{\texorpdfstring{T-13 · Vế phải là đa thức thì \(\beta=1\);
viết lại lũy
thừa}{T-13 · Vế phải là đa thức thì \textbackslash beta=1; viết lại lũy thừa}}\label{T-13}}

\begin{itemize}
\tightlist
\item
  \textbf{Tình huống:} coi \(12n+8\) là ``không có \(\beta\)'' và quên
  xét \(\beta=1\) có là nghiệm PTĐT; với \(2\cdot5^{n+1}\) lấy hệ số 2.
\item
  \textbf{Cách sửa:} \(P(n)=1^nP(n)\); \(2\cdot5^{n+1}=10\cdot5^n\);
  \((2-n)2^{n-2}=\frac{2-n}{4}2^n\). \textbf{Liên quan:}
  \protect\hyperlink{REC-03}{REC-03}. \textbf{Căn cứ:} S2-p48.
\end{itemize}

\hypertarget{T-14}{%
\subsection{\texorpdfstring{T-14 · Tìm hằng số \(C_1,C_2\) trước khi
cộng nghiệm
riêng}{T-14 · Tìm hằng số C\_1,C\_2 trước khi cộng nghiệm riêng}}\label{T-14}}

\begin{itemize}
\tightlist
\item
  \textbf{Vì sao sai:} điều kiện đầu áp cho nghiệm của hệ \textbf{không
  thuần nhất}.
\item
  \textbf{Cách sửa:} \(x_n=\bar x_n+x_n^*\) rồi mới thế \(x_0,x_1\).
  \textbf{Đối chiếu:} \protect\hyperlink{REC-03-E02}{REC-03-E02}: dùng
  sai thứ tự cho \(C_2=-2\) thay vì \(-5\). \textbf{Căn cứ:} S2-p47,
  p59.
\end{itemize}

\hypertarget{T-15}{%
\subsection{\texorpdfstring{T-15 · Bội của \(\beta\) và dạng nghiệm
kép}{T-15 · Bội của \textbackslash beta và dạng nghiệm kép}}\label{T-15}}

\begin{itemize}
\tightlist
\item
  \textbf{Tình huống:} chọn \(x_n^*=\beta^nQ_r(n)\) khi \(\beta\) là
  nghiệm (hệ vô nghiệm); với nghiệm kép viết
  \(C_1\lambda^n+C_2\lambda^n\).
\item
  \textbf{Cách sửa:} nhân \(n^m\) (\(m\) = bội của \(\beta\), xét
  \textbf{riêng từng} \(\beta\)); nghiệm kép: \((C_1+nC_2)\lambda_0^n\).
  \textbf{Liên quan:} \protect\hyperlink{REC-02}{REC-02},
  \protect\hyperlink{REC-03}{REC-03},
  \protect\hyperlink{REC-04}{REC-04}. \textbf{Căn cứ:} S2-p43, p49--50.
\end{itemize}

\hypertarget{T-16}{%
\subsection{T-16 · Tính chất quan hệ: chứng minh bằng ví dụ, bỏ sót
cặp}\label{T-16}}

\begin{itemize}
\tightlist
\item
  \textbf{Tình huống:} thử vài số rồi kết luận ``bắc cầu''; quên cặp
  \((a,b),(b,a)\Rightarrow(a,a)\).
\item
  \textbf{Cách sửa:} ``có'' → chứng minh tổng quát; ``không'' → một phản
  ví dụ. Kiểm tra bắc cầu với \textbf{mọi} cặp nối nhau. \textbf{Liên
  quan:} \protect\hyperlink{REL-02}{REL-02}. \textbf{Căn cứ:} S3-p7--11.
\end{itemize}

\hypertarget{T-17}{%
\subsection{\texorpdfstring{T-17 · Phản xứng ≠ ``không đối xứng'';
\((\mathbb Z,\mid)\) không là
poset}{T-17 · Phản xứng ≠ ``không đối xứng''; (\textbackslash mathbb Z,\textbackslash mid) không là poset}}\label{T-17}}

\begin{itemize}
\tightlist
\item
  \textbf{Tình huống:} cho rằng quan hệ không đối xứng thì phản xứng
  (hoặc ngược lại không thể vừa đối xứng vừa phản xứng).
\item
  \textbf{Vì sao sai:} ``='' vừa đối xứng vừa phản xứng;
  \(\{(1,2),(2,1),(1,3)\}\) không đối xứng, không phản xứng. Trên
  \(\mathbb Z\): \(3\mid-3\), \(-3\mid3\).
\item
  \textbf{Liên quan:} \protect\hyperlink{REL-02}{REL-02},
  \protect\hyperlink{POS-01}{POS-01}. \textbf{Căn cứ:} S3-p9, p32.
\end{itemize}

\hypertarget{T-18}{%
\subsection{T-18 · Hasse: vẽ cung bắc cầu; tối đại ≠ lớn
nhất}\label{T-18}}

\begin{itemize}
\tightlist
\item
  \textbf{Cách sửa:} chỉ vẽ trội trực tiếp; ``lớn nhất'' phải so sánh
  được và lớn hơn \textbf{mọi} phần tử --- có nhiều tối đại thì không có
  lớn nhất. \textbf{Đối chiếu:}
  \protect\hyperlink{POS-03-E01}{POS-03-E01}. \textbf{Căn cứ:}
  S3-p42--49.
\end{itemize}

\hypertarget{T-19}{%
\subsection{T-19 · Chặn trên/dưới và sup/inf}\label{T-19}}

\begin{itemize}
\tightlist
\item
  \textbf{Tình huống:} đòi chặn phải thuộc \(A\); chọn sup là một chặn
  trên bất kỳ; khẳng định sup luôn tồn tại.
\item
  \textbf{Cách sửa:} chặn thuộc \(S\); sup = chặn trên \textbf{nhỏ nhất}
  (so sánh được với mọi chặn trên); có thể không tồn tại. \textbf{Đối
  chiếu:} \protect\hyperlink{POS-04-E01}{POS-04-E01}. \textbf{Căn cứ:}
  S3-p50--53.
\end{itemize}

\hypertarget{T-20}{%
\subsection{T-20 · Karnaugh: quên quấn biên và 4 góc}\label{T-20}}

\begin{itemize}
\tightlist
\item
  \textbf{Cách sửa:} cột \(x\bar y\) kề cột \(\bar x\bar y\); hàng
  \(z\bar t\) kề hàng \(\bar z\bar t\); 4 góc là một tế bào
  (\(\bar y\bar t\)). \textbf{Liên quan:}
  \protect\hyperlink{BOO-02}{BOO-02},
  \protect\hyperlink{BOO-03}{BOO-03}. \textbf{Căn cứ:} S4-p25, p30.
\end{itemize}

\hypertarget{T-21}{%
\subsection{T-21 · Karnaugh: tế bào chưa ``lớn''; không so sánh độ đơn
giản}\label{T-21}}

\begin{itemize}
\tightlist
\item
  \textbf{Tình huống:} dùng hai tế bào 4 ô thay cho một tế bào 8 ô; liệt
  kê mọi phủ là tối tiểu.
\item
  \textbf{Cách sửa:} kiểm tra từng tế bào có mở rộng gấp đôi được không;
  Bước 5 loại công thức có công thức khác thực sự đơn giản hơn.
  \textbf{Đối chiếu:} \protect\hyperlink{BOO-03-E03}{BOO-03-E03}.
  \textbf{Căn cứ:} S4-p33, p40, p65.
\end{itemize}

\hypertarget{T-22}{%
\subsection{T-22 · Bù của tích ≠ tích các bù; mạch hằng}\label{T-22}}

\begin{itemize}
\tightlist
\item
  \textbf{Tình huống:} \(\overline{xy}\) (NAND) nhầm với
  \(\bar x\bar y\); không nhận ra
  \((x\vee y\vee z)\bar x\bar y\bar z\equiv0\).
\item
  \textbf{Cách sửa:} De Morgan: \(\overline{xy}=\bar x\vee\bar y\),
  \(\overline{x\vee y}=\bar x\bar y\). \textbf{Liên quan:}
  \protect\hyperlink{BOO-04}{BOO-04}. \textbf{Căn cứ:} S4-p70--71, p74.
\end{itemize}

\hypertarget{T-23}{%
\subsection{T-23 · Bậc và ma trận kề: khuyên, cạnh song
song}\label{T-23}}

\begin{itemize}
\tightlist
\item
  \textbf{Cách sửa:} khuyên cộng 2 vào bậc (vô hướng); cạnh song song:
  \(a_{ij}\) = số cạnh. \textbf{Kiểm tra:} \(\sum\deg=2m\). \textbf{Liên
  quan:} \protect\hyperlink{GRA-01}{GRA-01},
  \protect\hyperlink{GRA-02}{GRA-02}. \textbf{Căn cứ:} S5-p19, p26--28.
\end{itemize}

\hypertarget{T-24}{%
\subsection{T-24 · Đẳng cấu: bất biến chỉ là điều kiện cần}\label{T-24}}

\begin{itemize}
\tightlist
\item
  \textbf{Cách sửa:} bất biến khác ⇒ không đẳng cấu; bất biến giống ⇒
  phải dựng song ánh và kiểm tra \textbf{mọi} cạnh. \textbf{Đối chiếu:}
  \(C_6\) và 2 tam giác (\protect\hyperlink{GRA-03-E01}{GRA-03-E01}).
  \textbf{Căn cứ:} S5-p29--30.
\end{itemize}

\hypertarget{T-25}{%
\subsection{T-25 · Euler: quên liên thông, đếm sai đỉnh lẻ, đi qua cầu
sớm}\label{T-25}}

\begin{itemize}
\tightlist
\item
  \textbf{Cách sửa:} \protect\hyperlink{M-25}{M-25}: liên thông → đếm
  đỉnh lẻ (0: chu trình, 2: đường đi bắt đầu ở đỉnh lẻ) → Fleury tránh
  cầu. \textbf{Liên quan:} \protect\hyperlink{GRA-05}{GRA-05}.
  \textbf{Căn cứ:} S5-p49--50.
\end{itemize}

\hypertarget{T-26}{%
\subsection{T-26 · Dùng Dijkstra khi có trọng số âm}\label{T-26}}

\begin{itemize}
\tightlist
\item
  \textbf{Vì sao sai:} Dijkstra cố định đỉnh dựa trên giả thiết trọng số
  dương (S5-p56). \textbf{Đối chiếu:} \(a\to b\) (2), \(a\to c\) (3),
  \(c\to b\) (−2): Dijkstra cho 2, đúng là 1. \textbf{Cách sửa:} dùng
  Ford--Bellman. \textbf{Liên quan:} \protect\hyperlink{GRA-06}{GRA-06},
  \protect\hyperlink{GRA-07}{GRA-07}.
\end{itemize}

\hypertarget{T-27}{%
\subsection{T-27 · Dijkstra: cập nhật đỉnh đã cố định; chọn min sai phạm
vi; dừng sớm}\label{T-27}}

\begin{itemize}
\tightlist
\item
  \textbf{Cách sửa:} chỉ cập nhật đỉnh \textbf{chưa cố định} kề với đỉnh
  vừa cố định; chọn min trên \textbf{toàn bộ} đỉnh chưa cố định; với bài
  \(u_0\to z\) chỉ dừng khi \(z\) được cố định. \textbf{Liên quan:}
  \protect\hyperlink{GRA-06}{GRA-06}. \textbf{Căn cứ:} S5-p59.
\end{itemize}

\hypertarget{T-28}{%
\subsection{T-28 · Ford--Bellman: dùng giá trị cùng dòng; điều kiện
dừng}\label{T-28}}

\begin{itemize}
\tightlist
\item
  \textbf{Cách sửa:} \(L_k\) chỉ dùng giá trị \(L_{k-1}\); dừng khi
  \textbf{cả dòng} không đổi hoặc \(k=n\) (mạch âm). \textbf{Liên quan:}
  \protect\hyperlink{GRA-07}{GRA-07}. \textbf{Căn cứ:} S5-p73--74.
\end{itemize}

\hypertarget{T-29}{%
\subsection{T-29 · Thứ tự từ điển: so độ dài trước}\label{T-29}}

\begin{itemize}
\tightlist
\item
  \textbf{Cách sửa:} so ký tự tại vị trí khác nhau đầu tiên; tiền tố nhỏ
  hơn (\(0110\prec10\), \(0110\prec01100\)). \textbf{Liên quan:}
  \protect\hyperlink{POS-02}{POS-02}. \textbf{Căn cứ:} S3-p39--41.
\end{itemize}

\hypertarget{T-30}{%
\subsection{T-30 · Ma trận kề/trọng số: đọc sai hướng; không nêu quy ước
khuyên}\label{T-30}}

\begin{itemize}
\tightlist
\item
  \textbf{Cách sửa:} hàng = đỉnh đầu, cột = đỉnh cuối; ma trận không đối
  xứng ⇒ có hướng; ghi rõ quy ước khuyên (CON-01). \textbf{Liên quan:}
  \protect\hyperlink{GRA-02}{GRA-02},
  \protect\hyperlink{GRA-06-E02}{GRA-06-E02}. \textbf{Căn cứ:}
  S5-p26--28, p53.
\end{itemize}

\hypertarget{T-31}{%
\subsection{T-31 · Áp dụng công thức ngoài phạm vi slide}\label{T-31}}

\begin{itemize}
\tightlist
\item
  \textbf{Tình huống:} PTĐT có nghiệm phức; hệ thức cấp 3; Floyd;
  Quine--McCluskey.
\item
  \textbf{Cách sửa:} slide chỉ có nghiệm thực, cấp \(\le2\),
  Dijkstra/Ford--Bellman, Karnaugh (GAP-05/06/07). Nếu gặp trong đề,
  trình bày bằng kiến thức bổ sung và ghi rõ là ngoài slide.
  \textbf{Liên quan:} \protect\hyperlink{REC-02}{REC-02},
  \protect\hyperlink{GRA-06}{GRA-06},
  \protect\hyperlink{BOO-03}{BOO-03}.
\end{itemize}

\hypertarget{quickref}{%
\chapter{Tra cứu nhanh (Quick Reference)}\label{quickref}}

\hypertarget{muxe3-dux1ea1ng-buxe0i-mux1ed9t-duxf2ng}{%
\section{Mã dạng bài --- một
dòng}\label{muxe3-dux1ea1ng-buxe0i-mux1ed9t-duxf2ng}}

\begin{longtable}[]{@{}
  >{\raggedright\arraybackslash}p{(\columnwidth - 6\tabcolsep) * \real{0.2500}}
  >{\raggedright\arraybackslash}p{(\columnwidth - 6\tabcolsep) * \real{0.2500}}
  >{\raggedright\arraybackslash}p{(\columnwidth - 6\tabcolsep) * \real{0.2500}}
  >{\raggedright\arraybackslash}p{(\columnwidth - 6\tabcolsep) * \real{0.2500}}@{}}
\toprule\noalign{}
\begin{minipage}[b]{\linewidth}\raggedright
Mã
\end{minipage} & \begin{minipage}[b]{\linewidth}\raggedright
Dạng
\end{minipage} & \begin{minipage}[b]{\linewidth}\raggedright
Mã
\end{minipage} & \begin{minipage}[b]{\linewidth}\raggedright
Dạng
\end{minipage} \\
\midrule\noalign{}
\endhead
\bottomrule\noalign{}
\endlastfoot
\protect\hyperlink{LOG-01}{LOG-01} & Mệnh đề? sơ cấp/phức hợp &
\protect\hyperlink{CNT-05}{CNT-05} & Hoán vị lặp \\
\protect\hyperlink{LOG-02}{LOG-02} & Bảng chân trị, hằng đúng/sai &
\protect\hyperlink{CNT-06}{CNT-06} & Tổ hợp lặp, nghiệm nguyên \\
\protect\hyperlink{LOG-03}{LOG-03} & Tương đương / rút gọn &
\protect\hyperlink{REC-01}{REC-01} & Lập hệ thức đệ qui \\
\protect\hyperlink{LOG-04}{LOG-04} & Chứng minh suy luận &
\protect\hyperlink{REC-02}{REC-02} & Giải thuần nhất \\
\protect\hyperlink{LOG-05}{LOG-05} & Suy luận sai, phản ví dụ &
\protect\hyperlink{REC-03}{REC-03} & Vế phải \(\beta^nP_r(n)\) \\
\protect\hyperlink{QNT-01}{QNT-01} & Chân trị \(\forall,\exists\) &
\protect\hyperlink{REC-04}{REC-04} & Vế phải tổng nhiều dạng \\
\protect\hyperlink{QNT-02}{QNT-02} & Phủ định lượng từ &
\protect\hyperlink{REL-01}{REL-01} & Liệt kê, ma trận \(M_R\) \\
\protect\hyperlink{SET-01}{SET-01} & Đẳng thức tập hợp &
\protect\hyperlink{REL-02}{REL-02} & 4 tính chất quan hệ \\
\protect\hyperlink{SET-02}{SET-02} & \(P(X)\), \(A\times B\) &
\protect\hyperlink{REL-03}{REL-03} & Tương đương, lớp, đồng dư \\
\protect\hyperlink{FUN-01}{FUN-01} & Ảnh, ảnh ngược &
\protect\hyperlink{POS-01}{POS-01} & Poset, toàn phần \\
\protect\hyperlink{FUN-02}{FUN-02} & Đơn/toàn/song ánh, \(f^{-1}\) &
\protect\hyperlink{POS-02}{POS-02} & Thứ tự từ điển \\
\protect\hyperlink{FUN-03}{FUN-03} & \(g\circ f\) &
\protect\hyperlink{POS-03}{POS-03} & Hasse, tối đại/tối tiểu \\
\protect\hyperlink{IND-01}{IND-01} & Quy nạp &
\protect\hyperlink{POS-04}{POS-04} & Chặn, sup, inf \\
\protect\hyperlink{CNT-01}{CNT-01} & Cộng/nhân, lập số &
\protect\hyperlink{POS-05}{POS-05} & Sắp xếp topo \\
\protect\hyperlink{CNT-02}{CNT-02} & Chuồng bồ câu &
\protect\hyperlink{BOO-01}{BOO-01} & Hàm Bool, dạng chính tắc \\
\protect\hyperlink{CNT-03}{CNT-03} & Bù trừ &
\protect\hyperlink{BOO-02}{BOO-02} & Karnaugh ↔ đơn thức \\
\protect\hyperlink{CNT-04}{CNT-04} & \(P_n\), \(A_n^k\), \(C_n^k\) &
\protect\hyperlink{BOO-03}{BOO-03} & Đa thức tối tiểu \\
\protect\hyperlink{GRA-01}{GRA-01} & Bậc, bắt tay &
\protect\hyperlink{BOO-04}{BOO-04} & Mạch logic \\
\protect\hyperlink{GRA-02}{GRA-02} & Ma trận kề &
\protect\hyperlink{GRA-05}{GRA-05} & Euler, Fleury \\
\protect\hyperlink{GRA-03}{GRA-03} & Đẳng cấu &
\protect\hyperlink{GRA-06}{GRA-06} & Dijkstra \\
\protect\hyperlink{GRA-04}{GRA-04} & Liên thông, khớp, cầu &
\protect\hyperlink{GRA-07}{GRA-07} & Ford--Bellman \\
\end{longtable}

\hypertarget{cuxf4ng-thux1ee9c-trux1ecdng-yux1ebfu}{%
\section{Công thức trọng
yếu}\label{cuxf4ng-thux1ee9c-trux1ecdng-yux1ebfu}}

\begin{itemize}
\tightlist
\item
  \(p\to q\Leftrightarrow\neg p\vee q\Leftrightarrow\neg q\to\neg p\);
  \(\ \neg(p\to q)\Leftrightarrow p\wedge\neg q\); De Morgan
  \(\neg(p\wedge q)\Leftrightarrow\neg p\vee\neg q\).
\item
  \(\overline{\forall x\,p(x)}\Leftrightarrow\exists x\,\overline{p(x)}\);
  \(\ \exists x\forall y\Rightarrow\forall y\exists x\) (không ngược
  lại).
\item
  \(|P(X)|=2^{|X|}\), \(|A\times B|=|A||B|\);
  \(|A\cup B|=|A|+|B|-|A\cap B|\); 3 tập: \(+\)đơn \(-\)đôi \(+\)ba.
\item
  \(P_n=n!\), \(A_n^k=\frac{n!}{(n-k)!}\),
  \(C_n^k=\frac{n!}{k!(n-k)!}\), \(K_n^k=C_{n+k-1}^k\), hoán vị lặp
  \(\frac{n!}{n_1!\cdots n_k!}\); Dirichlet \(\lceil n/k\rceil\).
\item
  Đệ qui cấp 2: phân biệt \(C_1\lambda_1^n+C_2\lambda_2^n\); kép
  \((C_1+nC_2)\lambda_0^n\); nghiệm riêng \(n^m\beta^nQ_r(n)\), \(m\) =
  bội của \(\beta\).
\item
  Đồ thị: \(\sum\deg=2m\); có hướng \(\sum\deg^-=\sum\deg^+=m\); Euler:
  liên thông + 0 (chu trình) / 2 (đường đi) đỉnh lẻ.
\item
  Karnaugh 4 biến: cột \(x\bar y|xy|\bar xy|\bar x\bar y\), hàng
  \(z\bar t|zt|\bar zt|\bar z\bar t\); tế bào \(2^{n-k}\) ô ↔ đơn thức
  \(k\) từ đơn.
\end{itemize}

\hypertarget{quy-truxecnh-giux1ea3i-nhanh}{%
\section{Quy trình giải nhanh}\label{quy-truxecnh-giux1ea3i-nhanh}}

\begin{enumerate}
\def\labelenumi{\arabic{enumi}.}
\tightlist
\item
  \textbf{Đọc đề → gạch từ khóa} → tra
  \protect\hyperlink{keyword-index}{Mục lục từ khóa} → mã dạng bài.
\item
  Mở dạng bài trong \protect\hyperlink{taxonomy}{taxonomy}: đọc
  \textbf{Kiểm tra trước}.
\item
  Mở công thức (F-) và phương pháp (M-) được dẫn.
\item
  Mở ví dụ cùng mã, đối chiếu biến thể gần nhất với đề.
\item
  Giải theo các bước của phương pháp, trình bày theo mục ``Trình bày''.
\item
  Chạy \textbf{kiểm tra} (thế lại, tính \(x_2\), đếm bằng cách bù, thử
  phản ví dụ) và soát \protect\hyperlink{traps}{bẫy}.
\end{enumerate}

\hypertarget{kiux1ec3m-tra-trux1b0ux1edbc-khi-kux1ebft-luux1eadn}{%
\section{Kiểm tra trước khi kết
luận}\label{kiux1ec3m-tra-trux1b0ux1edbc-khi-kux1ebft-luux1eadn}}

\begin{itemize}
\tightlist
\item
  Logic: đã xét đủ \(2^n\) dòng / phản ví dụ đã \textbf{thử lại} mọi
  tiền đề? Lượng từ giữ đúng thứ tự và miền?
\item
  Ánh xạ: đã ghi rõ miền \(X\), tập đích \(Y\)? Nghiệm của \(f(x)=y\)
  thuộc \(X\)?
\item
  Quy nạp: \(N_0\) đúng đề? Bước quy nạp có dùng giả thiết?
\item
  Đếm: thứ tự? lặp? chữ số đầu \(\ne0\)? ràng buộc ``\(>\)'' đổi thành
  ``\(\ge\)'' đúng? kiểm tra bằng cách đếm phần bù?
\item
  Đệ qui: bội của \textbf{từng} \(\beta\)? \(C_1,C_2\) tính \textbf{sau}
  khi cộng nghiệm riêng? \(x_2\) từ hệ thức = từ công thức?
\item
  Quan hệ: ``không'' có phản ví dụ cụ thể? tương đương = PX+ĐX+BC; thứ
  tự = PX+phản xứng+BC.
\item
  Poset: Hasse chỉ có cung trực tiếp? sup/inf có thể không tồn tại?
\item
  Karnaugh: đủ tế bào lớn (quấn biên, 4 góc)? đã so sánh độ đơn giản các
  phủ?
\item
  Đồ thị: liên thông? trọng số âm (→ Ford--Bellman)? Dijkstra chỉ dừng
  khi đỉnh đích \textbf{được cố định}?
\end{itemize}

\hypertarget{lux1ed7i-dux1ec5-mux1ea5t-ux111iux1ec3m-nhux1ea5t}{%
\section{Lỗi dễ mất điểm
nhất}\label{lux1ed7i-dux1ec5-mux1ea5t-ux111iux1ec3m-nhux1ea5t}}

\protect\hyperlink{T-03}{T-03} đảo/phủ định kéo theo ·
\protect\hyperlink{T-06}{T-06} thứ tự lượng từ ·
\protect\hyperlink{T-08}{T-08} miền ánh xạ ·
\protect\hyperlink{T-10}{T-10} chữ số 0 · \protect\hyperlink{T-12}{T-12}
tổ hợp lặp · \protect\hyperlink{T-14}{T-14} thứ tự tìm hằng số ·
\protect\hyperlink{T-15}{T-15} bội của \(\beta\) ·
\protect\hyperlink{T-16}{T-16} chứng minh bằng ví dụ ·
\protect\hyperlink{T-18}{T-18} tối đại ≠ lớn nhất ·
\protect\hyperlink{T-20}{T-20} quấn biên Karnaugh ·
\protect\hyperlink{T-24}{T-24} đẳng cấu · \protect\hyperlink{T-26}{T-26}
Dijkstra với trọng số âm.

\hypertarget{phux1ee5-lux1ee5c-gap-conflict-report-khoux1ea3ng-trux1ed1ng-muxe2u-thuux1eabn-lux1ed7i-nguux1ed3n}{%
\chapter{Phụ lục --- Gap \& Conflict Report (Khoảng trống, mâu thuẫn,
lỗi
nguồn)}\label{phux1ee5-lux1ee5c-gap-conflict-report-khoux1ea3ng-trux1ed1ng-muxe2u-thuux1eabn-lux1ed7i-nguux1ed3n}}

Nhãn: \textbf{{[}SOURCE-VERIFIED{]}} có trong nguồn ·
\textbf{{[}CROSS-CHECKED{]}} kiểm tra độc lập · \textbf{{[}INFERRED{]}}
suy luận · \textbf{{[}UNVERIFIED{]}} chưa đủ căn cứ ·
\textbf{{[}CONFLICT{]}} mâu thuẫn.

\hypertarget{a.-khoux1ea3ng-trux1ed1ng-dux1eef-liux1ec7u-gap}{%
\section{A. Khoảng trống dữ liệu
(GAP)}\label{a.-khoux1ea3ng-trux1ed1ng-dux1eef-liux1ec7u-gap}}

\begin{longtable}[]{@{}
  >{\raggedright\arraybackslash}p{(\columnwidth - 8\tabcolsep) * \real{0.2000}}
  >{\raggedright\arraybackslash}p{(\columnwidth - 8\tabcolsep) * \real{0.2000}}
  >{\raggedright\arraybackslash}p{(\columnwidth - 8\tabcolsep) * \real{0.2000}}
  >{\raggedright\arraybackslash}p{(\columnwidth - 8\tabcolsep) * \real{0.2000}}
  >{\raggedright\arraybackslash}p{(\columnwidth - 8\tabcolsep) * \real{0.2000}}@{}}
\toprule\noalign{}
\begin{minipage}[b]{\linewidth}\raggedright
Mã
\end{minipage} & \begin{minipage}[b]{\linewidth}\raggedright
Khoảng trống
\end{minipage} & \begin{minipage}[b]{\linewidth}\raggedright
Ảnh hưởng
\end{minipage} & \begin{minipage}[b]{\linewidth}\raggedright
Cách xử lý hiện tại
\end{minipage} & \begin{minipage}[b]{\linewidth}\raggedright
Cần người dùng cung cấp
\end{minipage} \\
\midrule\noalign{}
\endhead
\bottomrule\noalign{}
\endlastfoot
GAP-01 & \textbf{Không có đề cương / phạm vi thi chính thức} & Không
khẳng định được chương nào thi, trọng số & Phạm vi = toàn bộ 5 chương
slide, nhãn \textbf{{[}INFERRED{]}} & Đề cương học phần hoặc thông báo
phạm vi thi \\
GAP-02 & \textbf{Không có đề thi cũ / bài tập giảng viên giao riêng} &
Không thống kê được tần suất dạng bài; dạng bài chỉ dựa vào ví dụ \&
``Bài tập'' trong slide & Không ghi tần suất; ví dụ bổ sung gắn
\textbf{{[}AI-GENERATED PRACTICE{]}} & Đề thi các năm trước (kèm đáp án
nếu có) \\
GAP-03 & Không có giáo trình (4 tài liệu tham khảo ở S1-p3 không được
cung cấp) & Một số định nghĩa chỉ có ở mức slide; không có số trang giáo
trình & Chỉ dẫn nguồn tới trang slide & (Tùy chọn) PDF giáo trình Nguyễn
Duy Nhất/Trần Ngọc Hội \\
GAP-04 & S2-p26 ``Hệ thức đệ qui'': slide chỉ có tiêu đề (nội dung không
trích được, không có ảnh) & Không ảnh hưởng nội dung chính & Bỏ qua &
--- \\
GAP-05 & S2 nói tới ``nghiệm phức'' không có; slide chỉ dạy PTĐT
\textbf{nghiệm thực phân biệt / nghiệm kép} (S2-p43) & Trường hợp nghiệm
phức \textbf{ngoài phạm vi slide} & Không đưa vào phần chính thức (ghi
chú ở Traps) & Xác nhận có thi nghiệm phức không \\
GAP-06 & S4-p79 chỉ dẫn link Quine--McCluskey, không có nội dung & Không
biết có thi Q--M không & Chỉ biên soạn phương pháp Karnaugh & Xác
nhận \\
GAP-07 & S5-p55 liệt kê ``Vét cạn, Dijkstra, Ford--Bellman,
\textbf{Floyd}'' nhưng slide chỉ trình bày Dijkstra \& Ford--Bellman &
Floyd có thể ngoài phạm vi & Không biên soạn Floyd & Xác nhận \\
GAP-08 & S5: không có cây, đồ thị Hamilton, tô màu, đồ thị phẳng,
BFS/DFS & --- & \textbf{Không} mở rộng phạm vi (Quy tắc 14) & --- \\
GAP-09 & S1 có ``Bài tập'' không lời giải (S1-p22, p50 lời giải ngắn,
p100 quy nạp; S2-p25 bài chia táo; S2-p67 6 hệ thức đệ qui; S5-p21, p35)
& --- & Đã giải đầy đủ trong Example Bank, gắn nhãn nguồn +
{[}CROSS-CHECKED{]} & --- \\
GAP-10 & Hình thức thi open-book không xuất hiện trong nguồn (chỉ có
trong S0) & --- & Theo S0 & --- \\
\end{longtable}

\hypertarget{b.-lux1ed7i-ux111iux1ec3m-cuxf3-khux1ea3-nux103ng-sai-trong-nguux1ed3n-khuxf4ng-sux1eeda-uxe2m-thux1ea7m}{%
\section{B. Lỗi / điểm có khả năng sai trong nguồn (không sửa âm
thầm)}\label{b.-lux1ed7i-ux111iux1ec3m-cuxf3-khux1ea3-nux103ng-sai-trong-nguux1ed3n-khuxf4ng-sux1eeda-uxe2m-thux1ea7m}}

\begin{longtable}[]{@{}
  >{\raggedright\arraybackslash}p{(\columnwidth - 8\tabcolsep) * \real{0.2000}}
  >{\raggedright\arraybackslash}p{(\columnwidth - 8\tabcolsep) * \real{0.2000}}
  >{\raggedright\arraybackslash}p{(\columnwidth - 8\tabcolsep) * \real{0.2000}}
  >{\raggedright\arraybackslash}p{(\columnwidth - 8\tabcolsep) * \real{0.2000}}
  >{\raggedright\arraybackslash}p{(\columnwidth - 8\tabcolsep) * \real{0.2000}}@{}}
\toprule\noalign{}
\begin{minipage}[b]{\linewidth}\raggedright
Mã
\end{minipage} & \begin{minipage}[b]{\linewidth}\raggedright
Vị trí
\end{minipage} & \begin{minipage}[b]{\linewidth}\raggedright
Nội dung nguồn
\end{minipage} & \begin{minipage}[b]{\linewidth}\raggedright
Vấn đề
\end{minipage} & \begin{minipage}[b]{\linewidth}\raggedright
Quyết định trong toolkit
\end{minipage} \\
\midrule\noalign{}
\endhead
\bottomrule\noalign{}
\endlastfoot
ERR-01 & S1-p94 & ``f(x)=2x+1 \ldots{} Khi đó \(f^{-1}(x)=(y-1)/2\)'' &
Lẫn biến: vế trái theo \(x\), vế phải theo \(y\) & Ghi đúng:
\(f^{-1}(y)=\frac{y-1}{2}\) (hay \(f^{-1}(x)=\frac{x-1}{2}\)).
\textbf{{[}CROSS-CHECKED{]}} \\
ERR-02 & S2-p29 & Định nghĩa hệ thức cấp k:
``\(a_0\neq 0, a_1,\dots,a_n\) là các hệ số thực'' & Chỉ số phải là
\(a_1,\dots,a_k\) (và để thực sự cấp \(k\) cần \(a_k\neq0\)) & Ghi
\(a_1,\dots,a_k\); nêu điều kiện \(a_k\ne0\) là
\textbf{{[}INFERRED{]}} \\
ERR-03 & S3-p25 & ``(i) a R b nếu \([a]_R=[b]_R\); (ii)
\([a]_R\ne[b]_R\) nếu \([a]_R\cap[b]_R=\emptyset\)'' & Phát biểu chỉ một
chiều; định lý chuẩn là ``⇔'' và (ii) phải là
``\([a]\ne[b] \Leftrightarrow [a]\cap[b]=\emptyset\)'' & Toolkit ghi
dạng ⇔, có chứng minh ngắn. \textbf{{[}CROSS-CHECKED{]}} \\
ERR-04 & S3-p37 & Thứ tự từ điển trên \(A\times B\):
``\((a_1,b_1)\preceq(a_2,b_2)\) nếu \(a_1<a_2\) hay (\(a_1=a_2\) và
\(b_1<' b_2\))'' & Thiếu trường hợp bằng nhau → định nghĩa như vậy là
thứ tự \textbf{chặt}, không phản xạ & Toolkit dùng ``\(b_1 \le' b_2\)''
và ghi chú. \textbf{{[}INFERRED{]}} \\
ERR-05 & S1-p85 (ví dụ 2 ánh xạ bằng nhau) & ``\(f(x)=x\) và
\(g(x)=(x+5)\bmod 5\) từ \([1,3]\to[1,3]\)'' & Nếu \([1,3]\) là đoạn số
thực, \((x+5)\bmod 5\) với số thực không chuẩn; nếu hiểu \(\{1,2,3\}\)
thì \(f=g\) & Không dùng làm ví dụ chính thức.
\textbf{{[}UNVERIFIED{]}} \\
ERR-06 & S4-p28 & Bản đồ Karnaugh ``Ví dụ 1'' vẽ 5 hàng & Lỗi vẽ (bản đồ
4 biến có 4 hàng) & Bỏ qua, không ảnh hưởng kết quả \\
ERR-07 & S4-p77 & ``Thiết kế mạch 3 cầu dao \ldots{} biến \(x,y\)
\ldots{} \(F(x,y)\)'' & Ghi nhầm \(F(x,y)\) thay vì \(F(x,y,z)\) &
Toolkit ghi \(F(x,y,z)\) \\
ERR-08 & S4-p74 & Sơ đồ cổng cho
\(f=(x\vee y\vee z)\,\bar x\bar y\bar z\) & Biểu thức đọc đúng, nhưng
\textbf{\(f\equiv 0\)} (hằng sai) --- slide không nêu & Đưa thành bẫy
\textbf{TRAP-BOO-04}. \textbf{{[}CROSS-CHECKED{]}} \\
ERR-09 & S2-p65 & ``\(f_n\) \ldots{} có dạng ở \textbf{Trường hợp 4}'' &
Slide chỉ định nghĩa ``Dạng 1, Dạng 2''; ``Trường hợp 4'' phải là
\textbf{Dạng 2} & Toolkit gọi là Dạng 2 \\
ERR-10 & S2-p40 & Ước lượng thời gian tháp Hà Nội dùng ``400 ngày/năm ×
25 giờ/ngày × 4000 giây/giờ'' & Không phải số thực tế; dùng làm
\textbf{chặn trên mẫu số} để có bất đẳng thức ``≥'' --- hợp lệ nhưng dễ
gây hiểu nhầm & Không đưa vào phần chính thức \\
ERR-11 & S1-p10, S1-p69--70, S3, S4 & Phủ định/bù được viết bằng gạch
trên (\(\bar p\)) và cả \(\neg p\) & Không phải lỗi; là \textbf{hai ký
hiệu tương đương} & Bảng ánh xạ ký hiệu trong Formula Handbook \\
ERR-13 & S1-p21 & Liệt kê \(C(p,q)=p\wedge\overline{(p\to q)}\) ngay sau
định nghĩa hằng sai & Như đã viết, \(C\Leftrightarrow p\wedge\neg q\)
\textbf{không} hằng sai (có thể ý định là
\(p\wedge\overline{(q\to p)}\)) & Toolkit nêu rõ trong LOG-02-E02
\textbf{{[}CROSS-CHECKED{]}} \\
ERR-12 & S1-p68 & Định lý hoán vị lượng từ ghi ``3) ∃x∀y ⇒ ∀y∃x'' &
Đúng; chiều ngược nói chung sai --- nguồn đã nêu & Dùng nguyên văn \\
\end{longtable}

\hypertarget{c.-muxe2u-thuux1eabn-quy-ux1b0ux1edbc-giux1eefa-cuxe1c-phux1ea7n-nguux1ed3n-conflict}{%
\section{C. Mâu thuẫn quy ước giữa các phần nguồn
(CONFLICT)}\label{c.-muxe2u-thuux1eabn-quy-ux1b0ux1edbc-giux1eefa-cuxe1c-phux1ea7n-nguux1ed3n-conflict}}

\begin{longtable}[]{@{}
  >{\raggedright\arraybackslash}p{(\columnwidth - 6\tabcolsep) * \real{0.2500}}
  >{\raggedright\arraybackslash}p{(\columnwidth - 6\tabcolsep) * \real{0.2500}}
  >{\raggedright\arraybackslash}p{(\columnwidth - 6\tabcolsep) * \real{0.2500}}
  >{\raggedright\arraybackslash}p{(\columnwidth - 6\tabcolsep) * \real{0.2500}}@{}}
\toprule\noalign{}
\begin{minipage}[b]{\linewidth}\raggedright
Mã
\end{minipage} & \begin{minipage}[b]{\linewidth}\raggedright
Vị trí
\end{minipage} & \begin{minipage}[b]{\linewidth}\raggedright
Mâu thuẫn
\end{minipage} & \begin{minipage}[b]{\linewidth}\raggedright
Quyết định
\end{minipage} \\
\midrule\noalign{}
\endhead
\bottomrule\noalign{}
\endlastfoot
CON-01 & S5-p26 vs S5-p28 & Định nghĩa: ``\(a_{ij}\) = số cạnh (cung) đi
từ \(i\) đến \(j\)'' --- theo định nghĩa này một khuyên tại \(e\) cho
\(a_{ee}=1\); nhưng ma trận ví dụ S5-p28 ghi \(a_{ee}=2\) &
\textbf{{[}CONFLICT{]}} Toolkit ghi cả hai quy ước, mặc định theo
\textbf{ví dụ S5-p28 (khuyên vô hướng đếm 2)} vì khớp với quy ước bậc
(khuyên đếm 2, S5-p19) và tổng hàng = bậc. \textbf{Cần xác nhận của
giảng viên} \\
CON-02 & S1-p58 vs thông lệ & \(\mathbb N=\{0,1,2,\dots\}\) (có số 0) &
Dùng theo nguồn: \(0\in\mathbb N\) \\
CON-03 & S2-p6 & Chuồng bồ câu dùng \(\lceil n/k\rceil\) --- ký hiệu bị
mất khi trích, đọc bằng ảnh & Theo nguồn \\
CON-04 & S5-p22 & ``Đỉnh bậc 0 gọi là đỉnh cô lập, đỉnh bậc 1 gọi là
đỉnh treo'' đặt trong mục đồ thị có hướng & Áp dụng cho cả vô hướng
(nghĩa phổ biến) \textbf{{[}INFERRED{]}} \\
CON-05 & S5-p49 & ``Nếu G có 2k đỉnh bậc lẻ thì vẽ được bằng k nét'' ---
đòi hỏi G liên thông, \(k\ge1\) (slide không nêu điều kiện) & Thêm điều
kiện liên thông \textbf{{[}INFERRED{]}} \\
CON-06 & S3-p33/S3-p30 & Ký hiệu thứ tự ``≼'' không hiển thị trong PDF
(ô trống) & Dùng \(\preceq\) \\
\end{longtable}

\hypertarget{d.-cuxe2u-hux1ecfi-cux1ea7n-ngux1b0ux1eddi-duxf9ng-xuxe1c-nhux1eadn-gom-mux1ed9t-lux1ea7n}{%
\section{D. Câu hỏi cần người dùng xác nhận (gom một
lần)}\label{d.-cuxe2u-hux1ecfi-cux1ea7n-ngux1b0ux1eddi-duxf9ng-xuxe1c-nhux1eadn-gom-mux1ed9t-lux1ea7n}}

\begin{enumerate}
\def\labelenumi{\arabic{enumi}.}
\tightlist
\item
  Có đề cương/phạm vi thi và đề thi cũ không? (GAP-01, GAP-02) --- nếu
  có, tôi sẽ thêm thống kê tần suất và điều chỉnh ưu tiên.
\item
  Có thi \textbf{nghiệm phức của PTĐT}, \textbf{Quine--McCluskey},
  \textbf{Floyd} không? (GAP-05/06/07)
\item
  Quy ước ma trận kề với khuyên: đếm 1 hay 2? (CON-01)
\end{enumerate}

Các phần không bị ảnh hưởng bởi các câu hỏi trên đã được biên soạn tiếp;
phần phụ thuộc được đánh dấu.

\hypertarget{phux1ee5-lux1ee5c-source-inventory-kiux1ec3m-kuxea-nguux1ed3n}{%
\chapter{Phụ lục --- Source Inventory (Kiểm kê
nguồn)}\label{phux1ee5-lux1ee5c-source-inventory-kiux1ec3m-kuxea-nguux1ed3n}}

Ngày kiểm kê: 2026-11-12. Tất cả vị trí nguồn trong bộ tài liệu dùng
dạng \texttt{S\textless{}k\textgreater{}-p\textless{}n\textgreater{}} =
\textbf{trang PDF thứ n} của tệp nguồn S (với hầu hết slide, trang PDF =
số slide in ở chân trang).

\begin{longtable}[]{@{}
  >{\raggedright\arraybackslash}p{(\columnwidth - 12\tabcolsep) * \real{0.1429}}
  >{\raggedright\arraybackslash}p{(\columnwidth - 12\tabcolsep) * \real{0.1429}}
  >{\raggedright\arraybackslash}p{(\columnwidth - 12\tabcolsep) * \real{0.1429}}
  >{\raggedright\arraybackslash}p{(\columnwidth - 12\tabcolsep) * \real{0.1429}}
  >{\raggedright\arraybackslash}p{(\columnwidth - 12\tabcolsep) * \real{0.1429}}
  >{\raggedright\arraybackslash}p{(\columnwidth - 12\tabcolsep) * \real{0.1429}}
  >{\raggedright\arraybackslash}p{(\columnwidth - 12\tabcolsep) * \real{0.1429}}@{}}
\toprule\noalign{}
\begin{minipage}[b]{\linewidth}\raggedright
Mã
\end{minipage} & \begin{minipage}[b]{\linewidth}\raggedright
Tệp gốc
\end{minipage} & \begin{minipage}[b]{\linewidth}\raggedright
Loại
\end{minipage} & \begin{minipage}[b]{\linewidth}\raggedright
Số trang
\end{minipage} & \begin{minipage}[b]{\linewidth}\raggedright
Nội dung chính
\end{minipage} & \begin{minipage}[b]{\linewidth}\raggedright
Chất lượng trích xuất
\end{minipage} & \begin{minipage}[b]{\linewidth}\raggedright
Vai trò
\end{minipage} \\
\midrule\noalign{}
\endhead
\bottomrule\noalign{}
\endlastfoot
S0 & Tệp prompt \emph{Claude Discrete Mathematics Open Book Toolkit
Prompt} (.md; 2 bản tải lên giống hệt nhau, 33 980 byte) & Project
Specification & --- & Kiến trúc sản phẩm, quy chuẩn, tiêu chí nghiệm thu
& Tốt & \textbf{Specification} --- không phải nguồn kiến thức \\
S1 & \texttt{01.pdf} --- ``CẤU TRÚC RỜI RẠC'', Phạm Thế Bảo (SGU),
Chương I & Slide bài giảng (PowerPoint, xuất PDF) & 100 & Mệnh đề, dạng
mệnh đề, luật logic, quy tắc suy diễn, vị từ--lượng từ, tập hợp, ánh xạ,
quy nạp & Tốt; \textbf{dấu gạch trên (phủ định/bù) mất khi trích văn
bản} → đã đối chiếu bằng ảnh các trang 21, 52--56, 69--70, 80 & Nguồn
kiến thức chính \\
S2 & \texttt{02.pdf} --- Chương II: Phép đếm & Slide & 67 & Nguyên lý
cộng, nhân, Dirichlet, bù trừ; hoán vị, chỉnh hợp, tổ hợp; hoán vị lặp,
tổ hợp lặp; hệ thức đệ qui tuyến tính & Tốt; trang 57--59 (ví dụ b)
\textbf{chỉ có ảnh} → đã đọc bằng ảnh; trang 51--66 dùng font VNI (mã
hóa lỗi khi trích) → đọc bằng ảnh khi cần & Nguồn kiến thức chính \\
S3 & \texttt{03.pdf} --- Chương 3: Quan hệ & Slide & 56 & Định nghĩa \&
tính chất quan hệ, ma trận biểu diễn, quan hệ tương đương, đồng dư, lớp
tương đương, quan hệ thứ tự, thứ tự từ điển, Hasse, tối đại/tối tiểu,
chặn trên/dưới, sup/inf, sắp xếp topo & Tốt; ký hiệu ``≼'' bị mất (hiển
thị trống); trang 30 có ảnh lỗi ``This image cannot currently be
displayed'' & Nguồn kiến thức chính \\
S4 & \texttt{04.pdf} --- Chương IV: Đại số Bool & Slide & 79 & Đại số
Bool, hàm Bool, dạng nối rời chính tắc, công thức đa thức tối tiểu, biểu
đồ Karnaugh (3--4 biến), tế bào lớn, thuật toán phủ, cổng logic, thiết
kế mạch & \textbf{Kém khi trích văn bản} (mất toàn bộ dấu bù; bản đồ
Karnaugh là hình) → toàn bộ ví dụ đã đọc lại bằng ảnh 90--200 dpi &
Nguồn kiến thức chính \\
S5 & \texttt{05.pdf} --- Chương 5: Đồ thị & Slide & 86 & Đồ thị vô/có
hướng, đơn/đa/giả đồ thị, bậc, định lý bắt tay, ma trận kề, đẳng cấu,
đường đi--chu trình--liên thông, đỉnh khớp, cầu, Euler (điều kiện,
Fleury), đường đi ngắn nhất (Dijkstra, Ford--Bellman) & Văn bản tốt;
\textbf{hình đồ thị chỉ có trong ảnh} → đã đọc bằng ảnh các trang
20--35, 42, 51, 60, 71, 75, 84, 86 & Nguồn kiến thức chính \\
\end{longtable}

\hypertarget{phuxe2n-loux1ea1i-nguux1ed3n}{%
\section{Phân loại nguồn}\label{phuxe2n-loux1ea1i-nguux1ed3n}}

\begin{itemize}
\tightlist
\item
  \textbf{Nguồn chính thức (theo thứ tự ưu tiên của S0 §5.3):} chưa có
  đề cương/phạm vi thi (mức 1), chưa có giáo trình (mức 2). Nguồn cao
  nhất hiện có là \textbf{slide giảng dạy S1--S5 (mức 3)}.
\item
  \textbf{Tài liệu tham khảo được slide nhắc tới (S1-p3) nhưng KHÔNG
  được cung cấp:} Nguyễn Duy Nhất -- Nguyễn Văn Phong -- Đinh Ngọc
  Thanh, \emph{Toán rời rạc}; Trần Ngọc Hội, \emph{Toán rời rạc}; Nguyễn
  Hữu Anh, \emph{Toán rời rạc} (NXB Giáo dục); Rosen, \emph{Discrete
  Mathematics and Its Applications}, 6th ed., 2007.
\item
  \textbf{Trùng lặp:} chỉ S0 (2 bản giống hệt). S1--S5 không trùng nhau.
\item
  \textbf{Thông tin đánh giá (S1-p4):} Quá trình 30\%, cuối kỳ 70\%,
  điểm thưởng 5--10\%. Không có thông tin hình thức thi (open-book được
  giả định từ S0).
\end{itemize}

\hypertarget{trang-chux1ec9-cuxf3-ux1ea3nh-khuxf3-truxedch-ux111uxe3-xux1eed-luxfd-bux1eb1ng-ux111ux1ecdc-ux1ea3nh}{%
\section{Trang chỉ có ảnh / khó trích (đã xử lý bằng đọc
ảnh)}\label{trang-chux1ec9-cuxf3-ux1ea3nh-khuxf3-truxedch-ux111uxe3-xux1eed-luxfd-bux1eb1ng-ux111ux1ecdc-ux1ea3nh}}

S1: 21, 52--56, 73, 76, 84. S2: 26--27, 40, 57--59. S3: 2, 4--5, 16--20,
29, 39, 45. S4: 9, 19, 22--37, 42--49, 53--65, 72--78. S5: 4, 6, 8, 13,
15--16, 20--35, 37, 39, 42--47, 51, 54, 60--86 (hình trong bảng).

\hypertarget{cuxf4ng-cux1ee5-ux111uxe3-duxf9ng-ux111ux1ec3-ux111ux1ecdc-nguux1ed3n}{%
\section{Công cụ đã dùng để đọc
nguồn}\label{cuxf4ng-cux1ee5-ux111uxe3-duxf9ng-ux111ux1ec3-ux111ux1ecdc-nguux1ed3n}}

\texttt{pdftotext} (văn bản), \texttt{pdftoppm} (ảnh 60--200 dpi) và đọc
ảnh trực tiếp. Mọi kết quả số/logic trong slide đã được kiểm tra lại
bằng \texttt{10\_Verification/verify\_source\_claims.py} (42/42 đạt).

\backmatter
\end{document}
